% Options for packages loaded elsewhere
\PassOptionsToPackage{unicode}{hyperref}
\PassOptionsToPackage{hyphens}{url}
\documentclass[
  11pt,
  openany]{report}
\usepackage{xcolor}
\usepackage[margin=1in]{geometry}
\usepackage{amsmath,amssymb}
\setcounter{secnumdepth}{5}
\usepackage{iftex}
\ifPDFTeX
  \usepackage[T1]{fontenc}
  \usepackage[utf8]{inputenc}
  \usepackage{textcomp} % provide euro and other symbols
\else % if luatex or xetex
  \usepackage{unicode-math} % this also loads fontspec
  \defaultfontfeatures{Scale=MatchLowercase}
  \defaultfontfeatures[\rmfamily]{Ligatures=TeX,Scale=1}
\fi
\usepackage{lmodern}
\ifPDFTeX\else
  % xetex/luatex font selection
    \setmainfont[]{STIX Two Text}
    \setmonofont[]{Consolas}
  \setmathfont[]{STIX Two Math}
\fi
% Use upquote if available, for straight quotes in verbatim environments
\IfFileExists{upquote.sty}{\usepackage{upquote}}{}
\IfFileExists{microtype.sty}{% use microtype if available
  \usepackage[]{microtype}
  \UseMicrotypeSet[protrusion]{basicmath} % disable protrusion for tt fonts
}{}
\usepackage{setspace}
\makeatletter
\@ifundefined{KOMAClassName}{% if non-KOMA class
  \IfFileExists{parskip.sty}{%
    \usepackage{parskip}
  }{% else
    \setlength{\parindent}{0pt}
    \setlength{\parskip}{6pt plus 2pt minus 1pt}}
}{% if KOMA class
  \KOMAoptions{parskip=half}}
\makeatother
\usepackage{color}
\usepackage{fancyvrb}
\newcommand{\VerbBar}{|}
\newcommand{\VERB}{\Verb[commandchars=\\\{\}]}
\DefineVerbatimEnvironment{Highlighting}{Verbatim}{commandchars=\\\{\}}
% Add ',fontsize=\small' for more characters per line
\usepackage{framed}
\definecolor{shadecolor}{RGB}{24,28,46}
\newenvironment{Shaded}{\begin{snugshade}}{\end{snugshade}}
\newcommand{\AlertTok}[1]{\textcolor[rgb]{1.00,0.54,0.54}{\textbf{#1}}}
\newcommand{\AnnotationTok}[1]{\textcolor[rgb]{0.50,0.54,0.62}{\textit{#1}}}
\newcommand{\AttributeTok}[1]{\textcolor[rgb]{0.75,0.60,1.00}{#1}}
\newcommand{\BaseNTok}[1]{\textcolor[rgb]{0.91,0.66,0.43}{#1}}
\newcommand{\BuiltInTok}[1]{\textcolor[rgb]{0.49,0.78,0.89}{#1}}
\newcommand{\CharTok}[1]{\textcolor[rgb]{0.66,0.83,0.63}{#1}}
\newcommand{\CommentTok}[1]{\textcolor[rgb]{0.50,0.54,0.62}{\textit{#1}}}
\newcommand{\CommentVarTok}[1]{\textcolor[rgb]{0.50,0.54,0.62}{\textit{#1}}}
\newcommand{\ConstantTok}[1]{\textcolor[rgb]{0.91,0.66,0.43}{#1}}
\newcommand{\ControlFlowTok}[1]{\textcolor[rgb]{0.62,0.65,1.00}{\textbf{#1}}}
\newcommand{\DataTypeTok}[1]{\textcolor[rgb]{0.75,0.60,1.00}{#1}}
\newcommand{\DecValTok}[1]{\textcolor[rgb]{0.91,0.66,0.43}{#1}}
\newcommand{\DocumentationTok}[1]{\textcolor[rgb]{0.50,0.54,0.62}{\textit{#1}}}
\newcommand{\ErrorTok}[1]{\textcolor[rgb]{1.00,0.54,0.54}{\textbf{#1}}}
\newcommand{\ExtensionTok}[1]{\textcolor[rgb]{0.49,0.78,0.89}{#1}}
\newcommand{\FloatTok}[1]{\textcolor[rgb]{0.91,0.66,0.43}{#1}}
\newcommand{\FunctionTok}[1]{\textcolor[rgb]{0.49,0.78,0.89}{#1}}
\newcommand{\ImportTok}[1]{\textcolor[rgb]{0.66,0.83,0.63}{#1}}
\newcommand{\InformationTok}[1]{\textcolor[rgb]{0.49,0.78,0.89}{#1}}
\newcommand{\KeywordTok}[1]{\textcolor[rgb]{0.62,0.65,1.00}{\textbf{#1}}}
\newcommand{\NormalTok}[1]{\textcolor[rgb]{0.78,0.80,0.85}{#1}}
\newcommand{\OperatorTok}[1]{\textcolor[rgb]{0.78,0.80,0.85}{#1}}
\newcommand{\OtherTok}[1]{\textcolor[rgb]{0.49,0.78,0.89}{#1}}
\newcommand{\PreprocessorTok}[1]{\textcolor[rgb]{0.62,0.65,1.00}{#1}}
\newcommand{\RegionMarkerTok}[1]{\textcolor[rgb]{0.50,0.54,0.62}{#1}}
\newcommand{\SpecialCharTok}[1]{\textcolor[rgb]{0.91,0.66,0.43}{#1}}
\newcommand{\SpecialStringTok}[1]{\textcolor[rgb]{0.66,0.83,0.63}{#1}}
\newcommand{\StringTok}[1]{\textcolor[rgb]{0.66,0.83,0.63}{#1}}
\newcommand{\VariableTok}[1]{\textcolor[rgb]{0.78,0.80,0.85}{#1}}
\newcommand{\VerbatimStringTok}[1]{\textcolor[rgb]{0.66,0.83,0.63}{#1}}
\newcommand{\WarningTok}[1]{\textcolor[rgb]{0.91,0.66,0.43}{#1}}
\usepackage{longtable,booktabs,array}
\usepackage{calc} % for calculating minipage widths
% Correct order of tables after \paragraph or \subparagraph
\usepackage{etoolbox}
\makeatletter
\patchcmd\longtable{\par}{\if@noskipsec\mbox{}\fi\par}{}{}
\makeatother
% Allow footnotes in longtable head/foot
\IfFileExists{footnotehyper.sty}{\usepackage{footnotehyper}}{\usepackage{footnote}}
\makesavenoteenv{longtable}
\setlength{\emergencystretch}{3em} % prevent overfull lines
\providecommand{\tightlist}{%
  \setlength{\itemsep}{0pt}\setlength{\parskip}{0pt}}
% ═══════════════════════════════════════════════════════════════════════════
%  The Seven-Layer Magic Trick — LaTeX Preamble
%  Design: "Midnight Network" — Dark-Mode Identity
% ═══════════════════════════════════════════════════════════════════════════

% ── Section Numbering ────────────────────────────────────────────────────
% Number chapters (0) and sections (1) in the TOC and body text.
% Subsections (2) are unnumbered. This matches --number-sections in pandoc.
\setcounter{secnumdepth}{2}

% ── Packages ─────────────────────────────────────────────────────────────
% microtype may already be loaded by pandoc; pass options safely
\PassOptionsToPackage{protrusion=true,tracking=true}{microtype}
\usepackage{xcolor}
\usepackage{tikz}
\usepackage{tikzpagenodes}
\usepackage{titlesec}
\usepackage{titletoc}
\usepackage{fancyhdr}
\usepackage{enumitem}
\usepackage{longtable,booktabs,array}
\usepackage{colortbl}
\usepackage[most]{tcolorbox}
\usepackage{etoolbox}
\usepackage{lettrine}
\usepackage{multicol}
\usepackage{graphicx}
\usepackage{needspace}

% ── Font Setup (Outfit from project assets) ─────────────────────────────
% Outfit is loaded from the assets/ directory so it works without
% system-wide font installation. fontspec is loaded by pandoc; we
% override the sansfont here with an explicit path.
\AtBeginDocument{%
  \setsansfont{Outfit}[
    Path = assets/,
    Extension = .ttf,
    UprightFont = *-Regular,
    BoldFont = *-Bold,
    FontFace = {sb}{n}{*-SemiBold},
    FontFace = {m}{n}{*-Medium},
  ]%
}

% ── Micro-typography ─────────────────────────────────────────────────────
\frenchspacing
\widowpenalty=10000
\clubpenalty=10000
\raggedbottom
\hyphenpenalty=300
\exhyphenpenalty=300
\tolerance=1414
\hbadness=1414
\emergencystretch=2em
\hfuzz=0.3pt
\setlength{\parskip}{2pt plus 1pt minus 0.5pt}

% Lettrine defaults
\setcounter{DefaultLines}{3}
\renewcommand{\DefaultLoversize}{0.1}
\renewcommand{\DefaultLhang}{0}
\setlength{\DefaultFindent}{3pt}
\setlength{\DefaultNindent}{0pt}

% Table spacing
\renewcommand{\arraystretch}{1.3}
\setlength{\tabcolsep}{8pt}

% ═══════════════════════════════════════════════════════════════════════════
%  COLOR DEFINITIONS — Midnight Dark Mode
% ═══════════════════════════════════════════════════════════════════════════

% Print mode toggle (documented, not active)
\newif\ifprintmode
\printmodefalse

% ── Brand Colors ───────────────────────────────────────────────────────
\definecolor{MidnightBg}{HTML}{0b0e1a}
\definecolor{MidnightText}{HTML}{DEDEDE}        % 87% white (Material "high emphasis")
\definecolor{MidnightBold}{HTML}{F5F5F5}       % 95% white (bold brightness delta)
\definecolor{MidnightWhite}{HTML}{ffffff}
\definecolor{MidnightHeaderBg}{HTML}{1e2236}   % Table header row background
\definecolor{MidnightBlue}{HTML}{0000fe}
\definecolor{MidnightPurple}{HTML}{8756ff}
\definecolor{MidnightMuted}{HTML}{a0a4b8}
\definecolor{MidnightSurface}{HTML}{141726}
\definecolor{MidnightBorder}{HTML}{252940}
\definecolor{MidnightGrey}{HTML}{6b7084}
\definecolor{MidnightStripe}{HTML}{101320}
\definecolor{MidnightCode}{HTML}{181c2e}

% ── Semantic Box Colors ────────────────────────────────────────────────
\definecolor{BoxTechnical}{HTML}{3b82f6}
\definecolor{BoxTechnicalBg}{HTML}{0f1a2e}
\definecolor{BoxInsight}{HTML}{8756ff}
\definecolor{BoxInsightBg}{HTML}{150f2e}
\definecolor{BoxCaution}{HTML}{f59e0b}
\definecolor{BoxCautionBg}{HTML}{1a160e}
\definecolor{BoxExample}{HTML}{0000fe}
\definecolor{BoxExampleBg}{HTML}{0f0e20}

% ═══════════════════════════════════════════════════════════════════════════
%  DARK PAGE BACKGROUND
% ═══════════════════════════════════════════════════════════════════════════

% Paint every page dark.
% \pagecolor ensures even blank/empty pages (from \clearpage) are dark.
% The TikZ hook provides a second layer for reliable coverage.
\pagecolor{MidnightBg}
\AddToHook{shipout/background}{%
  \begin{tikzpicture}[remember picture, overlay]
    \fill[MidnightBg] (current page.south west) rectangle (current page.north east);
  \end{tikzpicture}%
}
\AtBeginDocument{\color{MidnightText}}

% ═══════════════════════════════════════════════════════════════════════════
%  GRADIENT RULE HELPER
% ═══════════════════════════════════════════════════════════════════════════

\newcommand{\midnightgradientrule}[2]{% #1=width fraction, #2=height
  \begin{tikzpicture}
    \shade[left color=MidnightBlue, right color=MidnightPurple]
      (0,0) rectangle (#1\textwidth, #2);
  \end{tikzpicture}%
}

% ═══════════════════════════════════════════════════════════════════════════
%  HEADING HIERARCHY
% ═══════════════════════════════════════════════════════════════════════════

% Chapter (H1): Gradient rule above + border rule + large bold white title
\titleformat{\chapter}[display]
  {\normalfont\sffamily}
  {}
  {0pt}
  {\vspace*{-20pt}%
   \midnightgradientrule{1}{3pt}\\[8pt]%
   {\color{MidnightBorder}\rule{\textwidth}{0.5pt}}\\[2.5cm]%
   \fontsize{32pt}{38pt}\selectfont\bfseries\color{MidnightWhite}}
\titlespacing*{\chapter}{0pt}{0pt}{30pt}

% Section (H2): Centered bold sans, section number in blue, blue underrule
\titleformat{\section}[block]
  {\needspace{5\baselineskip}\normalfont\sffamily\fontsize{16pt}{20pt}\selectfont\bfseries\filcenter\color{MidnightWhite}}
  {{\color{MidnightBlue}\thesection}\enspace}
  {0pt}
  {}
  [\vspace{2pt}{\color{MidnightBlue}\titlerule[0.75pt]}]
\titlespacing*{\section}{0pt}{36pt plus 6pt minus 4pt}{12pt plus 2pt}

% Subsection (H3): Centered bold sans, muted
\titleformat{\subsection}[block]
  {\needspace{4\baselineskip}\normalfont\sffamily\fontsize{13pt}{17pt}\selectfont\bfseries\filcenter\color{MidnightMuted}}
  {\thesubsection\enspace}
  {0pt}
  {}
\titlespacing*{\subsection}{0pt}{18pt plus 3pt minus 1pt}{6pt plus 1pt}

% Sub-subsection (H4): Bold italic sans, muted
\titleformat{\subsubsection}
  {\needspace{3\baselineskip}\normalfont\sffamily\fontsize{11pt}{15pt}\selectfont\bfseries\itshape\color{MidnightMuted}}
  {}
  {0pt}
  {}
\titlespacing*{\subsubsection}{0pt}{14pt plus 2pt minus 1pt}{4pt plus 1pt}

% ═══════════════════════════════════════════════════════════════════════════
%  PART DIVIDER
% ═══════════════════════════════════════════════════════════════════════════

\newcommand{\partdivider}[1]{%
  \clearpage\thispagestyle{empty}%
  \vspace*{\fill}%
  \begin{center}
  \midnightgradientrule{0.4}{3pt}\\[2cm]
  {\fontsize{28pt}{34pt}\selectfont\sffamily\bfseries\color{MidnightWhite}#1}\\[2cm]
  \midnightgradientrule{0.4}{3pt}
  \end{center}
  \vspace*{\fill}\clearpage
}

% ═══════════════════════════════════════════════════════════════════════════
%  TABLE OF CONTENTS STYLING
% ═══════════════════════════════════════════════════════════════════════════

\titlecontents{chapter}
  [0pt]
  {\addvspace{12pt}\bfseries\large\sffamily\color{MidnightBlue}}
  {}
  {}
  {\titlerule*[8pt]{}\contentspage}
  [\addvspace{2pt}]

\titlecontents{section}
  [18pt]
  {\addvspace{2pt}\normalsize\sffamily\color{MidnightText}}
  {{\color{MidnightBlue}\thecontentslabel}\enspace}
  {}
  {{\color{MidnightGrey}\titlerule*[8pt]{.}}\contentspage}

\titlecontents{subsection}
  [36pt]
  {\small\color{MidnightGrey}}
  {\thecontentslabel\enspace}
  {}
  {\titlerule*[8pt]{.}\contentspage}

% ═══════════════════════════════════════════════════════════════════════════
%  RUNNING HEADERS & FOOTERS
% ═══════════════════════════════════════════════════════════════════════════

\pagestyle{fancy}
\fancyhf{}
\fancyhead[L]{\small\sffamily\color{MidnightGrey}\nouppercase{\leftmark}}
\fancyhead[R]{\small\sffamily\color{MidnightGrey}\thepage}
\fancyfoot[C]{\scriptsize\color{MidnightGrey}\textnormal{Midnight Network}}
\renewcommand{\headrulewidth}{0.5pt}
\renewcommand{\headrule}{\hbox to\headwidth{%
  \color{MidnightBlue}\leaders\hrule height \headrulewidth\hfill}}
\renewcommand{\footrulewidth}{0pt}

% Plain pages (chapter openings): just page number
\fancypagestyle{plain}{%
  \fancyhf{}%
  \fancyfoot[C]{\small\sffamily\color{MidnightGrey}\thepage}%
  \renewcommand{\headrulewidth}{0pt}%
}

% ═══════════════════════════════════════════════════════════════════════════
%  TABLE STYLING
% ═══════════════════════════════════════════════════════════════════════════

% Alternating row colors (applies to longtable)
\AtBeginEnvironment{longtable}{%
  \rowcolors{2}{MidnightBg}{MidnightStripe}%
  \arrayrulecolor{MidnightBorder}%
}

% ═══════════════════════════════════════════════════════════════════════════
%  BULLET LIST STYLING
% ═══════════════════════════════════════════════════════════════════════════

\setlist[itemize,1]{label={\color{MidnightBlue}\textbullet}, itemsep=3pt, parsep=1pt}
\setlist[itemize,2]{label={\color{MidnightPurple}\textendash}, itemsep=2pt, parsep=0pt}
\setlist[itemize,3]{label={\color{MidnightGrey}\textperiodcentered}, itemsep=1pt, parsep=0pt}

% ═══════════════════════════════════════════════════════════════════════════
%  PULL QUOTE & EPIGRAPH
% ═══════════════════════════════════════════════════════════════════════════

% Pull quote command
\newcommand{\pullquote}[1]{%
  \par\vspace{18pt}%
  \begin{center}
  \begin{minipage}{0.75\textwidth}
  \centering
  \midnightgradientrule{0.15}{1.5pt}\\[8pt]
  {\fontsize{14pt}{20pt}\selectfont\itshape\color{MidnightPurple} #1}\\[8pt]
  \midnightgradientrule{0.15}{1.5pt}
  \end{minipage}
  \end{center}
  \par\vspace{18pt}%
}

% Epigraph command for chapter-opening quotes
\newcommand{\zkepigraph}[2]{%
  \vspace{6pt}%
  \begin{center}
  \begin{minipage}{0.78\textwidth}
  \raggedleft
  {\fontsize{11pt}{16pt}\selectfont\itshape\color{MidnightMuted}#1}\\[4pt]
  {\footnotesize\color{MidnightGrey}--- #2}
  \end{minipage}
  \end{center}
  \vspace{12pt}%
}

% ═══════════════════════════════════════════════════════════════════════════
%  CONTENT-TYPE BOXES (tcolorbox)
% ═══════════════════════════════════════════════════════════════════════════

% Technical Detail box
\newtcolorbox{technicalbox}[1][Technical Detail]{%
  enhanced,
  breakable,
  colback=BoxTechnicalBg,
  colframe=BoxTechnical,
  coltext=MidnightText,
  fonttitle=\sffamily\bfseries\small,
  title={#1},
  coltitle=white,
  leftrule=4pt,
  toprule=0pt,
  bottomrule=0pt,
  rightrule=0pt,
  arc=0pt,
  outer arc=0pt,
  left=10pt, right=10pt, top=8pt, bottom=8pt,
  attach boxed title to top left={yshift=-\tcboxedtitleheight/2, xshift=10pt},
  boxed title style={colback=BoxTechnical, arc=2pt, boxrule=0pt},
}

% Insight box (replaces patientbox)
\newtcolorbox{insightbox}[1][Key Insight]{%
  enhanced,
  breakable,
  colback=BoxInsightBg,
  colframe=BoxInsight,
  coltext=MidnightText,
  fonttitle=\sffamily\bfseries\small,
  title={#1},
  coltitle=white,
  leftrule=4pt,
  toprule=0pt,
  bottomrule=0pt,
  rightrule=0pt,
  arc=0pt,
  outer arc=0pt,
  left=10pt, right=10pt, top=8pt, bottom=8pt,
  attach boxed title to top left={yshift=-\tcboxedtitleheight/2, xshift=10pt},
  boxed title style={colback=BoxInsight, arc=2pt, boxrule=0pt},
}

% Caution box
\newtcolorbox{cautionbox}[1][Important Caveat]{%
  enhanced,
  breakable,
  colback=BoxCautionBg,
  colframe=BoxCaution,
  coltext=MidnightText,
  fonttitle=\sffamily\bfseries\small,
  title={#1},
  coltitle=MidnightBg,
  leftrule=4pt,
  toprule=0pt,
  bottomrule=0pt,
  rightrule=0pt,
  arc=0pt,
  outer arc=0pt,
  left=10pt, right=10pt, top=8pt, bottom=8pt,
  attach boxed title to top left={yshift=-\tcboxedtitleheight/2, xshift=10pt},
  boxed title style={colback=BoxCaution, arc=2pt, boxrule=0pt},
}

% Example box (replaces regulatorybox)
\newtcolorbox{examplebox}[1][Example]{%
  enhanced,
  breakable,
  colback=BoxExampleBg,
  colframe=BoxExample,
  coltext=MidnightText,
  fonttitle=\sffamily\bfseries\small,
  title={#1},
  coltitle=white,
  leftrule=4pt,
  toprule=0pt,
  bottomrule=0pt,
  rightrule=0pt,
  arc=0pt,
  outer arc=0pt,
  left=10pt, right=10pt, top=8pt, bottom=8pt,
  attach boxed title to top left={yshift=-\tcboxedtitleheight/2, xshift=10pt},
  boxed title style={colback=BoxExample, arc=2pt, boxrule=0pt},
}

% ═══════════════════════════════════════════════════════════════════════════
%  MATH CONVENIENCES
% ═══════════════════════════════════════════════════════════════════════════

% Display math spacing (generous for dark backgrounds)
\setlength{\abovedisplayskip}{13pt plus 3pt minus 3pt}
\setlength{\belowdisplayskip}{13pt plus 3pt minus 3pt}
\setlength{\abovedisplayshortskip}{7pt plus 3pt}
\setlength{\belowdisplayshortskip}{7pt plus 3pt minus 3pt}

% ═══════════════════════════════════════════════════════════════════════════
%  BOLD TEXT — Brightness delta for dark-mode emphasis
% ═══════════════════════════════════════════════════════════════════════════

% On dark backgrounds, bold at the same color as body text is barely visible.
% Material Design / Apple HIG / GitHub Primer all use a brightness delta:
% body at 87% white, bold at 95% white. Color is reserved for interactive elements.
\let\oldtextbf\textbf
\renewcommand{\textbf}[1]{{\bfseries\color{MidnightBold}#1}}

% ═══════════════════════════════════════════════════════════════════════════
%  CODE BLOCKS — Styled container for pandoc's Shaded environment
% ═══════════════════════════════════════════════════════════════════════════

% Pandoc generates \begin{Shaded}\begin{Highlighting}...\end{Highlighting}\end{Shaded}
% for fenced code blocks. We override Shaded with a styled tcolorbox.
\renewenvironment{Shaded}{%
  \begin{tcolorbox}[
    colback=MidnightCode,
    colframe=MidnightBorder,
    arc=4pt,
    boxrule=0.5pt,
    left=8pt, right=8pt, top=6pt, bottom=6pt,
    breakable, enhanced
  ]%
}{\end{tcolorbox}}

% ═══════════════════════════════════════════════════════════════════════════
%  HYPERLINKS (deferred — hyperref is loaded later by pandoc)
% ═══════════════════════════════════════════════════════════════════════════

\AtBeginDocument{%
  \hypersetup{
    colorlinks=true,
    linkcolor=MidnightPurple,
    urlcolor=MidnightPurple,
    citecolor=MidnightPurple,
  }%
}
\usepackage{bookmark}
\IfFileExists{xurl.sty}{\usepackage{xurl}}{} % add URL line breaks if available
\urlstyle{same}
\hypersetup{
  hidelinks,
  pdfcreator={LaTeX via pandoc}}

\author{}
\date{}

\begin{document}

% The Seven-Layer Magic Trick — Midnight Network Branded Title Page
% Design: Midnight Dark Mode

\begin{titlepage}
\newgeometry{margin=1in}

% Background is handled by the global shipout hook (MidnightBg)

% Top gradient rule — 3pt blue-to-purple, full width
\noindent\midnightgradientrule{1}{3pt}

\vspace{3cm}

\begin{center}

% Midnight logo — text-based rendering
% (SVG logo available in assets/ for future vector inclusion)
{\fontsize{28pt}{34pt}\selectfont\sffamily\bfseries\color{MidnightWhite}%
\textls[100]{MIDNIGHT}}

\vspace{2cm}

% Main title — large, bold, white
{\fontsize{32pt}{38pt}\selectfont\sffamily\color{MidnightWhite}%
\textbf{The Seven-Layer\\Magic Trick}}

\vspace{1cm}

% Subtitle — medium, muted
{\fontsize{14pt}{20pt}\selectfont\sffamily\color{MidnightMuted}%
A Complete Guide to\\
Zero-Knowledge Proof Systems}

\vspace{1.2cm}

% Decorative gradient rule — 40% width
\midnightgradientrule{0.4}{1.5pt}

\vspace{0.3cm}

% Edition label — letterspaced, grey
{\fontsize{9pt}{12pt}\selectfont\sffamily\color{MidnightGrey}%
\textls[150]{FIRST EDITION}}

\vspace{0.8cm}

% Author
{\fontsize{14pt}{18pt}\selectfont\sffamily\color{MidnightText}%
Charles Hoskinson}

\vspace{0.8cm}

% Date
{\fontsize{12pt}{16pt}\selectfont\sffamily\color{MidnightMuted}%
March 2026}

\vspace{1.5cm}

% Organization name — letterspaced
{\fontsize{10pt}{14pt}\selectfont\sffamily\color{MidnightGrey}%
\textls[200]{MIDNIGHT NETWORK}}

\end{center}

\restoregeometry
\end{titlepage}

% The Seven-Layer Magic Trick — Copyright and License Page
% Placed on verso of title page (standard book convention)

\newpage
\thispagestyle{empty}

\vspace*{\fill}

\begin{center}
\begin{minipage}{0.75\textwidth}

{\small\color{MidnightMuted}

\noindent Copyright \textcopyright\ 2026 Charles Hoskinson.\\[1.5em]

\noindent This work is licensed under the Creative Commons Attribution 4.0
International License (CC BY 4.0).\\[1em]

\noindent You are free to share and adapt this material for any purpose,
including commercial use, as long as you give appropriate credit, provide
a link to the license, and indicate if changes were made.\\[1em]

\noindent Full license text:\\
{\color{MidnightPurple}\url{https://creativecommons.org/licenses/by/4.0/legalcode}}\\[2em]

\noindent First edition, March 2026.\\[0.5em]
\noindent Published by Midnight Network.\\[2em]

{\color{MidnightGrey}\rule{0.3\textwidth}{0.5pt}}\\[1em]

\noindent\textit{The Seven-Layer Magic Trick: A Complete Guide to Zero-Knowledge Proof Systems.}

}

\end{minipage}
\end{center}

\vspace*{\fill}

\newpage

\renewcommand*\contentsname{Table of Contents}
{
\setcounter{tocdepth}{1}
\tableofcontents
}
\setstretch{1.4}
title: ``The Seven-Layer Magic Trick'' subtitle: ``A Complete Guide to
Zero-Knowledge Proof Systems'' author: ``Charles Hoskinson'' edition: ``First
Edition'' date: ``March 2026'' ---

\chapter*{Glossary of Key Terms}\label{glossary-of-key-terms}
\addcontentsline{toc}{chapter}{Glossary of Key Terms}

\textbf{AIR (Algebraic Intermediate Representation)} -- A way of encoding
computation as a table of numbers governed by polynomial rules that must hold
between consecutive rows. Used by STARK proof systems.

\textbf{Arithmetization} -- The conversion of a computation into a system of
polynomial equations that can be verified mathematically.

\textbf{BabyBear} -- A 31-bit prime field (\(2^{31} - 2^{27} + 1\)) used by SP1
and Plonky3. Enables SIMD-friendly arithmetic on modern CPUs.

\textbf{BLS12-381} -- A specific elliptic curve providing approximately 128-bit
security. Used by Zcash, Ethereum's blob commitments, and Midnight. Vulnerable
to quantum computers.

\textbf{BN254 (alt\_bn128)} -- An older elliptic curve hardcoded into Ethereum's
built-in operations. Its estimated security has eroded from 128 bits to
approximately 100 bits.

\textbf{CCS (Customizable Constraint Systems)} -- A unified mathematical
framework that includes R1CS, AIR, and PLONKish as special cases.

\textbf{Circle STARKs} -- A STARK variant operating over the circle group of a
Mersenne prime field, enabling efficient FFTs without requiring smooth-order
multiplicative subgroups.

\textbf{Circuit} -- A representation of a computation as a network of
mathematical operations over a finite field. Not an electrical circuit -- a
mathematical one.

\textbf{Commitment scheme} -- A cryptographic method that lets you ``seal'' a
value in an envelope, later proving properties about it without opening it. The
four main families are KZG, FRI, IPA, and lattice-based (Ajtai).

\textbf{CycleFold} -- An engineering optimization that delegates the non-native
scalar multiplication in folding schemes to a co-processor circuit on a
secondary curve, reducing overhead.

\textbf{Data availability (DA)} -- The guarantee that transaction data is
published and accessible so that anyone can reconstruct state and verify proofs
independently.

\textbf{Discrete logarithm problem (DLP)} -- Given points \(P\) and \(Q\) on an
elliptic curve where \(Q = nP\), finding the integer \(n\). Believed hard
classically; broken by Shor's algorithm on a quantum computer.

\textbf{Elliptic curve} -- A mathematical curve whose point-addition operation
is easy to perform but hard to reverse. The foundation of most modern public-key
cryptography.

\textbf{Execution trace} -- The complete record of every step in a computation:
register values, memory accesses, intermediate results. The raw material from
which a witness is derived.

\textbf{Extension field} -- A larger field constructed from a base field by
adjoining roots of an irreducible polynomial, used to achieve sufficient
security when working with small base fields.

\textbf{FHE (Fully Homomorphic Encryption)} -- Encryption that allows
computation on encrypted data without decrypting it first. Currently 10,000x to
1,000,000x slower than plaintext computation.

\textbf{Fiat-Shamir transform} -- A technique for converting an interactive
proof into a non-interactive one using a hash function. Every blockchain-based
zero-knowledge proof uses this.

\textbf{Finite field} -- A set of numbers where arithmetic ``wraps around'' like
a clock. All arithmetic in zero-knowledge proofs happens in finite fields.

\textbf{Flash loan} -- A DeFi mechanism allowing uncollateralized borrowing
within a single transaction. Used in governance attacks when combined with
token-weighted voting.

\textbf{Folding} -- A technique for combining two mathematical claims into one
using random challenges. Dramatically reduces the cost of proving long
computations.

\textbf{FRI (Fast Reed-Solomon Interactive Oracle Proof of Proximity)} -- A
method for verifying polynomial proximity using only hash functions. The
commitment scheme inside every STARK. Transparent and plausibly post-quantum.

\textbf{Gas} -- The unit of computational cost on Ethereum. A simple transfer
costs roughly 21,000 gas; verifying a Groth16 proof costs roughly 250,000 gas.

\textbf{Goldilocks field} -- A specific 64-bit prime (\(2^{64} - 2^{32} + 1\)).
Called ``Goldilocks'' because it fits in a single machine word for fast
arithmetic.

\textbf{Groth16} -- The most compact zero-knowledge proof system: proofs are
exactly 192 bytes. Requires a per-circuit trusted setup ceremony.

\textbf{Halo 2} -- A PLONK-family proof system developed by the Electric Coin
Company. The original Halo used IPA commitments (no ceremony); Halo 2 as
deployed typically uses KZG commitments.

\textbf{HyperNova} -- A folding scheme by Kothapalli and Setty that generalizes
Nova from R1CS to CCS using the sumcheck protocol, enabling folding of any
constraint system.

\textbf{ISA (Instruction Set Architecture)} -- The fundamental language a
processor understands. RISC-V is the dominant ISA for zero-knowledge virtual
machines.

\textbf{IVC (Incrementally Verifiable Computation)} -- A framework for proving
that a sequence of computation steps was performed correctly, where each step's
proof implicitly covers all previous steps.

\textbf{KZG (Kate-Zaverucha-Goldberg)} -- A polynomial commitment scheme
producing constant-size proofs (\textasciitilde48 bytes) using elliptic curve
pairings. Requires a trusted setup. Not quantum-resistant.

\textbf{Lattice} -- A regular grid of points in high-dimensional space. Finding
a short vector in such a lattice is believed to be hard even for quantum
computers.

\textbf{M31 / Mersenne-31} -- A 31-bit prime field (\(2^{31} - 1\)) used by
Circle STARKs and Stwo. Its circle group has order \(2^{31}\), enabling
efficient radix-2 FFTs.

\textbf{Merkle tree} -- A pyramid-shaped data structure where a single hash at
the top summarizes an entire collection, and membership can be proven with a
short path of hashes.

\textbf{Module-LWE (Module Learning With Errors)} -- A lattice problem believed
hard even for quantum computers. The foundation for NIST's post-quantum
encryption standard (FIPS 203) alongside Module-SIS.

\textbf{Module-SIS (Module Short Integer Solution)} -- A mathematical problem
believed hard even for quantum computers. The foundation for lattice-based ZK
proof systems and NIST's post-quantum standards.

\textbf{MPC (Secure Multi-Party Computation)} -- Protocols that let multiple
parties jointly compute a function without any party revealing its input to the
others.

\textbf{MSM (Multi-Scalar Multiplication)} -- A computation where many elliptic
curve points are multiplied by different scalars and summed. Highly
parallelizable on GPUs.

\textbf{Neo} -- A lattice-based folding scheme by Nguyen and Setty that adapts
HyperNova's CCS folding to small fields (Goldilocks) with Ajtai commitments,
introducing pay-per-bit commitment costs.

\textbf{Nova} -- The first practical folding scheme, by Kothapalli, Setty, and
Tzialla (2022). Introduced relaxed R1CS and reduced the per-step overhead of IVC
by orders of magnitude.

\textbf{NTT (Number Theoretic Transform)} -- The finite-field version of the
Fast Fourier Transform. A critical bottleneck in modern proof systems.

\textbf{Nullifier} -- A cryptographic value derived from a secret key and a
commitment, used to prevent double-spending in privacy-preserving systems
without revealing the spender's identity.

\textbf{Pedersen commitment} -- A cryptographic commitment using elliptic curve
arithmetic. Not quantum-resistant.

\textbf{PLONK} -- A universal zero-knowledge proof system (2019) that works with
any circuit up to a fixed size using a single trusted setup ceremony.

\textbf{PLONKish} -- The arithmetization format used by PLONK and variants. Uses
selector columns for different gate types plus copy constraints for wiring.

\textbf{Post-quantum (PQ)} -- Resistant to attacks by quantum computers.
Hash-based and lattice-based cryptography are believed to be PQ. Elliptic curve
cryptography is not.

\textbf{Proof aggregation} -- Combining multiple proofs into a single proof that
is cheaper to verify than checking each proof individually. A key technique for
amortizing on-chain verification costs.

\textbf{Prover} -- The entity that generates a zero-knowledge proof. In the
magician metaphor, the prover is the magician who performs the trick.

\textbf{R1CS (Rank-1 Constraint System)} -- The simplest constraint system
format. The ``assembly language'' of arithmetization. Used natively by Groth16.

\textbf{Relaxed R1CS} -- A generalization of R1CS
(\(\mathbf{A}\mathbf{z} \circ \mathbf{B}\mathbf{z} = u \cdot \mathbf{C}\mathbf{z} + \mathbf{E}\))
that introduces a scalar \(u\) and error vector \(\mathbf{E}\), enabling two
instances to be combined via random linear combination. The key enabler of
folding.

\textbf{RISC-V} -- An open, royalty-free instruction set architecture. Eight of
ten major zkVMs use RISC-V as their target.

\textbf{Rollup} -- A blockchain scaling technique that executes transactions off
the main chain, posting a compact proof back for verification.

\textbf{Shor's algorithm} -- A quantum algorithm (1994) that breaks all elliptic
curve cryptography in polynomial time.

\textbf{SNARK (Succinct Non-interactive ARgument of Knowledge)} -- A proof
system producing compact proofs that can be verified quickly without
interaction. Groth16, PLONK, and Marlin are SNARKs.

\textbf{Soundness} -- The guarantee that a cheating prover cannot create a valid
proof for a false statement, except with negligible probability.

\textbf{SRS (Structured Reference String)} -- The public mathematical parameters
from a trusted setup ceremony. The stage on which the magician performs.

\textbf{STARK (Scalable Transparent ARgument of Knowledge)} -- A proof system
using hash-based commitments (FRI) with no trusted setup. Transparent, plausibly
post-quantum, but with larger proofs than SNARKs.

\textbf{Stwo} -- StarkWare's production implementation of Circle STARKs,
deployed on Starknet mainnet. Achieved 940x throughput improvement over the
previous Stone prover.

\textbf{Sumcheck protocol} -- An interactive proof protocol (1992) that reduces
verifying a sum over an exponentially large domain to checking a single
evaluation. The backbone of modern ZK proof systems.

\textbf{TEE (Trusted Execution Environment)} -- A hardware-isolated secure
enclave (e.g., Intel SGX, ARM TrustZone) that processes data in a protected
region even the hardware operator cannot inspect.

\textbf{UTXO (Unspent Transaction Output)} -- A model where each digital coin is
a discrete object created by one transaction and consumed by another, like
physical bills in a wallet.

\textbf{Verifier} -- The entity that checks a zero-knowledge proof. In the
magician metaphor, the verifier is the audience that renders the verdict.

\textbf{Witness} -- The private data and execution trace that the prover uses to
generate a proof. The verifier never sees the witness.

\textbf{Zero-knowledge} -- The property that a proof reveals nothing beyond the
truth of the statement being proved.

\textbf{ZKIR} -- Midnight's Zero-Knowledge Intermediate Representation. A
24-instruction typed bytecode that the Compact compiler produces, which a
backend lowers to PLONKish constraints.

\textbf{zkVM (Zero-Knowledge Virtual Machine)} -- A virtual processor that
executes programs and automatically generates a zero-knowledge proof of correct
execution.

\partdivider{Part I: The Invitation}

\chapter{Chapter 1: The Promise of Provable and Programmable
Secrets}\label{chapter-1-the-promise-of-provable-and-programmable-secrets}

\begin{quote}
\emph{``Any sufficiently advanced technology is indistinguishable from magic.''}
-- Arthur C. Clarke, \emph{Profiles of the Future} (1962)
\end{quote}

\section{The Trick}\label{the-trick}

Every civilization that has ever existed has faced the same problem. You know
something. I need to verify it. And the only method anyone has ever devised for
verification is \emph{disclosure} -- you open the books, you show the document,
you reveal the source code, you hand over the key. Knowledge flows from the
prover to the checker, and some of it inevitably spills.

For ten thousand years, this was not a problem. It was a law of nature, as
inescapable as gravity. To prove is to show. To show is to reveal. To reveal is
to lose control.

Then, in 1985, three researchers at MIT wrote a paper and broke the law.

Shafi Goldwasser, Silvio Micali, and Charles Rackoff demonstrated something that
sounds, on first hearing, like a contradiction: it is possible to prove a
statement is true while revealing \emph{nothing} about \emph{why} it is true.
Not approximately nothing. Not mostly nothing. Nothing -- in a mathematically
rigorous sense that can be stated as a theorem and verified by anyone who cares
to check. The proof convinces. It does not inform. The audience sees the trick
succeed and learns nothing about how it was performed.

Clarke was half right. This technology is indistinguishable from magic -- until
you understand it. Then it is more astonishing than magic, because magic relies
on deception while this relies on its opposite. The magician hides the
mechanism. The zero-knowledge proof hides the data \emph{and lets you verify the
mechanism is honest}. That distinction will occupy us for the rest of this book.

Here is the thesis, stated plainly so you can argue with it for the next three
hundred pages.

Zero-knowledge proofs do not eliminate trust. They \emph{decompose} it. They
take one monolithic act of faith -- trust the bank, trust the platform, trust
the government -- and shatter it into seven independent, weaker assumptions.
Each one is testable. Each one is replaceable. What remains after the
mathematics has done its work is not zero trust but \emph{less} trust,
distributed across more points of failure, each auditable on its own terms.

The word ``trustless'' is marketing. The accurate word is ``trust-minimized.''
This book will show you exactly how much trust remains, where it lives, and what
breaks when each assumption fails.

There are only two characters in this story. One is the magician. The other is
the audience. In the technical literature, they are called the \emph{prover} and
the \emph{verifier}. The prover knows the secret and wants to convince someone
of a fact without revealing private information. The verifier checks the proof
and renders a verdict: accept or reject. Every zero-knowledge system ever built
-- every billion-dollar rollup, every privacy protocol, every identity
credential -- reduces to this elemental exchange.

We will use the magician-and-audience framing throughout, not because it is
charming but because it is precise. When the analogy works, we will ride it.
When it breaks -- and it will, most severely at Layer 4, where the mathematics
of arithmetization overwhelms any theatrical metaphor -- we will say so. A
strained metaphor is worse than no metaphor at all. By Chapter 10, we will
replace the seven-layer stack entirely with a directed graph. We will tell you
when each transition happens.

\section{The Proof at the Door}\label{the-proof-at-the-door}

Before naming the properties that make this possible, let us watch the trick at
the scale of a single human interaction.

You are twenty years old. You walk into a bar. The bouncer asks for your ID. You
reach into your wallet and hand over a government-issued card containing your
full legal name, your date of birth, your home address, your photograph, your
driver's license number, and -- depending on your jurisdiction -- your organ
donor status and corrective lens prescription. The bouncer glances at the birth
year, confirms you are of legal age, and hands it back.

Pause on this. In a three-second interaction, you have disclosed six pieces of
personally identifiable information to solve a problem that requires exactly one
bit of answer: \emph{are you twenty-one or older? Yes or no.}

The card was designed in an era when the only way to prove a fact about yourself
was to reveal the document containing that fact. We have been living with this
design for so long that it feels inevitable. It is not.

Now imagine a different exchange. You carry a digital credential on your phone
-- issued by the same government authority, carrying the same legal weight. The
bouncer's terminal displays a challenge: \emph{prove you are 21 or older.} Your
phone computes for a fraction of a second and transmits a proof. The green light
appears.

What did the bouncer learn? That you are old enough. Nothing else. Not your
name. Not your address. Not your exact age. Not even which government issued the
credential. The proof is a single bit of verified truth, wrapped in mathematics.

Three properties make this work. They are worth naming now because they recur at
every layer of the story -- and because, if you watch closely, you have already
seen all three in action.

The first: \emph{the honest succeed}. If you genuinely are twenty-one, the proof
will always verify. No glitch, no false rejection, no edge case where valid
credentials fail. Cryptographers call this \textbf{completeness}. Without it,
honest people get turned away, and the technology dies on contact with reality.

The second: \emph{the dishonest fail}. A nineteen-year-old cannot forge this
proof. They cannot borrow your credential and make it work -- the proof is bound
to a secret key only you possess. They cannot manipulate the computation to make
nineteen appear as twenty-one -- the polynomial constraints will not verify. The
mathematics enforces the rule on the bouncer's behalf. Cryptographers call this
\textbf{soundness}, and the word ``cannot'' is doing real work: the probability
of a successful forgery is roughly one in \(2^{128}\), a number with thirty-nine
digits. The sun will burn out first.

The third: \emph{nothing leaks}. Not a partial hint about your birth month. Not
a statistical correlation that narrows your age range. An adversary who
intercepts the proof learns exactly what the bouncer learned -- you are old
enough -- and nothing more. The proof is, in a precise and formal sense,
\emph{simulatable}: anyone could generate something that looks identical to it
without knowing your date of birth at all. Cryptographers call this
\textbf{zero-knowledge}. It is the property that makes the trick feel
impossible, and the one this book exists to explain.

\section{The Phenomenon}\label{the-phenomenon}

Completeness, soundness, zero-knowledge. These three properties were the subject
of Goldwasser, Micali, and Rackoff's 1985 paper, ``The Knowledge Complexity of
Interactive Proof Systems'' -- one of the most consequential in the history of
computer science, and the work that earned Goldwasser and Micali the Turing
Award. For roughly twenty years after its publication, zero-knowledge proofs
remained almost entirely theoretical. Beautiful mathematics, waiting for the
world to catch up.

The world caught up.

What does a proof actually look like? Not a page of equations. Not an argument
in English. A Groth16 proof -- named after the Danish cryptographer Jens Groth,
who published the scheme in 2016 -- is exactly three groups of numbers, each
about 48 bytes long. The entire proof fits in 192 bytes. That is smaller than a
tweet, shorter than this paragraph's first sentence. Three numbers that took a
GPU cluster seconds to compute, that a smart contract on Ethereum can verify for
roughly 250,000 gas, and that reveal absolutely nothing about the secret they
certify.

A STARK proof is larger -- typically 50 to 200 kilobytes -- but still small
compared to the computation it certifies. And unlike Groth16, a STARK requires
no trusted setup ceremony: its security rests on hash functions alone. We will
spend all of Chapter 2 on the difference, because the choice between these two
families is the first real decision any system designer faces, and its
consequences cascade through every subsequent layer.

The bar scenario is vivid, but it understates the stakes. The real power emerges
when you chain zero-knowledge proofs together.

A bank must prove to a regulator that it holds sufficient reserves to cover all
deposits. Today, this requires opening the books -- revealing individual account
balances, transaction histories, counterparty relationships. The bank's
customers never consented to this exposure. The regulator receives more
information than necessary. With a zero-knowledge proof, the bank proves a
single statement: ``The sum of all deposits is less than or equal to the sum of
all reserves.'' The regulator verifies in milliseconds. The proof reveals
nothing about any individual account. The regulator is convinced -- not because
they trust the bank, but because the mathematics makes it impossible for the
bank to have produced a valid proof of a false statement.

Or take the supply chain. A pharmaceutical company must prove to the FDA that a
batch of medication was manufactured in a GMP-certified facility, that the
active ingredient was sourced from an approved supplier, and that the cold chain
was maintained throughout transit. Today, this audit requires revealing
proprietary manufacturing processes, supplier contracts, and logistics routes --
competitive intelligence that companies spend millions to protect. With
zero-knowledge proofs, each attestation becomes a sealed certificate: the proof
says the constraints are satisfied, the FDA verifies, and the proprietary
process stays proprietary.

Or take the scenario that scales furthest. You prove you are over twenty-one
\emph{and} that you hold a valid medical prescription \emph{and} that your
insurance covers the medication -- all without revealing your name, your
diagnosis, or your policy number to anyone except the parties who strictly need
each fact. Each proof is independent. Each reveals only its single bit of truth.
The pharmacy fills your prescription. Your insurer processes the claim. Your
medical history remains yours.

The pattern is always the same. Someone must verify a fact. Verification has
always required disclosure. Disclosure leaks information irrelevant to the fact
being verified. Zero-knowledge proofs sever the coupling between verification
and disclosure. They let you prove the fact and \emph{only} the fact.

This is privacy as an engineering constraint -- enforced by polynomial
equations, verified in milliseconds, falsifiable by anyone who cares to check.

\section{Three Converging Forces}\label{three-converging-forces}

For twenty years after Goldwasser, Micali, and Rackoff, zero-knowledge proofs
were a theoretical marvel and a practical impossibility. Three forces converged
to change that.

\textbf{The privacy crisis became quantifiable.} Proving you are over eighteen
requires revealing your full date of birth. Proving your creditworthiness
requires revealing your financial history. Proving your vaccination status
requires revealing your medical records. Each is a case where a zero-knowledge
proof could replace full disclosure with a single verified bit: yes or no,
nothing more. The European Union's eIDAS 2.0 regulation, taking effect in 2026,
mandates that digital identity wallets for 450 million EU citizens support
selective disclosure -- revealing only the attributes necessary for a given
interaction. That mandate, at continental scale, requires zero-knowledge
infrastructure. The regulatory pull is no longer hypothetical. It is law.

\textbf{The scaling problem acquired a price tag.} Ethereum, the world's largest
smart-contract platform, processes roughly 15 transactions per second on its
base layer. ZK rollups -- Layer 2 systems that batch hundreds or thousands of
transactions into a single proof verified on-chain -- process orders of
magnitude more. The rollup inherits the security of the main chain without
requiring the main chain to re-execute every transaction. The proof is
sufficient. By early 2026, the total value locked in ZK rollups exceeded \$20
billion. These are not research projects. They are financial infrastructure
securing real capital, and every dollar locked is a dollar betting that the
proof system works.

\textbf{The cost curve broke loose from history.} In December 2023, generating a
single zero-knowledge proof of a meaningful computation cost approximately \$80.
By December 2025, it cost \$0.04. A 2,000-fold reduction in twenty-four months.
To appreciate how unusual this is: solar energy costs dropped 99\% over
\emph{four decades}. DNA sequencing outran Moore's Law, but it took a decade to
achieve its steepest plunge. Zero-knowledge proving achieved a comparable
collapse in two years. Nothing in recent technological history matches this
pace.

What drove it? Four independent teams -- SP1 Hypercube, RISC Zero, Stwo, and the
Ethereum Foundation's zkEVM effort -- achieved real-time proving of Ethereum
blocks by late 2025. SP1 proved a block in 6.9 seconds on 16 GPUs, handling
99.7\% of Ethereum L1 blocks in under 12 seconds. The improvements compounded
across the entire stack: smaller fields reduced arithmetic cost, Circle STARKs
enabled efficient proving over those fields, better arithmetization reduced
constraint counts, GPU parallelization accelerated the proving pipeline, and
aggregation services amortized verification. The seven-layer model we will build
in this book is not merely an analytical framework. It is a map of where the
cost reductions came from.

The speed race reached its finish line. In December 2025, the Ethereum
Foundation declared it won and pivoted the competition axis to provable
security: 128-bit security proofs by end of 2026. The performance frontier has
been crossed. The security frontier is now.

At \$0.04 per proof, the economics undergo a phase transition -- the way water
does not gradually become steam but shifts states at a threshold. At \$80, only
the most valuable financial settlements justify the expense. At four cents, you
can prove \emph{everything}: identity checks, game-state transitions, AI model
inferences, supply-chain attestations, compliance audits. Projections for 2027
suggest sub-cent proofs for computations exceeding one billion cycles. The
question stops being ``can we afford to prove this?'' and becomes ``why would we
\emph{not} prove everything?''

These three forces -- privacy demand, scaling need, and cost collapse -- form a
flywheel. Cheaper proofs make more applications viable. More applications create
more demand. More demand funds more engineering. More engineering produces
cheaper proofs. The cycle is self-reinforcing.

\section{The Seven Layers at a Glance}\label{the-seven-layers-at-a-glance}

Every zero-knowledge system, from the simplest proof to the most complex rollup,
decomposes into seven layers. These are not independent floors in a building.
They are organs in a body -- deeply interdependent, each shaping what the others
can do.

Here they are, one sentence each, in the language of our two characters.

\textbf{Layer 1 -- The Setup (Building the Stage).} Before the magician can
perform, someone must construct the mathematical parameters that make proving
and verifying possible -- and who builds them, and whether you must trust them,
determines the system's deepest guarantees.

\textbf{Layer 2 -- The Language (Choreographing the Act).} The magician needs a
script: a programming language in which to express the computation she wants to
prove, where her choice of language determines what bugs she can make and what
assurances the compiler can enforce.

\textbf{Layer 3 -- The Witness (The Secret Backstage Preparation).} The magician
goes backstage to run the computation with her private data, recording every
step in an execution trace that no one else will ever see -- and this
preparation, not the proof itself, is the most underestimated bottleneck in the
entire stack.

\textbf{Layer 4 -- The Arithmetization (Encoding the Performance).} The
backstage recording is transformed into a system of polynomial equations: a
mathematical puzzle that can be checked in seconds, even though solving it took
hours.

\textbf{Layer 5 -- The Proof System (Sealing the Certificate).} The magician
compresses the entire puzzle into a compact, tamper-proof certificate -- the
proof -- using cryptographic machinery that guarantees soundness and
zero-knowledge simultaneously.

\textbf{Layer 6 -- The Primitives (The Deep Craft).} Beneath the proof system
lie the fundamental cryptographic building blocks: elliptic curves, hash
functions, and polynomial commitment schemes, whose mathematical hardness
assumptions are the bedrock on which everything above rests.

\textbf{Layer 7 -- The Verification (The Audience's Verdict).} The audience
checks the proof on a public stage -- a blockchain, a smart contract, a verifier
endpoint -- and the economics, governance, and data availability of that stage
determine whether the trick actually matters in the world outside mathematics.

A warning. We present these layers in order, from 1 to 7, because a book must be
linear. The dependencies do not share this constraint. The choice of
cryptographic primitive (Layer 6) determines which proof systems are available
(Layer 5). That determines which arithmetizations work (Layer 4). That shapes
which languages are efficient (Layer 2). That constrains the setup (Layer 1). A
single decision -- say, choosing the Goldilocks field over BLS12-381 -- can
cascade through all seven layers and reshape the entire architecture.

In practice, some layers fuse. Jolt merges witness generation and
arithmetization into a single lookup step. Cairo co-designs its language around
its constraint system. The proof core -- Layers 4, 5, and 6 -- behaves as one
inseparable design unit in every production system. We will flag each collapse
when we reach it. By Chapter 10, we will redraw the model in its honest final
form: a directed acyclic graph with fourteen causal edges, not seven tidy
floors. The map is provisional. The territory is more tangled. We will be honest
about where the map fails.

You do not need to memorize this list. We will build each layer one at a time,
and by the end, this overview will feel like a map you drew yourself.

\section{The Deepest Question}\label{the-deepest-question}

Each of the seven layers introduces a trust assumption -- a bet you are making
that something is true, without the ability to verify it from first principles
in real time.

At Layer 1, you trust that the setup ceremony was conducted honestly, or you
choose a transparent setup that requires no ceremony at all, accepting a cost in
proof size or verification time. The largest ceremony ever conducted --
Ethereum's KZG Summoning -- drew 141,416 participants. Among them, there is
almost certainly someone whose operational security was impeccable. Almost
certainly. Chapter 2 will make that ``almost'' precise.

At Layer 2, you trust that the program was written correctly. This is where the
deepest irony of the field lives: 67\% of real-world zero-knowledge
vulnerabilities are under-constrained circuits -- bugs in the \emph{program},
not in the cryptography {[}Chaliasos et al., USENIX Security 2024{]}. The most
sophisticated proof system in the world will faithfully prove a false statement
if the circuit asks it to. One character -- literally, a single \texttt{=} where
\texttt{\textless{}==} was needed -- caused a complete soundness break in
Tornado Cash. Chapter 3 tells that story.

At Layer 3, you trust that the hardware generating the witness does not leak
secrets through side channels. The proof is zero-knowledge, but the
\emph{process of generating the proof} may not be. Privacy, it turns out, is
partly a luxury good: the architecture that maximizes it (client-side proving)
demands hardware most people cannot afford. Chapter 4 confronts this
uncomfortable truth.

At Layers 4 and 5, you trust that the arithmetization faithfully encodes the
computation and that the proof system's security reduction is tight. The
overhead tax for this encoding is staggering -- 10,000x to 50,000x over native
execution, falling to 1,000x-5,000x with current optimizations. You are paying
for verifiability in compute. Chapters 5 and 6 show where that tax comes from
and how the field is reducing it.

At Layer 6, you trust that the underlying mathematical problems -- discrete
logarithms, hash collision resistance, lattice problems -- are genuinely hard.
These are conjectures, not theorems. BN254, a widely deployed elliptic curve,
has already seen its security erode from an estimated 128 bits to roughly 100.
NIST targets 2035 for retiring all pre-quantum algorithms. If you are building
infrastructure meant to last decades, this layer's assumption has an expiration
date. Chapter 7 explains why, and what comes after.

At Layer 7, you trust the governance of the verification layer. Most ZK rollups
today operate at Stage 0 or Stage 1 on the L2Beat maturity framework, meaning a
small group holding multisig keys can replace the verifier contract. Six layers
of mathematical elegance, and the seventh is a committee. The Beanstalk protocol
lost \$182 million in 13 seconds when an attacker used a flash loan to capture
governance power and drain the treasury. The governance mechanism worked exactly
as designed. The system was not broken. It was \emph{used}. Chapter 8 tells that
story, and it is the emotional nadir of the book.

No single layer requires trusting one entity with everything. Each assumption is
independently falsifiable, independently auditable, independently replaceable --
in principle. In practice, the coupling between layers makes this harder than it
sounds, and that tension between principle and practice is the spine of the next
thirteen chapters.

There is a deep symmetry worth noting early. Layer 1 (setup) and Layer 7
(verification) are both \emph{social} trust -- human decisions about who to
trust with power. The mathematical layers (2 through 6) are sandwiched between
them. The system converts social trust into mathematical certainty and then
converts mathematical certainty back into social trust. The cryptography is a
bridge between two shores of human judgment. Understanding this symmetry -- and
the vulnerability it creates at both ends -- is the difference between
understanding zero-knowledge proofs as mathematics and understanding them as
systems that operate in the real world.

\section{How to Read This Guide}\label{how-to-read-this-guide}

This book has three audiences.

\textbf{The Executive Path (\textasciitilde45 minutes).} Chapter 1, the opening
of Chapter 2, the opening of Chapter 8, the Chapter 11 landscape table, and
Chapter 14. You will emerge able to evaluate vendor claims, explain ZK proofs to
a board, and ask the right questions about governance and quantum risk.

\textbf{The Engineer Path (\textasciitilde2 hours).} Chapters 1-2 in full,
Chapters 3-5 through their core sections, Chapter 6 through ``Folding: The
Snowball,'' Chapter 7 through ``Four Families of Commitment Schemes,'' Chapter 8
through ``Governance: The Achilles Heel,'' and Chapters 10, 11, and 13.

\textbf{The Researcher Path (\textasciitilde4+ hours).} Read everything. Then
work through the seven open questions in Chapter 14. The active research fronts
are the folding genealogy (Chapter 6), the lattice revolution (Chapter 7), the
CCS unification (Section 5.3), and the proof core triad (Chapters 6, 10, and
11).

\textbf{Regardless of path:} do not skip Chapter 2 or the trust decomposition in
Chapter 10. The first is where the deepest trust decisions live. The second is
the book's thesis in its most complete form -- seven assumptions, fourteen
causal edges, three architectural paths, each with a different failure profile.
Everything between here and there is building the case.

The field is crossing three sequential frontiers. \emph{Performance} (2023-2025)
is largely crossed: real-time proving is achieved, costs are sub-cent.
\emph{Security} (2026-2028) is the current frontier: formal verification,
128-bit provable security, post-quantum readiness. \emph{Privacy} (2027+) is
approaching: compiler-enforced disclosure boundaries, constant-time
implementations, metadata protection. Where you enter this story depends on
which frontier matters most to your work.

\section{The First Decision}\label{the-first-decision}

We have described the trick. We have named its three properties. We have
sketched the seven layers, previewed the trust assumptions hiding inside each
one, and measured the cost trajectory that is making all of this economically
inevitable. A running example -- a 4x4 Sudoku puzzle -- will ground every
abstraction from Chapter 3 onward. It is the computation we prove at each layer:
first as a program, then as a witness, then as 72 polynomial constraints, then
as a sealed certificate, and finally as a verified proof. The same small puzzle,
viewed through seven different lenses.

Now we need to build the stage.

Every magic show requires one. In zero-knowledge proof systems, ``building the
stage'' means creating the mathematical parameters that both characters will use
-- the prover to generate proofs, the verifier to check them. This step happens
before any proof is ever created. And the most important question about the
stage is not how it is built but \emph{who builds it, and whether you have to
trust them}.

In 2016, the Zcash ceremony involved six participants, each generating a share
of the secret parameters and then destroying their share. If even one of the six
was honest, the system is secure. In 2023, Ethereum's KZG Summoning Ceremony
drew 141,416 participants. The growth from six to 141,416 tells you something
about how seriously the field takes this problem. The alternative -- transparent
setups that require no ceremony at all -- tells you something about how
seriously the field wishes it could avoid the problem entirely.

That question -- trusted setup or transparent setup, ceremony or glass stage,
hidden trapdoors or none -- is the subject of Chapter 2. It is the first fork in
the road. Its consequences echo through every layer that follows.

The stage is where the trust story begins. Let us build it.

\chapter{Chapter 2: Layer 1 -- Building the
Stage}\label{chapter-2-layer-1-building-the-stage}

\begin{quote}
*``If the toxic waste is destroyed, how does anyone* know?'' -- A question that
sounds naive until you realize nobody has a clean answer.
\end{quote}

\section{The Fair Shuffle Problem}\label{the-fair-shuffle-problem}

Here is a puzzle that has nothing to do with cryptography -- at first.

You and a group of friends want to play a card game. Nobody trusts anyone to
shuffle the deck fairly. The dealer might stack the cards. She might glimpse the
order. She might memorize a few key positions and exploit them later. So you
need a method -- a protocol -- that produces a provably fair shuffle, even
though someone must do the physical shuffling.

One solution: get 141,416 people to each add a bit of randomness to the shuffle,
one after another. Each person takes the deck, makes some unpredictable change
to the ordering, and passes it on. Then each person destroys their memory of
what they did. At the end, the deck is fair as long as \emph{even one} of those
141,416 people acted honestly -- genuinely random and genuinely forgetful.

That number -- 141,416 -- is not hypothetical. It is the exact count of
participants in the Ethereum KZG Summoning Ceremony, completed in 2023. The
largest cryptographic ceremony in history. And the thing they were ``shuffling''
was not a deck of cards but a set of mathematical parameters called a Structured
Reference String, or SRS -- the stage on which every proof in the system would
subsequently be performed.

Now consider a second solution: skip the dealer entirely. Design a game where no
shuffling is needed at all. Derive every rule from publicly known mathematics --
hash functions, no secrets, no waste. Anyone can verify the rules are fair by
reading them. No trust required. But the game runs slower, and the scorecards
take up more space.

These two solutions -- the ceremony and the transparent alternative -- define
the fundamental choice at Layer 1. Every zero-knowledge system ever deployed
sits on one side of this divide or the other. And increasingly, the most
sophisticated systems use both: a transparent inner proof wrapped in a
ceremony-derived outer shell.

Understanding this choice, and its consequences, is where the trust story
begins.

\section{The Structured Reference String}\label{the-structured-reference-string}

Before we discuss who builds the stage, we need to understand what it is.

A \emph{Structured Reference String} (SRS) is a list of specially constructed
numbers. These numbers are points on an elliptic curve -- a mathematical curve
with a remarkable property. You can ``add'' points on the curve together easily.
Given the result, however, it is extremely hard to figure out how many additions
were performed. We will explain elliptic curves more fully in Chapter 7. For
now, think of the SRS as a ruler with very precise markings that everyone uses
to measure, but that no one can reverse-engineer to discover how the markings
were made. The markings let you measure certain things (verify polynomial
evaluations). They do not let you reconstruct the manufacturing process (recover
the secret value from the markings). The scheme that makes this possible was
invented by Kate, Zaverucha, and Goldberg in 2010 {[}Kate, Zaverucha, Goldberg,
2010{]} and is universally known as KZG, after the authors' initials.

The secret value used to generate the SRS is called the \emph{trapdoor}. In the
zero-knowledge community, it has acquired a more evocative name: \emph{toxic
waste}. The metaphor is precise. Like radioactive material, the trapdoor is
dangerous if it persists and must be destroyed after it has served its purpose.
If anyone -- anyone at all -- retains knowledge of the trapdoor, they can forge
proofs of false statements. They can prove that \(2 + 2 = 5\), that an empty
bank account holds a billion dollars, that an invalid transaction is valid. No
one would be able to tell the forged proof from a real one.

The SRS is the stage. The toxic waste is the spare key to the stage's trapdoor.
And the ceremony is the manufacturing process designed to ensure that the spare
key is destroyed.

\section{Two Ways to Build a Stage}\label{two-ways-to-build-a-stage}

\subsection{The Trusted Stage: Ceremonies and the 1-of-N
Model}\label{the-trusted-stage-ceremonies-and-the-1-of-n-model}

The first generation of zero-knowledge systems relied on what is called a
\emph{trusted setup}. One person, or a small group, would generate the SRS using
secret randomness, produce the public parameters, and then -- hopefully --
destroy the secret.

The problem is obvious: you have to trust them. The original Zcash Sprout
ceremony in 2016 involved six pre-selected participants {[}Bowe, Gabizon, Miers,
2017{]}. (Note: the BGM17 paper describes the Sapling ceremony protocol, not
Sprout. The Sprout ceremony used the earlier BCTV14 protocol.) They used
air-gapped computers -- machines deliberately disconnected from the internet.
After generating their contributions, they physically destroyed the hardware.
One participant incinerated their computer. Another dissolved theirs in acid.
The ceremony took two days. It was as much a ritual as a protocol.

Six people. Six chances for a mistake, a compromise, or a lie. The Zcash team
needed the community to trust that at least one of these six acted honestly.
This is the \emph{1-of-N trust model}: security holds if even one participant
out of N genuinely destroyed their contribution.

Why, then, would you want six-figure participation instead of six? The
mathematics says one honest participant suffices. But the sociology says
something different. With six participants, you can name them, investigate them,
pressure them, bribe them, subpoena them. With over a hundred thousand anonymous
contributors from around the world, many submitting their randomness through a
simple web browser, collusion becomes operationally infeasible. The security is
sociological as much as cryptographic: you are not trusting a specific person.
You are trusting that it is impossible to corrupt \emph{everyone}.

The evolution of ceremonies tells the story of this insight pushed to its limit:

\begin{itemize}
\tightlist
\item
  \textbf{2016 -- Zcash Sprout.} Six pre-selected participants. Air-gapped
  machines physically destroyed. Two days. The first trusted setup ceremony in
  production.
\item
  \textbf{2018 -- Zcash Sapling.} Approximately ninety participants across two
  phases (87 in phase 1, 91 in phase 2) via the BGM17 ``MMORPG'' framework
  {[}Bowe, Gabizon, Miers, 2017{]}. 2.5 hours per contribution. The first
  scalable ceremony protocol.
\item
  \textbf{2019-2022 -- Proliferation.} Dozens of projects (Tornado Cash, Hermez,
  Aztec, Loopring) ran their own ceremonies, mostly for Groth16 circuits on the
  BN254 curve.
\item
  \textbf{2023 -- Ethereum KZG Summoning.} 141,416 contributors. Permissionless.
  Web-based. Anyone could participate. The ceremony ran for months. The
  ``more-the-merrier'' model at maximum scale.
\item
  \textbf{2025+ -- On-chain ceremonies.} Smart-contract-mediated setups where
  contributions are verified on-chain and no coordinator is needed
  {[}Nikolaenko, Ragsdale, Bonneau, Boneh, 2022{]}.
\end{itemize}

The SRS that emerged from the Ethereum KZG ceremony is \emph{universal}: it
works for any circuit up to a certain size. This distinction matters. The older
Groth16 system {[}Groth, 2016{]} required a \emph{circuit-specific} setup --
every new program needed a new ceremony. PLONK {[}Gabizon, Williamson,
Ciobotaru, 2019{]} and its descendants changed this: one ceremony, one universal
SRS, and then deterministic, public key derivation for each specific circuit.
The ceremony is the capital expenditure. Everything after it is operating
expense.

\subsection{What It Felt Like: The Human Interior of a
Ceremony}\label{what-it-felt-like-the-human-interior-of-a-ceremony}

The timeline above tells the institutional story. It says nothing about what it
was like to be there -- to be one of those 141,416 people, sitting at a laptop,
waiting for a progress bar to crawl across a browser tab, knowing that you were
participating in something whose purpose you might only half understand.

Consider the six who gathered for the Zcash Sprout ceremony in 2016. They came
from different countries. They did not all know each other. Some were
cryptographers by training; others were engineers or entrepreneurs who happened
to believe that financial privacy was a human right worth defending with
mathematics. They met in locations that were, by design, not disclosed in
advance. Each brought a computer -- purchased new, never connected to the
internet, destined to be destroyed within hours of its first and only use.

Peter Todd, one of the six, later described his experience. He purchased a
laptop with cash from a store he chose at random, drove to a location he had
selected from a map, and generated his contribution in a room with no wireless
signals. When he was done, he filmed himself incinerating the machine. The video
exists. You can watch a man feed a laptop into a furnace and understand that
what you are witnessing is not performance art. It is a protocol step. The
computer had touched a secret -- a single random number, a fragment of toxic
waste -- and the protocol required that the computer cease to exist. Todd
performed the ceremony alone, with no audience except the camera, in a
deliberate act of informational suicide: the knowing destruction of something
known, so that the knowing would end.

There is something ancient in this. Cultures have always understood that certain
acts of creation require corresponding acts of destruction. The Navajo sand
painting is swept away upon completion. The Tibetan mandala is dissolved into a
river. The meaning of the creation is inseparable from the certainty of its
erasure. What the Zcash participants were doing -- generating a secret, using it
exactly once, and annihilating every trace of it -- belongs to this lineage. The
difference is that here, the stakes were not symbolic. A single surviving photon
of information, a single bit of toxic waste persisting in a swap file or a
memory register, would compromise the soundness of every proof the system would
ever produce.

Andrew Miller, another Sprout participant, took a different approach. He used a
CD-ROM-based operating system that ran entirely in volatile memory -- nothing
was ever written to permanent storage. When he powered down the machine, the
secret evaporated with the electrical charge in the RAM chips. The computer
itself survived. The information did not. This was, in its way, more elegant
than incineration: rather than destroying the vessel, you ensured the vessel
never held the secret in any durable form. The secret existed for perhaps ninety
seconds as a pattern of electrical charges in silicon, performed its
mathematical function, contributed its randomness to the growing SRS, and then
returned to thermal noise. If you had frozen the RAM chips in liquid nitrogen
within those ninety seconds, you might have recovered something. Nobody did.

These are not colorful anecdotes. They are the actual security analysis. When a
cryptographer evaluates the Sprout ceremony, they must reason about furnaces,
cash purchases, volatile memory, and the thermal decay time of DRAM cells. The
mathematics of the 1-of-N trust model is impeccable. The implementation of that
model passes through the physical world, where ``destruction'' is an engineering
judgment, not a mathematical theorem.

The Sapling ceremony in 2018 changed the texture of participation entirely. Sean
Bowe, Ariel Gabizon, and Ian Miers designed the BGM17 protocol specifically to
lower the barrier. You no longer needed to buy a disposable laptop or own a
furnace. You needed a computer, an internet connection, and approximately 2.5
hours of patience. The protocol was sequential -- each participant received the
current state of the SRS, multiplied it by their own secret randomness, and
passed the result to the next participant -- but the infrastructure was
centralized enough to be manageable and decentralized enough to be credible.

Ninety people participated. They came from the cryptography community, from the
Zcash ecosystem, from academia, from the broader blockchain world. Some ran the
computation on high-end workstations. Others used modest laptops. One
participant, reportedly, ran their contribution on a Raspberry Pi -- the \$35
single-board computer popular with hobbyists -- just to prove it could be done.
The contribution took longer on the Pi. The mathematics did not care.

But ninety people is still a small enough group to name, to investigate, to
subpoena. The operational security improvement over six participants was real
but bounded. A nation-state adversary with sufficient motivation could, in
principle, identify and compromise all ninety. The names were not all public,
but the ceremony's sequential nature meant that the coordinating infrastructure
knew who was participating and when. The 1-of-N model was theoretically sound.
The N was not yet sociologically overwhelming.

The KZG Summoning Ceremony that Ethereum conducted in 2023 was something
qualitatively different. It was not a ceremony in the way the Zcash events were
ceremonies -- a small group of known individuals performing careful, bespoke
acts of cryptographic hygiene. It was a mass participatory event. A ritual
scaled to the size of a small city.

The interface was a web page. You navigated to ceremony.ethereum.org, connected
a wallet or signed in with an Ethereum account, and waited in a queue. When your
turn came, your browser generated a random number, performed the necessary
elliptic curve multiplication, submitted the result, and discarded the secret.
The entire contribution took between twenty seconds and two minutes, depending
on your hardware. You did not need to understand polynomial commitments or KZG
or elliptic curves. You needed to click a button and wait.

One hundred and forty-one thousand, four hundred and sixteen people did this.

They did it from bedrooms and offices and coffee shops. They did it on phones
and tablets and gaming rigs and ancient ThinkPads. They did it in Tokyo and
Lagos and Sao Paulo and Reykjavik and places that no one will ever catalog
because the ceremony did not require identification beyond an Ethereum address.
Some understood exactly what they were contributing and why. Others participated
because someone they followed on social media said it mattered. Others did it
out of curiosity, or boredom, or a vague sense of civic duty toward an
infrastructure they used but did not fully comprehend.

And here is what matters: the ceremony did not need any of them to understand.
The mathematics does not require informed consent. It requires honest
randomness. A participant who clicks the button because a friend told them to,
generating a random number from hardware entropy they could not explain,
contributes exactly as much security as a participant who has read every paper
on polynomial commitments since 2010. The secret is not in the understanding.
The secret is in the entropy.

This is the sense in which the KZG ceremony was a ritual in the anthropological
meaning of the word. A ritual is a formalized action whose efficacy does not
depend on the participant's comprehension of the underlying mechanism. A child
who is baptized does not understand the theology. A voter who marks a ballot
does not understand constitutional law. A ceremony participant who clicks a
button does not understand the discrete logarithm problem. In each case, the act
is valid because the protocol says it is valid, not because the actor grasps the
full chain of reasoning that connects the act to its consequences.

Anthropologists who study ritual observe that large-scale ceremonies serve
functions beyond their stated purpose. They create shared experience. They
generate social cohesion. They mark boundaries between in-groups and out-groups.
They transform abstract commitments into embodied actions. The Ethereum KZG
ceremony did all of these things. Participants shared screenshots of their
contribution receipts on social media. They discussed queue times and hardware
performance. They created a shared memory of participation in something that
mattered, even if they could not articulate exactly what it was or why.

The ceremony also served a function that its designers may not have fully
anticipated: it created a constituency. One hundred and forty-one thousand
people who personally contributed to the SRS have a stake in its legitimacy that
goes beyond economic interest. They performed an act -- a small act, a few
seconds of computation -- and that act binds them to the outcome in a way that
merely reading about a ceremony does not. The SRS is not just a mathematical
object to them. It is something they helped build. This psychological investment
is not captured in any security model, but it is real, and it matters for the
long-term governance and social resilience of the system.

\subsection{Ritual as Protocol, Protocol as
Ritual}\label{ritual-as-protocol-protocol-as-ritual}

The progression from Sprout to the KZG Summoning traces an arc that any
historian of technology would recognize: the democratization of a priestly
function.

In 2016, generating ceremony contributions required expertise, specialized
equipment, and operational security practices borrowed from intelligence
agencies. The participants were, in effect, priests -- a small caste entrusted
with a dangerous sacrament, selected for their competence and their willingness
to undergo the destruction ritual. The community watched from outside. They
could verify the transcript after the fact, but they could not participate.
Trust flowed downward, from the priests to the congregation.

By 2023, the priestly function had been automated and distributed. The browser
handled the cryptography. The hardware random number generator provided the
entropy. The protocol enforced the sequencing. The participant's role was
reduced to a single act of will (the decision to show up) and a single act of
trust (the belief that their browser was not compromised). The ceremony had
become, in the language of distributed systems, \emph{permissionless}. Anyone
could join. No one could be excluded. The trust model had inverted: instead of
the congregation trusting the priests, each participant trusted only themselves.

This inversion changes how we should think about the foundations of
cryptographic systems. The traditional model of security assumes a small number
of trusted parties who perform critical operations on behalf of a larger
community. Certificate authorities sign TLS certificates. Key escrow agents hold
recovery keys. Hardware security module operators maintain root secrets. The
ceremony model -- especially at the scale of the KZG Summoning -- replaces this
with something closer to a vote. Not a vote on a proposition, but a vote of
presence. By participating, each contributor cast a ballot that said: ``I exist,
I generated honest randomness, and I destroyed my secret.'' The SRS is the
tally.

The algebra is worth seeing. The SRS is a sequence of elliptic curve points:
\([s]G, [s^2]G, [s^3]G, \ldots, [s^d]G\), where \(s\) is the toxic waste and
\(G\) is the generator point. During the ceremony, each participant \(i\) takes
the current SRS, multiplies every point by their own secret \(\tau_i\), and
passes the result forward. After all \(N\) participants have contributed, the
effective secret is \(s = \tau_1 \cdot \tau_2 \cdot \tau_3 \cdots \tau_N\). To
recover \(s\), an adversary must know \emph{every} \(\tau_i\) -- every
participant's secret. If even one participant's \(\tau_i\) is truly random and
truly destroyed, the product \(s\) is irrecoverable. The secrets do not add.
They multiply. And a product of unknowns is unknown.

This multiplicative structure is why the 1-of-N model works, and it is also why
the human story matters. Each \(\tau_i\) is not just a number. It is a moment in
a person's life: the instant their hardware random number generator sampled
thermal noise, the few seconds of computation, the silent deletion that
followed. The SRS encodes 141,416 such moments, multiplied together, fused into
a mathematical object that no participant can reverse and no adversary can
deconstruct. The ceremony is a protocol that transforms individual acts of trust
into a collective mathematical artifact. It is, in this precise sense, an act of
collective authorship -- a document written by 141,416 hands, each contributing
a single character in a language none of them can read, producing a text that
none of them can forge.

Whether this is beautiful or terrifying depends on your disposition. From the
perspective of pure mathematics, the construction is elegant: a simple algebraic
property (the hardness of factoring a product when the factors are destroyed)
scaled to sociological proportions. From the perspective of operational
security, it is a source of permanent anxiety: the entire edifice rests on the
assumption that physical destruction is final, that no forensic technique will
ever recover a \(\tau_i\) from a formatted hard drive or a power-cycled RAM chip
or the electromagnetic emanations of a laptop that existed for three hours in
2016 and then met a furnace.

Both perspectives are correct simultaneously. The ceremony is the most robust
trust-minimization technique we know how to build. It is also a construction
whose security guarantee is, at bottom, a claim about the physical world --
about the irreversibility of certain physical processes -- and the physical
world is full of surprises. Nobody has ever recovered toxic waste from a
destroyed ceremony machine. But nobody has ever proved, in any rigorous sense,
that it cannot be done.

\subsection{The Transparent Stage: No Secrets, No Waste, No
Ceremony}\label{the-transparent-stage-no-secrets-no-waste-no-ceremony}

The alternative is to build the stage from glass.

STARKs -- Scalable Transparent ARguments of Knowledge -- were introduced by
Ben-Sasson, Bentov, Horesh, and Riabzev in 2018 {[}Ben-Sasson et al., 2018{]}.
The word ``transparent'' is the operative one. A STARK requires no trusted setup
at all. The only ``setup'' is agreeing on a hash function -- a publicly known
algorithm like SHA-256 or BLAKE3.

A note on the arithmetic that underlies everything in this book: all the numbers
in a zero-knowledge system live in a ``finite field'' -- a set of numbers (say,
0 through some large prime minus 1) where arithmetic wraps around, like a clock.
On a 12-hour clock, 7 + 8 gives 3, not 15. Field arithmetic works the same way,
but the ``clock'' has billions or trillions of hours. Every time you see ``field
element,'' ``field arithmetic,'' or ``field size'' in the intervening chapters,
this is what we mean: clock arithmetic with a very large clock. There is no
secret. There is no toxic waste. There is no ceremony. There is nothing to
destroy because nothing dangerous was ever created.

The security of a STARK rests on a different mathematical assumption: collision
resistance of the hash function. It must be computationally infeasible to find
two inputs that produce the same output. This is a weaker trust assumption than
the discrete logarithm problem underlying KZG, and it has a property worth
remembering: collision-resistant hash functions are believed to resist quantum
computers. We will return to this point shortly, because it shapes every
architectural decision in the ZK ecosystem.

A third family deserves mention: Bulletproofs {[}Bunz et al., 2017{]}.
Bulletproofs are transparent (no ceremony needed) and produce logarithmic-size
proofs, but they require linear verification time. Monero adopted Bulletproofs
for its confidential transactions. The Inner Product Argument (IPA) at the core
of Bulletproofs inspired the Halo approach to recursion without pairings. But
Bulletproofs' linear verification cost makes them unsuitable for on-chain
verification of large circuits, and -- importantly -- they rely on the discrete
logarithm assumption. They are \emph{not} quantum-resistant.

The setup spectrum, ordered from most to least trust required, looks like this:

\begin{longtable}[]{@{}
  >{\raggedright\arraybackslash}p{(\linewidth - 8\tabcolsep) * \real{0.2000}}
  >{\raggedright\arraybackslash}p{(\linewidth - 8\tabcolsep) * \real{0.2000}}
  >{\raggedright\arraybackslash}p{(\linewidth - 8\tabcolsep) * \real{0.2000}}
  >{\raggedright\arraybackslash}p{(\linewidth - 8\tabcolsep) * \real{0.2000}}
  >{\raggedright\arraybackslash}p{(\linewidth - 8\tabcolsep) * \real{0.2000}}@{}}
\toprule\noalign{}
\begin{minipage}[b]{\linewidth}\raggedright
Setup Type
\end{minipage} & \begin{minipage}[b]{\linewidth}\raggedright
Trust Assumption
\end{minipage} & \begin{minipage}[b]{\linewidth}\raggedright
Example
\end{minipage} & \begin{minipage}[b]{\linewidth}\raggedright
Proof Size
\end{minipage} & \begin{minipage}[b]{\linewidth}\raggedright
PQ Secure?
\end{minipage} \\
\midrule\noalign{}
\arrayrulecolor{MidnightBlue}\midrule\arrayrulecolor{MidnightBorder}
\endhead
\bottomrule\noalign{}
\endlastfoot
Circuit-specific trusted & 1-of-N honest in ceremony & Groth16 & 192 bytes &
No \\
Universal trusted & 1-of-N honest in ceremony & PLONK, Marlin &
\textasciitilde880 bytes & No \\
Transparent (DL-based) & Discrete log hard & Bulletproofs & \textasciitilde700
bytes & No \\
Transparent (hash-based) & Collision resistance & STARKs & \textasciitilde100 KB
& Yes \\
\end{longtable}

Notice the pattern: as you move down the table, trust decreases but proof size
increases. No known construction breaks this tradeoff. The SoK paper on trusted
setups {[}Wang, Cohney, Bonneau, 2025{]} identifies the ``ideal polynomial
commitment scheme'' -- transparent setup, constant-size proofs, constant-time
verification -- as the central open problem in the field. Nobody has built it.
Nobody has proved it impossible. It is the holy grail.

\section{The Capex/Opex Framework}\label{the-capexopex-framework}

The most useful lens for understanding setup economics is the capital
expenditure / operating expenditure distinction.

\textbf{Ceremony costs are capex.} The Ethereum KZG ceremony required
approximately \$2-5 million in coordination, engineering, and security auditing.
This is a one-time cost. Because the resulting SRS is universal, every system
that uses it amortizes that cost. With fifty or more rollups sharing the same
SRS, the per-rollup cost converges to approximately \$60,000 -- trivial relative
to the annual operating budget of any serious blockchain project.

\textbf{Per-proof costs are opex.} Once the stage is built, the cost that
matters is the cost of each proof. Here the gap between trusted and transparent
setups opens wide:

\begin{itemize}
\tightlist
\item
  \textbf{Groth16 on-chain verification}: approximately 200,000-300,000 gas,
  which at typical Ethereum gas prices translates to about \$0.50-\$1.00. This
  is the cheapest on-chain verification available. Groth16 proofs are exactly
  192 bytes -- three elliptic curve group elements. Smaller than a tweet.
\item
  \textbf{Raw STARK on-chain verification}: approximately 2-5 million gas,
  translating to \$5-\$25. STARK proofs run roughly 100 KB -- about 500 times
  larger than Groth16 proofs.
\end{itemize}

These numbers explain why trusted setups persist despite their trust
assumptions. The argument is not philosophical. It is economic. And it explains
why the dominant production pattern in 2026 is neither pure trusted nor pure
transparent, but \emph{hybrid}: use a transparent STARK as the inner proof (no
ceremony required, post-quantum security for the computation), then wrap it in a
Groth16 or KZG-based outer proof for cheap on-chain verification. You get the
transparency of STARKs and the economics of SNARKs. The cost is complexity --
you now maintain two proof systems -- and the outer wrapper still requires a
trusted setup.

\section{The 141,416-Person Question}\label{the-141416-person-question}

Let us dwell on a question that sounds naive but turns out to be fundamental.

If the 1-of-N trust model says you only need one honest participant, why did
141,416 people participate in the Ethereum KZG ceremony? What is the marginal
value of participant number 141,417?

The mathematical answer: zero. One honest participant is sufficient. The 141,415
others add nothing to the mathematical security guarantee.

The operational answer: everything. Security is not a binary state -- it is a
confidence level. With six participants, a well-resourced adversary could
conceivably investigate, coerce, or compromise all six. With 141,416 anonymous
participants from around the world, many contributing through ephemeral browser
sessions, the attack surface becomes unmanageable. You would need to compromise
not just the people but their hardware, their network connections, their
randomness sources, their memories.

There is also a signaling dimension. A ceremony with 141,416 participants is a
public demonstration that the community takes the trust assumption seriously. It
is \emph{social consensus made visible}. The ceremony functions as both a
cryptographic protocol and a coordination event.

But a deeper question hides beneath this one.

``If the toxic waste is destroyed, how does anyone \emph{know} it was actually
destroyed?''

The epigraph that opened this chapter asked it. Here is the honest answer:
nobody knows. ``Destruction'' is a physical claim about a digital artifact. You
cannot prove you deleted something from your own computer. You can only promise
you did. The 1-of-N model rests on: ``I trust that at least one person is
honest, \emph{and} that their computer was not compromised, \emph{and} that no
backup exists anywhere, \emph{and} that no side-channel leaked the secret during
generation.''

This is weaker than trusting a single entity. But it is not zero trust. The word
``trustless'' does not apply here. The accurate word is ``trust-minimized.''
This distinction matters, and it will recur at every layer.

\subsection{The Sociology of Trust at
Scale}\label{the-sociology-of-trust-at-scale}

The 141,416-person question admits a deeper analysis than the one just given,
one that touches the foundations of how human societies construct shared belief.

Consider what it means to ``trust'' a cryptographic ceremony. You are not
trusting that a specific person told the truth. You are not trusting that a
specific institution acted in good faith. You are trusting a \emph{statistical
claim about a population}: that among 141,416 independent actors, at least one
behaved honestly. This is a fundamentally different kind of trust than anything
that existed before networked cryptography. It is closer to the trust you place
in thermodynamic laws -- not trust in any particular molecule, but trust in the
aggregate behavior of an astronomically large ensemble.

The sociologist Niklas Luhmann distinguished between \emph{confidence} and
\emph{trust}. Confidence is the passive expectation that things will continue to
work as they have -- you are confident the sun will rise, confident your bank
will not vanish overnight. Trust is the active decision to make yourself
vulnerable to another agent's choices -- you trust a surgeon, trust a business
partner, trust a babysitter. The 1-of-N ceremony model sits at the boundary
between these two concepts. You are not trusting any specific ceremony
participant. You are placing confidence in a system designed so that trust in
any specific individual is unnecessary. The confidence is in the design, not in
the people.

This distinction has practical consequences. When trust fails -- when a specific
person or institution betrays you -- the natural response is blame,
accountability, punishment. When confidence fails -- when the sun does not rise,
or when a financial system collapses -- the natural response is existential
shock, a rewriting of your model of the world. If the Ethereum KZG ceremony were
ever found to be compromised, the response would be closer to the second kind.
There would be no individual to blame, no villain to prosecute. There would only
be the discovery that a system designed to be trustworthy was not -- that the
statistical claim about the population was wrong.

This is why the ceremony's social dimensions matter as much as its cryptographic
ones. The ceremony must not only \emph{be} secure. It must be \emph{perceived}
as secure by a community large and diverse enough to sustain confidence over
decades. Perception is not a frivolous concern. It is an engineering
requirement. A ceremony that is cryptographically perfect but that the community
does not believe in provides no security in practice, because projects will not
build on an SRS the community doubts, and users will not deposit funds into
systems built on questioned foundations.

The Ethereum KZG ceremony achieved its social credibility through three
mechanisms that future ceremony designers should study.

First, \emph{radical openness}. The ceremony's source code was published. The
coordination protocol was specified in public documents. The queue was visible
in real time. Every contribution was logged in a publicly verifiable transcript.
This transparency did not make the ceremony secure -- security came from the
1-of-N assumption -- but it made the ceremony \emph{auditable}, which is a
precondition for sustained confidence.

Second, \emph{permissionless participation}. The decision to let anyone with an
Ethereum address contribute was not merely a design choice. It was a statement
about who the SRS belongs to. A ceremony restricted to credentialed
cryptographers would have been more efficient and arguably more secure in a
narrow technical sense -- each contribution would have been generated with
better operational hygiene. But it would have been a ceremony \emph{for} the
community, not \emph{by} the community. The permissionless design traded some
individual contribution quality for a qualitative shift in collective ownership.

Third, \emph{diversity of entropy sources}. Participants contributed from
different hardware, different operating systems, different network
configurations, different physical locations, different jurisdictions. An
adversary attempting to compromise the ceremony would need to mount simultaneous
attacks across an impossibly heterogeneous set of targets. This is the
cryptographic equivalent of biodiversity: monocultures are fragile, but
ecosystems with many species are resilient. The heterogeneity of 141,416
independent computing environments is itself a security property, one that no
formal model fully captures.

There is a fourth mechanism, less often discussed: the \emph{unforgeable cost of
participation}. Each contributor spent real time -- waiting in the queue,
running the computation, verifying the receipt. This time cannot be faked or
automated away entirely (anti-Sybil measures ensured that each Ethereum address
could contribute only once). The aggregate time investment of 141,416 people,
even at a few minutes each, represents thousands of person-hours of attention.
This attention is a scarce resource, and its expenditure signals genuine
commitment in a way that no amount of on-chain capital can replicate. A ceremony
with 141,416 participants who each spent three minutes is harder to dismiss than
a ceremony with six participants who each spent three days, even though the
per-person investment is radically different, because the \emph{collective}
investment is distributed across a population too large to coordinate against.

The game-theoretic analysis reinforces this intuition. In a ceremony with six
participants, the cost of corrupting the ceremony is the maximum of the costs of
corrupting each participant -- you need all six. With 141,416 participants, the
cost is still the maximum of individual corruption costs (since you need all of
them), but the \emph{minimum} of those individual costs is now determined by the
most resilient participant in a population of 141,416. Among that many people,
there is almost certainly someone whose operational security is excellent, whose
hardware is air-gapped, whose randomness source is physical, and whose
motivation is ideological rather than financial. The ceremony's security is, in
a precise sense, determined by its strongest link rather than its weakest.

This is the inverse of the chain metaphor that dominates security thinking. A
chain is only as strong as its weakest link. A ceremony is only as
\emph{insecure} as its \emph{strongest} link. This inversion -- this structural
optimism built into the mathematics of the 1-of-N model -- is what makes
mass-participation ceremonies viable. You do not need to ensure that every
participant is careful. You need to ensure that the population is large enough
and diverse enough that at least one participant, somewhere, is careful enough.

The question of how large is ``large enough'' has no clean mathematical answer,
because it depends on the adversary model. Against a lone hacker, six
participants may suffice. Against a corporate adversary, ninety. Against a
nation-state with global surveillance capabilities, perhaps tens of thousands.
Against a coalition of nation-states -- the most paranoid threat model, but not
an absurd one for infrastructure securing tens of billions of dollars -- perhaps
hundreds of thousands. The Ethereum ceremony's 141,416 participants do not
guarantee security against any specific adversary. They guarantee that the
\emph{sociological} bar for a successful attack is extraordinarily high, higher
than any previous cryptographic ceremony has set, and plausibly higher than any
adversary can clear.

Plausibly. Not provably. And that gap -- between plausible and provable -- is
where the honest assessment of ceremony security must live.

\section{The Bug That Was Not a Ceremony
Failure}\label{the-bug-that-was-not-a-ceremony-failure}

In 2019, a vulnerability designated CVE-2019-7167 surfaced in the BCTV14
construction used by Zcash {[}BCTV14{]}. The bug was devastating: it would have
allowed unlimited counterfeiting of Zcash tokens. And it had nothing to do with
the trusted setup ceremony.

The flaw lived in the cryptographic construction itself. The BCTV14 protocol
included an extra element in the SRS that was unnecessary for the proof system
but that an attacker could exploit to forge proofs. The Zcash ceremony could
have been conducted perfectly -- every participant honest, every piece of toxic
waste destroyed -- and the system would still have been vulnerable.

This case study carries a lesson that extends far beyond Layer 1: \emph{ceremony
integrity is necessary but not sufficient}. You can build a perfect stage and
still get the show wrong. The construction must be correct independently of the
ceremony. The program must be correct independently of the proof system.
Security is not a chain where one strong link compensates for a weak one. It is
a conjunction: every link must hold simultaneously.

Sean Bowe and Ariel Gabizon discovered the BCTV14 bug before anyone exploited
it. The team deployed a fix transparently. But the episode is a permanent
reminder that the most dangerous vulnerabilities are not the ones in the
ceremony. They are the ones in the mathematics that the ceremony exists to
protect.

\section{Universal versus Circuit-Specific
Setups}\label{universal-versus-circuit-specific-setups}

The transition from circuit-specific to universal setups changed how
zero-knowledge systems are built. Understanding it requires understanding what
``circuit-specific'' means in practice.

Groth16 {[}Groth, 2016{]} produces the most compact proofs ever constructed:
exactly 192 bytes, three group elements, verified with three pairing operations.
No other system comes close. But Groth16's setup is circuit-specific. The SRS
encodes the structure of a particular circuit: the specific polynomial
constraints, the wiring, the gates. Change the circuit -- add a feature, fix a
bug, upgrade the protocol -- and you need a new ceremony. New coordination. New
participants. New toxic waste to destroy. At \$2-5 million per ceremony, this is
unsustainable for any system that evolves.

PLONK {[}Gabizon, Williamson, Ciobotaru, 2019{]} and Marlin {[}Chiesa et al.,
2019{]} solved this by making the SRS \emph{universal}. The SRS encodes a
generic mathematical structure (powers of a secret value on an elliptic curve,
the same structure as KZG commitments). Any circuit up to a maximum size can use
the same SRS. The ``derived'' per-circuit setup -- selecting the evaluation
domain, computing selector polynomials, generating proving and verification keys
-- is entirely \emph{deterministic and public}. No new secrets. No new ceremony.
No new toxic waste.

Midnight, the privacy-focused blockchain, adopted this model. Midnight uses a
PLONK-family proof system (specifically, a variant of Halo2) with a universal
SRS on the BLS12-381 elliptic curve {[}Midnight Developer Guide, 2026{]}. The
Compact compiler takes each smart contract, compiles it to a circuit
intermediate representation (ZKIR), and derives per-circuit proving and
verification keys from the universal SRS. One ceremony serves all contracts. The
per-contract compilation runs in seconds, produces deterministic output (same
source code yields same keys), and requires no trust beyond the original
ceremony.

The architecture diagram from Midnight's developer documentation shows this
cleanly:

\begin{itemize}
\tightlist
\item
  \texttt{midnight-trusted-setup} -- the ceremony, producing the universal SRS
\item
  \texttt{midnight-zk} -- the proof system (PLONK / Halo2 / BLS12-381),
  consuming the SRS
\item
  \texttt{compactc\ compile} -- the compiler, deriving per-circuit keys
  deterministically
\end{itemize}

Each compiled contract produces three artifact types: TypeScript bindings for
the dApp frontend, ZKIR circuit files for the proof server, and
proving/verification key pairs derived from the universal SRS. The proving key
for a simple counter contract is 13.7 KB. The verification key is 1.3 KB. These
keys inherit their security from the original ceremony -- no additional trust
assumptions enter during compilation.

The universal setup model has a consequence for the capex/opex analysis that is
easy to miss: the capital expenditure of the ceremony is paid once and amortized
indefinitely. Every new contract, every new circuit, every upgrade deploys under
the same SRS. The marginal cost of adding a new application to the ecosystem
converges to the cost of compilation -- effectively zero.

\section{The Quantum Shelf Life}\label{the-quantum-shelf-life}

Every pairing-based proof system -- Groth16, PLONK, Marlin, Sonic, KZG -- rests
on the hardness of the discrete logarithm problem on elliptic curves. Shor's
algorithm, running on a sufficiently powerful quantum computer, solves the
discrete logarithm problem in polynomial time.

This is not ``might break.'' This is ``does break, given sufficient hardware.''
The question is when, not whether.

The implications for Layer 1 are stark. A quantum computer would not need to
compromise any ceremony participant. It would not need to break into anyone's
hardware or bribe anyone. It would simply take the \emph{public SRS} -- the list
of elliptic curve points that everyone can see -- and extract the original
trapdoor from it. Mathematically, irreversibly, without detection.

Once the trapdoor is extracted, every proof ever generated under that SRS
becomes suspect. The attacker can forge proofs of false statements
indistinguishable from legitimate ones. The six-figure ceremony becomes
irrelevant. The toxic waste was not in anyone's memory or on anyone's hard drive
-- it was encoded in the public parameters, waiting for a quantum computer
powerful enough to read it.

Conservative estimates place cryptographically relevant quantum computers --
machines capable of breaking 256-bit elliptic curve cryptography -- in the
2032-2035 timeframe. The National Institute of Standards and Technology (NIST)
finalized its first post-quantum cryptography standards (FIPS 203, 204, 205) in
August 2024 and published NIST IR 8547, a deprecation roadmap that targets 2035
for the retirement of pre-quantum cryptographic algorithms.

This creates a concept we will call the \emph{quantum shelf life} of a trusted
setup. A KZG ceremony conducted in 2023 produces an SRS that is secure against
classical computers indefinitely but has a finite lifespan against quantum
adversaries. If that SRS secures a blockchain expected to operate for 20 years,
the setup's shelf life may expire before the system's intended lifetime ends.

STARKs, by contrast, rely on collision-resistant hash functions, which are
believed to resist quantum computers. (The caveat ``believed to be'' is doing
real work here: Grover's algorithm provides a quadratic speedup for hash
preimage search, which means security parameters need to be doubled -- 128-bit
classical security becomes 64-bit quantum security -- but this is a quantitative
adjustment, not a qualitative break.) A transparent STARK-based setup has no
quantum shelf life problem because there is no trapdoor to extract.

This brings us to the concept of \emph{option value}. A transparent setup
preserves the option to migrate to post-quantum cryptography without
re-ceremony. A trusted setup on a pairing-friendly curve burns this option: if
the post-quantum deadline arrives and your SRS lives on BLS12-381, you need a
new setup from scratch. Given the NIST 2035 timeline, systems designed today
with 10+ year lifespans should seriously consider whether the performance
advantages of pairing-based setups justify the loss of post-quantum migration
flexibility.

Midnight's choice of BLS12-381 illustrates this tension concretely. BLS12-381
provides approximately 128-bit classical security but offers zero post-quantum
security. Every component of Midnight's proof system -- the SRS, the polynomial
commitments, the proof verification -- rests on assumptions that quantum
computers break. Midnight originally experimented with Pluto-Eris curves (which
enabled recursive proof composition), but switched back to BLS12-381 for
pragmatic reasons: faster proof generation, wider ecosystem compatibility, and
better tooling support. The switch was a deliberate choice to optimize for
present-day deployment at the cost of future quantum vulnerability.

The alternative approach -- exemplified by systems like Neo {[}Nguyen, Setty,
2025{]} -- uses lattice-based cryptography (Module-SIS and Module-LWE
assumptions) that is conjectured to resist quantum computers. The tradeoff:
larger proofs, more expensive verification, and a transparent setup that
requires no ceremony at all. The trust assumption shrinks from ``1-of-N ceremony
participants were honest'' to ``the lattice problem is hard'' -- a purely
mathematical assumption with no sociological component.

Neither choice is unambiguously correct. They represent different bets on the
future: Midnight bets that pairing-based cryptography will remain secure long
enough for the system to establish itself and migrate later. Lattice-based
systems bet that the quantum threat justifies paying the performance cost now.
Nobody knows which bet is right, because nobody knows when quantum computers
will arrive at cryptographic scale.

\section{BN254's Eroding Security Margin}\label{bn254s-eroding-security-margin}

There is a more immediate security concern than quantum computers, and it
affects the most widely deployed curve in the ZK ecosystem.

BN254 (also called alt\_bn128 or BN128) is the curve hardcoded into Ethereum's
elliptic curve precompile opcodes. Every Groth16 proof verified on Ethereum's
base layer uses BN254. But the estimated security of BN254 has been eroding. The
Tower Number Field Sieve (Tower NFS) -- a family of discrete log algorithms that
exploits the tower structure of extension fields -- has reduced BN254's
estimated security from 128 bits to approximately 100 bits {[}Kim, Barbulescu,
2016; Menezes, Sarkar, Singh, 2016{]}.

One hundred bits of security is not broken. But it sits below the 128-bit
threshold that NIST mandates for new cryptographic deployments. The ZK ecosystem
has been slowly migrating from BN254 to BLS12-381, which provides a comfortable
128-bit security margin even under Tower NFS analysis. Migration is slow,
though: existing smart contracts reference the BN254 precompile directly, and
changing the curve means changing the verification logic, which means upgrading
every contract that verifies proofs.

The setup layer casts a long shadow. Choices made at Layer 1 -- which curve,
which parameters, which ceremony -- propagate forward through years or decades
of deployment. The curve choice that was state-of-the-art in 2018 shows its age
by 2026. The ceremony that was sufficient in 2023 may be insufficient by 2035.
Building the stage is not a one-time act. It is a commitment with a time
horizon.

\section{The ADOPT Framework}\label{the-adopt-framework}

The most comprehensive survey of trusted setup ceremonies, the SoK paper by
Wang, Cohney, and Bonneau {[}Wang, Cohney, Bonneau, 2025{]}, catalogs over forty
real-world ceremonies and evaluates them against five properties (the ADOPT
framework):

\begin{itemize}
\tightlist
\item
  \textbf{A}vailable: Can anyone access and verify the SRS?
\item
  \textbf{D}ecentralized: Is there a single coordinator who could manipulate the
  ceremony?
\item
  \textbf{O}pen: Can anyone participate without permission?
\item
  \textbf{P}ersistent: Will the ceremony data survive long-term for auditing?
\item
  \textbf{T}ransparent: Is every step of the ceremony publicly observable?
\end{itemize}

Their finding: \emph{no existing ceremony satisfies all five properties}. The
Ethereum KZG ceremony scores well on openness and availability but still relied
on a coordinating entity (the Ethereum Foundation). Many older ceremonies fail
on persistence -- the intermediate transcript data for projects like Hermez has
already become unrecoverable just a few years later.

This is not a failure of any specific project. It is a structural limitation of
the ceremony model. Ceremonies are social events, and social events are messy.
The gap between the abstract ``1-of-N honest participant'' security model and
the operational reality of running a ceremony with hundreds of thousands of
participants -- each using different hardware, different software, different
randomness sources, across different jurisdictions -- is wide. The protocol can
be mathematically perfect. The ceremony is always imperfect.

The transparent alternative -- STARKs, hash-based commitments, no ceremony at
all -- avoids this entire problem class. There is no ceremony to evaluate, no
transcript to preserve, no coordinator to trust. The cost is measured in proof
size and verification time, not in coordination complexity and social trust.

\section{Midnight's BLS12-381 Stage}\label{midnights-bls12-381-stage}

Let us trace the full Layer 1 story through a real system.

Midnight is a privacy-focused blockchain built on the Cardano ecosystem. It uses
BLS12-381 with a PLONK-family proof system. Here is how its Layer 1 choices
cascade through the stack:

\textbf{The ceremony.} Midnight operates a \texttt{midnight-trusted-setup}
repository for conducting its Powers-of-Tau ceremony. The ceremony produces a
universal SRS on BLS12-381 -- a set of elliptic curve points encoding powers of
a secret trapdoor. The SRS bounds the maximum circuit size: you cannot prove
statements about circuits larger than the SRS supports.

\textbf{Per-circuit keys.} The Compact compiler (\texttt{compactc\ compile})
takes each contract's circuit description (ZKIR) and the universal SRS to
produce per-circuit proving and verification keys. This is deterministic: same
source plus same compiler yields same keys. No new trust assumption enters. A
simple counter contract produces a 13.7 KB proving key and a 1.3 KB verification
key.

\textbf{Deployment flow.} The proving key goes to the proof server (running
locally at localhost:6300 -- witnesses never cross the network). The
verification key reaches the blockchain via \texttt{submitInsertVerifierKeyTx}.
Every node that validates transactions uses this key to check proofs.

\textbf{The Pluto-Eris detour.} Midnight originally adopted Pluto-Eris, a cycle
of curves that enabled recursive proof composition (KZG on Pluto, IPA on Eris).
In April 2025, the team announced a switch back to BLS12-381. The reason: faster
proof generation, better ecosystem tooling, and support for higher transaction
volumes. Theoretical optimality yielded to engineering pragmatism. This is a
pattern we will see repeatedly: the choices that win in production are not
always the choices that win in papers.

\textbf{The quantum exposure.} BLS12-381 provides approximately 128-bit
classical security but zero post-quantum security. Every ZK proof in Midnight --
state transitions, shielded transfers, Zswap privacy operations -- becomes
forgeable if a quantum computer extracts the trapdoor from the SRS. Midnight's
entire proof system would need replacement in a post-quantum world. The
migration path is unclear.

The Midnight case study illustrates a broader truth about Layer 1: the setup
choice is a bet on the future. Midnight bet on BLS12-381 because it is the most
mature, best-tooled, most widely deployed pairing-friendly curve in the
ecosystem. That bet optimizes for today's performance and ecosystem
compatibility. Whether it survives the quantum transition is an open question --
one that the Midnight team will face, along with every other pairing-based
system, before the decade is out.

\section{Option-Value Analysis}\label{option-value-analysis}

We close the economic analysis with a concept from financial options theory that
the ZK ecosystem has not yet internalized: \emph{option value}.

A trusted setup on BLS12-381 buys performance today but forecloses certain
futures. When -- not if -- post-quantum migration becomes necessary, every
system on BLS12-381 will need a new setup. New ceremony. New coordination. New
toxic waste. If the system has accumulated ten years of state and ten million
users, the migration cost is not just the \$2-5 million of the ceremony. It is
the coordination cost of upgrading every node, every contract, every verifier
key.

A transparent setup on hash-based primitives costs more per proof today but
preserves the option to migrate without re-ceremony. The hash function can be
swapped. The security parameters can be adjusted. No toxic waste needs to be
re-destroyed because none was ever created.

What is this option worth? The answer depends on your estimate of the quantum
timeline. If you believe cryptographically relevant quantum computers are 20+
years away, the option has low present value, and the performance advantage of
pairing-based setups dominates. If you believe the timeline is 10-15 years --
consistent with NIST's 2035 deprecation target -- the option is substantial. And
if the ``Harvest Now, Decrypt Later'' threat model applies (where adversaries
record encrypted data today to decrypt it when quantum computers arrive), the
option may be critical even for systems launched tomorrow.

The Federal Reserve's FEDS 2025-093 working paper on quantum threats to
financial infrastructure uses exactly this option-value framework. Their
conclusion: systems with 10+ year operational lifespans should use post-quantum
or post-quantum-ready primitives. Transparent setups satisfy this criterion by
default. Trusted setups on pairing-friendly curves do not.

\section{The Setup Tradeoff}\label{the-setup-tradeoff}

The stage is built. Let us take stock.

Layer 1 is where the trust story begins. The choice between trusted and
transparent setups is not a philosophical preference -- it is an economic,
operational, and temporal decision with consequences that persist for the
lifetime of the system. Trusted setups produce smaller proofs and cheaper
verification but require a ceremony, create toxic waste, and carry a quantum
shelf life. Transparent setups avoid all ceremony-related risks but pay a cost
in proof size and verification expense. The hybrid approach -- transparent inner
proof, trusted outer wrapper -- dominates production because it captures most of
the benefits of both, at the price of added complexity.

The 1-of-N trust model is a genuine advance over trusting a single entity, but
it is trust-minimized, not trustless. The BCTV14 bug proves that ceremony
integrity alone is insufficient -- the construction must be independently
correct. The quantum shelf life means that no pairing-based setup is permanent.
And the ADOPT framework reveals that no ceremony conducted to date achieves the
ideal of full availability, decentralization, openness, persistence, and
transparency.

We have introduced our two characters -- the magician (prover) and the audience
(verifier) -- and we have seen that even the first layer of their interaction
involves non-trivial trust decisions. The magician cannot perform without a
stage. The audience cannot judge without a ruler. And someone, somewhere, had to
build both.

Now, the stage is ready. The mathematical parameters are in place. The SRS is
published. The toxic waste is (we hope) destroyed. The verifier keys are derived
and deployed.

The magician needs a script. She needs to express her computation -- the thing
she wants to prove -- in a language the proof system can understand. That
language choice turns out to be the single most consequential decision for the
security of the entire system. Not because of the cryptography. Because of the
bugs. Sixty-seven percent of real-world SNARK vulnerabilities trace back to the
script, not the stage -- to the program layer, not the ceremony or the proof
system. The next chapter explains why, and what the field is doing about it.

The stage is built. Now, who writes the script?

\partdivider{Part II: The Craft}

\emph{The audience has seen the stage. They know the rules. Now the curtain
rises, and the real work begins -- backstage, in the mathematics, where every
step must be recorded, encoded, sealed, and verified. The next six chapters
trace the magician's craft from the first line of code to the final proof on the
blockchain.}

\chapter{Chapter 3: Choreographing the
Act}\label{chapter-3-choreographing-the-act}

\emph{Layer 2 -- Language}

\section{RISC-V Won. Why Taxonomy Still
Matters.}\label{risc-v-won.-why-taxonomy-still-matters.}

If RISC-V has won -- if eight of ten major zero-knowledge virtual machines now
target the same general-purpose instruction set -- why bother presenting a
taxonomy of competing philosophies at all? Why not just say ``write Rust,
compile to RISC-V, the proof system handles the rest''?

Because the taxonomy is not about which instruction set the machine runs. It is
about what \emph{the developer sees}. And what the developer sees determines
what bugs the developer makes. Those bugs -- not cryptographic breaks, not
quantum computers, not governance attacks -- are the single largest source of
real-world failures in zero-knowledge systems. Sixty-seven percent of all known
SNARK vulnerabilities are under-constrained circuits: programs where the
developer said less than they meant, and the proof system happily proved false
statements as a result.

The language layer is where the magician writes the choreography for the act.
The choice of notation determines what mistakes are possible. Some notations let
the performer accidentally step into the audience's view. Others physically
prevent it. That distinction is worth understanding, even in a world where
RISC-V has won the instruction set war.

\begin{quote}
\textbf{The Running Example: A Sudoku Proof}

To ground the abstractions that follow, we will trace a single computation
through every layer of the stack. The computation: proving you know the solution
to a 4x4 Sudoku puzzle without revealing it.

The puzzle has these givens:

\begin{verbatim}
+---+---+---+---+
| 1 |   |   | 4 |
+---+---+---+---+
|   | 4 | 1 |   |
+---+---+---+---+
|   | 1 |   |   |
+---+---+---+---+
| 4 |   |   | 1 |
+---+---+---+---+
\end{verbatim}

The program checks: every row contains \{1,2,3,4\} with no repeats, every column
contains \{1,2,3,4\}, every 2x2 box contains \{1,2,3,4\}, and every filled cell
matches the given clue. The prover knows the solution. The verifier knows only
the puzzle. At Layer 2, this is a program. We follow it through every layer to
come.
\end{quote}

\section{From Circuits to Virtual Machines: A Brief
Evolution}\label{from-circuits-to-virtual-machines-a-brief-evolution}

To understand where we are, we need to understand where we came from.

The first generation of zero-knowledge programming looked nothing like
programming. In 2018, if you wanted to prove a computation in zero knowledge,
you wrote \emph{circuits} -- not programs, but direct descriptions of
mathematical relationships. The dominant tool was Circom, a domain-specific
language created by Jordi Baylina and the iden3 team. In Circom, you did not
write \texttt{if\ balance\ \textgreater{}=\ amount\ then\ approve}. You wrote
constraint templates: mathematical equations that the proof system would verify.
The developer was simultaneously writing two programs in one file -- one that
computed the witness (the private data), and one that generated the constraints
(the mathematical rules). These two programs had to agree perfectly on every
input. When they did not, the result was an under-constrained circuit: a proof
system that would certify false statements as true.

Imagine asking a playwright to write both the script and the stage directions in
a single document, using two different notations that had to be perfectly
synchronized, with no compiler to check whether they matched. That is what early
ZK development felt like.

Circom was powerful. It gave developers complete control over the constraint
system. But that control came at a cost. The Chaliasos SoK confirmed the scale
of the problem: 95 of 141 catalogued vulnerabilities were under-constrained
circuits -- the epidemic described at the opening of this chapter. The Tornado
Cash bug was a single character: \texttt{=} where \texttt{\textless{}==} was
needed. One character. Complete soundness break. A malicious prover could
generate proofs for false statements, and the verifier would accept them.

The second generation asked a different question: what if the developer never
saw the constraints at all? What if they wrote a program in a language they
already knew, and a compiler handled the translation to mathematics?

Cairo, created by StarkWare in 2021, pioneered this approach. Cairo defined a
new instruction set architecture -- a virtual CPU designed from the ground up so
that every instruction's execution could be efficiently encoded as polynomial
constraints. The developer wrote programs. The compiler generated constraints.
The proof system verified the constraints. The developer never touched a
mathematical equation.

But Cairo required learning a new language, a new toolchain, a new way of
thinking about computation. The programs you had already written -- in Rust, in
C++, in Python -- could not run on Cairo. You had to rewrite everything.

The third generation asked the obvious follow-up: what if we proved a processor
that developers already targeted? What if the instruction set was not some
exotic ZK-native design, but plain RISC-V -- the open standard that Rust, C, and
C++ compilers already produce code for?

This is the generation we live in now. SP1 (Succinct), RISC Zero, Airbender
(ZKsync), ZisK (the team formerly known as Polygon Hermez), and Pico Prism all
prove RISC-V execution. The developer writes standard Rust. The compiler targets
standard RISC-V. The proof system proves the execution trace. The developer may
never know they are working with zero-knowledge proofs at all.

But this triumphant narrative -- circuits to custom VMs to RISC-V -- leaves out
a fourth thread. There is another approach, one that does not fit the
evolutionary story. It does not prove a processor at all. It proves \emph{state
transitions}. And its compiler does something no instruction-set-based approach
can do: it prevents privacy leaks at compile time.

\section{The Four Philosophies}\label{the-four-philosophies}

Every approach to Layer 2 answers the same question differently: how should a
developer express a computation that will be proven in zero knowledge?

Four distinct philosophies have emerged. Each makes a different bet about what
matters most.

\textbf{Philosophy A: EVM-Compatible.} Prove the Ethereum Virtual Machine
itself. The developer writes Solidity, the same language used for ordinary
Ethereum smart contracts. The proof system executes the EVM bytecode and proves
that execution was correct. The advantage is obvious: every existing Ethereum
application ``just works.'' The Solidity ecosystem -- its tooling, its auditors,
its libraries, its millions of lines of battle-tested code -- transfers
wholesale.

The leading examples are Scroll, with \$748 million in total value locked, and
Linea (ConsenSys), with over \$2 billion. Both prove EVM execution using
different strategies: Scroll uses a custom zkEVM circuit, while Linea has been
converging toward a RISC-V backend with EVM compatibility layered on top.

The disadvantage is equally obvious: the EVM was not designed for proving. It
has 256-bit arithmetic (proof systems prefer 31-bit or 64-bit fields). It has
dynamic gas metering. It has a complex memory model. Proving every quirk of the
EVM is computationally expensive -- like translating a novel into a language
that has no word for half the concepts.

\textbf{Philosophy B: ZK-Native ISA.} Design a new instruction set from scratch,
optimized for proving. The developer learns a new language, but every
instruction translates efficiently into polynomial constraints. Cairo is the
canonical example. Its instruction set was co-designed with StarkWare's STARK
proof system. Layer 2 was literally shaped by Layer 4 -- the language exists
\emph{because} of the arithmetization.

Cairo's execution trace model -- columns for program counter, flags, operands,
with polynomial constraints between adjacent rows -- became the template for
every subsequent zkVM. The Goldilocks prime field (\(2^{64} - 2^{32} + 1\)) that
Cairo adopted has since been used by RISC Zero, SP1, and others. Cairo bet on
algebraic efficiency over developer familiarity, and for the Starknet ecosystem,
that bet paid off: it is the most battle-tested ZK system after Circom, handling
real assets on mainnet since 2020.

The cost is ecosystem isolation. Cairo programs do not run anywhere except
Starknet. The tooling, libraries, and developer community are Cairo-specific.
Every line of code is a bet on one ecosystem. But that bet has paid returns:
Starknet processes real assets, real DeFi, real NFTs. Cairo is not a research
project. It is production infrastructure.

The evolution from Cairo 0 to Cairo 1.0 is itself instructive. The original
Cairo was barely recognizable as a programming language -- it looked more like
assembly with syntactic sugar. Cairo 1.0 aligned with mainstream language
features: Rust-like syntax, a borrow checker, generic types. The lesson for the
field: even a language designed for proof efficiency eventually converges toward
familiar developer ergonomics, because adoption requires accessibility.

\textbf{Philosophy C: General-Purpose ISA.} Prove a standard processor. RISC-V
has won this category so decisively that it is no longer a competition. SP1
(Succinct), RISC Zero (with its R0VM 2.0 reducing Ethereum block proving from 35
minutes to 44 seconds), Airbender (ZKsync, achieving 21.8 million cycles per
second on a single H100 GPU), ZisK (the ex-Polygon team), and Pico Prism all
target RISC-V.

The appeal is obvious: the developer writes Rust, C, or C++. The compiler
produces RISC-V machine code. The zkVM executes that code and generates a proof.
Standard toolchains, standard debuggers, standard testing frameworks. The
ZK-specific complexity hides behind the compilation boundary.

Even projects that started with EVM compatibility are converging on RISC-V.
ZKsync's Airbender proves RISC-V execution and layers EVM compatibility on top.
The trend is clear: RISC-V is the assembly language of the zero-knowledge world.

\textbf{Philosophy D: Application-Specific DSL.} This is the philosophy that
does not fit the evolutionary narrative. Instead of proving a processor, these
languages prove \emph{state transitions}. Instead of hiding ZK complexity behind
a compilation boundary, they make privacy a first-class language concept.

Three languages define this category.

\emph{Compact} (Midnight/IOG) is a TypeScript-like language for
privacy-preserving smart contracts. Its compiler produces three artifacts from a
single source file: a ZKIR circuit (the constraint system), TypeScript bindings
(the application interface), and proving keys (the cryptographic setup
material). No other ZK language generates all three from one source. More
importantly, Compact's compiler includes a \emph{disclosure analysis} pass that
statically rejects any program where a private value might leak to a public
surface without explicit developer consent. We will return to this in detail.

\emph{Noir} (Aztec Labs) is a Rust-inspired, backend-agnostic ZK language. It
compiles to an Abstract Circuit Intermediate Representation (ACIR). Noir does
not target a specific proof system -- it targets multiple backends. This makes
it the closest thing the ZK world has to a ``write once, prove anywhere''
language.

Noir 1.0 was pre-released in late 2025. It became an officially recognized
language on GitHub. NoirCon conferences have been held (NoirCon0 in November
2024). The ecosystem includes over 600 projects and 900 GitHub stars. Key
adopters include zkEmail, zkPassport, and zkLogin. Aztec's Ignition Chain went
live in November 2025 as the first decentralized L2 on Ethereum, with 185+
operators across 5 continents and 3,400+ sequencers -- and its core cryptography
was rewritten in Noir.

Noir breaks the three-philosophy taxonomy because it is neither ISA-based nor
chain-specific: it is a universal circuit language. Its privacy model is
straightforward -- all inputs are private unless explicitly declared
\texttt{pub} -- but it lacks the compile-time disclosure analysis that Compact
provides. The developer bears responsibility for correctly managing the
public/private boundary. In exchange, Noir programs can target any proving
backend that accepts ACIR, making them portable across proof systems in a way
that no other ZK language achieves.

\emph{Leo} (Aleo) is a privacy-first language for the Aleo blockchain, with
syntax borrowing from Rust and TypeScript. Leo targets Aleo's record-based
(UTXO-like -- a UTXO, or Unspent Transaction Output, is a model where each
digital coin is a discrete object created by one transaction and consumed by
another, like physical bills in a wallet) privacy model and includes hooks for
formal verification. With over 400,000 CLI downloads, Leo represents the
privacy-specialized variant of the application DSL approach.

To understand why these three languages matter -- why they are not merely
syntactic alternatives to writing Rust and compiling to RISC-V -- you need to
see what development actually looks like in each one. The differences are not
cosmetic. They are structural. Each language imposes a different mental model on
the developer, and that mental model determines what classes of bug are
possible, what classes of privacy leak are preventable, and what the compiler
can guarantee before a single proof is generated.

\subsection{Compact: Privacy by
Compilation}\label{compact-privacy-by-compilation}

Consider a simple scenario: a private token transfer. The sender wants to prove
they have sufficient balance without revealing what that balance is. In Compact,
the developer writes something that looks, at first glance, like an ordinary
TypeScript function:

\begin{Shaded}
\begin{Highlighting}[]
\ImportTok{export}\NormalTok{ circuit }\FunctionTok{transfer}\NormalTok{(}
\NormalTok{  recipient}\OperatorTok{:}\NormalTok{ Bytes}\OperatorTok{\textless{}}\DecValTok{32}\OperatorTok{\textgreater{},}
\NormalTok{  amount}\OperatorTok{:}\NormalTok{ Unsigned Integer}
\NormalTok{)}\OperatorTok{:}\NormalTok{ [] \{}
  \KeywordTok{const}\NormalTok{ my\_balance }\OperatorTok{=} \FunctionTok{disclose}\NormalTok{(}\FunctionTok{get\_balance}\NormalTok{())}\OperatorTok{;}
  \FunctionTok{assert}\NormalTok{(my\_balance }\OperatorTok{\textgreater{}=}\NormalTok{ amount}\OperatorTok{,} \StringTok{"insufficient funds"}\NormalTok{)}\OperatorTok{;}

  \KeywordTok{const}\NormalTok{ new\_sender\_balance }\OperatorTok{=}\NormalTok{ my\_balance }\OperatorTok{{-}}\NormalTok{ amount}\OperatorTok{;}
  \KeywordTok{const}\NormalTok{ new\_recipient\_balance }\OperatorTok{=} \FunctionTok{get\_recipient\_balance}\NormalTok{(recipient) }\OperatorTok{+}\NormalTok{ amount}\OperatorTok{;}

\NormalTok{  ledger}\OperatorTok{.}\AttributeTok{sender\_balances}\NormalTok{[}\FunctionTok{sender}\NormalTok{()] }\OperatorTok{=}\NormalTok{ new\_sender\_balance}\OperatorTok{;}
\NormalTok{  ledger}\OperatorTok{.}\AttributeTok{recipient\_balances}\NormalTok{[recipient] }\OperatorTok{=}\NormalTok{ new\_recipient\_balance}\OperatorTok{;}
\NormalTok{\}}
\end{Highlighting}
\end{Shaded}

The keyword \texttt{circuit} is the first departure from ordinary programming.
This function will not execute on a server. It will be compiled into a
zero-knowledge circuit -- a set of polynomial constraints that a prover can
satisfy and a verifier can check. Every variable inside this function will
become a wire in that circuit. Every operation will become a gate.

The keyword \texttt{disclose} is the second departure, and the more
consequential one. The function \texttt{get\_balance()} is a witness function --
it retrieves the sender's private balance from off-chain storage. That value is,
by default, invisible to anyone but the prover. The \texttt{disclose()} call is
the developer's explicit declaration: ``I acknowledge that this private value
will influence public state.'' Without it, the compiler rejects the program. Not
at runtime. Not during testing. At compile time, before any proof is generated,
before any key material is created, before any circuit is emitted.

The practical consequence is that a developer cannot accidentally write a
transfer function that leaks the sender's balance. They can intentionally leak
it -- \texttt{disclose()} is an explicit consent mechanism, not a prohibition.
But the accidental case, which accounts for the majority of real-world privacy
bugs, is eliminated by the compiler's disclosure analysis pass.

The compilation itself is worth understanding. Compact's 26-pass nanopass
pipeline transforms the source through a sequence of increasingly specialized
intermediate languages. The program begins as \texttt{Lsrc} -- essentially the
developer's TypeScript-like code. It passes through type inference
(\texttt{Ltypes}), where the compiler determines the bit-width of every value.
It passes through disclosure analysis (\texttt{Lnodisclose}), where the compiler
traces every data-flow path from witness inputs to public outputs. It passes
through loop unrolling (\texttt{Lunrolled}), where bounded loops are expanded
into straight-line code -- necessary because ZK circuits have no concept of
iteration. It passes through circuit flattening (\texttt{Lflattened}), where
nested expressions are decomposed into individual gates. And it emerges as ZKIR:
a JSON-formatted circuit description ready for the proof system.

At no point does the developer interact with constraints, gates, or polynomials.
The entire mathematical substrate is hidden behind the compilation boundary. The
developer writes TypeScript-like code. The compiler produces a zero-knowledge
circuit. The gap between intent and implementation -- the gap where
under-constrained bugs live -- is bridged by the compiler, not by the developer.

\subsection{Noir: Write Once, Prove
Anywhere}\label{noir-write-once-prove-anywhere}

Noir takes a different approach to the same problem. Where Compact is
chain-specific and privacy-first, Noir is backend-agnostic and
correctness-first. A Noir program compiles not to a specific proof system's
constraint format, but to ACIR -- Abstract Circuit Intermediate Representation
-- which can then be lowered to any compatible backend.

Consider the same token transfer scenario in Noir:

\begin{Shaded}
\begin{Highlighting}[]
\KeywordTok{fn}\NormalTok{ main(}
\NormalTok{    sender\_balance}\OperatorTok{:}\NormalTok{ Field}\OperatorTok{,}
\NormalTok{    amount}\OperatorTok{:} \KeywordTok{pub}\NormalTok{ Field}\OperatorTok{,}
\NormalTok{    recipient\_balance}\OperatorTok{:}\NormalTok{ Field}\OperatorTok{,}
\NormalTok{) }\OperatorTok{{-}\textgreater{}} \KeywordTok{pub}\NormalTok{ Field }\OperatorTok{\{}
\NormalTok{    assert(sender\_balance }\OperatorTok{\textgreater{}=}\NormalTok{ amount)}\OperatorTok{;}
    \KeywordTok{let}\NormalTok{ new\_recipient\_balance }\OperatorTok{=}\NormalTok{ recipient\_balance }\OperatorTok{+}\NormalTok{ amount}\OperatorTok{;}
\NormalTok{    new\_recipient\_balance}
\OperatorTok{\}}
\end{Highlighting}
\end{Shaded}

The privacy model is visible in the function signature. Parameters without the
\texttt{pub} keyword are private -- they are part of the witness, known only to
the prover. Parameters marked \texttt{pub} are public inputs, visible to the
verifier. The return value, also marked \texttt{pub}, is the public output.

This is simpler than Compact's disclosure analysis. There is no
\texttt{disclose()} mechanism, no 26-pass pipeline tracing data-flow paths. The
developer declares privacy at the function boundary: this input is private, that
input is public, and the compiler enforces the declaration. It is the
developer's responsibility to get the declaration right. Noir will not warn you
if you accidentally mark a sensitive value as \texttt{pub}. But it will
guarantee that every private input remains invisible to the verifier -- that the
proof reveals nothing about private inputs beyond what the public outputs
logically imply.

Where Noir stands apart is composability. A Noir developer can write a library
of circuit components -- hash functions, signature verifiers, Merkle tree
checkers -- and reuse them across projects targeting different proof systems.
The same Noir code that runs on Aztec's Barretenberg backend today could, in
principle, run on a PLONK backend, a Groth16 backend, or a future proof system
that does not yet exist. This is not theoretical: the Noir ecosystem already
includes standard libraries for common cryptographic primitives, and projects
like zkEmail, zkPassport, and zkLogin use these libraries to build real
applications.

A more realistic Noir program demonstrates the language's expressiveness. Here
is a simplified credential verification -- proving you are over 18 without
revealing your birthdate:

\begin{Shaded}
\begin{Highlighting}[]
\KeywordTok{use} \PreprocessorTok{std::hash::}\NormalTok{poseidon}\OperatorTok{;}

\KeywordTok{fn}\NormalTok{ main(}
\NormalTok{    birth\_year}\OperatorTok{:}\NormalTok{ Field}\OperatorTok{,}
\NormalTok{    birth\_month}\OperatorTok{:}\NormalTok{ Field}\OperatorTok{,}
\NormalTok{    birth\_day}\OperatorTok{:}\NormalTok{ Field}\OperatorTok{,}
\NormalTok{    credential\_hash}\OperatorTok{:} \KeywordTok{pub}\NormalTok{ Field}\OperatorTok{,}
\NormalTok{    current\_year}\OperatorTok{:} \KeywordTok{pub}\NormalTok{ Field}\OperatorTok{,}
\NormalTok{    threshold\_age}\OperatorTok{:} \KeywordTok{pub}\NormalTok{ Field}\OperatorTok{,}
\NormalTok{) }\OperatorTok{\{}
    \CommentTok{// Verify the credential hash matches the private birthdate}
    \KeywordTok{let}\NormalTok{ computed\_hash }\OperatorTok{=} \PreprocessorTok{poseidon::bn254::}\NormalTok{hash\_3([birth\_year}\OperatorTok{,}\NormalTok{ birth\_month}\OperatorTok{,}\NormalTok{ birth\_day])}\OperatorTok{;}
\NormalTok{    assert(computed\_hash }\OperatorTok{==}\NormalTok{ credential\_hash)}\OperatorTok{;}

    \CommentTok{// Verify age threshold without revealing exact birthdate}
    \KeywordTok{let}\NormalTok{ age }\OperatorTok{=}\NormalTok{ current\_year }\OperatorTok{{-}}\NormalTok{ birth\_year}\OperatorTok{;}
\NormalTok{    assert(age }\OperatorTok{\textgreater{}=}\NormalTok{ threshold\_age)}\OperatorTok{;}
\OperatorTok{\}}
\end{Highlighting}
\end{Shaded}

The developer writes Rust-like code. The standard library provides cryptographic
primitives. The compiler handles the translation to constraints. The program is
readable, auditable, and portable across proof backends. What the program cannot
do -- what no Noir program can do -- is enforce at compile time that the
developer has not accidentally marked \texttt{birth\_year} as \texttt{pub}. That
responsibility rests with the developer, not the compiler.

\subsection{Leo: Privacy as a Record
System}\label{leo-privacy-as-a-record-system}

Leo takes a third approach, one rooted in Aleo's record-based execution model.
Where Compact models privacy as a property of data flow and Noir models it as a
property of function signatures, Leo models privacy as a property of
\emph{records} -- discrete objects that are created, consumed, and transferred,
much like physical banknotes.

A Leo token transfer looks different from both Compact and Noir:

\begin{Shaded}
\begin{Highlighting}[]
\NormalTok{program token}\OperatorTok{.}\NormalTok{aleo }\OperatorTok{\{}
\NormalTok{    record Token }\OperatorTok{\{}
\NormalTok{        owner}\OperatorTok{:}\NormalTok{ address}\OperatorTok{,}
\NormalTok{        amount}\OperatorTok{:} \DataTypeTok{u64}\OperatorTok{,}
    \OperatorTok{\}}

\NormalTok{    transition transfer(}
\NormalTok{        input}\OperatorTok{:}\NormalTok{ Token}\OperatorTok{,}
\NormalTok{        recipient}\OperatorTok{:}\NormalTok{ address}\OperatorTok{,}
\NormalTok{        amount}\OperatorTok{:} \DataTypeTok{u64}\OperatorTok{,}
\NormalTok{    ) }\OperatorTok{{-}\textgreater{}}\NormalTok{ (Token}\OperatorTok{,}\NormalTok{ Token) }\OperatorTok{\{}
        \KeywordTok{let}\NormalTok{ remaining}\OperatorTok{:} \DataTypeTok{u64} \OperatorTok{=}\NormalTok{ input}\OperatorTok{.}\NormalTok{amount }\OperatorTok{{-}}\NormalTok{ amount}\OperatorTok{;}

        \KeywordTok{let}\NormalTok{ sender\_token}\OperatorTok{:}\NormalTok{ Token }\OperatorTok{=}\NormalTok{ Token }\OperatorTok{\{}
\NormalTok{            owner}\OperatorTok{:} \KeywordTok{self}\OperatorTok{.}\NormalTok{caller}\OperatorTok{,}
\NormalTok{            amount}\OperatorTok{:}\NormalTok{ remaining}\OperatorTok{,}
        \OperatorTok{\};}

        \KeywordTok{let}\NormalTok{ recipient\_token}\OperatorTok{:}\NormalTok{ Token }\OperatorTok{=}\NormalTok{ Token }\OperatorTok{\{}
\NormalTok{            owner}\OperatorTok{:}\NormalTok{ recipient}\OperatorTok{,}
\NormalTok{            amount}\OperatorTok{:}\NormalTok{ amount}\OperatorTok{,}
        \OperatorTok{\};}

        \ControlFlowTok{return}\NormalTok{ (sender\_token}\OperatorTok{,}\NormalTok{ recipient\_token)}\OperatorTok{;}
    \OperatorTok{\}}
\OperatorTok{\}}
\end{Highlighting}
\end{Shaded}

The keyword \texttt{record} defines a private data structure. Records in Leo are
encrypted on-chain -- only the owner can decrypt them. When a
\texttt{transition} consumes a record and produces new records, the old record
is nullified (marked as spent) and the new records are encrypted for their
respective owners. The entire UTXO lifecycle -- creation, transfer, consumption
-- is expressed in the language itself.

The keyword \texttt{transition} is Leo's equivalent of Compact's
\texttt{circuit}. It defines a function that will be proven in zero knowledge.
But where Compact's circuits operate on abstract state and Noir's functions
operate on field elements, Leo's transitions operate on records -- typed, owned,
encrypted objects with a lifecycle managed by the Aleo runtime.

Leo's privacy model is structural rather than analytical. Privacy does not
emerge from a disclosure analysis pass or from \texttt{pub} annotations on
function parameters. It emerges from the record model itself: records are
encrypted, transitions consume and produce records, and the only public artifact
is a nullifier (proving a record was spent) and a commitment (proving a new
record was created). The developer does not choose what is private. Everything
inside a record is private by default. Publicity is the exception, declared
through explicit \texttt{public} annotations on transition inputs.

The tradeoff is the same one that appears throughout Philosophy D: Leo programs
run only on Aleo. The record model, the transition semantics, the encryption
scheme -- all are Aleo-specific. But for developers building on Aleo, Leo
provides something that general-purpose languages cannot: a programming model
where privacy is not a layer of annotation on top of ordinary computation, but
the fundamental unit of state.

\subsection{The Philosophy D Synthesis}\label{the-philosophy-d-synthesis}

What unites these three languages is a conviction that privacy cannot be an
afterthought. In Philosophies A, B, and C, the proof system guarantees
computational integrity -- it proves the computation was done correctly. But it
says nothing about what information the computation reveals. Privacy depends on
the developer correctly managing which values are public and which are private.
A mistake does not produce a compiler error. It produces a privacy leak.

Philosophy D languages treat privacy as a compiler concern. The language itself
knows the difference between public and private. And in Compact's case, the
compiler physically prevents the developer from accidentally crossing that
boundary.

What unites Philosophies A, B, and C is a shared assumption: the developer does
not need to think about privacy. The proof system guarantees computational
integrity -- it proves the computation was done correctly. But it says nothing
about what information the computation reveals. Whether to encrypt inputs, hide
outputs, or shield metadata is left to the application layer, if it is
considered at all.

Philosophy D breaks this assumption. Privacy is not a layer above the language
-- it is embedded in the language itself.

The following table summarizes the four philosophies:

\begin{longtable}[]{@{}
  >{\raggedright\arraybackslash}p{(\linewidth - 10\tabcolsep) * \real{0.1667}}
  >{\raggedright\arraybackslash}p{(\linewidth - 10\tabcolsep) * \real{0.1667}}
  >{\raggedright\arraybackslash}p{(\linewidth - 10\tabcolsep) * \real{0.1667}}
  >{\raggedright\arraybackslash}p{(\linewidth - 10\tabcolsep) * \real{0.1667}}
  >{\raggedright\arraybackslash}p{(\linewidth - 10\tabcolsep) * \real{0.1667}}
  >{\raggedright\arraybackslash}p{(\linewidth - 10\tabcolsep) * \real{0.1667}}@{}}
\toprule\noalign{}
\begin{minipage}[b]{\linewidth}\raggedright
Philosophy
\end{minipage} & \begin{minipage}[b]{\linewidth}\raggedright
Representative
\end{minipage} & \begin{minipage}[b]{\linewidth}\raggedright
What It Proves
\end{minipage} & \begin{minipage}[b]{\linewidth}\raggedright
Privacy Model
\end{minipage} & \begin{minipage}[b]{\linewidth}\raggedright
Developer Experience
\end{minipage} & \begin{minipage}[b]{\linewidth}\raggedright
Tradeoff
\end{minipage} \\
\midrule\noalign{}
\arrayrulecolor{MidnightBlue}\midrule\arrayrulecolor{MidnightBorder}
\endhead
\bottomrule\noalign{}
\endlastfoot
\textbf{A: EVM-Compatible} & Scroll, Linea & EVM execution & None (transparent)
& Familiar (Solidity) & Proving the EVM is expensive \\
\textbf{B: ZK-Native ISA} & Cairo (Starknet) & Custom CPU trace & None
(transparent) & New language required & Locked to one ecosystem \\
\textbf{C: General-Purpose ISA} & SP1, RISC Zero, Airbender & RISC-V execution &
None (transparent) & Familiar (Rust, C++) & Arithmetization overhead \\
\textbf{D: Application DSL} & Compact, Noir, Leo & State transitions or circuits
& Compiler-enforced & Domain-specific syntax & Locked to one chain (Compact/Leo)
or one IR (Noir) \\
\end{longtable}

The taxonomy is not a ranking. Each philosophy serves a different constituency.
The following decision guide distills each philosophy's sweet spot:

\begin{longtable}[]{@{}
  >{\raggedright\arraybackslash}p{(\linewidth - 4\tabcolsep) * \real{0.2750}}
  >{\raggedright\arraybackslash}p{(\linewidth - 4\tabcolsep) * \real{0.4250}}
  >{\raggedright\arraybackslash}p{(\linewidth - 4\tabcolsep) * \real{0.3000}}@{}}
\toprule\noalign{}
\begin{minipage}[b]{\linewidth}\raggedright
Philosophy
\end{minipage} & \begin{minipage}[b]{\linewidth}\raggedright
Choose This When
\end{minipage} & \begin{minipage}[b]{\linewidth}\raggedright
Avoid When
\end{minipage} \\
\midrule\noalign{}
\arrayrulecolor{MidnightBlue}\midrule\arrayrulecolor{MidnightBorder}
\endhead
\bottomrule\noalign{}
\endlastfoot
\textbf{A: EVM-Compatible} & Existing Solidity codebase; need L1 security
inheritance; team knows EVM tooling & Building from scratch;
performance-critical new system; proving cost is binding constraint \\
\textbf{B: ZK-Native ISA} & Maximum proving efficiency; willing to learn Cairo;
building within Starknet ecosystem & Need ecosystem portability; team cannot
invest in new language; multi-chain deployment required \\
\textbf{C: General-Purpose ISA} & Standard Rust/C++ codebase; want ecosystem
portability; broadest developer pool; zkVM as proving backend & Need
compile-time privacy enforcement; need lowest possible proof latency for simple
circuits \\
\textbf{D: Application DSL} & Privacy-preserving smart contracts; want
compiler-enforced disclosure rules; domain-specific state model & Need
general-purpose computation; team wants familiar language; multi-backend
portability (except Noir) \\
\end{longtable}

To make these four philosophies concrete: \textbf{Philosophy A} is like
translating a novel into a language that has no word for half the concepts --
faithful but expensive. \textbf{Philosophy B} is like building a custom theater
for a specific play -- the acoustics are perfect, but the theater can only stage
that one production. \textbf{Philosophy C} is like staging the play in any
theater in the world by writing it for a universal stage: RISC-V is the
universal stage, and the play (your Rust program) works anywhere.
\textbf{Philosophy D} is like writing a play where the script physically
prevents the actors from breaking character -- in Compact, the compiler will not
let a private value reach a public surface without explicit consent. Philosophy
A serves Ethereum developers who want to reuse existing code. Philosophy B
serves ecosystems willing to invest in a custom stack for maximum proof
efficiency. Philosophy C serves the broadest developer community with the least
friction. Philosophy D serves applications where privacy is not optional --
where the language must prevent the developer from accidentally revealing what
should stay hidden.

\section{The Polygon zkEVM Cautionary
Tale}\label{the-polygon-zkevm-cautionary-tale}

Before we proceed, a cautionary tale.

The original version of this book listed Polygon zkEVM as the flagship example
of Philosophy A -- the EVM-compatible approach. Five independent reviewers
flagged the same problem: Polygon zkEVM was shut down.

Polygon acquired Hermez, the team behind one of the most ambitious zkEVM
implementations, for approximately \$250 million. The project aimed to prove
every Ethereum opcode, faithfully reproducing the EVM's behavior inside a
zero-knowledge circuit. It was technically impressive. It was also, in the end,
commercially unviable.

In 2025, Polygon announced the sunsetting of zkEVM Mainnet Beta. The Hermez
team, led by co-founder Jordi Baylina, spun off to form ZisK -- and pivoted to
RISC-V. The very team that had spent years building the most faithful
EVM-compatible proof system concluded that proving the EVM directly was not the
path forward. Better to prove a clean, efficient instruction set (RISC-V) and
layer EVM compatibility on top.

The lesson is not that EVM compatibility is wrong. Scroll and Linea are
thriving. The lesson is that \emph{how} you achieve compatibility matters
enormously. Polygon bet on faithfully reproducing every EVM quirk at the circuit
level. The cost -- in engineering complexity, in proving time, in the sheer
number of constraints required to model 256-bit arithmetic -- proved
prohibitive. Scroll and Linea survived because they found more efficient paths
to the same goal.

For a system architect evaluating ZK infrastructure, the takeaway is concrete: a
\$250 million investment with a world-class team is not sufficient if the
architectural approach creates an exponential constraint-generation problem.
Philosophy A is viable, but only if you resist the temptation to prove the EVM
literally. Abstract where you can. Approximate where you must. And keep one eye
on the RISC-V projects that may make your compatibility layer unnecessary.

The cautionary tale has a coda. Jordi Baylina, who led Polygon Hermez and
co-created Circom, took the team to ZisK and chose RISC-V. The person who knew
more about EVM-compatible ZK proving than almost anyone else alive concluded
that the direct approach was not the path forward. When the world's leading
practitioner abandons a technique, the rest of the field should pay attention.

\section{The Developer's Actual
Experience}\label{the-developers-actual-experience}

The taxonomy of philosophies describes how \emph{architects} think about Layer
2. But what does a \emph{developer} actually do?

The original paper described choosing a language but never described using one.
A smart reader would ask: ``Fine, I picked a language. Now what?'' The answer
involves a lifecycle that most ZK documentation glosses over.

\textbf{Step 1: Write.} The developer writes source code. In SP1, this is
standard Rust. In Circom, this is a template-based constraint description
language. In Compact, this is TypeScript-like code with \texttt{witness}
declarations and \texttt{circuit} exports. The writing experience varies
enormously. An SP1 developer uses familiar tools -- VS Code, cargo, clippy. A
Circom developer works in a specialized IDE with no debugger, no step-through
execution, and error messages that refer to constraint indices rather than
variable names.

\textbf{Step 2: Compile.} The compiler translates source code into a form the
proof system can work with. For SP1, this means Rust to RISC-V machine code via
the standard LLVM backend. For Circom, this means templates to R1CS constraint
systems. For Compact, this means a 26-pass nanopass compilation pipeline that
transforms the source through 26 intermediate languages -- from \texttt{Lsrc}
through type checking (\texttt{Ltypes}), disclosure analysis
(\texttt{Lnodisclose}), loop unrolling (\texttt{Lunrolled}), circuit flattening
(\texttt{Lflattened}), and finally ZKIR output.

A critical finding from recent research: standard LLVM optimization passes (-O3)
yield over 40\% improvement when targeting zkVMs, compared to much larger gains
on traditional CPUs. This is because LLVM's optimization heuristics are tuned
for hardware features -- cache locality, branch prediction, instruction-level
parallelism -- that do not exist in a zkVM. By refining a small set of LLVM
passes to use a ZK-aware cost model, researchers achieved up to 45\% on
individual benchmarks (with average gains of 1-4\%). The compiler is an
underexplored optimization surface.

\textbf{Step 3: Test.} The developer tests their program. In the RISC-V world,
this means running the program natively (without proof generation) and checking
outputs. SP1 supports this directly: you can execute your Rust program on a
standard CPU and verify correctness before paying the cost of proof generation.
In Compact, the SDK provides a local execution oracle that simulates the
blockchain environment -- block time, token balances, contract state -- without
a running chain.

\textbf{Step 4: Prove.} The developer generates a proof. This is where the cost
hits. Proof generation for a Compact circuit takes seconds to tens of seconds on
local development hardware (detailed timings appear in Chapter 6). For an SP1
program, proving time depends on the number of RISC-V cycles executed -- a
simple computation might take seconds; an Ethereum block might take minutes even
on GPU clusters.

What does ``prove'' actually feel like? The experience is unlike anything else
in software development. There is no analogy in web development, in systems
programming, in machine learning training. It is its own thing, and it deserves
honest description.

PLACEHOLDER\_PROVE\_REST

\textbf{Step 5: Deploy.} The proven program is deployed. For rollup-based
systems, this means posting the proof and public inputs to Ethereum. For
Compact, this means deploying the ZKIR circuit, TypeScript bindings, and proving
keys to Midnight's network. Contract deployment on Midnight's devnet is
dominated by proof generation for the constructor circuit (see Chapter 6 for
measured latencies).

\textbf{Step 6: Monitor.} In production, the developer monitors for correctness,
performance, and security. This step is almost entirely undocumented in the ZK
ecosystem. There are no standard monitoring tools for ZK deployments. No
dashboards for constraint utilization. No alerting for proof generation
failures. The gap between ``deploy'' and ``done'' is where real-world systems
fail.

One emerging bright spot: LLM-assisted ZK development. The ZK-Coder system
improved Circom circuit generation success rates from 20\% (baseline large
language model) to 88\%. This suggests that the developer experience barrier --
the steep learning curve, the unfamiliar constraint semantics, the cryptic error
messages -- may be partially addressable through AI tooling. But 88\% is not
100\%, and the 12\% failure cases may be precisely the subtle
under-constrainedness bugs that are hardest to detect. An LLM that generates a
circuit with a missing constraint is more dangerous than an LLM that fails to
generate a circuit at all.

The developer workflow reveals something about Layer 2: the \emph{language} is
the part the developer sees, but the \emph{compiler} is the part that matters. A
language with beautiful syntax and a buggy compiler is worse than an ugly
language with a correct compiler. The field is beginning to understand this.
CirC, a unifying compiler infrastructure from Stanford, demonstrated that the
compilation problem for ZK circuits shares fundamental structure with SMT
solving and software verification. The same optimizations -- constant folding,
dead code elimination, common subexpression elimination -- apply across all
targets. This suggests that investment in ZK compiler infrastructure could pay
off disproportionately, improving every language simultaneously rather than
optimizing each one independently.

\section{Under-Constrained Circuits: The Dominant Failure
Mode}\label{under-constrained-circuits-the-dominant-failure-mode}

Here is the fact that should keep every ZK developer awake at night: the most
common way a zero-knowledge system fails in practice is not a cryptographic
break. It is not a quantum computer. It is not a governance attack. It is a bug
in the program.

The under-constrained vulnerability epidemic documented in Section 3.1 -- 95 of
141 real-world bugs -- is not an abstraction. To understand why it happens at
this scale, you need to understand the dual-track problem.

In Circom -- the most widely deployed ZK language by project count -- every line
of code simultaneously describes two things: how to \emph{compute} a value
(witness generation), and how to \emph{constrain} that value (the mathematical
rule the proof system enforces). These two descriptions use different operators.
The arrow \texttt{\textless{}-\/-} computes a value. The triple-equals
\texttt{===} constrains it. The combined operator \texttt{\textless{}==} does
both.

The Tornado Cash bug was this: a developer used \texttt{=} (JavaScript
assignment) where \texttt{\textless{}==} (constrained assignment) was needed.
The witness generator computed the correct value. The constraint system did not
enforce it. A malicious prover could substitute any value, and the proof would
still verify. One character. Complete soundness break.

This is not a rare edge case. The ZKAP static analysis framework, which
introduced a Circuit Dependence Graph abstraction to detect vulnerabilities in
Circom circuits, found 34 previously unknown vulnerabilities across 258 circuits
in 15 open-source projects. Its analysis identified three root causes:

First, \emph{nondeterministic signals} -- circuit outputs that can take multiple
values for a given input because constraints are missing. Twenty-four percent of
ZKAP's findings fell in this category.

Second, \emph{unsafe component usage} -- sub-circuits invoked without properly
constraining their inputs or outputs. This created gaps where values passed
between components were computed but not verified.

Third, \emph{constraint-computation discrepancies} -- places where the witness
generator and the constraint system diverged. Division is the classic example:
in witness generation, division computes a quotient. In constraints, division is
expressed as multiplication (if \(a/b = c\), the constraint is
\(b \cdot c = a\)). When the divisor can be zero, these two formulations behave
differently. Division-by-zero was the single most common vulnerability class in
ZKAP's findings, accounting for 41\% of all bugs.

The defensive tooling is evolving. Picus (also called QED\^{}2) uses SMT-based
techniques to automatically detect under-constrained circuits in R1CS, reducing
the under-constrainedness problem to queries on systems of polynomial equations
over finite fields. ZKAP introduced the Circuit Dependence Graph abstraction --
combining data flow edges (witness computation) with constraint edges (R1CS) --
and achieved an F1 score of 0.82, compared to 0.64 for the earlier Circomspect
tool. zkFuzz, a fuzz-testing framework for ZK circuits, found 66 bugs including
38 zero-days. MTZK discovered 21 bugs across 4 different ZK compilers.

But these tools share a limitation: they work primarily on Circom circuits. The
Rust-based systems that dominate production -- halo2 (used by Scroll and ZK
Bridge projects), Plonky3 (used by SP1 and Stwo), and custom constraint systems
in RISC Zero and Jolt -- are not covered. The most common ZK bug class has
automated detection for the oldest ZK language but not for the systems where new
code is being written.

NAVe, a formal verification tool for Noir programs announced in 2025, begins to
close this gap. It formalizes Noir's ACIR intermediate representation and uses
the cvc5 SMT solver to verify program properties. But formal verification at
scale -- for circuits with millions of constraints -- remains beyond current
tools. The combination of compile-time prevention (refinement types, disclosure
analysis) and post-hoc verification (static analysis, formal methods) could
provide comprehensive coverage, but no system achieves both today.

The evidence is clear: the most common failure mode in zero-knowledge systems is
not a failure of cryptography. It is a failure of software engineering. The
magician's choreography has a typo, and the proof system performs the typo
faithfully.

\section{Compact's Disclosure Analysis}\label{compacts-disclosure-analysis}

Compact, Midnight's smart contract language, takes a fundamentally different
approach to the bug problem. Rather than giving developers direct access to
constraints and hoping they get it right, Compact's compiler enforces a
\emph{disclosure rule}: witness values are private by default, and any attempt
to use a private value in a public context without explicit consent is a
compile-time error.

This is not a style guide. It is not a best practice recommendation. It is a
hard compiler rejection.

Here is how it works. In Compact, private data enters the circuit through
\emph{witness} functions -- declared in Compact, implemented in TypeScript:

\begin{Shaded}
\begin{Highlighting}[]
\NormalTok{witness }\FunctionTok{get\_secret}\NormalTok{()}\OperatorTok{:}\NormalTok{ Bytes}\OperatorTok{\textless{}}\DecValTok{32}\OperatorTok{\textgreater{};}
\end{Highlighting}
\end{Shaded}

The implementation runs off-chain, in the user's browser or node:

\begin{Shaded}
\begin{Highlighting}[]
\KeywordTok{const}\NormalTok{ witnesses }\OperatorTok{=}\NormalTok{ \{}
  \DataTypeTok{get\_secret}\OperatorTok{:}\NormalTok{ (ctx) }\KeywordTok{=\textgreater{}}\NormalTok{ [ctx}\OperatorTok{.}\AttributeTok{privateState}\OperatorTok{,} \FunctionTok{hexToBytes}\NormalTok{(}\StringTok{"aabb..."}\NormalTok{)]}
\NormalTok{\}}\OperatorTok{;}
\end{Highlighting}
\end{Shaded}

The witness value is private. It exists only on the user's device. To use it
inside a circuit -- where it will be constrained and proven -- the developer
must explicitly call \texttt{disclose()}:

\begin{Shaded}
\begin{Highlighting}[]
\ImportTok{export}\NormalTok{ circuit }\FunctionTok{verify}\NormalTok{()}\OperatorTok{:}\NormalTok{ [] \{}
  \KeywordTok{const}\NormalTok{ sk }\OperatorTok{=} \FunctionTok{disclose}\NormalTok{(}\FunctionTok{get\_secret}\NormalTok{())}\OperatorTok{;}
  \KeywordTok{const}\NormalTok{ my\_hash }\OperatorTok{=} \FunctionTok{persistentHash}\NormalTok{([}\FunctionTok{pad}\NormalTok{(}\DecValTok{32}\OperatorTok{,} \StringTok{"auth:"}\NormalTok{)}\OperatorTok{,}\NormalTok{ sk])}\OperatorTok{;}
  \FunctionTok{assert}\NormalTok{(my\_hash }\OperatorTok{==}\NormalTok{ stored\_hash}\OperatorTok{,} \StringTok{"not authorized"}\NormalTok{)}\OperatorTok{;}
\NormalTok{\}}
\end{Highlighting}
\end{Shaded}

Without \texttt{disclose()}, the compiler rejects the program. The error message
traces the complete path from the witness value to the public surface:

\begin{verbatim}
potential witness-value disclosure must be declared but is not:
  witness value potentially disclosed:
    the return value of witness get_amount at line 8 char 1
  nature of the disclosure:
    ledger operation might disclose the witness value
  via this path through the program:
    the argument to increment at line 11 char 8
\end{verbatim}

The compiler catches three categories of accidental leakage: witness values used
in ledger operations (writing to on-chain state), witness values returned from
circuits (visible in the proof's public outputs), and witness values passed to
kernel operations (token transfers, balance queries).

The disclosure analysis pass runs as part of the 26-stage nanopass compilation
pipeline, at the \texttt{Lnodisclose} intermediate language stage. The privacy
boundary is verified before any code generation occurs. The compiler has already
type-checked the program, analyzed data flow paths, and confirmed that every
potential disclosure is explicitly marked -- or the compilation fails.

Consider a concrete case. The Midnight developer guide documents that the first
attempt at implementing a private voting contract -- using naive if/else
branching on witness values -- was rejected by the compiler with 11 disclosure
errors. Each error traced the path from a witness value to a public surface. The
compiler forced a fundamental redesign: Merkle trees replaced per-slot
branching, nullifiers replaced voted-flags, and arithmetic tallying replaced
conditional increments. The resulting design was not just compiler-compliant --
it was architecturally superior. The naive approach would have leaked which
candidate each voter chose through the pattern of ledger writes. The
compiler-forced redesign made this impossible.

No other ZK language provides this guarantee. In Circom, Noir, and Cairo,
privacy depends on the developer correctly managing which values are public and
which are private. A mistake does not produce a compiler error. It produces a
privacy leak that may not be discovered until an attacker exploits it. Compact
makes privacy a compiler guarantee rather than a developer responsibility.

The tradeoff is clear: Compact contracts are locked to Midnight's proof system
(PLONK on BLS12-381), token model (Zswap), and ledger architecture. They cannot
be deployed on another chain. In exchange, the developer gets something no
general-purpose approach can offer: the compiler will not let you accidentally
show the audience what is behind the curtain.

\section{Midnight: Compiler, IR, Circuit}\label{midnight-compiler-ir-circuit}

Compact's three-part compilation architecture illustrates a principle that
applies across all of Layer 2: the compilation target shapes the developer's
world.

From a single \texttt{.compact} source file, the compiler produces three
distinct artifacts:

\textbf{Artifact 1: ZKIR circuits.} ZKIR (Zero-Knowledge Intermediate
Representation) is a JSON-formatted circuit description with 24 typed
instructions organized into eight categories: arithmetic (\texttt{add},
\texttt{mul}, \texttt{neg}), constraints (\texttt{assert},
\texttt{constrain\_bits}, \texttt{constrain\_eq}), control flow
(\texttt{cond\_select}, \texttt{copy}), type encoding (\texttt{decode},
\texttt{encode}), cryptographic operations (\texttt{ec\_mul},
\texttt{hash\_to\_curve}, \texttt{persistent\_hash}), and I/O (\texttt{impact},
\texttt{output}, \texttt{private\_input}, \texttt{public\_input}).

Every ZKIR circuit has two transcript channels. The \texttt{publicTranscript}
records ledger operations -- reads, writes, comparisons -- visible to the
on-chain verifier. The \texttt{privateTranscriptOutputs} contains
witness-derived values visible only to the prover. The ZKIR checker verifies
that the serialized public transcript matches exactly what the circuit computed.
Tampering with either transcript causes rejection with specific, diagnostic
error messages.

\textbf{Artifact 2: TypeScript bindings.} The compiler generates type-safe
JavaScript/TypeScript API code that handles contract interaction from the dApp
frontend. This includes witness provider interfaces (the functions that supply
private inputs), serialization between TypeScript types and field elements, and
the transaction construction and submission pipeline. The witness functions run
off-chain with access to private state, external APIs, and databases --
computation that could never run inside a ZK circuit.

\textbf{Artifact 3: Proving keys.} Contract-specific cryptographic material for
the PLONK-based proof system on BLS12-381. Different circuits produce different
keys. The proof server requires these keys to generate proofs; validators
require them to verify proofs.

This three-part split is not an implementation detail. It reflects the
fundamental architecture of privacy-preserving computation: what can be proven
(ZKIR), what runs privately (TypeScript), and what makes proofs possible (keys).
No other ZK language produces all three from a single source file with
first-class blockchain integration. Circom produces R1CS plus a witness
generator but no blockchain API. Noir produces ACIR but no TypeScript bindings
and no blockchain integration. Cairo produces execution traces but no privacy
separation. Compact unifies the entire dApp development pipeline.

For a system architect, the lesson generalizes beyond Midnight: the choice of
Layer 2 language is not just a choice of syntax. It is a choice of compilation
target, developer tooling, privacy model, and deployment pipeline. The language
shapes everything downstream.

\par\vspace{18pt}

The choreography is written. The developer has expressed their computation in a
language -- whether Rust targeting RISC-V, Cairo targeting Starknet, or Compact
targeting Midnight's ZKIR. The compiler has translated the program into a form
the proof system can work with.

But the program is just the \emph{plan}. It describes what the computation
should do. It does not contain the private data. It does not contain the
execution trace. It does not contain the secret.

Now the magician goes backstage. The curtain closes. The audience waits. Behind
the curtain, the magician will run the computation with real data -- real bank
balances, real identity credentials, real votes -- and record every step. This
recording is the witness: the complete execution trace that later layers will
prove properties about without ever revealing.

The recording is where the real cost lives. And it is where the real
vulnerabilities hide.

\chapter{Chapter 4: The Secret
Performance}\label{chapter-4-the-secret-performance}

\emph{Layer 3 -- Witness Generation}

\section{The Hidden Bottleneck}\label{the-hidden-bottleneck}

Here is a number that should have been a scandal: 50-70\%.

That is the fraction of total proving time consumed by witness generation in
modern GPU-accelerated systems. Not 10-25\%, as the field commonly claimed until
recently. The outright majority. If witness generation accounts for more than
half the cost of producing a zero-knowledge proof, why has the field treated it
as a minor backstage interlude? Why have billions of dollars in optimization
effort focused on the cryptographic proving step while the dominant bottleneck
hid in plain sight?

The answer reveals something important about how technology communities deceive
themselves. When GPU acceleration made the cryptographic proving step 10 times
faster, witness generation did not get faster. It stayed the same speed. But its
\emph{share} of total time climbed from a modest 20\% to a dominant 67\%. The
better the proof system got, the worse the witness gap became. The field
celebrated its proving breakthroughs while the actual bottleneck quietly grew.

This chapter is about what happens backstage. The curtain has closed. The
audience (the verifier) cannot see what the magician does next. She takes her
private data -- your bank balance, your identity, your vote -- and runs the
computation, recording every step. This recording is the witness: the complete
execution trace. Later layers will prove properties about it without revealing
its contents.

But three problems lurk behind that curtain. The recording is expensive to make.
The recording room has thin walls. And if the recording is wrong, the entire
system breaks.

\section{Execution Traces}\label{execution-traces}

Start with the simplest possible example. You want to prove you know a number
whose square is 25. The witness is that number: 5. The proof convinces the
verifier that you know such a number, without revealing that the number is 5.
The verifier learns ``this person knows a square root of 25.'' The verifier does
not learn ``the number is 5.'' (In fact, the verifier does not even learn
whether you chose 5 or -5 -- both are valid witnesses for the same statement.)

Now scale up. You want to prove that an Ethereum state transition is valid --
that a block of transactions, when applied to the current state, produces the
claimed new state. The witness is not a single number. It is the \emph{entire
execution trace}: every memory access, every register value, every instruction
execution, every intermediate hash computation, every storage read and write.
For a complex Ethereum block, this means millions of steps, each with dozens of
columns of data. The execution trace for a block proved by SP1 or RISC Zero
might contain billions of field elements.

The witness is, in the most literal sense, a complete recording of everything
that happened during the computation. Think of it as a security camera that
films the magician backstage: it captures not just the final result but every
intermediate movement, every prop placement, every sleight of hand. The
recording exists so that the proof system can later verify that every step was
consistent with the rules -- without the verifier ever watching the footage.

Here is what a witness looks like for a trivial computation: proving that
\(x^2 + x = 12\) where \(x = 3\).

\begin{longtable}[]{@{}lll@{}}
\toprule\noalign{}
Step & Operation & Result \\
\midrule\noalign{}
\arrayrulecolor{MidnightBlue}\midrule\arrayrulecolor{MidnightBorder}
\endhead
\bottomrule\noalign{}
\endlastfoot
1 & Load x & 3 \\
2 & \(t_1 = x \times x\) & 9 \\
3 & \(t_2 = t_1 + x\) & 12 \\
4 & Assert \(t_2 = 12\) & ✓ \\
\end{longtable}

Four field elements. Four rows in the trace table. Each intermediate value is a
cell the prover must fill in and the constraint system must verify. For a real
zkVM executing a RISC-V program, the trace has columns for every register (32
registers), the program counter, the current instruction, memory addresses
accessed, and intermediate ALU results -- thousands of columns, millions of
rows, but the same fundamental structure: a table where each row is one clock
cycle of execution.

This distinction between the \emph{computation} (what the magician does) and the
\emph{recording} (the witness) matters more than it looks. The computation might
take milliseconds. The recording is vastly larger because it includes every
intermediate value. And generating the recording -- running the computation
while capturing every detail -- is the expensive part.

\begin{quote}
\textbf{The Running Example: The Sudoku Witness}

For our 4x4 Sudoku, the witness is the completed grid -- sixteen field elements:

\begin{verbatim}
+---+---+---+---+
| 1 | 2 | 3 | 4 |
+---+---+---+---+
| 3 | 4 | 1 | 2 |
+---+---+---+---+
| 2 | 1 | 4 | 3 |
+---+---+---+---+
| 4 | 3 | 2 | 1 |
+---+---+---+---+
\end{verbatim}

The execution trace records every check: ``cell (0,0) = 1, matches given? yes.
Row 0 sum = 10, all distinct? yes. Column 0 = \{1,3,2,4\}, all distinct? yes.
Box (0,0) = \{1,2,3,4\}, all distinct? yes.'' Sixteen values, plus every
intermediate comparison and boolean result -- roughly 80 field elements in
total. The verifier never sees this grid. The verifier sees only the original
puzzle (the public input) and, eventually, the proof.
\end{quote}

In a zkVM like SP1 or RISC Zero, witness generation means \emph{emulating the
entire RISC-V processor}. Every instruction is fetched, decoded, and executed.
Every register update is recorded. Every memory access is logged. This is full
virtual machine emulation, step by step, sequentially. It cannot be easily
parallelized because each instruction depends on the state left by the previous
instruction. The program counter moves forward one step at a time, and the
witness generator must follow.

Here is why witness generation is CPU-bound. Polynomial arithmetic -- the core
of the proving step -- is naturally parallel. Number-theoretic transforms,
multi-scalar multiplications, and polynomial evaluations split naturally across
thousands of GPU cores. But VM emulation is inherently sequential. The next
instruction depends on the current instruction's result. You cannot execute
instruction 1,000 before you know the outcome of instruction 999.

\section{Witness Generation Costs}\label{witness-generation-costs}

The paper that this book revises claimed witness generation accounted for
10-25\% of total proving time. That figure was approximately correct in 2023,
when the proving step was slow enough to dwarf everything else. It is no longer
correct.

Modern GPU-accelerated provers have transformed the cost structure. When the
cryptographic proving step runs on an NVIDIA H100 GPU -- or a cluster of them --
it becomes 10 to 100 times faster than on a CPU. Multi-scalar multiplications,
the traditional bottleneck, have been optimized to the point where
number-theoretic transforms (NTTs) now account for up to 90\% of GPU proving
time. And NTTs themselves have been accelerated through pipelining, with systems
like BatchZK achieving over 3,000 times speedup over CPU baselines for Merkle
tree commitment operations.

But witness generation did not ride this wave. It remains CPU-bound. Sequential.
Memory-intensive. Accelerate the proving step by 10x and leave witness
generation unchanged, and watch the proportions shift:

Before GPU acceleration: witness generation takes 2 seconds, proving takes 8
seconds. Witness share: 20\%.

After GPU acceleration: witness generation still takes 2 seconds, proving takes
0.8 seconds. Witness share: 71\%.

Welcome to the Witness Gap -- and it is growing, not shrinking. Every
improvement to the proving step makes the witness gap worse by proportion. The
field has been optimizing the fast part and ignoring the slow part.

The numbers from recent profiling studies confirm this. ZKPOG, the first
end-to-end GPU acceleration system that treats witness generation as a
first-class optimization target, demonstrated that moving witness generation to
the GPU can yield 3-10x speedups. But this requires different parallelization
strategies than proving. Proving parallelizes naturally because polynomial
arithmetic is regular and data-independent. Witness generation requires
analyzing the circuit's dependency graph, topologically sorting gates to
identify independent clusters, and mapping irregular computation patterns onto
GPU hardware. The parallelism is there, but extracting it is harder.

BatchZK took a different approach: pipelining. Instead of generating the entire
witness first and then proving, BatchZK overlaps witness generation with proof
computation. The witness is generated in chunks and fed into the prover as a
stream. This significantly improves GPU utilization throughout the pipeline
compared to sequential approaches where the GPU sits idle during witness
generation.

The most radical proposal comes from Nair, Thaler, and Zhu, who showed that the
Jolt zkVM can be implemented with \emph{streaming witness generation} -- the
prover never materializes the full witness in memory. Instead, it generates
witness chunks on the fly, consumes them immediately in the sum-check protocol,
and discards them. Checkpoints at regular intervals enable parallel
regeneration. The space requirement drops from linear in the trace length to the
square root of the trace length, with less than 2x time overhead. For a
computation with \(2^{35}\) cycles, this means roughly 100 GB of working memory
instead of terabytes.

A third approach attacks the problem from the constraint side rather than the
computation side. Ozdemir, Laufer, and Boneh developed new algebraic interactive
proofs for RAM consistency checking -- the process of verifying that memory
reads and writes in the execution trace are consistent. Memory checking
dominates zkVM witness generation cost because the standard approach (Merkle
tree commitments) requires approximately \(600 \cdot A \cdot \log(N)\)
constraints for A accesses to N-sized memory. Each Poseidon hash in the Merkle
tree costs roughly 300 field multiplications. Their approach reduces this to
\(3N + 2A + O(1)\) constraints -- up to 51.3x fewer for persistent memory and
even more for sparse memory. By shrinking the constraint count for memory
operations, this work proportionally shrinks the witness that must be generated,
since memory-related witness entries often dominate the total witness size.

These are not theoretical proposals. They represent the frontier of a field that
has finally recognized where the real bottleneck lives.

\subsection{The Bottleneck That Flipped}\label{the-bottleneck-that-flipped}

The before-and-after illustrates a pattern that recurs throughout engineering:
the phenomenon where solving one problem promotes the next problem to dominance.

In 2023, a typical zero-knowledge proof for a moderately complex computation --
say, verifying a batch of Ethereum transactions -- took approximately 10 seconds
end to end. Of those 10 seconds, approximately 2 seconds were spent on witness
generation: emulating the virtual machine, recording every register value,
logging every memory access, building the complete execution trace. The
remaining 8 seconds were spent on the cryptographic proving step: computing
multi-scalar exponentiations, performing number-theoretic transforms, building
polynomial commitments. The witness generation share was 20\%. It was a minor
cost. Nobody optimized it because the proving step was the obvious target.

The field poured billions of dollars into making the proving step faster. GPU
implementations replaced CPU implementations. Custom NTT kernels exploited the
butterfly structure of the transform. Batched MSM algorithms amortized curve
operations across thousands of scalar multiplications. Pipeline architectures
overlapped memory transfers with computation. It worked. By 2025, GPU
acceleration had reduced the proving step from 8 seconds to 0.8 seconds -- a 10x
improvement, and some systems achieved 100x.

Witness generation was still 2 seconds.

The arithmetic is pitiless. Before: 2 / (2 + 8) = 20\%. After: 2 / (2 + 0.8) =
71\%. The proving step had become a solved problem. The witness step had become
the problem. And the worse the proving step got (in the sense of ``faster''),
the more dominant the witness step became. At 100x GPU acceleration, the proving
step drops to 0.08 seconds, and witness generation accounts for 96\% of total
time. At that point, further GPU optimization yields approximately zero
improvement to the end-to-end latency. You could make the proving step
infinitely fast and the proof would still take 2 seconds.

This is Amdahl's Law applied to zero-knowledge proofs. Amdahl's Law states that
the speedup of a program is limited by the fraction of the program that cannot
be parallelized. If witness generation accounts for 71\% of total time and
cannot be parallelized (because VM emulation is inherently sequential), then
even infinite parallelization of the remaining 29\% yields at most a 3.4x
overall speedup. The bottleneck is not a bottleneck you can throw hardware at.
It is a bottleneck in the structure of the computation itself.

\subsection{Why Witness Generation Resists
Parallelization}\label{why-witness-generation-resists-parallelization}

The reason witness generation is sequential is not an accident of
implementation. It is a consequence of what witness generation \emph{is}.

Consider a RISC-V program executing inside a zkVM. Each instruction reads from
registers, performs an operation, and writes the result to a register. The next
instruction reads from registers that may have been written by the previous
instruction. The program counter advances by one. Branches depend on comparison
results that were just computed. Memory loads depend on addresses that were just
calculated.

This is a dependency chain. Instruction 1000 cannot execute before instruction
999 because instruction 1000 might read a register that instruction 999 wrote.
Instruction 999 cannot execute before instruction 998 for the same reason. The
entire execution trace is a single sequential thread of dependencies, from the
first instruction to the last.

Compare this to the proving step. A number-theoretic transform operates on a
vector of field elements. Each butterfly operation in the transform reads two
elements, computes two outputs, and writes them back. The butterfly operations
within a single stage of the transform are \emph{independent} -- they read and
write disjoint memory locations. This means they can execute in parallel across
thousands of GPU cores. The parallelism is not extracted by clever engineering;
it is inherent in the mathematical structure of the transform.

Witness generation has no such structure. The ``butterfly'' of VM emulation is a
single instruction, and each instruction depends on the previous one. There is
no way to execute instruction 1000 before instruction 999 completes, because you
do not know what instruction 1000 \emph{is} until instruction 999 has updated
the program counter. A branch instruction at position 999 could send execution
to position 1000 or to position 5000 or anywhere else. You cannot know until the
branch condition is evaluated.

This is why three different research groups attacked the problem from three
different angles -- each targeting a different dimension of the sequential
bottleneck.

\textbf{Pipelining (BatchZK).} Instead of generating the entire witness first
and then proving, BatchZK overlaps the two phases. The witness is generated in
chunks: the first chunk of the execution trace is produced, then immediately fed
to the GPU prover while the CPU generates the next chunk. The GPU is never idle.
The CPU is never idle. The total wall-clock time drops because the two phases
execute concurrently. But the witness generation itself is not faster -- only
the overlap is new. Pipelining does not solve the sequential bottleneck; it
hides it behind the proving step.

\textbf{Streaming (Nair, Thaler, Zhu).} The streaming approach attacks the
memory dimension. Instead of materializing the full witness -- which for large
computations can require hundreds of gigabytes of RAM -- the streaming prover
generates witness chunks on demand, feeds them into the sum-check protocol, and
discards them. Checkpoints at regular intervals (every \(\sqrt{T}\) steps) allow
the prover to restart from any checkpoint, enabling a form of parallelism:
different proving threads can work on different segments of the trace, each
regenerating its segment from the nearest checkpoint. The space requirement
drops from \(O(T)\) to \(O(\sqrt{T})\), and the time overhead is less than 2x.
For a computation with \(2^{35}\) cycles, this means roughly 100 GB of working
memory instead of terabytes.

\textbf{Algebraic RAM reduction (Ozdemir, Laufer, Boneh).} The most radical
approach does not try to make witness generation faster or smaller. It tries to
make the witness \emph{simpler}. Memory checking -- verifying that memory reads
and writes in the execution trace are consistent -- is the dominant cost in zkVM
witness generation because the standard approach uses Merkle trees. Each memory
access requires a Poseidon hash, and each Poseidon hash costs approximately 300
field multiplications. For a program with A memory accesses to N-sized memory,
the standard approach requires approximately \(600 \cdot A \cdot \log(N)\)
constraints. Ozdemir et al.~replace Merkle tree checking with algebraic
interactive proofs that require only \(3N + 2A + O(1)\) constraints -- up to
51.3x fewer. Fewer constraints means a smaller witness for the memory-checking
portion, which often dominates the total witness.

Each approach has trade-offs. Pipelining requires careful synchronization
between the CPU witness generator and the GPU prover. Streaming requires
checkpoint management and sacrifices some proving speed for memory savings.
Algebraic RAM reduction requires new interactive proof protocols that are not
yet implemented in production systems. But together, they represent the first
serious engineering effort to close the witness gap -- the gap that the field
spent years ignoring because the proving step was the shinier problem.

The key performance numbers tell the story:

\begin{longtable}[]{@{}lll@{}}
\toprule\noalign{}
Metric & Value & Source \\
\midrule\noalign{}
\arrayrulecolor{MidnightBlue}\midrule\arrayrulecolor{MidnightBorder}
\endhead
\bottomrule\noalign{}
\endlastfoot
Witness generation share (with GPU proving) & 50-70\% of total time & ZKPOG \\
NTT share of GPU proving time & up to 90\% & ZKProphet \\
GPU pipeline speedup over CPU (sum-check protocol) & 3,040x & BatchZK \\
GPU pipeline speedup over CPU (Merkle tree) & 793x & BatchZK \\
GPU pipeline memory per proof & 0.08-0.44 GB & BatchZK \\
Streaming prover space reduction & \(O(\sqrt{KT})\) & Nair et al. \\
RAM constraint reduction & up to 51.3x & Ozdemir et al. \\
ZKPOG end-to-end GPU speedup & 22.8x average & ZKPOG \\
\end{longtable}

\section{Memory: The Binding Constraint}\label{memory-the-binding-constraint}

The Witness Gap is also a memory problem.

GPU proving requires a minimum of 24 GB of VRAM -- which excludes every consumer
GPU below the NVIDIA RTX 4090 (approximately \$2,000). Large computations demand
far more. Jolt and ZKM can require 128 GB of system RAM. Groth16 for circuits
with \(2^{25}\) constraints needs approximately 200 GB of RAM. An Ethereum block
execution trace, with its millions of state accesses and storage operations, can
require specialized hardware with 512 GB or more.

These are not abstract numbers. They determine who can generate proofs and who
cannot.

If you are a rollup operator running a proving cluster in a data center with 16
NVIDIA H100 GPUs and dual-socket servers with 512 GB of RAM, the memory
requirements are a budgeted operating expense. You buy the hardware, you run the
provers, you amortize the cost across millions of transactions.

But if you are an individual user generating proofs on your own device -- the
scenario that provides maximum privacy, because your private data never leaves
your machine -- the memory requirements become a barrier to entry. A laptop with
16 GB of RAM cannot generate proofs for non-trivial computations. A phone cannot
even attempt it.

This leads to an uncomfortable conclusion, and it deserves to land with full
force: \emph{privacy is, in part, a luxury good}. The privacy-maximizing
architecture (client-side proving, where your secrets never leave your device)
requires hardware that most people do not own. The privacy-minimizing
architecture (delegated proving, where you send your private data to a proving
service) works on any device but requires trusting the service with your
secrets.

Read that again. The architecture that protects your data the most demands
hardware that costs the most. The architecture that exposes your data to third
parties is the one available to everyone. This is not a theoretical concern. It
is the economic structure of privacy in 2026, and it should disturb anyone who
believes privacy is a right rather than a commodity.

\subsection{The Hardware Ladder}\label{the-hardware-ladder}

The abstract claim becomes concrete when you map it to specific hardware tiers.
Here is what zero-knowledge proving looks like at each rung of the hardware
ladder, from the bottom up.

\textbf{Tier 1: A laptop with 16 GB of RAM (\textasciitilde\$800-1,500).} You
can prove trivial computations -- a few thousand constraints, the kind found in
simple identity attestations or basic Merkle membership proofs. Anything
resembling a real application (a token transfer with privacy, a complex smart
contract execution, a state transition proof) exceeds your memory budget. The
witness alone may require more RAM than you have. At this tier, you must
delegate proving to a remote service. Your private data -- the witness, which
contains every secret the proof is supposed to protect -- leaves your machine.
You are trusting someone else's hardware with your secrets.

\textbf{Tier 2: An NVIDIA RTX 4090 with 24 GB VRAM (\textasciitilde\$2,000 for
the GPU alone, \textasciitilde\$4,000 for a workstation that can host it).} You
can prove moderately complex circuits locally. This is the minimum hardware for
client-side privacy-preserving proofs of meaningful complexity. ZKPOG targets
this tier explicitly, arguing that it represents the frontier of
``democratized'' proving. Midnight's proof server targets similar hardware. At
this tier, your secrets stay on your machine. You have genuine privacy. But you
have also spent \$4,000 on a workstation, which places you in the top 5\% of
computing hardware ownership globally.

\textbf{Tier 3: A data center server with an NVIDIA H100 GPU, 80 GB of HBM3
memory (\textasciitilde\$30,000 for the GPU, \textasciitilde\$50,000-80,000 for
a server).} You can prove Ethereum blocks in real time. You can run a rollup's
proving infrastructure. You can handle the most complex circuits that production
systems generate today. This is the rollup operator tier -- the hardware that
Succinct, RISC Zero, and other proving services deploy. The H100's HBM3 provides
3.35 TB/s of memory bandwidth, which is the binding constraint for NTT
performance at large polynomial sizes. No consumer GPU comes close.

\textbf{Tier 4: A cluster of 16 NVIDIA H100 GPUs (\textasciitilde\$500,000 for
the GPUs alone, \textasciitilde\$1 million or more for the full cluster with
networking, cooling, and redundancy).} You can prove the most complex
computations that exist in the ZK ecosystem: large zkVM programs with billions
of cycles, full Ethereum block proofs with all precompiles, recursive proof
compositions with deep nesting. This is the proving-as-a-service tier. Companies
like Succinct and Gevulot operate at this level. A single proof that would take
an RTX 4090 several minutes completes in seconds when the computation is sharded
across 16 H100s with high-bandwidth interconnect.

\textbf{Tier 0: A smartphone with 4-8 GB of RAM (\textasciitilde\$200-800).}
This is what most of the world actually owns. At this tier, zero-knowledge
proving of any meaningful complexity is not slow -- it is impossible. The memory
is insufficient. The compute is insufficient. The thermal envelope is
insufficient. A phone cannot generate a ZK proof for a shielded transaction, a
private vote, or a verifiable computation of any significant size. Not ``cannot
generate it quickly'' -- cannot generate it at all.

The uncomfortable math: approximately 5.5 billion people on Earth own a
smartphone. Approximately 200 million own a desktop or laptop with a discrete
GPU capable of ZK proving. Of those, perhaps 10-20 million own hardware at Tier
2 or above. That means roughly 95\% of the world's population -- including
nearly all smartphone-only users in developing economies -- cannot perform
client-side ZK proving. They must delegate. They must trust. The cryptographic
guarantee of privacy is available to them only through the intermediation of
someone else's hardware.

This is not a bug in the technology. It is a structural feature of the cost
curve. Moore's Law may eventually bring proving hardware to lower price points.
Algorithmic improvements (streaming provers, algebraic RAM reduction) may lower
the hardware floor. But in 2026, the privacy hierarchy is clear: the richer your
hardware, the more private your computation. The poorer your hardware, the more
you must trust others with your secrets.

There is a historical parallel. In the 1990s, strong encryption was classified
as a munition by the United States government. Export-grade encryption was
deliberately weakened to 40-bit keys, ensuring that only domestic users (and the
NSA) had access to real cryptographic security. The rest of the world got a
pantomime of privacy. The ``Crypto Wars'' ended when the government relented and
strong encryption became universally available. The current ZK hardware barrier
is not a government restriction -- it is an economic one. But the effect is
similar: strong privacy for the few, weak privacy (or no privacy) for the many.
Whether this barrier will fall, as the export restrictions did, depends on
whether the field can make proving cheap enough to run on the hardware that
people actually own.

The field is aware of this tension. ZKPOG specifically targets the NVIDIA RTX
4090 (24 GB VRAM, approximately \$2,000) as its hardware platform, arguing that
democratizing ZK proving requires targeting accessible consumer hardware rather
than data center GPUs. Streaming witness generation reduces memory requirements
at the cost of additional computation. And proof delegation with trusted
execution environments (TEEs) offers a middle path -- your data is processed
inside a secure enclave that even the hardware operator cannot inspect. But TEEs
have their own vulnerability history (Foreshadow, AEPIC Leak, Downfall), and
Intel deprecated SGX on consumer processors in 2021.

The memory constraint also interacts with NTT performance in a way that matters
for system design. NTT -- the number-theoretic transform used in polynomial
commitment schemes -- is memory-bandwidth-limited at large sizes. The butterfly
structure of the NTT requires global memory accesses with poor locality in later
stages, meaning that the speed of proving is ultimately limited not by compute
(FLOPS) but by memory bandwidth (GB/s). HBM (High Bandwidth Memory) on data
center GPUs provides the bandwidth; consumer GPUs do not.

For a system architect making infrastructure decisions, the takeaway is this:
when evaluating ZK proving solutions, ask about memory, not just speed. A system
that generates proofs in 3 seconds but requires 256 GB of RAM is not the same as
a system that generates proofs in 10 seconds but runs in 32 GB. The first
requires a data center. The second runs on a workstation. A proving service that
advertises ``sub-second proofs'' but requires an H100 cluster is making a
different claim than one that advertises ``ten-second proofs'' on consumer
hardware.

The memory question also intersects with the privacy question. If client-side
proving requires 24 GB of VRAM, and only data center GPUs and the NVIDIA RTX
4090 have that much, then privacy-preserving client-side proving is available to
a narrow slice of users. Everyone else must delegate proving to a service, which
means sending their private witness to someone else's hardware. The Midnight
architecture partially addresses this by running the proof server locally
alongside the dApp -- but ``locally'' still means ``on hardware with sufficient
resources.'' The 18-24 second proof generation time observed on Midnight's
development environment reflects desktop-class hardware. On a mobile device, the
same computation might be infeasible.

\section{Side-Channel Attacks: When the Walls
Leak}\label{side-channel-attacks-when-the-walls-leak}

The mathematical definition of zero-knowledge is precise: the proof reveals
nothing about the witness beyond the truth of the statement being proved. But
the mathematical definition describes the \emph{proof}. It does not describe the
\emph{process of generating the proof}.

The process can leak.

The most vivid demonstration came in a 2020 USENIX Security paper by Tramer,
Boneh, and Paterson. They showed that Zcash's Groth16 prover leaked information
about transaction amounts through proof generation time. The attack was simple
in concept: the prover's multi-scalar exponentiation (MSM) implementation
optimized away terms where the witness coefficient was zero. More zeros in the
binary representation of the transaction amount meant fewer curve
multiplications, which meant faster proof generation. By measuring how long
proof generation took -- remotely, across the network -- an attacker could
estimate the Hamming weight of the transaction amount.

The correlation coefficient was R = 0.57. Not perfect, but far from zero. A
timing measurement that should have been meaningless -- how long did the proof
take? -- revealed information about the secret that the proof was supposed to
protect.

Monero's Bulletproofs implementation was safe. The correlation was R = 0.04 --
essentially noise. The difference was architectural: Bulletproofs operates on
both the binary decomposition of the amount \emph{and its complement}, making
the number of curve operations constant regardless of the value. The proof
generation time was the same whether the amount was 1 or 1,000,000.
Constant-time implementation is not an optimization. It is a security
requirement.

\subsection{The Zcash Timing Attack: A Detective
Story}\label{the-zcash-timing-attack-a-detective-story}

The Zcash attack deserves to be told as the detective story it was, because the
investigative method reveals how side-channel analysis works in practice -- and
why it is so difficult to defend against.

The researchers began with a hypothesis: if the Groth16 prover's multi-scalar
exponentiation skips zero-valued scalar multiplications, then the proof
generation time should correlate with the number of nonzero bits in the witness.
They did not need to break any cryptography. They did not need to find a flaw in
the mathematics. They needed a stopwatch.

They set up a Zcash node and generated shielded transactions with known amounts.
For each transaction, they measured the wall-clock time of the
\texttt{r1cs\_gg\_ppzksnark\_prover} function -- the core Groth16 proving
routine in libsnark. They varied the transaction amount systematically, from
round numbers like 1.000 ZEC (which has a simple binary representation with many
trailing zeros) to numbers with many nonzero digits like 1.337 ZEC (which has a
denser binary representation).

The pattern emerged immediately. Round numbers produced faster proofs. Dense
numbers produced slower proofs. The relationship was not subtle. A transaction
of exactly 1.000 ZEC generated a proof measurably faster than a transaction of
1.337 ZEC, because the scalar representation of 1.000 ZEC in the finite field
contained more zero coefficients, and each zero coefficient allowed the MSM
algorithm to skip a point addition.

The correlation coefficient was R = 0.57 -- not strong enough to determine the
exact amount, but strong enough to distinguish broad categories. An attacker
could not tell whether you sent 1.337 ZEC or 1.338 ZEC. But the attacker could
distinguish ``approximately 1 ZEC'' from ``approximately 100 ZEC'' with
meaningful confidence. The ``zero-knowledge'' proof leaked the order of
magnitude of the transaction amount through the duration of its generation.

What makes this attack memorable is not its severity -- R = 0.57 is a partial
leak, not a catastrophic one -- but what it reveals about the gap between
mathematical proof and physical implementation. The Groth16 proof system is
provably zero-knowledge. The mathematical proof of this property is correct. The
implementation of the prover was not constant-time, and so the \emph{process} of
generating the proof leaked information that the \emph{proof itself} was
mathematically guaranteed to conceal. The proof was zero-knowledge. The prover
was not.

The fix was straightforward: replace the variable-time MSM with a constant-time
implementation that processes all scalar coefficients identically, whether they
are zero or nonzero. The constant-time version is slower -- it performs
multiplications that the variable-time version skips -- but it eliminates the
timing channel. Zcash implemented this fix. The performance cost was real. The
privacy gain was essential.

\subsection{The Poseidon Cache-Timing
Attack}\label{the-poseidon-cache-timing-attack}

The second attack class is subtler, and in some ways more disturbing, because it
targets a component that was specifically designed for zero-knowledge systems.

Poseidon is an ``algebraic'' hash function. Unlike SHA-256, which was designed
for general-purpose hashing and happens to be usable (expensively) inside ZK
circuits, Poseidon was built from the ground up to minimize the number of
constraints required to express its computation as a polynomial relation. Where
SHA-256 requires approximately 25,000 constraints per hash in a Groth16 circuit,
Poseidon requires approximately 300. This 80x reduction in constraint count
translates directly into faster proving and smaller proofs. Poseidon is, by
design, the ideal hash function for zero-knowledge systems. Every major ZK
project uses it or a close variant.

But Poseidon's S-box -- the nonlinear component that provides cryptographic
security -- involves computing \(x^5\) (or \(x^7\), depending on the variant)
over a large prime field. In software, this is typically implemented using
lookup tables that map input values to output values. The S-box computation
accesses these tables at indices determined by the internal state of the hash,
which depends on the secret input being hashed.

This is where cache timing enters. Modern CPUs use a hierarchy of caches (L1,
L2, L3) to speed up memory access. When a program accesses a memory location,
the CPU loads the surrounding cache line (typically 64 bytes) into the L1 cache.
Subsequent accesses to the same cache line are fast (a few cycles). Accesses to
different cache lines that map to the same cache set can evict earlier entries,
making them slow again (hundreds of cycles).

An attacker sharing the same physical CPU -- a realistic scenario in cloud
computing environments where virtual machines share hardware -- can observe
which cache lines the victim's Poseidon computation accesses. The attacker
primes the cache (fills it with known data), waits for the victim's hash
computation to execute, then probes the cache to see which of the attacker's
entries were evicted. The eviction pattern reveals which table entries the
victim accessed, which reveals information about the internal state of the hash,
which reveals information about the secret input.

The hash function was designed to be ZK-friendly in algebra. It was not designed
to be constant-time in hardware. The algebraic design and the implementation
security were treated as separate concerns, and the gap between them is
exploitable.

This is not hypothetical. Cache-timing attacks are a well-studied attack class
with decades of published results against AES, RSA, and other cryptographic
primitives. The novelty of Mukherjee et al.'s work is showing that the same
attack class applies to ZK-specific constructions -- and that the ZK community's
emphasis on algebraic efficiency has, in some cases, actively increased
vulnerability by encouraging table-based designs.

\subsection{The Electromagnetic Channel}\label{the-electromagnetic-channel}

The third attack class operates at a physical level that most software engineers
never consider.

Every electronic circuit, when it operates, produces electromagnetic emanations.
A transistor switching from 0 to 1 consumes a different amount of current than a
transistor remaining at 0. This current difference creates a magnetic field that
propagates outward from the chip. The field is weak -- microwatts -- but it is
measurable with equipment that costs a few hundred dollars: a near-field
electromagnetic probe, a low-noise amplifier, and a digital oscilloscope.

Field operations in elliptic curve arithmetic are particularly vulnerable. When
a prover computes a point addition on an elliptic curve, the specific operations
performed (and their power consumption) depend on the coordinates of the points
being added, which depend on the witness values. A modular multiplication where
both operands are large consumes more power than one where an operand is small.
A conditional branch (reduce or do not reduce after multiplication) produces a
different electromagnetic signature depending on which path is taken.

Measuring these emanations from a few centimeters away -- close enough to touch
the device but not close enough to require physical modification -- can
reconstruct the scalar values used in multi-scalar exponentiation. This means
reconstructing the witness coefficients. This means reconstructing the private
inputs to the proof.

This is not science fiction. Electromagnetic side-channel attacks are a
published, demonstrated attack class. They have been used to extract AES keys
from smartcards, RSA private keys from laptops, and ECDSA signing keys from
hardware security modules. The equipment required is modest: a near-field probe
(\$50-200), an amplifier (\$100-500), and an oscilloscope (\$500-5,000). A
university research lab can mount this attack. A well-funded adversary can mount
it from across a room using more sensitive antennas.

For zero-knowledge provers running on commodity hardware -- laptops, desktops,
even data center servers without electromagnetic shielding -- the EM channel is
an open question. No major ZK implementation has published an electromagnetic
side-channel analysis. The attack surface is real, the equipment is cheap, and
the countermeasures (electromagnetic shielding, randomized execution ordering,
amplitude-flattening power regulation) are not part of any ZK prover's design
requirements.

The three attack channels -- timing, cache, electromagnetic -- form a hierarchy
of increasing physical intimacy. Timing attacks can be mounted remotely, across
a network. Cache attacks require co-location on the same physical machine.
Electromagnetic attacks require physical proximity to the hardware. But all
three extract information from the same fundamental source: the fact that
computation is a physical process, and physical processes leave physical traces.
The mathematical abstraction of a zero-knowledge proof exists in a world of pure
information. The implementation exists in a world of transistors, cache lines,
and electromagnetic fields. The gap between those worlds is where privacy leaks.

The attack extended beyond timing. Mukherjee, Rechberger, and Schofnegger
published the first systematic study of cache timing leakages in zero-knowledge
protocols in 2024. They examined ZK-friendly hash functions (Poseidon,
Reinforced Concrete, Tip5, Monolith) and popular proof systems (Groth16,
Plonky2, Plonky3, halo2, Circle STARKs). Here is what they found.

ZK-friendly hash functions were designed for \emph{algebraic} efficiency -- they
minimize the number of constraints required to express a hash computation inside
a circuit. But nobody designed them for \emph{implementation} security.
Reinforced Concrete uses large lookup tables (256 KB for its Bars function)
indexed by secret-dependent data. These table lookups create cache access
patterns that vary with the secret. In a shared cloud environment, where the
attacker runs on the same physical machine as the prover, these cache patterns
are observable.

The irony cuts deep: the move toward lookup-based designs in ZK hash functions
-- motivated by the algebraic efficiency gains that Lasso and Jolt demonstrated
-- actively increases the side-channel attack surface. The very optimization
that makes proving faster makes the proving process less private.

Field arithmetic itself leaks. The Goldilocks field (\(2^{64} - 2^{32} + 1\))
uses conditional reductions after arithmetic operations. If the result exceeds
the modulus, a reduction step is needed; if not, it is skipped. This conditional
branch creates a timing signal. Assembly implementations with branch-free code
mitigate this, but many deployed field arithmetic libraries use branch-dependent
paths for performance.

The accurate assessment: ``zero-knowledge'' is a mathematical property of the
\emph{proof}. Implementation zero-knowledge depends on the hardware, the
operating system, the runtime, and the network. The gap between ``the proof
reveals nothing'' and ``the timing reveals everything'' is the gap between
theory and practice. The backstage walls are thinner than the magician thinks.

For defensive implementation, the standard is clear:

\begin{itemize}
\tightlist
\item
  No zero-skipping optimizations in multi-scalar exponentiation.
\item
  No secret-dependent table lookups in hash functions.
\item
  Branch-free field arithmetic, especially for conditional reduction steps.
\item
  Cache-aligned memory access patterns.
\item
  Constant-time prover implementations throughout the pipeline.
\end{itemize}

These requirements conflict with performance optimization at almost every turn.
Making a prover constant-time means foregoing shortcuts that can halve
computation time. The tension between proving speed and implementation privacy
is real and ongoing.

The interaction between side channels and GPU proving adds another dimension.
GPU architectures use SIMT (Single Instruction, Multiple Threads) execution,
where groups of 32 threads (warps on NVIDIA hardware) execute the same
instruction simultaneously. Constant-time code requires all threads in a warp to
follow the same execution path. When some threads need a conditional reduction
and others do not, the warp must execute both paths, with inactive threads
masking their results. This ``thread divergence'' reduces GPU utilization -- the
very parallelism that makes GPUs fast for proving works against the
constant-time requirement.

The open question is whether witness generation can be made fully constant-time
on GPUs without unacceptable performance loss. The answer is not yet clear. What
is clear is that any system claiming both GPU-accelerated proving and
zero-knowledge must address this tension explicitly. Most do not.

For the reader who wants a single mental model: the backstage walls are made of
different materials at different heights. The cryptographic walls (the proof
itself) are mathematically perfect -- no information passes through. The
implementation walls (the proving process) are made of timing signals, cache
patterns, and memory access traces. They leak. Not catastrophically, not in
every deployment, but measurably and exploitably in the wrong environment. The
system architect's job is not to eliminate all leakage -- that may be impossible
-- but to understand where the walls are thin and what information an attacker
on the other side could extract.

\section{Witness-Constraint Divergence}\label{witness-constraint-divergence}

The witness is not just the most expensive artifact in the ZK pipeline. It is
also the most dangerous place for bugs.

Remember the dual-track problem from Chapter 3: in many ZK systems, the witness
generator and the constraint system are two separate programs that must compute
identical functions on all inputs. When they disagree, the result is either a
soundness bug (the proof system accepts false statements) or a completeness bug
(the proof system rejects true statements). Both are bad. Soundness bugs are
worse.

Two real-world examples illustrate the stakes.

RISC Zero disclosed CVE-2025-52484: a missing constraint in the RISC-V circuit
that allowed confusion between the rs1 and rs2 register operands. The witness
generator computed the correct values for both registers. The constraint system
did not enforce that the registers were distinct. A malicious prover could
substitute one register for another, and the proof would still verify. The
computation would appear valid -- the proof would pass -- but the result would
be wrong.

The zkSync Era MemoryWriteQuery bug was worse. The struct that handled memory
write operations failed to call \texttt{lc.enforce\_zero(cs)} on the highest 128
bits of a 256-bit value, leaving those bits unconstrained. A malicious prover
could modify the highest 128 bits of any memory write -- including withdrawal
amounts. The proof system would accept the modification. An attacker could
change a withdrawal of 0.00002 ETH to a withdrawal of 100,000 ETH, and the proof
would verify.

These bugs share a common structure: the witness was correct, but the
constraints were insufficient. The magician performed the trick honestly
backstage, but the constraint system's description of what ``honest'' meant was
incomplete. The proof certified that the computation matched the constraints.
The constraints did not match the intended computation.

Call it the correctness gap: the distance between what the developer meant and
what the constraints actually enforce. It is measured not in bits of security
but in lines of code that were not written.

Multiple valid witnesses can exist for the same statement, and the proof system
does not care which one the prover uses. If you prove that you know a square
root of 25, the proof system accepts whether you use 5 or -5. This is by design
-- the proof system guarantees soundness (you cannot prove a false statement),
not uniqueness (there is only one valid witness).

This property -- that multiple valid witnesses exist -- is fundamental, not a
bug. The proof system guarantees that the prover knows \emph{some} valid
witness. It does not guarantee which one. For most applications, this is exactly
right: if you prove you know a valid password, it does not matter whether you
know the first password in the hash table or the last.

Non-deterministic hints exploit this deliberately. Instead of computing a square
root step by step -- which is expensive inside a circuit -- the prover
\emph{guesses} the answer (as a witness value) and the circuit verifies that the
square of the guess equals 25. The computation goes from expensive (sequential
square root algorithm) to cheap (one multiplication and one comparison). This
``guess and check'' pattern is a standard programming technique in ZK systems,
not a bug. But it requires that the checking constraints are complete. If the
check only verifies \(x \cdot x = 25\) without constraining that x is in the
correct range, a malicious prover might find a field element that satisfies the
equation but is not the intended square root.

\section{\texorpdfstring{The \texttt{disclose()} Boundary: Midnight's Witness
Architecture}{The disclose() Boundary: Midnight's Witness Architecture}}\label{the-disclose-boundary-midnights-witness-architecture}

Midnight's approach to witness generation illustrates both the power and the
subtlety of the witness/circuit separation.

In Compact, the two worlds are sharply delineated:

\begin{longtable}[]{@{}
  >{\raggedright\arraybackslash}p{(\linewidth - 4\tabcolsep) * \real{0.3333}}
  >{\raggedright\arraybackslash}p{(\linewidth - 4\tabcolsep) * \real{0.3333}}
  >{\raggedright\arraybackslash}p{(\linewidth - 4\tabcolsep) * \real{0.3333}}@{}}
\toprule\noalign{}
\begin{minipage}[b]{\linewidth}\raggedright
Property
\end{minipage} & \begin{minipage}[b]{\linewidth}\raggedright
Witnesses
\end{minipage} & \begin{minipage}[b]{\linewidth}\raggedright
Circuits
\end{minipage} \\
\midrule\noalign{}
\arrayrulecolor{MidnightBlue}\midrule\arrayrulecolor{MidnightBorder}
\endhead
\bottomrule\noalign{}
\endlastfoot
Where they run & Off-chain (user's browser or node) & Compiled to ZKIR, proven
locally, verified on-chain \\
What they can do & Arbitrary computation (JavaScript) & Only ZK-provable
computation (field arithmetic, hashing, comparisons) \\
What they access & Private state, external APIs, databases & Only circuit inputs
(public and disclosed private values) \\
Privacy & Completely private -- never leaves the device & Proven but not
revealed \\
\end{longtable}

Witness functions are \emph{declared} in Compact but \emph{implemented} in
TypeScript:

\begin{Shaded}
\begin{Highlighting}[]
\NormalTok{witness }\FunctionTok{get\_secret}\NormalTok{()}\OperatorTok{:}\NormalTok{ Bytes}\OperatorTok{\textless{}}\DecValTok{32}\OperatorTok{\textgreater{};}  \CommentTok{// declared in Compact}
\end{Highlighting}
\end{Shaded}

\begin{Shaded}
\begin{Highlighting}[]
\CommentTok{// implemented in TypeScript}
\KeywordTok{const}\NormalTok{ witnesses }\OperatorTok{=}\NormalTok{ \{}
  \DataTypeTok{get\_secret}\OperatorTok{:}\NormalTok{ (ctx) }\KeywordTok{=\textgreater{}}\NormalTok{ [ctx}\OperatorTok{.}\AttributeTok{privateState}\OperatorTok{,}\NormalTok{ secretKey]}
\NormalTok{\}}\OperatorTok{;}
\end{Highlighting}
\end{Shaded}

Each witness function receives the current context -- including a
\texttt{privateState} object that persists across invocations -- and returns a
tuple of the updated private state and the witness value. This means witnesses
can maintain state, query external services, read databases, and perform
arbitrary computation. They are full JavaScript programs, not circuit-compatible
snippets.

The \texttt{disclose()} operator is the \emph{sole gateway} from the witness
world to the circuit world. Without it, witness values are invisible to the
circuit, the proof, and the chain. With it, the value enters the circuit as a
\texttt{private\_input} in the ZKIR -- the verifier never sees the value, but
the constraints verify properties about it.

The compiled ZKIR circuit has two transcript channels that encode this
separation physically:

The \texttt{publicTranscript} records every ledger operation -- reads, writes,
comparisons -- as a sequence of VM operations. The on-chain verifier sees this
transcript and checks that it is consistent with the proof. If the developer's
circuit reads a counter value from the ledger, that read appears in the public
transcript. If the circuit increments the counter, that increment appears. The
public transcript is the complete audit trail of on-chain state changes.

The \texttt{privateTranscriptOutputs} contains the witness values that entered
the circuit through \texttt{disclose()}. Only the prover sees these. They are
consumed during proof generation by \texttt{private\_input} instructions in the
ZKIR and then discarded. The ZKIR checker verifies that both transcripts are
fully consumed -- no extra values, no missing values, no tampering.

This two-transcript model achieves something precise: the verifier knows
\emph{what changed} on the ledger (from the public transcript) and is convinced
that the changes are valid (from the ZK proof), but does not know \emph{why} the
changes were made (the private inputs that drove the computation). The ``what''
is public. The ``why'' is private.

In practice, the pipeline for a Midnight transaction involves four sequential
steps:

\begin{enumerate}
\def\labelenumi{\arabic{enumi}.}
\item
  \texttt{contracts.callTx()} reads current ledger state via the indexer,
  executes the circuit locally with current state and private witnesses, and
  produces an unproven transaction.
\item
  \texttt{proofProvider.proveTx()} generates the ZK proof -- the dominant
  latency step, as detailed in Chapter 6 -- producing a proven transaction.
\item
  \texttt{walletProvider.balanceTx()} binds the transaction, runs token
  balancing (unshielded, shielded, and dust), signs UTXO inputs with BIP-340
  Schnorr signatures, and merges the balancing transaction with the original.
\item
  \texttt{midnightProvider.submitTx()} submits to the blockchain, where the node
  verifies the ZK proof, checks that the public transcript matches the ledger
  state, and applies the state transition.
\end{enumerate}

At no point do witnesses cross the network. The developer guide states this
explicitly: ``Witnesses stay local. Never sent to chain.''

The side-channel implications matter here. Proof generation takes a fixed amount
of time regardless of witness values -- dominated by the cryptographic proving
step, not the witness computation. This provides natural but unintentional
timing uniformity: the fixed cost of proof generation drowns out any timing
variation in witness computation. But the documentation does not address cache
timing, network timing (when a user queries the indexer immediately before
submitting a transaction, the timing correlation reveals which contract state
they are acting on), or transaction structure analysis (the number of segments
in a transaction could reveal which circuit was called).

Privacy on Midnight is genuine at the cryptographic level. It is unexamined at
the implementation level. This is not a criticism unique to Midnight -- it
applies to every privacy-preserving system in production. But it is worth noting
that the same project that provides the most rigorous compile-time privacy
guarantees (disclosure analysis) has the least documented runtime privacy
analysis.

One accidental privacy benefit: the fixed cost of proof generation -- dominated
by the PLONK proving step -- provides natural timing padding. Whether the
witness contains a trivial secret or a complex multi-step computation, the proof
generation time is approximately the same. An observer measuring transaction
timing cannot easily distinguish between different witness computations. The
padding is natural but not designed. A dedicated attacker measuring sub-second
timing variations in the witness computation phase (which occurs before proof
generation) might still extract information. But the 18-second proving step
provides a large, fixed-duration buffer that dominates the total transaction
time.

The developer guide's demonstration of the private voting dApp provides a
concrete end-to-end example. The off-chain witness construction computes
Poseidon hashes using the same \texttt{persistentHash} function available in
both the Compact circuit and the JavaScript SDK. The voter provides their secret
key, vote choice, Merkle sibling hash, and direction flag as witnesses. The
circuit reconstructs the Merkle root from these witnesses and asserts it matches
the on-chain root. The circuit also computes a nullifier --
\texttt{persistentHash({[}pad(32,\ "nullf:"),\ secret\_key{]})} -- using domain
separation so that the nullifier hash and the voter's attestation hash are
cryptographically unlinkable. The nullifier is disclosed (it must appear
on-chain to prevent double-voting), but it cannot be traced back to the voter's
identity.

What the on-chain verifier sees per vote: a nullifier hash and updated vote
tallies. What the verifier does not see: which voter, their exact identity, or
which Merkle leaf they occupy. The privacy boundary is sharp, and it is enforced
at every level: by the compiler (disclosure analysis), by the ZKIR checker
(transcript integrity), and by the cryptographic construction (domain-separated
hashing).

\section{The Witness as a Multi-Dimensional
Problem}\label{the-witness-as-a-multi-dimensional-problem}

The research literature reveals that the ``Witness Gap'' is not a single
bottleneck but a convergence of four distinct challenges, each requiring
different solutions.

\textbf{The Performance Gap.} Witness generation has been neglected relative to
MSM and NTT optimization because it does not parallelize in the same way. ZKPOG
shows that GPU-accelerated witness generation is possible (3-10x speedups) but
requires analyzing the circuit's dependency graph and topologically sorting
gates to identify independent clusters. The parallelism exists, but extracting
it is harder than parallelizing polynomial arithmetic.

\textbf{The Memory Gap.} The full witness for a large computation can require
hundreds of gigabytes of RAM. Streaming witness generation -- never
materializing the full witness, instead generating chunks on the fly and
consuming them immediately -- is the path forward. Nair, Thaler, and Zhu showed
this can be achieved with \(O(\sqrt{T})\) space and less than 2x time overhead,
using checkpoints at regular intervals for parallel regeneration.

\textbf{The Security Gap.} The witness is the most sensitive artifact in the
system. It contains the private inputs. Side-channel attacks can leak witness
information through timing (R=0.57 in Zcash), cache patterns (Mukherjee et
al.~2024), and network metadata. Constant-time implementation is a security
requirement, not a performance optimization.

\textbf{The Correctness Gap.} The witness generator and the constraint system
must compute identical functions. When they disagree, the result is a soundness
bug. Static analysis tools like ZKAP (F1 score 0.82, 34 previously unknown
vulnerabilities discovered) can detect divergence, but they currently work only
on Circom. Extending them to Rust-based systems (halo2, Plonky3) remains an open
problem.

These four gaps interact. Solving the performance gap (GPU acceleration) can
worsen the security gap (GPU thread divergence from constant-time code reduces
SIMT utilization). Solving the memory gap (streaming) changes the architecture
in ways that affect the correctness gap (streaming provers must handle state
differently than batch provers). There is no single fix. The witness problem is
systemic.

The analogy holds: the magician's backstage is not just dark -- it is expensive,
fragile, and surveilled. The recording equipment costs a fortune. The walls have
cracks. And if the recording is wrong, the audience will believe a lie. Layer 3
is where the practical reality of zero-knowledge systems diverges most sharply
from the elegant theory. The mathematics is beautiful. The engineering is
brutal.

For the system architect, Layer 3 generates the most concrete questions in any
ZK evaluation:

\begin{itemize}
\tightlist
\item
  What is the witness generation time for your target workload, and how does it
  compare to the proving time? If it exceeds 50\%, your proving GPU is idle most
  of the time.
\item
  What are the memory requirements? Can the witness fit in the VRAM of your
  target GPU, or must it be streamed from system RAM?
\item
  Is the prover constant-time? If not, what information does the timing profile
  reveal about private inputs?
\item
  Is client-side proving feasible on your users' hardware? If not, what is the
  trust model for delegated proving?
\item
  How does the witness generator handle the correctness gap? Is the constraint
  system formally verified, statically analyzed, or tested against the witness
  generator?
\end{itemize}

These are not theoretical questions. They determine whether a ZK system provides
the properties it claims. A system with fast proving but slow witness
generation, insufficient memory, timing leaks, and unverified constraints is a
system that looks good on benchmarks and fails in production.

\par\vspace{18pt}

The recording is made. It is expensive to produce. It is vulnerable to side
channels. It is the most common site of implementation bugs.

But the witness is just a recording. It proves nothing by itself. Anyone could
fabricate a recording. The question that Layer 4 must answer is: how do we turn
this recording into a mathematical puzzle -- a system of polynomial equations --
such that checking the puzzle is vastly cheaper than re-doing the computation?
How do we encode a million steps of execution into a form where a few random
spot-checks are enough to guarantee that every step was correct?

That transformation is the subject of Layer 4: the most technically demanding
layer in the stack, and the one where the magic trick metaphor will finally
strain to its breaking point.

\chapter{Chapter 5: Encoding the
Performance}\label{chapter-5-encoding-the-performance}

\section{Layer 4 -- Arithmetization}\label{layer-4-arithmetization}

In the previous chapter, we watched the magician go backstage and produce a
detailed recording of every step in the computation -- the witness. That
recording is private. Now comes the hardest transformation in the entire stack:
converting that recording into a form that can be checked mathematically,
without re-doing the computation, and without revealing the private data.

Arithmetization: the art of turning a computer program into a system of
polynomial equations.

If the previous chapter was about what the magician does backstage, this chapter
is about the notation system used to write down what happened. The notation must
be precise enough to catch any error, compact enough to be checked quickly, and
structured enough to reveal nothing about the performance except its
correctness. Finding such a notation -- and making it efficient enough for
practical use -- has been the central technical challenge of the zero-knowledge
field for the past decade.

\par\vspace{18pt}

\emph{The Sudoku analogy implies a unique solution. Does a ZK proof have one
solution or many?}

The answer is: many. There are typically many valid witnesses for a given public
statement. If you are proving you know a number whose square is 25, both 5 and
-5 work. If you are proving you have a valid passport, any valid passport will
do. The Sudoku comparison, popular in introductory ZK writing, misleads
precisely because it implies uniqueness -- one grid, one solution, one truth. A
zero-knowledge proof is closer to proving you hold \emph{a} winning lottery
ticket without showing which one. The distinction matters because the entire
machinery of this chapter exists to handle a richer, messier reality than any
single-solution puzzle can capture.

\par\vspace{18pt}

Let us be honest at the outset: this is where the magic trick analogy strains
hardest. A sealed scorecard, a crossword puzzle, a spreadsheet with rules --
every metaphor we reach for captures one aspect and distorts another. So we will
do what Feynman recommended when analogies fail: state what is actually
happening in plain language, and trust the reader to follow.

This chapter is longer and more technical than the others. That is because
arithmetization is where the conceptual rubber meets the mathematical road. The
ideas here -- constraint systems, polynomial identities, lookup arguments, the
sumcheck protocol -- are the load-bearing structures of every ZK system in
existence. A reader who understands this chapter understands why zero-knowledge
proofs work. A reader who skips it must take the rest of the book on faith.

What is actually happening is this. The computation -- every addition, every
comparison, every memory access -- gets encoded as relationships between numbers
in a finite field. These relationships take the form of polynomial equations. If
the computation was performed correctly, all the equations are satisfied
simultaneously. If the prover cheated at any step, at least one equation is
violated. And here is the key insight that makes the entire field of
zero-knowledge proofs possible: checking whether all these polynomial equations
hold can be done by evaluating them at a few random points, which is vastly
faster than re-executing the original computation.

This chapter tells the story of how the encoding schemes evolved, from the rigid
first attempts to the unified framework that powers every modern proof system.
It is also, unavoidably, a story about the overhead this encoding imposes -- and
whether that overhead is an immutable tax or a temporary engineering constraint.

The story has five acts. First, we establish the spreadsheet metaphor that makes
constraint systems intuitive. Second, we trace the evolution from R1CS to AIR to
PLONKish, with concrete worked examples showing how each system encodes
computation differently. Third, we encounter CCS -- the unifying grammar that
reveals all three systems as dialects of the same language -- and the sumcheck
protocol that powers verification. Fourth, we follow the lookup revolution from
Plookup through Jolt, watching as table lookups replace polynomial constraints
as the primary computation paradigm. Fifth, we confront the overhead tax
honestly, with concrete numbers showing what the encoding costs in practice and
where those costs are falling.

\section{The Spreadsheet Metaphor (And Where It
Works)}\label{the-spreadsheet-metaphor-and-where-it-works}

Before we encounter the formal constraint systems, we need a mental model. The
best available one is the spreadsheet.

The spreadsheet metaphor is not perfect -- we will say where it breaks down --
but it is the most productive starting point. Every major constraint system
(R1CS, AIR, PLONKish, CCS) can be understood as a particular way of organizing a
spreadsheet and writing rules for its cells. The differences between the systems
are differences in how the rules are structured, not in the underlying idea.

Imagine a computation with a thousand steps. You create a giant table. Each row
represents one moment in time -- one step of the computation. Each column
represents a variable: a processor register, a memory value, a boolean flag.
Every cell contains a number drawn from a finite field (think: integers modulo a
large prime).

Consider a tiny computation: \(x = 3\), \(y = 4\), \(z = x + y = 7\), then
\(w = z \times 2 = 14\).

\begin{longtable}[]{@{}lllll@{}}
\toprule\noalign{}
Row & A (input 1) & B (input 2) & C (result) & Rule \\
\midrule\noalign{}
\arrayrulecolor{MidnightBlue}\midrule\arrayrulecolor{MidnightBorder}
\endhead
\bottomrule\noalign{}
\endlastfoot
1 & 3 & 4 & 7 & C = A + B \\
2 & 7 & 2 & 14 & C = A * B \\
\end{longtable}

If every rule holds across every row, the spreadsheet faithfully records the
computation. Change any cell, and at least one rule breaks. Suppose a cheating
prover changes the result in row 1 from 7 to 9. Row 1 now violates its own rule
(3 + 4 is not 9), and row 2 also breaks (because row 2 expects to read 7 from
row 1's output, not 9). Errors propagate. This is a two-row example, but the
principle scales to millions of rows: one wrong cell poisons the entire
spreadsheet.

Now you write rules. ``The value in column B at row 5 must equal the value in
column A at row 4 plus the value in column C at row 4.'' ``If the opcode column
at row 7 says `multiply,' then column D at row 7 must equal column B at row 7
times column C at row 7.'' These rules are polynomial equations that relate
cells to one another.

Notice that the rules are not arbitrary. They are polynomial equations --
expressions built from addition, subtraction, and multiplication of cell values.
This restriction is fundamental. A polynomial rule like ``\(A \cdot B = C\)'' is
checkable by the proof system. A non-polynomial rule like ``if A \textgreater{}
B then C = 1 else C = 0'' cannot be directly encoded as a polynomial equation
because comparison is not a polynomial operation. (It can be encoded
\emph{indirectly}, by decomposing A and B into bits and constraining the
bit-level comparison, but this adds many auxiliary constraints.) The polynomial
restriction is the price of admittance to the proof system. Only relationships
expressible as polynomial equations over finite fields can be directly verified.
Everything else must be translated into polynomial form first.

If every rule holds across every row, the spreadsheet is \emph{consistent} -- it
faithfully records a valid computation. If any rule is violated, the computation
was not performed correctly. The prover's job is to fill in the spreadsheet
(this is the witness from the previous chapter) and then convince the verifier
that all the rules hold. The verifier's job is to check -- but not by examining
every cell. Instead, the verifier picks random evaluation points and checks
whether the polynomial equations are satisfied there. By the Schwartz-Zippel
lemma, a polynomial that is not identically zero will be nonzero at a random
point with high probability. The Schwartz-Zippel lemma is the mathematical fact
that makes this work: a nonzero polynomial of degree \(d\), evaluated at a
random point from a field of size \(q\), is zero with probability at most
\(d/q\). For the fields used in ZK (where \(q\) is astronomically large), this
probability is negligible. One random check is almost as good as checking
everywhere. So if the equations check out at the random points, the spreadsheet
is almost certainly correct everywhere.

That is the core idea of arithmetization. Every constraint system in this
chapter is a different way of organizing the spreadsheet, choosing the rules,
and encoding the computation.

\begin{quote}
\textbf{The Running Example: The Sudoku Constraints}

Our Sudoku witness becomes a 16-row constraint system. Each cell must satisfy:

\begin{itemize}
\tightlist
\item
  \textbf{Range constraint}:
  \((\text{cell} - 1)(\text{cell} - 2)(\text{cell} - 3)(\text{cell} - 4) = 0\).
  This polynomial evaluates to zero only when the cell contains a valid value.
  Four values, one degree-\(4\) polynomial per cell.
\item
  \textbf{Given-cell constraint}: For each clue,
  \(\text{cell}_i = \text{given}_i\). Eight equalities for our 8-given puzzle.
\item
  \textbf{Uniqueness constraint}: For each row, column, and 2x2 box, the product
  \((a - b)\) for all pairs must be nonzero. Equivalently: the polynomial
  product over all pairs of \((a - b)\) must be nonzero for each group. Eight
  groups, \(\binom{4}{2} = 6\) pairs each, yielding 48 pair checks.
\end{itemize}

Total: 16 range constraints + 8 given-cell constraints + 48 uniqueness checks =
72 constraints over 16 witness variables. In R1CS form, each degree-\(4\) range
constraint decomposes into intermediate multiplications, expanding to roughly
120 R1CS constraints. In CCS form, the higher-degree constraints can be
expressed directly. The witness (the completed grid) satisfies all 72
constraints. A wrong value in any cell makes at least one polynomial nonzero,
and the Schwartz-Zippel lemma catches it with overwhelming probability at a
random evaluation point.
\end{quote}

Why polynomials? Because of a fact about polynomials: a polynomial of degree
\(d\) is completely determined by its values at any \(d+1\) points. If you know
a line (degree \(1\)), two points fix it exactly. If you know a cubic (degree
\(3\)), four points fix it exactly. This means that if a polynomial
``misbehaves'' at even a single point, it must be the wrong polynomial -- and
checking it at a random point catches this misbehavior with near certainty. A
polynomial commitment scheme exploits this: the prover seals a polynomial into a
short commitment, and the verifier can spot-check it at random points to confirm
it is correct -- without ever seeing the full polynomial.

The metaphor is imperfect in one important respect: a real spreadsheet has rows
and columns with human-readable labels. A constraint system is a set of abstract
polynomial equations over vectors of field elements. The ``rows'' and
``columns'' are a convenient way to think about structure, but the mathematics
does not require a rectangular layout. Keep this in mind as we move through the
specific systems.

With the spreadsheet image in hand, we can state the central question of this
chapter: What is the best way to organize the rules? Should each row have its
own custom rule (like R1CS)? Should all rows share the same rule (like AIR)?
Should rows have switchable rules controlled by flags (like PLONKish)? Or should
all these approaches be unified under a single framework (like CCS)? The history
of arithmetization is the history of answering this question, and the answer
keeps changing as proof systems evolve and new mathematical tools become
available. The constraint systems in the next sections are not mere notation.
They are architectural decisions that determine the performance, flexibility,
and security of every zero-knowledge proof system built on top of them.

\section{The Constraint System Evolution: R1CS, AIR,
PLONKish}\label{the-constraint-system-evolution-r1cs-air-plonkish}

The three major constraint systems -- R1CS, AIR, and PLONKish -- emerged in a
span of just seven years (2012-2019). Each was designed to solve a specific
limitation of its predecessor, but each also introduced new trade-offs.
Understanding this evolution is not optional for understanding modern ZK: every
proof system, every zkVM, and every privacy protocol in production today is
built on one of these three foundations (or, increasingly, on CCS, which unifies
all three).

The history of arithmetization is a history of increasing expressiveness. Each
new constraint system solved a specific limitation of its predecessor.
Understanding this genealogy is essential because the constraint system you
choose determines which proof systems you can use, which fields are efficient,
and how much overhead the encoding imposes.

The genealogy also reveals how rapidly the field moves. R1CS was introduced in
2012. By 2023, it was already the ``legacy'' format -- still deployed in
production (Groth16 is not going away), but superseded by more expressive
systems for new development. The eleven-year span from R1CS to CCS saw more
architectural innovation in constraint system design than the previous four
decades of theoretical computer science produced. This acceleration is driven by
practical pressure: the economic value of efficient zero-knowledge proofs
creates strong incentives for better arithmetization.

\subsection{R1CS: The Assembly Language
(2012)}\label{r1cs-the-assembly-language-2012}

The first practical arithmetization emerged from the work of Gennaro, Gentry,
Parno, and Raykova (GGPR) in 2012, who introduced the QAP (Quadratic Arithmetic
Program) framework that R1CS later formalized. The name ``Rank-1 Constraint
System'' describes the mathematical structure precisely: each constraint has
rank 1 (it is the product of two linear functions), and the system is a
collection of such constraints. The ``rank-1'' designation means each constraint
captures exactly one multiplication -- a bilinear relationship between
variables. Addition is free (it does not require a constraint, because linear
combinations can be folded into the matrix entries). Only multiplication
generates constraints. This is why the number of R1CS constraints for a circuit
equals the number of multiplication gates, not the total number of gates.

Before the mathematical notation, let us ground it in the spreadsheet from the
previous section. Imagine your spreadsheet has three special columns -- call
them A, B, and C. For each row, column A and column B each contain a combination
of the variables, and column C contains the result. The rule for every row is:
(what is in column A) multiplied by (what is in column B) must equal (what is in
column C). That is all R1CS is: a spreadsheet where every row enforces one
multiplication rule. The mathematical notation below says exactly this, just
more precisely.

R1CS encodes a computation as a list of constraints, each of the form:

\emph{(linear combination of variables) times (linear combination of variables)
equals (linear combination of variables)}

Or, in mathematical notation:
\((\mathbf{A} \cdot \mathbf{z}) \circ (\mathbf{B} \cdot \mathbf{z}) = \mathbf{C} \cdot \mathbf{z}\),
where \(\mathbf{A}\), \(\mathbf{B}\), and \(\mathbf{C}\) are sparse matrices and
\(\mathbf{z}\) is the vector of all variables (public inputs, private witness,
and intermediate values). Each row of the matrices defines one constraint. Each
constraint captures one multiplication gate.

For the tiny spreadsheet example (3 + 4 = 7, then 7 * 2 = 14), the witness
vector is \(\mathbf{z} = (1, 3, 4, 7, 2, 14)\) -- the constant 1 followed by the
variables x, y, z, w, and the final result. The first row of A selects ``x''
(entry 1 in position corresponding to x), the first row of B selects ``1''
(indicating the addition is encoded as (x + y) * 1 = z after reformulation), and
the first row of C selects ``z''. The second row of A selects ``z'' (value 7),
the second row of B selects ``w'' (value 2), and the second row of C selects
``result'' (value 14). The matrices are mostly zeros -- only a handful of
entries are nonzero. This sparsity is typical: a circuit with millions of
constraints has matrices with millions of rows but only a few nonzero entries
per row.

R1CS is the assembly language of constraint systems -- simple, well-understood,
and directly amenable to proof systems like Groth16 and Spartan. Groth16,
deployed across most of the Ethereum ecosystem, works natively with R1CS and
produces the smallest possible proofs -- three elliptic curve group elements,
constant-time verification.

But R1CS has a fundamental limitation: each constraint is bilinear -- degree
\(2\). You can encode a multiplication (\(a \cdot b = c\)) directly. But what
about a computation that requires checking a hash function, where a single
invocation might need thousands of multiplications? You encode each
multiplication as a separate constraint, one after another. There is no way to
express a higher-degree relationship in a single constraint, and there is no
notion of ``this constraint applies uniformly across all time steps.'' Every
gate gets its own row.

For small circuits, R1CS works beautifully. For large, repetitive computations
-- like executing millions of processor instructions in a zkVM -- the lack of
structure becomes a liability.

To see this concretely, consider encoding a computation with 10 million
multiplication gates in R1CS. You need 10 million rows in the matrices A, B, and
C. Each matrix is sparse (most entries are zero), but the total number of
nonzero entries -- and hence the prover's work -- scales linearly with the gate
count. There is no compression possible: R1CS treats each gate independently,
with no awareness that gate 5,000,001 might be performing exactly the same
operation as gate 1. Compare this to a function that repeats the same 1,000-gate
circuit 10,000 times: R1CS still requires 10,000,000 separate constraints, while
a structured constraint system could potentially describe the repeating pattern
once and instantiate it 10,000 times. This structural blindness is what
motivated the search for richer constraint formats.

Yet R1CS persists. Groth16 proofs (which require R1CS) remain the gold standard
for on-chain verification because of their unmatched proof size: 3 group
elements, roughly 128 bytes, verifiable in constant time. No other proof system
achieves this compactness. When Ethereum smart contracts verify ZK proofs, the
gas cost of verification is proportional to proof size -- and Groth16's tiny
proofs mean minimal gas costs. Many production systems (Zcash, Tornado Cash,
Worldcoin) use Groth16 for precisely this reason, accepting the constraint
system's limitations in exchange for the proof system's efficiency. R1CS is the
assembly language: nobody wants to write it directly, but the machine code it
produces is unbeatable.

\subsection{AIR: The State Machine (2018)}\label{air-the-state-machine-2018}

The Algebraic Intermediate Representation arrived alongside STARKs, introduced
by Ben-Sasson, Bentov, Horesh, and Riabzev in 2018. The name is precise:
``Algebraic'' because the constraints are polynomial equations over fields (not
boolean circuits or SAT formulas). ``Intermediate'' because AIR sits between the
high-level computation and the low-level proof system -- it is the ``compiled''
form of the computation, analogous to LLVM IR in a compiler toolchain.
``Representation'' because it is a way of representing the computation, not the
computation itself. AIR solves the structure problem by embracing the
spreadsheet metaphor directly.

An AIR consists of two things: an execution trace (a 2D matrix where rows are
time steps and columns are algebraic registers) and transition constraints
(polynomial equations that must hold between consecutive rows). The constraints
are \emph{uniform} -- the same polynomial equations apply at every row.

This uniformity is AIR's great strength and its great limitation. It is a
perfect fit for sequential computations where the same operation repeats: hash
chains, state machine execution, virtual machine instruction cycles. The
constraint ``if the opcode is ADD, then the output register equals the sum of
the two input registers'' applies identically at every step. You write it once;
it is enforced everywhere.

The word ``uniform'' deserves emphasis. In R1CS, each constraint row can encode
a completely different relationship between variables. Row 1 might enforce
\(a + b = c\); row 2 might enforce \(d \cdot e = f\); row 3 might enforce
\(g = h\). Each row is independent. In AIR, every row obeys the same set of
transition polynomials. If the transition constraint says
``\(\text{column\_3}[\text{next}] = \text{column\_1}[\text{current}] + \text{column\_2}[\text{current}]\),''
then this relationship holds at every pair of consecutive rows. The prover
cannot make exceptions. This rigidity is what enables compression: a single
polynomial equation describes the entire computation, regardless of how many
steps it contains.

The limitation is that AIR cannot natively handle non-uniform computation. If
your program has different instruction types -- additions, multiplications, hash
invocations, memory accesses -- you need tricks to encode the selection logic
(``which instruction is executing at this row?'') within the uniform framework.
This is possible but adds complexity and overhead.

In practice, real STARK-based systems handle non-uniformity by using multiple
AIR traces -- one per ``sub-machine.'' Cairo's architecture, for example,
decomposes the VM into separate traces for the CPU, memory, range checks, and
each built-in operation (Pedersen hash, ECDSA, bitwise operations). Each trace
is a separate AIR with its own transition constraints. Cross-trace consistency
is enforced through permutation arguments and lookup arguments that connect the
traces: when the CPU trace records ``hash instruction at step 1000,'' the hash
trace must contain a corresponding row with matching inputs and outputs. This
multi-trace design preserves AIR's uniformity within each trace while handling
the non-uniformity of a full instruction set across traces. The cost is the
cross-trace connection overhead, but for computations dominated by a single
operation type (as hash-heavy applications often are), the overhead is
manageable.

AIR became the foundation of the STARK ecosystem: StarkWare's Stone and Stwo
provers, Polygon Miden, and others. Its tight coupling with FRI (the hash-based
polynomial commitment scheme) means AIR-based systems are transparent (no
trusted setup) and plausibly post-quantum secure.

The AIR-FRI coupling deserves a moment of attention because it illustrates how
Layers 4 and 5 (arithmetization and proof system) become inseparable. FRI works
by repeatedly folding a polynomial in half -- reducing its degree by a factor of
2 at each step -- and checking consistency at random points. This folding
requires the polynomial to be evaluated on a domain with a specific
multiplicative structure (a ``coset'' of a subgroup of the field). AIR traces
are naturally expressed as polynomials on such domains because the trace rows
correspond to consecutive powers of a group generator. The uniformity of AIR's
transition constraints means the constraint polynomial has the same degree
structure as the trace polynomial, which is exactly what FRI needs. Try to use
FRI with a non-uniform constraint system (like PLONKish), and you need
additional machinery (permutation polynomials, selector commitments) that adds
overhead. AIR and FRI were born for each other.

\subsubsection{A Tiny AIR: The Counter}\label{a-tiny-air-the-counter}

To make AIR concrete, consider the simplest possible state machine: a counter
that starts at 0 and increments by 1 each step. The execution trace has two
columns -- \texttt{counter} (the current value) and \texttt{flag} (whether to
increment) -- and three rows:

\begin{longtable}[]{@{}lll@{}}
\toprule\noalign{}
Row & counter & flag \\
\midrule\noalign{}
\arrayrulecolor{MidnightBlue}\midrule\arrayrulecolor{MidnightBorder}
\endhead
\bottomrule\noalign{}
\endlastfoot
0 & 0 & 1 \\
1 & 1 & 1 \\
2 & 2 & 1 \\
\end{longtable}

The transition constraint is a single polynomial equation that must hold between
every pair of consecutive rows:

\(\text{counter}[i+1] = \text{counter}[i] + \text{flag}[i]\)

Check it. Between rows 0 and 1: does 1 = 0 + 1? Yes. Between rows 1 and 2: does
2 = 1 + 1? Yes. The trace is valid.

Now suppose a cheating prover submits a trace where row 2 claims
\texttt{counter\ =\ 5}. The transition constraint between rows 1 and 2 becomes:
does 5 = 1 + 1? No.~The constraint is violated, and the proof fails.

There is a subtlety that the transition constraint alone does not capture: the
starting value. The transition constraint says ``each row follows correctly from
the previous row,'' but it says nothing about where the counter begins. A trace
starting at counter = 1000 with flag = 1 at every row would satisfy the
transition constraint perfectly -- 1000, 1001, 1002 -- even though the counter
was supposed to start at 0. This is where \textbf{boundary constraints} enter. A
boundary constraint pins a specific cell to a specific value: ``counter at row 0
must equal 0.'' In an AIR, you typically have transition constraints (which
apply uniformly between consecutive rows) and boundary constraints (which apply
at specific rows, usually the first and last). The transition constraints ensure
the computation proceeds correctly; the boundary constraints ensure it starts
and ends in the right place.

The full AIR for this counter is therefore:

\begin{itemize}
\tightlist
\item
  Transition constraint:
  \(\text{counter}[i+1] = \text{counter}[i] + \text{flag}[i]\) (for all
  consecutive row pairs)
\item
  Boundary constraint: \(\text{counter}[0] = 0\) (the counter starts at zero)
\item
  Boundary constraint: \(\text{flag}[i] \cdot (1 - \text{flag}[i]) = 0\) (each
  flag is boolean -- either 0 or 1)
\end{itemize}

The boolean constraint on the flag deserves attention. Without it, a cheating
prover could set flag = 7 at some row, making the counter jump by 7 instead of 0
or 1. The constraint \(\text{flag} \cdot (1 - \text{flag}) = 0\) is satisfied
only when flag is 0 or 1 (plug in either value and one factor is zero). This is
a polynomial constraint of degree \(2\) -- exactly the kind of equation that AIR
handles naturally.

The critical observation is what the prover \emph{wrote down} to define this
constraint system. Not three separate rules -- one for each row -- but a small
set of polynomial equations applied uniformly across the entire trace. One
transition rule and a few boundary conditions, enforced everywhere. If the
counter had a million rows instead of three, the constraint description would be
identical: the same equations. Only the trace grows; the constraints stay fixed.

This is the structural difference from R1CS. In R1CS, you would write a separate
constraint for each row: ``row 0's output equals row 0's input plus row 0's
flag,'' then ``row 1's output equals row 1's input plus row 1's flag,'' and so
on -- one constraint per step. For a million-step computation, you need a
million constraints. In AIR, you write the rule once. The prover fills in the
trace; the polynomial machinery checks the rule everywhere simultaneously.

For repetitive computations -- hash chains where the same compression function
executes thousands of times, virtual machines where the same instruction cycle
repeats for every step -- AIR's uniformity is not just convenient. It is a
compression of the constraint description itself, from linear in the number of
steps to constant in the number of distinct transition rules. That compression
is what made STARKs practical for proving large computations.

One further detail illuminates how the polynomial machinery works behind the
scenes. The prover does not submit the raw trace table to the verifier. Instead,
the prover interpolates each column of the trace as a polynomial. For the
counter column with values (0, 1, 2), the prover finds a polynomial \(P(x)\)
such that \(P(0) = 0\), \(P(1) = 1\), \(P(2) = 2\) -- in this case, simply
\(P(x) = x\). For the flag column with values (1, 1, 1), the polynomial is
\(F(x) = 1\). The transition constraint ``\(P(x+1) = P(x) + F(x)\)'' becomes a
polynomial identity that must hold at \(x = 0\) and \(x = 1\) (every pair of
consecutive rows). The proof system checks this identity at a random evaluation
point -- not at \(x = 0\) or \(x = 1\), but at some random \(r\) chosen by the
verifier -- and the Schwartz-Zippel lemma guarantees that a false identity will
fail this random check with overwhelming probability.

The trace, the constraints, and the verification all reduce to polynomials. The
trace columns are polynomials. The transition constraint is a polynomial
identity. The verification is a polynomial evaluation. This is what
``arithmetization'' means in practice: every aspect of the computation becomes a
polynomial, and every check becomes a polynomial evaluation.

The polynomial encoding also reveals the source of the overhead. The counter
trace has 3 rows and 2 columns -- 6 values. But the polynomials that interpolate
these columns have degree \(2\) (you need a degree-\(2\) polynomial to pass
through 3 points in general). The transition constraint, when expressed as a
polynomial, produces a ``constraint polynomial'' whose degree is the sum of the
degrees of the trace polynomials it involves. If the trace polynomial has degree
\(d\) and the transition constraint has algebraic degree \(k\), the constraint
polynomial has degree roughly \(d \cdot k\). For a trace with \(n\) rows, \(d\)
is roughly \(n\), so the constraint polynomial has degree roughly \(n \cdot k\).
This polynomial must be shown to vanish on all consecutive-row pairs, which
means it is divisible by a ``vanishing polynomial'' \(Z(x)\) that has roots at
the evaluation domain. The quotient \(T(x) = \text{constraint}(x) / Z(x)\) is
the polynomial the prover commits to; if the division is exact (no remainder),
the constraints are satisfied. If the prover cheated, the division leaves a
remainder, and the random evaluation check catches it.

This is where the NTTs come in. Converting between coefficient and evaluation
representations of these high-degree polynomials requires the Number Theoretic
Transform -- the finite-field version of the FFT -- which dominates the prover's
computation time.

The limitation is equally visible. Suppose you want some rows to add and other
rows to multiply. With a single uniform constraint, you cannot express ``do
addition at row 5 and multiplication at row 6'' without encoding the selection
logic into the polynomial itself -- adding flag columns, conditional terms, and
degree overhead. The counter example is clean because every row does the same
thing. Real programs do not.

\subsection{PLONKish: The Custom Workshop
(2019)}\label{plonkish-the-custom-workshop-2019}

PLONK, introduced by Gabizon, Williamson, and Ciobotaru in 2019, took a
different approach. Instead of uniform constraints, PLONKish arithmetization
uses \emph{selector columns} to enable non-uniform gates.

The key innovation separates the constraint system into two components:

\begin{enumerate}
\def\labelenumi{\arabic{enumi}.}
\item
  \textbf{Gate constraints}: polynomial equations controlled by selector
  polynomials. Different rows can have different gate types. If the selector for
  ``addition'' is active at row 5, the addition constraint is enforced there. If
  the selector for ``multiplication'' is active at row 6, the multiplication
  constraint is enforced there. You can define custom gates for any operation
  you need.
\item
  \textbf{Copy constraints}: a permutation argument that enforces wiring --
  ensuring that the output of one gate is correctly fed as input to another
  gate. This replaces R1CS's matrix-based variable assignment with a more
  flexible connection mechanism.
\end{enumerate}

PLONKish sits between R1CS and AIR in expressiveness. Like AIR, it uses a
structured trace with rows and columns. Unlike AIR, different rows can follow
different rules. Like R1CS, it can handle arbitrary circuits. Unlike R1CS, it
supports custom gates that capture complex operations in fewer constraints.

PLONKish became the dominant arithmetization in deployed systems. Halo2 (used by
Zcash and Scroll), Polygon zkEVM (before its shutdown), and numerous other
production systems chose PLONKish because its flexibility handles the diverse
instruction sets of real-world computations. The Halo2 library, originally
developed by the Electric Coin Company for the Zcash Orchard protocol, became
the de facto standard for PLONKish circuit development. Its ``region-based'' API
lets developers define gates, assign cells, and specify copy constraints in a
structured way that catches many common errors at compile time. Scroll's zkEVM
-- one of the most ambitious ZK projects ever attempted -- encoded the entire
Ethereum Virtual Machine instruction set as Halo2 PLONKish circuits, using
custom gates for EVM opcodes, lookup arguments for bytecode verification, and
copy constraints to wire the data path. The resulting circuit has millions of
constraints per block and requires GPU clusters to prove, but it works -- which
says something about PLONKish's flexibility.

\subsubsection{A Tiny PLONKish Circuit: Compute 3 + 4, Then Multiply, Then Add
Again}\label{a-tiny-plonkish-circuit-compute-3-4-then-multiply-then-add-again}

Here is a three-row PLONKish trace that computes (3 + 4) * 2 + 1 = 15. The trace
has three ``witness'' columns (a, b, c) and two ``selector'' columns (q\_add,
q\_mul):

\begin{longtable}[]{@{}llllll@{}}
\toprule\noalign{}
Row & a & b & c & q\_add & q\_mul \\
\midrule\noalign{}
\arrayrulecolor{MidnightBlue}\midrule\arrayrulecolor{MidnightBorder}
\endhead
\bottomrule\noalign{}
\endlastfoot
0 & 3 & 4 & 7 & 1 & 0 \\
1 & 7 & 2 & 14 & 0 & 1 \\
2 & 14 & 1 & 15 & 1 & 0 \\
\end{longtable}

The gate constraint is a single equation evaluated at every row:

\(q_{\text{add}} \cdot (a + b - c) + q_{\text{mul}} \cdot (a \cdot b - c) = 0\)

At row 0: \(q_{\text{add}} = 1\), \(q_{\text{mul}} = 0\), so the equation
becomes \(1 \cdot (3 + 4 - 7) + 0 \cdot (\ldots) = 0\). Check: \(0 = 0\). The
addition gate is active.

At row 1: \(q_{\text{add}} = 0\), \(q_{\text{mul}} = 1\), so the equation
becomes \(0 \cdot (\ldots) + 1 \cdot (7 \cdot 2 - 14) = 0\). Check: \(0 = 0\).
The multiplication gate is active.

At row 2: \(q_{\text{add}} = 1\), \(q_{\text{mul}} = 0\), so the equation
becomes \(1 \cdot (14 + 1 - 15) + 0 \cdot (\ldots) = 0\). Check: \(0 = 0\). The
addition gate is active again.

The selectors act as switches. When \(q_{\text{add}} = 1\) and
\(q_{\text{mul}} = 0\), only the addition constraint is ``on.'' When the
selectors flip, only the multiplication constraint is ``on.'' One polynomial
equation, evaluated identically at every row, enforces different gate types
depending on which selector is active.

But the gate constraint alone does not guarantee correctness. Look at the trace:
row 0 produces c = 7, and row 1 consumes a = 7. How does the proof system know
these two 7s are the \emph{same} value -- that the output of row 0 actually
flows into the input of row 1?

This is the job of the \textbf{copy constraint}. A permutation argument -- a
separate cryptographic mechanism outside the gate equation -- enforces that the
cell (row 0, column c) contains the same value as the cell (row 1, column a).
Similarly, (row 1, column c) must equal (row 2, column a). The permutation
argument works by proving that a certain set of values is a rearrangement of
another set, which can only be true if the ``wired'' cells agree. Without copy
constraints, a cheating prover could fill row 1 with a = 999 and the gate
equation would still pass (as long as 999 * 2 = c at row 1). The copy constraint
is what stitches the circuit together.

To see the copy constraint at work, consider what happens without it. A cheating
prover submits this trace:

\begin{longtable}[]{@{}llllll@{}}
\toprule\noalign{}
Row & a & b & c & q\_add & q\_mul \\
\midrule\noalign{}
\arrayrulecolor{MidnightBlue}\midrule\arrayrulecolor{MidnightBorder}
\endhead
\bottomrule\noalign{}
\endlastfoot
0 & 3 & 4 & 7 & 1 & 0 \\
1 & 99 & 2 & 198 & 0 & 1 \\
2 & 198 & 1 & 199 & 1 & 0 \\
\end{longtable}

Every gate constraint passes: 3 + 4 = 7, 99 * 2 = 198, 198 + 1 = 199. But the
computation is wrong -- row 1 should have used a = 7 (the output of row 0), not
a = 99. Without the copy constraint binding cell (row 0, c) to cell (row 1, a),
the prover is free to insert any value it likes. The copy constraint catches
this: it requires that position (row 0, column c) and position (row 1, column a)
hold the same value. Since 7 is not 99, the permutation check fails, and the
proof is rejected.

This reveals why PLONKish is more flexible than AIR. In the AIR counter example,
every row obeyed the same transition rule. Here, row 0 adds, row 1 multiplies,
and row 2 adds again -- three different operations in three rows, controlled by
selector values. You can define custom gates for any operation: a ``range
check'' gate, a ``Poseidon hash round'' gate, an ``elliptic curve addition''
gate. Each gets its own selector column, and the prover activates whichever gate
the computation requires at each row. The trace is a heterogeneous computation
log, not a uniform state machine.

The power of custom gates becomes clearer with a slightly more complex example.
Suppose you want to enforce that a value lies in the range \([0, 255]\) -- an
8-bit range check. In R1CS, you would decompose the value into 8 bits, constrain
each bit to be boolean (8 constraints), and constrain the sum to equal the
original value (1 constraint) -- 9 constraints total. In PLONKish, you can
define a single custom ``range gate'' that encodes the entire check in one row,
using a lookup argument or a specialized polynomial identity. One row, one gate,
one constraint. The circuit designer creates the gate once; the prover activates
it wherever a range check is needed.

The cost of this flexibility is the copy constraint machinery. The permutation
argument adds overhead -- both in proof size and in prover computation -- that
AIR avoids because AIR's uniform structure implicitly handles data flow between
consecutive rows. But for computations that mix many different operations (as
real programs do), the overhead is worth paying.

\subsubsection{The Same Computation, Three
Encodings}\label{the-same-computation-three-encodings}

To crystallize the differences, consider encoding the same simple computation --
``compute \(x \cdot (x + 1)\) where \(x = 3\), so the result is \(12\)'' -- in
all three systems.

\textbf{In R1CS:} You need two constraints. First, an addition: an intermediate
variable \(t = x + 1 = 4\). Then a multiplication:
\(\text{result} = x \cdot t = 3 \cdot 4 = 12\). Each constraint takes one row in
the R1CS matrix. The matrices A, B, C encode the variable wiring. Two rows, two
constraints, done.

\textbf{In AIR:} You set up a two-row trace. Row 0 holds \(x = 3\) and computes
\(t = x + 1 = 4\). Row 1 holds the multiplication
\(\text{result} = x \cdot t = 12\). The transition constraint relates
consecutive rows. But here is the awkwardness: the addition and the
multiplication are \emph{different} operations, and AIR wants uniform
constraints across all rows. You either need to encode both operations into a
single transition polynomial (using conditional logic with flag columns, which
increases the constraint degree), or you define a two-step cycle where even rows
add and odd rows multiply (which works but means half the trace structure is
``wasted'' on selection logic). For this tiny example, AIR is overkill.

\textbf{In PLONKish:} Row 0 uses the addition gate: \(a = 3\), \(b = 1\),
\(c = 4\) (\(q_{\text{add}} = 1\)). Row 1 uses the multiplication gate:
\(a = 3\), \(b = 4\), \(c = 12\) (\(q_{\text{mul}} = 1\)). A copy constraint
links (row 0, column c) to (row 1, column b), and another links the input
\(x = 3\) to (row 1, column a). Two rows, two different gate types, clean and
direct.

The comparison reveals each system's natural habitat. R1CS handles this
computation most directly -- two bilinear constraints, no overhead. PLONKish
handles it almost as directly, with slight overhead from the copy constraints.
AIR handles it least naturally, because the computation is not repetitive --
there is no pattern that repeats across many rows. For a computation that
\emph{is} repetitive (running the same hash compression 1000 times), the ranking
reverses: AIR wins by writing the constraint once, while R1CS and PLONKish must
either repeat the constraint description 1000 times or use recursion to simulate
repetition.

The constraint count for the same computation across different systems is
instructive:

\begin{longtable}[]{@{}
  >{\raggedright\arraybackslash}p{(\linewidth - 6\tabcolsep) * \real{0.1159}}
  >{\raggedright\arraybackslash}p{(\linewidth - 6\tabcolsep) * \real{0.3478}}
  >{\raggedright\arraybackslash}p{(\linewidth - 6\tabcolsep) * \real{0.4638}}
  >{\raggedright\arraybackslash}p{(\linewidth - 6\tabcolsep) * \real{0.0725}}@{}}
\toprule\noalign{}
\begin{minipage}[b]{\linewidth}\raggedright
System
\end{minipage} & \begin{minipage}[b]{\linewidth}\raggedright
Constraints for x*(x+1)
\end{minipage} & \begin{minipage}[b]{\linewidth}\raggedright
Constraints for 1000 hash rounds
\end{minipage} & \begin{minipage}[b]{\linewidth}\raggedright
Why
\end{minipage} \\
\midrule\noalign{}
\arrayrulecolor{MidnightBlue}\midrule\arrayrulecolor{MidnightBorder}
\endhead
\bottomrule\noalign{}
\endlastfoot
R1CS & 2 & \textasciitilde30,000,000 & 1 constraint per gate,
\textasciitilde30,000 per hash round \\
AIR & \textasciitilde4 (with padding) & \textasciitilde30,000 & One transition
polynomial, reused 1000 times \\
PLONKish & 2 (+ copy constraints) & \textasciitilde15,000,000 & Custom hash
gates cut per-round cost in half \\
\end{longtable}

The numbers are approximate, but the ratios are revealing. For the tiny
computation, all three systems are roughly comparable. For the hash chain, AIR's
constraint count is independent of the number of repetitions -- it depends only
on the number of distinct transition types. This is why STARKs dominate in
hash-heavy workloads (blockchain state verification, recursive proof
composition) while PLONKish dominates in mixed workloads (smart contract
execution, general-purpose circuits).

A clarification for the precise reader: AIR's ``\textasciitilde30,000'' for 1000
hash rounds refers to the \emph{constraint description} size -- the number of
distinct polynomial equations that must hold. The actual \emph{trace} still has
1000 * (rows per hash round) rows, each of which must satisfy the constraints.
The prover's work is proportional to the trace size, not the constraint
description size. But the constraint description size matters for the verifier
(who must check the polynomial identity, not every row) and for the proof size
(which depends on the degree of the constraint polynomial, not the number of
trace rows). The asymmetry between ``small constraint description, large trace''
is precisely what makes AIR efficient for repetitive computations: the
verifier's work grows with the constraint complexity, not with the number of
repetitions.

\subsection{Three Dialects, One Problem}\label{three-dialects-one-problem}

By 2022, the ZK ecosystem had three constraint system families, each with its
own proof systems, tooling, and community:

\begin{longtable}[]{@{}
  >{\raggedright\arraybackslash}p{(\linewidth - 8\tabcolsep) * \real{0.1250}}
  >{\raggedright\arraybackslash}p{(\linewidth - 8\tabcolsep) * \real{0.0938}}
  >{\raggedright\arraybackslash}p{(\linewidth - 8\tabcolsep) * \real{0.3281}}
  >{\raggedright\arraybackslash}p{(\linewidth - 8\tabcolsep) * \real{0.1562}}
  >{\raggedright\arraybackslash}p{(\linewidth - 8\tabcolsep) * \real{0.2969}}@{}}
\toprule\noalign{}
\begin{minipage}[b]{\linewidth}\raggedright
System
\end{minipage} & \begin{minipage}[b]{\linewidth}\raggedright
Year
\end{minipage} & \begin{minipage}[b]{\linewidth}\raggedright
Constraint Structure
\end{minipage} & \begin{minipage}[b]{\linewidth}\raggedright
Best For
\end{minipage} & \begin{minipage}[b]{\linewidth}\raggedright
Key Proof Systems
\end{minipage} \\
\midrule\noalign{}
\arrayrulecolor{MidnightBlue}\midrule\arrayrulecolor{MidnightBorder}
\endhead
\bottomrule\noalign{}
\endlastfoot
R1CS & 2012 & Bilinear (degree 2) & Small circuits, Groth16 & Groth16, Spartan,
Nova \\
AIR & 2018 & Uniform polynomial & VM traces, STARKs & STARKs (Stone, Stwo) \\
PLONKish & 2019 & Selector-gated, custom & Flexible circuits & PLONK, Halo2 \\
\end{longtable}

A folding scheme designed for R1CS could not accept AIR input. A proof system
built for AIR could not handle PLONKish circuits. A developer choosing an
arithmetization was simultaneously choosing a proof system ecosystem -- and
switching later meant rewriting everything.

The fragmentation had real costs. When the Polygon team decided to migrate from
Hermez (PLONK-based) to a STARK-based architecture, the circuit rewrite took
years. When Scroll built their zkEVM on Halo2 (PLONKish), they could not easily
adopt the newer sumcheck-based proof systems that emerged in 2023-2024 without
rewriting their entire constraint system. Research teams working on folding
schemes had to choose: target R1CS (like Nova did) and exclude the STARK
ecosystem, or target a custom format and exclude everyone else. The constraint
system choice was a one-way door -- enter through it, and you are locked into
the corresponding proof system family for the life of the project.

The field needed a unifier. Not a compromise format that sacrificed efficiency
for generality, but a mathematical framework that could express R1CS, AIR, and
PLONKish as special cases -- preserving the efficiency of each while providing a
single target for proof systems to implement.

\section{CCS: The Rosetta Stone}\label{ccs-the-rosetta-stone}

In 2023, Srinath Setty, Justin Thaler, and Riad Wahby published a paper that
changed constraint system design. Customizable Constraint Systems (CCS) unified
R1CS, AIR, and PLONKish into a single mathematical framework, without overhead.

The word ``without overhead'' is the point. Previous attempts at unification
existed -- for example, you can always convert AIR to R1CS by expanding every
transition into individual constraints, or convert R1CS to AIR by padding with
identity transitions. But these conversions incur blowup: the converted instance
is larger, sometimes much larger, than the original. CCS achieves something
different: it captures each constraint format in its native form, preserving the
sparsity and structure that makes each format efficient. An R1CS instance
becomes a CCS instance of exactly the same size. An AIR instance becomes a CCS
instance of exactly the same size. Nothing is wasted in translation.

\subsection{The Idea}\label{the-idea}

A CCS instance is defined by a set of sparse matrices \(M_1, \ldots, M_t\) over
a finite field, a collection of multisets \(S_1, \ldots, S_q\) (each specifying
which matrices to multiply together element-wise), and constants
\(c_1, \ldots, c_q\). The satisfying condition is:

\(\sum_{i=1}^{q} c_i \cdot \bigcirc_{j \in S_i} (M_j \cdot \mathbf{z}) = \mathbf{0}\)

(The Hadamard product is simply element-wise multiplication:
\([a, b, c] \circ [d, e, f] = [ad, be, cf]\). When the formula says ``Hadamard
product of \(M_j \cdot \mathbf{z}\),'' it means: compute each matrix-vector
product separately, then multiply the resulting vectors element by element.)

This looks abstract. Here is what it means concretely.

\textbf{R1CS is CCS with two terms.} Set \(q = 2\), \(S_1 = \{1, 2\}\),
\(S_2 = \{3\}\), \(c_1 = 1\), \(c_2 = -1\). The satisfying condition becomes
\((M_1 \cdot \mathbf{z}) \circ (M_2 \cdot \mathbf{z}) - (M_3 \cdot \mathbf{z}) = 0\),
which is exactly the R1CS equation
\(\mathbf{A}\mathbf{z} \circ \mathbf{B}\mathbf{z} = \mathbf{C}\mathbf{z}\) when
you identify \(M_1 = A\), \(M_2 = B\), \(M_3 = C\).

\textbf{AIR is CCS with matrices encoding shift relations.} The transition
constraints between consecutive rows become matrix-vector products with
appropriate shift structure.

\textbf{PLONKish is CCS with matrices encoding selector-weighted gate
equations.} The selector polynomials become entries in the matrices; the copy
constraints map to specific matrix structures.

The key insight is not that CCS enables new computations -- any NP statement can
already be expressed in R1CS. The insight is that CCS provides a \emph{uniform
interface}. A proof system that targets CCS automatically handles R1CS, AIR, and
PLONKish inputs without conversion overhead. Write one folding scheme for CCS,
and it works with every constraint format the industry has produced.

\subsection{Why CCS Matters Now}\label{why-ccs-matters-now}

CCS is the native constraint system for every major modern folding scheme:

\begin{itemize}
\tightlist
\item
  \textbf{HyperNova} (Kothapalli and Setty, 2023): multi-folding for CCS, using
  the sumcheck protocol to fold multiple CCS instances simultaneously.
\item
  \textbf{ProtoStar} and \textbf{ProtoGalaxy} (2023): folding schemes that
  generalize Nova to higher-degree constraint systems -- which CCS naturally
  supports.
\item
  \textbf{Neo} (Nguyen and Setty, 2025): the first lattice-based folding scheme
  for CCS, achieving post-quantum security with native small-field efficiency.
\item
  \textbf{LatticeFold+} (Boneh and Chen, 2025): extends LatticeFold with faster,
  simpler lattice-based folding and shorter proofs.
\end{itemize}

Without CCS, none of these systems could claim generality. Each would be locked
to R1CS (like Nova) or would need separate implementations for each constraint
format. CCS is the abstraction layer that made the folding revolution possible.

The gap between research and deployment is real, however. Production systems in
early 2026 still largely use PLONKish (Halo2, Scroll) or AIR (StarkWare, Stwo).
The CCS-native stack is approximately two to three years behind the research
frontier. But the trajectory is clear: as folding-based proof systems move from
research prototypes to production deployments, CCS will become the standard
target.

The parallel to programming language history is instructive. In the 1960s, each
computer had its own instruction set, its own assembler, and its own operating
conventions. Writing a program for an IBM 7090 required completely different
code than writing for a UNIVAC 1108. Then came C and UNIX, which provided a
common language and a common operating system interface. Programs written in C
could run on any machine with a C compiler. CCS plays the same role for
constraint systems: it provides a common mathematical interface that any proof
system can target. Write your constraints in CCS, and any CCS-compatible proof
system -- HyperNova, ProtoStar, Neo -- can prove them. The ``operating system''
for zero-knowledge proof systems is being standardized, even if the
``applications'' (production deployments) have not yet caught up.

\subsection{The Degree Parameter}\label{the-degree-parameter}

One subtle but important feature of CCS is the degree parameter \(d\), which
captures the maximum degree of the constraint polynomials. R1CS has \(d = 2\)
(bilinear constraints). PLONKish can have \(d = 2\) or higher, depending on the
custom gate design. CCS handles arbitrary degree without modification.

This matters because higher-degree constraints can capture more complex
operations in fewer constraints. A single degree-\(4\) constraint can express
relationships that would require multiple degree-\(2\) R1CS constraints. The
tradeoff is that higher-degree constraints require more sophisticated proof
techniques -- but the sumcheck protocol, which we turn to next, handles
arbitrary degrees naturally.

To see the degree parameter in action, return to the ``x * (x + 1)''
computation. In R1CS (\(d = 2\)), this requires two constraints: \(t = x + 1\)
(degree \(1\), but padded to the bilinear form as \((x + 1) \cdot 1 = t\)) and
\(\text{result} = x \cdot t\) (degree \(2\)). In a CCS instance with \(d = 3\),
you could express the entire computation in a single constraint:
\(x \cdot (x + 1) - \text{result} = 0\), which is a degree-\(2\) polynomial in
\(x\). With \(d = 4\), you could encode
\(x \cdot (x + 1) \cdot (x + 2) - \text{result} = 0\) in a single constraint --
a relationship that would require three R1CS constraints (one for each pairwise
multiplication). Higher degree means more computation packed into fewer
constraints, at the cost of more complex proof machinery.

\subsection{Three Dialects, One Grammar}\label{three-dialects-one-grammar}

Return to the three micro-examples we built in the previous sections and look at
them through the CCS lens. What CCS reveals is that R1CS, AIR, and PLONKish are
not three different formalisms. They are three dialects of the same language.

\textbf{R1CS is CCS with \(q = 2\).} You need exactly two multisets:
\(S_1 = \{1, 2\}\) (which multiplies the results of matrices \(M_1\) and \(M_2\)
element-wise) and \(S_2 = \{3\}\) (which provides \(M_3\)'s result). The CCS
equation becomes
\(c_1 \cdot (M_1 \cdot \mathbf{z} \circ M_2 \cdot \mathbf{z}) + c_2 \cdot (M_3 \cdot \mathbf{z}) = 0\),
with \(c_1 = 1\) and \(c_2 = -1\). This is exactly
\((\mathbf{A} \cdot \mathbf{z}) \circ (\mathbf{B} \cdot \mathbf{z}) - \mathbf{C} \cdot \mathbf{z} = 0\)
-- the R1CS equation, expressed in CCS notation. Two matrix-vector products, one
Hadamard product, one subtraction. That is the entire constraint system. Every
R1CS instance that has ever been deployed -- every Groth16 proof, every Spartan
verification -- is a CCS instance with \(q = 2\).

For the ``3 * 4 = 12'' multiplication from the spreadsheet example, the CCS
encoding has \(t = 3\) matrices (\(M_1 = A\), \(M_2 = B\), \(M_3 = C\)),
\(q = 2\) multisets (\(S_1 = \{1, 2\}\) and \(S_2 = \{3\}\)), and the witness
vector \(\mathbf{z} = (1, 3, 4, 12)\) where the first entry is the constant 1.
The matrix A selects the left operand (3), B selects the right operand (4), and
C selects the output (12). The Hadamard product
\((\mathbf{A} \cdot \mathbf{z}) \circ (\mathbf{B} \cdot \mathbf{z})\) computes
\(3 \cdot 4 = 12\) element-wise, and subtracting
\(\mathbf{C} \cdot \mathbf{z} = 12\) yields zero. One constraint, three
matrices, one Hadamard product.

\textbf{AIR is CCS with shift matrices.} The counter example from earlier had
the transition constraint
\(\text{counter}[i+1] = \text{counter}[i] + \text{flag}[i]\). In CCS, this
becomes a set of matrices where one matrix \(M_{\text{shift}}\) extracts the
``next row'' values and another \(M_{\text{current}}\) extracts the ``current
row'' values. For a 3-row trace, \(M_{\text{current}}\) might have ones on the
diagonal (selecting counter{[}0{]}, counter{[}1{]}, counter{[}2{]}) while
\(M_{\text{shift}}\) has ones on the superdiagonal (selecting counter{[}1{]},
counter{[}2{]}, counter{[}0{]} with wraparound). The shift operation -- looking
at row i + 1 instead of row i -- is encoded in the matrix structure itself. The
polynomial constraint between consecutive rows becomes a matrix-vector product
where the matrix has ones on a shifted diagonal. What looked like a
fundamentally different formalism (rules between consecutive rows, rather than
rules within a single row) turns out to be a specific matrix pattern within CCS.

\textbf{PLONKish is CCS with selector-weighted matrices.} The PLONKish trace
with its \(q_{\text{add}}\) and \(q_{\text{mul}}\) columns maps to CCS matrices
where the selector values are baked into the matrix entries. The gate equation
\(q_{\text{add}} \cdot (a + b - c) + q_{\text{mul}} \cdot (a \cdot b - c) = 0\)
becomes a CCS instance where one multiset captures the addition term (weighted
by \(q_{\text{add}}\)) and another captures the multiplication term (weighted by
\(q_{\text{mul}}\)). The copy constraints -- the permutation argument that wires
outputs to inputs -- map to additional matrix structure that enforces equality
between specific positions in the witness vector.

The visual is this: imagine three spreadsheets, each with different column
headers and different rules. The R1CS spreadsheet has columns A, B, C with the
rule ``A times B equals C.'' The AIR spreadsheet has columns for registers with
the rule ``next row relates to current row by this transition polynomial.'' The
PLONKish spreadsheet has witness columns and selector columns with the rule
``the active gate constraint must be satisfied.'' Three different layouts. Three
different conventions. But CCS says: they are all just matrices times a witness
vector, combined with Hadamard products and summed to zero. One grammar, three
dialects.

This is not a metaphor. It is a theorem. Any R1CS instance, any AIR instance,
any PLONKish instance can be mechanically translated into a CCS instance with no
increase in constraint count or witness size. The translation preserves
everything -- the structure, the sparsity, the degree. When a proof system like
HyperNova targets CCS, it is not accepting a lowest-common-denominator format.
It is accepting the universal format that contains every existing constraint
dialect as a special case.

The practical consequence is immediate. Before CCS, a developer choosing R1CS
was simultaneously choosing Groth16 or Spartan. A developer choosing AIR was
choosing STARKs. A developer choosing PLONKish was choosing Halo2 or PLONK.
Switching constraint systems meant rewriting the circuit and the proof system
integration. CCS breaks this coupling. Write your constraints in whichever
dialect is natural for your computation -- R1CS for simple circuits, AIR for VM
traces, PLONKish for mixed-gate workloads -- and any CCS-compatible proof system
will accept them without translation overhead. The grammar is universal; the
dialects are a matter of convenience.

There is a deeper mathematical point here, one that Penrose would appreciate.
The existence of a universal constraint grammar is not obvious. One might have
expected that the structural differences between R1CS (bilinear, flat), AIR
(uniform, sequential), and PLONKish (selector-gated, permutation-wired) would
require genuinely different proof techniques -- that no single algebraic
framework could capture all three without paying some conversion tax. CCS
demonstrates that the differences are shallow. At the level of sparse
matrix-vector products and Hadamard products, all three constraint systems are
doing the same thing. The ``three families'' narrative that dominated ZK from
2018 to 2022 was a historical artifact, not a mathematical necessity.

CCS provides the universal grammar. Sumcheck provides the universal verification
engine. The two are partners: CCS tells us \emph{what} the constraints look like
-- a sum of Hadamard products of matrix-vector pairs -- and sumcheck tells us
\emph{how to check} that sum without evaluating every term. The verifier does
not inspect every cell of the constraint spreadsheet. Instead, sumcheck reduces
the problem: ``does this multilinear polynomial sum to zero over the boolean
hypercube?'' becomes, after n rounds of interaction, ``does this polynomial
evaluate correctly at one random point?'' The reduction is exponential -- from
\(2^n\) checks to \(n\) rounds -- and it is the reason modern proof systems can
verify in time logarithmic in the computation size.

\section{The Sumcheck Protocol: The Hidden
Foundation}\label{the-sumcheck-protocol-the-hidden-foundation}

If there is one protocol that deserves to be called the backbone of modern
zero-knowledge proof systems, it is the sumcheck protocol. Lund, Fortnow,
Karloff, and Nisan introduced it in 1992 -- decades before practical ZK systems
existed. Sumcheck has since become the common thread running through every major
proof system of the current era.

\subsection{What Sumcheck Does}\label{what-sumcheck-does}

Before stating the problem, one definition. A multilinear polynomial is one
where no variable appears with degree higher than one -- for example,
\(g(x_1, x_2) = 3x_1 x_2 + 2x_1 + x_2 + 1\). Each variable is either present or
absent in any given term, but never squared or cubed. The sumcheck protocol
works naturally with multilinear polynomials because their structure matches the
binary hypercube over which the sum is taken: each variable takes the value 0 or
1, so higher powers would collapse anyway (since \(0^k = 0\) and \(1^k = 1\)).
Multilinear polynomials are the native representation for most sumcheck-based
proof systems, including Spartan, HyperNova, and Jolt.

The problem sumcheck solves is deceptively simple. You have a multivariate
polynomial \(g(x_1, \ldots, x_n)\) over a finite field, and you want to verify
that its sum over all binary inputs equals a claimed value:

\(\sum_{(x_1, \ldots, x_n) \in \{0,1\}^n} g(x_1, \ldots, x_n) = T\)

Naively, checking this requires evaluating \(g\) at all \(2^n\) binary inputs.
For a polynomial over 30 variables, that is a billion evaluations. Sumcheck
reduces this to \(n\) rounds of interaction (or, via the Fiat-Shamir transform,
\(n\) rounds of hash-based challenge generation), each involving a single
univariate polynomial of low degree.

The protocol is interactive: the prover and verifier exchange messages in
rounds. In practice, the Fiat-Shamir transform replaces the verifier's random
challenges with hash outputs, making the protocol non-interactive. But the
conceptual structure remains round-based.

Here is the intuition. In round 1, the prover sends a univariate polynomial
\(p_1(x_1)\) that claims to be the sum of \(g\) over all remaining variables:
\(p_1(x_1) = \sum_{x_2,\ldots,x_n} g(x_1, x_2, \ldots, x_n)\). The verifier
checks that \(p_1(0) + p_1(1) = T\) (this ensures consistency with the claimed
total), then sends a random challenge \(r_1\). In round 2, the prover sends
\(p_2(x_2) = \sum_{x_3,\ldots,x_n} g(r_1, x_2, x_3, \ldots, x_n)\). The verifier
checks \(p_2(0) + p_2(1) = p_1(r_1)\), then sends another random challenge. This
continues until all variables are bound to random values, at which point the
verifier checks one evaluation of \(g\) at the random point.

The result: verifying a sum over \(2^n\) inputs reduces to checking \(n\)
low-degree univariate polynomials and one evaluation of \(g\). For \(n = 30\),
that is 30 polynomial checks instead of a billion evaluations.

To make this concrete, suppose we want to verify that a polynomial
\(g(x_1, x_2)\) sums to \(T\) over all binary inputs. There are four inputs:
\(g(0,0) + g(0,1) + g(1,0) + g(1,1) = T\). Instead of checking all four, the
prover sends a univariate polynomial \(p_1(x_1)\) that claims to be the partial
sum over \(x_2\). The verifier checks: does \(p_1(0) + p_1(1) = T\)? If yes, the
verifier picks a random \(r_1\) and asks for the next round. Now the prover
sends \(p_2(x_2)\) claiming to sum \(g(r_1, x_2)\). The verifier checks
\(p_2(0) + p_2(1) = p_1(r_1)\). After two rounds, the verifier holds a single
point \(g(r_1, r_2)\) that can be checked directly. Two rounds replaced four
evaluations. For \(n\) variables, \(n\) rounds replace \(2^n\) evaluations.

This exponential compression is why sumcheck appears everywhere in modern ZK.
The reduction from \(2^n\) evaluations to \(n\) rounds is not merely a
constant-factor improvement. It is an exponential improvement -- the kind that
turns impossible problems into trivial ones. For a polynomial over 100
variables, naively verifying the sum would require \(2^{100}\) evaluations (more
than the number of atoms in the observable universe). Sumcheck reduces this to
100 rounds. The gap between ``impossible'' and ``trivial'' is the gap that
sumcheck bridges.

To see how sumcheck serves CCS specifically, consider a CCS instance with a
single constraint:
\(M_1 \cdot \mathbf{z} \circ M_2 \cdot \mathbf{z} = M_3 \cdot \mathbf{z}\)
(where \(\circ\) is the Hadamard product). The verifier needs to check that this
equation holds at every row. Equivalently, the verifier needs to check that the
polynomial
\(h(x) = (M_1 \cdot \mathbf{z})(x) \cdot (M_2 \cdot \mathbf{z})(x) - (M_3 \cdot \mathbf{z})(x)\)
sums to zero over all row indices. This is exactly a sumcheck instance. The
prover sends the univariate restriction of \(h\) in the first variable, the
verifier checks its degree and evaluates at a random point, and the process
recurses on the remaining variables. After \(\log_2(n)\) rounds, the verifier
holds a single evaluation claim that can be checked against the committed
polynomials. The CCS constraint -- which could be R1CS, AIR, or PLONKish in
disguise -- has been verified without the verifier ever touching the witness.

Spartan uses it for R1CS verification. HyperNova uses it for CCS folding. Jolt
and Lasso use it for lookup arguments. SP1 Hypercube builds its entire
polynomial stack on sumcheck. When the Ethereum Foundation's zkEVM effort
evaluated proof system designs, sumcheck-based architectures won -- not because
they are simplest to implement, but because the exponential reduction in
verifier work is too large to ignore.

\subsection{Why Sumcheck Is Everywhere}\label{why-sumcheck-is-everywhere}

The sumcheck protocol is the verification mechanism that makes polynomial-based
arithmetization practical. Every time a modern proof system needs to verify that
a polynomial identity holds over a large domain, sumcheck is how it does so.

\begin{itemize}
\tightlist
\item
  \textbf{Spartan} (Setty, 2019) uses sumcheck to verify R1CS satisfaction
  directly, without FFTs. The prover expresses the R1CS check as a multilinear
  polynomial sum and runs sumcheck to prove it holds.
\item
  \textbf{HyperNova} uses sumcheck as the core of its multi-folding protocol.
  Folding multiple CCS instances reduces to a sumcheck instance.
\item
  \textbf{Jolt and Lasso} reduce lookup verification to sumcheck instances.
  Every table lookup becomes a polynomial sum that sumcheck can verify.
\item
  \textbf{LogUp-GKR} combines the sumcheck protocol with the GKR interactive
  proof to verify lookup arguments with logarithmic overhead.
\item
  \textbf{SP1 Hypercube} (Succinct, 2025) uses sumcheck over the Boolean
  hypercube as its primary verification strategy.
\item
  \textbf{Binius} (Irreducible, 2025) applies sumcheck over binary tower fields.
\end{itemize}

The sumcheck protocol is to modern ZK proof systems what the internal combustion
engine was to early automobiles: the mechanism that makes the entire apparatus
work, even though the user never sees it directly. Understanding that sumcheck
exists, and that it reduces exponential verification to linear communication, is
essential for understanding why the overhead of arithmetization is not as
catastrophic as it might first appear.

A note on presentation order: this chapter covers arithmetization (Layer 4)
before the proof system (Layer 5) and cryptographic primitives (Layer 6). But in
practice, the causal arrow often runs the other way. The sumcheck protocol is a
proof technique that shaped which arithmetization formats became practical. The
field choice at Layer 6 determines which polynomial representations are
efficient at Layer 4. These three layers -- field, commitment scheme, polynomial
representation -- form an inseparable ``proof core.'' We present them in the
standard order, but the reader should understand that the dependency is
circular, not linear.

The sumcheck protocol also illustrates a recurring theme in this book: the most
important technical ideas are often invisible to the end user. A developer
writing a Compact smart contract on Midnight, or a Solidity developer deploying
a Groth16 verifier on Ethereum, will never interact with the sumcheck protocol
directly. They will never see a multilinear polynomial or check a partial sum.
But sumcheck is running underneath, silently reducing the verification cost from
exponential to linear, making the entire stack practical. The seven layers of
the magic trick include mechanisms that the audience never sees -- and sumcheck
is the most consequential of them all.

\section{Lookup Arguments}\label{lookup-arguments}

In the classical approach to arithmetization, every operation in a computation
is encoded as polynomial constraints. Addition and multiplication are natural:
they are already arithmetic operations over the field. But what about operations
that are \emph{not} naturally arithmetic?

Consider the problem from the perspective of a circuit designer building a zkVM.
The RISC-V instruction set contains 47 base instructions. Of these, roughly half
are ``arithmetic-friendly'' -- ADD, SUB, MUL, and their variants map naturally
to field operations. But the other half are ``arithmetic-hostile'': AND, OR,
XOR, SLL (shift left logical), SRL (shift right logical), SLT (set less than),
BEQ (branch if equal), and the memory load/store operations. Each of these
requires bitwise decomposition or comparison logic that does not map neatly to
field arithmetic.

Bitwise AND, comparison, range checks, hash functions -- these require
decomposing the values into individual bits, constraining each bit to be 0 or 1,
and then reconstructing the result. A single SHA-256 hash invocation can require
tens of thousands of constraints.

The cost is staggering when you trace through a concrete example. Consider XOR
-- the bitwise exclusive-or of two 8-bit numbers. On a CPU, this is one
instruction, one clock cycle, done. In a polynomial constraint system, you must
first decompose each 8-bit input into 8 individual bits (8 range-check
constraints per input, 16 total), then constrain each output bit to be the XOR
of the corresponding input bits (each bit-level XOR requires the polynomial
\(a + b - 2ab\), which is a degree-\(2\) constraint, so 8 more constraints), and
finally reconstruct the output from its bits (8 more constraints). That is
roughly 32 constraints for an operation that takes a single machine cycle.
SHA-256 calls XOR, AND, NOT, and rotation thousands of times. Multiply 32
constraints per operation by thousands of operations and you arrive at the tens
of thousands of constraints that a single hash invocation demands.

This is not a fixable inefficiency in the constraint system design. It is a
fundamental mismatch between the polynomial language (which speaks addition and
multiplication over large fields) and the bitwise language (which speaks AND,
OR, XOR over individual bits). No amount of cleverness in the constraint layout
will make polynomials natively express bit manipulation. The operations live in
different algebraic worlds.

The mismatch created a two-tier cost structure in early ZK systems.
``Arithmetic-friendly'' operations (Poseidon hash, MiMC, elliptic curve
arithmetic) -- operations designed from the ground up to minimize constraint
count -- were cheap. ``Arithmetic-hostile'' operations (SHA-256, Keccak, AES,
bitwise logic) -- operations from the traditional computing world -- were
expensive by comparison. This is why the ZK community designed entirely new hash
functions (Poseidon, Rescue, Neptune) that use only field additions and
multiplications, avoiding bitwise operations entirely. A Poseidon hash costs
roughly 300 constraints in R1CS. A SHA-256 hash costs roughly 25,000
constraints. Same security level. Hundred-fold difference in proving cost. The
constraint system penalizes any computation that crosses the boundary between
field arithmetic and bit arithmetic.

Lookup arguments offer a fundamentally different approach: instead of encoding
the operation as constraints, you look up the answer in a pre-computed table.
Instead of proving \emph{how} you computed XOR, you prove \emph{that} your
answer appears in a table of all correct XOR results. The philosophical shift is
from verification-by-recomputation to verification-by-membership.

\subsection{Plookup: The First Practical Lookup
(2020)}\label{plookup-the-first-practical-lookup-2020}

Gabizon and Williamson introduced Plookup in 2020. The idea: if you have a table
of pre-approved input-output pairs (for example, all possible 8-bit XOR
results), you can prove that a set of values appears in the table without
recomputing the operation.

Consider that table of 8-bit XOR results. It has 256 * 256 = 65,536 entries,
each of the form (a, b, a XOR b). If the prover claims that 0x3F XOR 0xA7 =
0x98, the verifier does not check the XOR. Instead, the verifier checks that the
triple (0x3F, 0xA7, 0x98) appears somewhere in the table. If it does, the answer
is correct -- because the table was constructed correctly, and membership in a
correct table implies correctness of the result.

Plookup works by sorting the lookup values and the table entries into a single
sorted sequence, then verifying the sorting through a grand product argument. If
every lookup value appears in the table, the sorted sequence has a specific
structure that the grand product captures. The prover merges the table and the
lookup values into one sorted list, then proves that the merged list is a valid
interleaving of the original table with the queried entries. A polynomial
identity -- checked via a grand product over the entire sorted sequence --
catches any entry that does not belong.

The catch: sorting costs \(O(n \log n)\), and the grand product requires
committing to the sorted sequence. For a circuit with \(n\) lookups into a table
of size \(T\), the prover must commit to a sorted list of length \(n + T\). This
makes Plookup a meaningful optimization for expensive operations (hashes, range
checks) but not a universal solution. The sorting overhead is the bottleneck,
and it resists parallelization -- you cannot sort half a list on one machine and
half on another without a merge step.

Despite its limitations, Plookup was immediately adopted. Lookup arguments for
range checks (proving a value is between \(0\) and \(2^N\)) replaced the naive
approach of decomposing into N bits and constraining each one. For a 16-bit
range check, the naive approach requires 16 boolean constraints plus a
reconstruction constraint (17 total). A Plookup-based range check requires one
lookup into a table of 65,536 entries. The lookup is more expensive in absolute
prover time (sorting overhead), but far cheaper in constraint count -- and for
systems where constraint count is the bottleneck, this tradeoff is worthwhile.
By 2021, every major PLONKish system used lookup arguments for range checks.

\subsection{LogUp: The Sorting-Free Revolution
(2022)}\label{logup-the-sorting-free-revolution-2022}

Ulrich Haboeck's LogUp paper in 2022 replaced Plookup's sorting with an
observation that is, in retrospect, elegant to the point of inevitability. The
name ``LogUp'' comes from ``logarithmic derivative'' -- the key mathematical
technique. If you have a polynomial P(X) = Product of (X - r\_i) whose roots are
exactly the lookup values, then the logarithmic derivative of P is P'(X)/P(X) =
Sum of 1/(X - r\_i). This transforms a product (which is hard to check
incrementally) into a sum (which is easy to accumulate and verify). The idea
comes from complex analysis, where logarithmic derivatives convert
multiplicative structures into additive ones. Haboeck's insight was to apply
this classical technique to the lookup problem.

Instead of sorting, LogUp observes that if every lookup value \(f_i\) appears in
the table \(t\), then a specific identity over rational functions must hold:

\(\sum_i \frac{1}{X - f_i} = \sum_j \frac{m_j}{X - t_j}\)

where \(m_j\) counts how many times table entry \(t_j\) is looked up. The left
side sums one term per lookup. The right side sums one term per table entry,
weighted by the number of times it was accessed. If every lookup value is in the
table, these two sums are equal as formal rational functions -- and therefore
they are equal at a random evaluation point with overwhelming probability.

A tiny example makes this concrete. Suppose the table contains \{1, 2, 3\} and
the prover claims to look up the values \{2, 3, 2\}. The left side (one term per
lookup) is: \(1/(X-2) + 1/(X-3) + 1/(X-2) = 2/(X-2) + 1/(X-3)\). The right side
(one term per table entry, weighted by frequency) is:
\(0/(X-1) + 2/(X-2) + 1/(X-3)\) -- because entry 1 is looked up 0 times, entry 2
is looked up twice, and entry 3 is looked up once. Both sides equal
\(2/(X-2) + 1/(X-3)\). The identity holds. Now suppose the prover cheats and
claims to look up \{2, 3, 5\}, where 5 is not in the table. The left side
becomes \(1/(X-2) + 1/(X-3) + 1/(X-5)\). No assignment of multiplicities to the
table entries \{1, 2, 3\} can produce a term \(1/(X-5)\) on the right side. The
identity fails at a random evaluation point with overwhelming probability.

This transforms the lookup argument from a product check to a \emph{sum check}
-- and sums, unlike products, compose beautifully. Why does this matter? Because
a product of n terms can be thrown off by a single corrupted factor (the product
becomes wrong, but localizing the error requires inspecting every factor). A sum
of n terms, by contrast, is naturally decomposable: you can split the sum into
batches, compute partial sums independently, and aggregate them. The algebraic
structure of addition is friendlier than the algebraic structure of
multiplication.

The advantages are substantial:

\begin{itemize}
\tightlist
\item
  \textbf{No sorting required.} Eliminates the \(O(n \log n)\) overhead
  entirely. The prover's cost drops to \(O(n + T)\), linear in the number of
  lookups plus the table size.
\item
  \textbf{Natural batching.} Multiple lookup tables can be handled
  simultaneously via random linear combinations. If your circuit uses a XOR
  table, an AND table, and a range-check table, LogUp handles all three in one
  pass.
\item
  \textbf{Parallelizable.} The summation structure means partial lookups can be
  computed independently on separate machines and aggregated with a single
  addition. This is exactly the property that GPU-based provers exploit.
\item
  \textbf{Sumcheck-compatible.} The rational function identity can be verified
  using the sumcheck protocol, connecting lookups directly to the same
  verification backbone that handles constraint checking.
\end{itemize}

LogUp became the production standard. It replaced Plookup in deployed systems
and enabled the next generation of lookup-based architectures. The transition
was rapid: by 2024, virtually every new proof system design used LogUp or a
LogUp variant for its lookup needs.

\subsection{LogUp-GKR: The Verifier Gets Faster
(2023)}\label{logup-gkr-the-verifier-gets-faster-2023}

LogUp made the prover efficient. But the verifier still had to check the
rational function identity, which naively requires work proportional to the
number of lookups. Papini and Haboeck combined LogUp with the GKR interactive
proof protocol to solve this.

The GKR protocol provides an efficient way to verify layered arithmetic circuits
-- circuits where the computation flows through layers, each layer depending
only on the one before it. The fractional sum computation in LogUp (adding up
all the 1/(X - f\_i) terms) has exactly this layered structure: it is a sum
reduction tree. LogUp-GKR applies the GKR protocol to this tree, reducing the
verifier's work from linear to logarithmic in the number of lookups.

The result: logarithmic proof size and verification time for lookup arguments.
The verifier does \(O(\log n)\) work regardless of how many lookups the prover
performed.

LogUp-GKR is now used in Polygon's Plonky3 framework and in SP1 Hypercube. It
makes lookups nearly free for the verifier, which is critical for recursive
proof composition -- where the verifier's circuit size directly affects the cost
of the next recursion step. If verifying a lookup takes \(O(n)\) work, then a
recursive verifier circuit must be \(O(n)\) in size, which is expensive to prove
in the next recursion layer. With LogUp-GKR, the recursive verifier circuit is
\(O(\log n)\), making deep recursion practical.

The combination of LogUp (efficient prover) and GKR (efficient verifier) made
lookups a first-class operation in the proof system -- no longer an optimization
to be applied selectively, but a general-purpose tool to be used wherever a
pre-computed table exists. This set the stage for the two results that completed
the lookup revolution: Lasso, which made table size irrelevant, and Jolt, which
made lookups the only computation paradigm needed.

\subsection{Lasso: The Table Size Disappears
(2023)}\label{lasso-the-table-size-disappears-2023}

The fundamental limitation of both Plookup and LogUp is that the prover must
somehow touch the entire table. For a table of \(2^{16}\) entries (65,536 rows),
this is manageable. For a table of all possible 64-bit operations -- \(2^{128}\)
entries -- it is physically impossible to even store the table, let alone commit
to it. The table of all 64-bit additions alone has \(2^{128}\) rows. At one byte
per row, that is \(10^{38}\) bytes -- more than the number of atoms in the
observable universe. No amount of hardware solves this.

Lasso, by Setty, Thaler, and Wahby (2023), solves this through
\emph{decomposition}. The insight is that most useful tables have internal
structure that can be exploited. Specifically, if the table's multilinear
extension (MLE) can be evaluated efficiently -- meaning you can compute the
table's value at any point without materializing the entire table -- then each
lookup can be decomposed into lookups into much smaller subtables.

The intuition is best seen through an analogy. Suppose you have a multiplication
table for two-digit numbers. Instead of storing all 90 * 90 = 8,100 entries (for
digits 10-99), you could decompose each two-digit number into its tens and units
digits, store separate multiplication tables for single digits (only 10 * 10 =
100 entries each), and reconstruct the full product from partial products. The
full table has 8,100 entries; the subtables have a combined 200 entries. You
traded one large lookup for several small lookups plus some arithmetic glue.
Lasso does exactly this, but for arbitrary structured tables over finite fields,
using the multilinear extension as the decomposition mechanism.

For a table of size \(2^{2W}\), Lasso decomposes each lookup index into \(c\)
chunks. Instead of one lookup into a table of size \(2^{2W}\), the prover
performs \(c\) lookups into subtables of size \(2^{2W/c}\). For 64-bit RISC-V
operations with \(c = 6\), each subtable has roughly \(2^{22}\) entries -- about
4 million rows. That fits comfortably in memory.

The prover's cost becomes \(O(n \cdot c \cdot \log(N)/c)\), which is
proportional to the number of lookups (n) and the number of chunks (c), but
\emph{independent of the table size} (N). You can look up values in a table of
\(2^{128}\) entries without ever materializing the table. The table exists as a
mathematical function -- its MLE -- not as a stored data structure. The prover
only pays for the subtable entries it actually accesses.

This is the kind of result that reshapes what is considered possible. Before
Lasso, ``table size'' was a hard constraint on lookup arguments. After Lasso,
table size is irrelevant -- only table structure matters. A structured table of
\(2^{128}\) entries is no harder to use than a structured table of \(2^{16}\)
entries. The prover's work scales with the number of lookups it performs, not
with the number of entries it could theoretically look up.

\subsection{Jolt: The Lookup Singularity Realized
(2023)}\label{jolt-the-lookup-singularity-realized-2023}

Arun, Setty, and Thaler's Jolt paper took Lasso's theoretical framework and
applied it to its logical conclusion: what if \emph{every} instruction in a
processor were a lookup?

The concept, originally proposed by Barry Whitehat as the ``lookup
singularity,'' posits that circuits should be expressed entirely as lookups into
pre-computed tables. Jolt demonstrates this is achievable for a complete RISC-V
instruction set:

\begin{enumerate}
\def\labelenumi{\arabic{enumi}.}
\tightlist
\item
  Every RISC-V instruction has an evaluation table mapping inputs to outputs.
\item
  All instruction tables are MLE-structured (their multilinear extensions can be
  evaluated efficiently).
\item
  Lasso's decomposition makes these lookups efficient regardless of the
  theoretical table size.
\item
  Memory consistency is verified through offline memory checking
  (fingerprint-based techniques), not Merkle trees.
\end{enumerate}

For each instruction, the prover decomposes the operands into chunks, performs
lookups into small subtables (typically around 4 million entries each), and
commits to roughly 18 field elements per instruction (3 per chunk, with
\(c = 6\) chunks).

The memory-checking component (point 4) is worth highlighting separately. In a
real processor, memory is read-write: the program loads and stores values
freely. Proving that every load returns the value of the most recent store to
the same address is the memory consistency problem discussed in the overhead
section. Jolt handles this through ``offline memory checking'' -- a technique
where the prover computes a cryptographic fingerprint of the sequence of all
reads and writes, and the verifier checks that the fingerprint is consistent
with a valid read-write memory. This avoids the per-access cost of Merkle tree
proofs and makes memory checking nearly as cheap as instruction checking. The
technique is not specific to Jolt; it was developed by Blum et al.~in the 1990s
and adapted for ZK by Setty (Spartan) and others. But Jolt's integration of
offline memory checking with Lasso-based instruction lookups produces a complete
zkVM architecture where every component -- instruction verification, memory
consistency, program counter management -- is handled by either a lookup or a
fingerprint check.

The result is a zkVM where the constraint system is almost entirely lookups,
with minimal arithmetic ``glue.'' This is not an optimization applied to an
existing constraint system -- it is a different paradigm for encoding
computation.

Pause and consider what Jolt achieved. A complete RISC-V instruction set -- ADD,
SUB, AND, OR, XOR, SLL, SRL, SRA, SLT, SLTU, BEQ, BNE, BLT, BGE, LW, SW, and
dozens more -- expressed without writing a single arithmetic constraint by hand.
Every instruction is a table lookup. The ``constraint system'' is a collection
of tables plus the Lasso machinery to prove that every instruction's result
appears in the correct table. No custom gates. No selector polynomials. No
hand-optimized constraint layouts. Just tables.

To see how this works for a specific instruction, trace through a 32-bit ADD.
The prover needs to prove that register\_a + register\_b = register\_c.~The
``addition table'' for 32-bit inputs has \(2^{64}\) entries -- impossibly large
to store. But Lasso decomposes each 32-bit operand into \(c = 4\) chunks of 8
bits each. Each chunk lookup goes into a subtable of size \(2^{16} = 65{,}536\)
entries (all possible 8-bit additions, accounting for carry). The prover
performs 4 small lookups instead of one impossible lookup, commits to the chunk
values and the carry bits, and uses the Lasso sumcheck machinery to prove that
the chunks reconstruct the full addition correctly.

Compare this to how a traditional zkVM would prove the same 32-bit ADD. In RISC
Zero's earlier architecture, the prover would encode the addition as a
polynomial constraint over the full 32-bit values, with range checks to ensure
the operands fit in 32 bits (costing roughly 32 constraints for bit
decomposition per operand), a constraint for the addition itself, and further
constraints for carry propagation and overflow detection. Roughly 70 to 100
constraints per ADD instruction. In Jolt, the same instruction costs
approximately 18 field element commitments (3 per chunk, 6 chunks for 64-bit
RISC-V) and a handful of sumcheck rounds. The constraint count per instruction
drops by roughly 4x.

This is genuinely surprising. For years, the ZK community assumed that building
a practical zkVM required painstaking constraint engineering -- hand-crafting
gate designs for each instruction type, optimizing selector layouts, minimizing
constraint counts through algebraic tricks. Jolt demonstrates that all of that
complexity can be replaced by a single, uniform mechanism: look up the answer.
The engineering effort shifts from ``design clever constraints'' to ``design
decomposable tables,'' and the latter turns out to be systematically easier.

\subsection{The Genealogy in Full}\label{the-genealogy-in-full}

The progression from auxiliary technique to primary computation paradigm took
just three years:

\begin{longtable}[]{@{}
  >{\raggedright\arraybackslash}p{(\linewidth - 10\tabcolsep) * \real{0.1143}}
  >{\raggedright\arraybackslash}p{(\linewidth - 10\tabcolsep) * \real{0.0857}}
  >{\raggedright\arraybackslash}p{(\linewidth - 10\tabcolsep) * \real{0.1571}}
  >{\raggedright\arraybackslash}p{(\linewidth - 10\tabcolsep) * \real{0.1429}}
  >{\raggedright\arraybackslash}p{(\linewidth - 10\tabcolsep) * \real{0.2714}}
  >{\raggedright\arraybackslash}p{(\linewidth - 10\tabcolsep) * \real{0.2286}}@{}}
\toprule\noalign{}
\begin{minipage}[b]{\linewidth}\raggedright
System
\end{minipage} & \begin{minipage}[b]{\linewidth}\raggedright
Year
\end{minipage} & \begin{minipage}[b]{\linewidth}\raggedright
Technique
\end{minipage} & \begin{minipage}[b]{\linewidth}\raggedright
Sorting?
\end{minipage} & \begin{minipage}[b]{\linewidth}\raggedright
Table Size Limit
\end{minipage} & \begin{minipage}[b]{\linewidth}\raggedright
Key Innovation
\end{minipage} \\
\midrule\noalign{}
\arrayrulecolor{MidnightBlue}\midrule\arrayrulecolor{MidnightBorder}
\endhead
\bottomrule\noalign{}
\endlastfoot
Plookup & 2020 & Grand product & Yes (\(O(n \log n)\)) & Fixed, materializable &
First practical lookup \\
LogUp & 2022 & Logarithmic derivatives & No & Fixed, materializable & Sum-based,
parallelizable \\
LogUp-GKR & 2023 & LogUp + GKR & No & Fixed, materializable & Logarithmic
verifier \\
Lasso & 2023 & Decomposition & No & Unlimited (structured) & Table-size
independent \\
Jolt & 2023 & Lasso for full ISA & No & Unlimited (structured) & Lookup
singularity realized \\
\end{longtable}

The three-year progression is worth pausing on. In 2020, lookups were an
auxiliary optimization -- useful for range checks and hash functions, but
secondary to the main constraint system. By 2023, lookups had become a complete
computation paradigm -- capable of replacing the constraint system entirely for
general-purpose ISA execution. Nothing else in ZK moved this fast.

The industry has not yet fully absorbed this shift. Jolt is in alpha
(open-sourced by a16z), with key missing features including full recursion
support and GPU-accelerated proving. Production systems in 2026 still largely
use LogUp or LogUp-GKR for specific operations (range checks, hash functions)
while relying on polynomial constraints for the core computation. But the
trajectory is clear: lookups are moving from a useful optimization to the
primary computation paradigm.

One qualification is important. The lookup singularity works well for ISA-level
computation -- adding two registers, comparing values, shifting bits. For
application-specific circuits with rich algebraic structure (elliptic curve
operations, pairing computations), direct polynomial constraints remain more
efficient. The lookup approach excels for \emph{general-purpose} computation;
specialized circuits may still prefer custom constraints.

The lesson of the lookup revolution is about the relationship between
computation and verification. The classical approach to ZK asks: ``How do I
re-express this computation as polynomial constraints?'' The lookup approach
asks a different question: ``How do I prove that the answer is correct, without
re-expressing the computation at all?'' The table is a certificate of
correctness. If the answer is in the table, the answer is correct. The proof
reduces to a membership test. This conceptual shift -- from ``prove the
computation'' to ``prove membership in a table of correct answers'' -- may be
the most consequential shift in arithmetization since CCS unified the constraint
systems. It is also, notably, the idea that connects most directly to the
sumcheck protocol: LogUp's rational function identity is verified by sumcheck,
Lasso's decomposition is verified by sumcheck, and Jolt's per-instruction
correctness reduces to sumcheck instances. Lookups and sumcheck are two sides of
the same coin.

Three advances -- CCS, sumcheck, and lookups -- form a complete verification
stack. CCS provides the universal constraint format: any polynomial relation,
any degree, any structure, expressed as sums of Hadamard products. Sumcheck
provides the universal verification engine: any polynomial sum over an
exponential domain, reduced to a single-point check in logarithmic rounds.
Lookups replace the most expensive constraint-by-constraint encoding with table
references: instead of proving that a value satisfies a complex relation through
dozens of polynomial constraints, the prover demonstrates that the value appears
in a precomputed table. Together, they answer a question that was open as
recently as 2022: can we build a proof system that handles any constraint
format, verifies in near-linear time, and avoids the worst overhead of
hand-crafted constraint engineering? By 2024, the answer was yes -- and the
combination of CCS + sumcheck + lookups is the engine inside every frontier
proof system built since.

\section{The Overhead Tax: 10,000x to
50,000x}\label{the-overhead-tax-10000x-to-50000x}

We have spent this chapter describing how computation is encoded as mathematics.
Now we must confront the cost.

A computation that runs natively in 1 millisecond -- executing instructions
directly on a processor at billions of operations per second -- takes 10 to 50
seconds to prove in a zkVM. That is an overhead of 10,000x to 50,000x. Where
does this multiplier come from?

The answer is not a single bottleneck but three interlocking sources of overhead
that multiply together. Each source has its own physics, its own improvement
trajectory, and its own fundamental limits. Understanding the decomposition is
essential because it determines which engineering improvements will matter most
for which applications.

\subsection{Source 1: Field Arithmetic
Encoding}\label{source-1-field-arithmetic-encoding}

Native computation uses 32-bit or 64-bit integers with hardware-accelerated
arithmetic. A single addition or multiplication takes one CPU cycle -- roughly
0.3 nanoseconds on a modern processor.

ZK computation uses finite field arithmetic. If the proof system requires a
254-bit prime field (as BLS12-381 and BN254 do), every ``addition'' becomes
multi-precision arithmetic over four 64-bit machine words. Each field
multiplication requires a Barrett or Montgomery reduction. The per-operation
cost is roughly 10 to 100 times higher than native integer arithmetic, depending
on the field size.

The small-field revolution -- BabyBear (31-bit), Mersenne-31 (31-bit),
Goldilocks (64-bit) -- addresses this directly. A 31-bit field element fits in a
single 32-bit register; a 64-bit Goldilocks element fits in a single machine
word. Arithmetic in these fields runs 10x to 100x faster than in 254-bit fields.
But even in the smallest fields, the overhead of field arithmetic versus native
integer operations remains significant.

The Mersenne-31 field is particularly elegant. Its modulus is \(2^{31} - 1\),
which is a Mersenne prime. Reduction modulo \(2^{31} - 1\) is exceptionally
fast: after computing the product of two 31-bit numbers (yielding a 62-bit
result), you split the result into a high part and a low part at bit position
31, and add them together. If the sum exceeds \(2^{31} - 1\), subtract once. Two
shifts and two additions -- faster than any other prime modular reduction. This
is why SP1 and Stwo chose Mersenne-31 as their base field: the per-operation
cost is nearly as fast as native 32-bit integer arithmetic, closing the gap
between ``native'' and ``field'' computation to a factor of roughly 3x to 5x.

\subsection{Source 2: Constraint Expansion}\label{source-2-constraint-expansion}

A single native instruction (say, a 64-bit addition) becomes multiple
constraints in the arithmetized form. The addition itself is one constraint, but
proving that the operands are within the correct range (range checks), that the
memory was read correctly (memory consistency), and that the instruction was
selected properly (opcode decoding) can require dozens to hundreds of additional
constraints. This multiplicative blowup is the most counter-intuitive aspect of
arithmetization: the ``interesting'' computation (the actual addition) accounts
for a tiny fraction of the total constraint count. The vast majority of
constraints are devoted to proving that the computational environment is
correctly maintained -- that registers hold the right values, that memory is
consistent, that the instruction pointer advanced correctly.

To make this tangible, here is a rough breakdown of the constraints required to
prove a single ADD instruction in a typical zkVM (based on published analyses of
RISC Zero and SP1 architectures):

\begin{itemize}
\tightlist
\item
  \textbf{Instruction decode}: the prover must prove that the opcode field of
  the instruction word equals the ADD opcode. This requires extracting bit
  fields and constraining them -- roughly 5 to 10 constraints.
\item
  \textbf{Register read}: the prover must prove it correctly read the values of
  the source registers (rs1 and rs2) from the register file. Memory consistency
  for each read requires 3 to 5 constraints (depending on the memory-checking
  technique).
\item
  \textbf{Arithmetic operation}: the actual addition is 1 constraint.
\item
  \textbf{Overflow handling}: the result must be reduced modulo \(2^{32}\) (for
  32-bit RISC-V). This requires proving that the result fits in 32 bits -- a
  range check costing 32 constraints (one per bit) in naive approaches, or fewer
  with lookup-based range checks.
\item
  \textbf{Register write}: the prover must prove it correctly wrote the result
  to the destination register (rd). Another 3 to 5 memory consistency
  constraints.
\item
  \textbf{Program counter update}: the prover must prove that the PC advanced by
  4 (the size of one instruction). 2 to 3 constraints.
\end{itemize}

Total: roughly 50 to 80 constraints for a single ADD instruction that, on a real
CPU, takes one clock cycle. The arithmetic itself (the actual addition) accounts
for exactly 1 of those constraints. The other 49 to 79 are bookkeeping: proving
that the right values were read from the right places, that the result was
stored correctly, and that the instruction was decoded properly. This
bookkeeping overhead is the dominant cost in constraint expansion.

Memory consistency is particularly expensive. In a native processor, reading
from memory is a single operation -- the cache or main memory returns the value,
and the hardware guarantees it is the value that was most recently written to
that address. The CPU does not need to \emph{prove} this; the hardware enforces
it physically. In a constraint system, there is no hardware. The prover claims
``I read value 42 from address 0x1000,'' and the verifier has no way to check
this claim without a mathematical argument.

Classical approaches use Merkle tree hashing -- roughly 300 multiplication
constraints per Poseidon hash invocation -- to authenticate every memory access.
The idea: maintain a Merkle tree over the entire memory state. Before each read,
prove that the value at the target address is consistent with the Merkle root.
After each write, update the Merkle tree and prove the new root is correct. Each
access requires a Merkle proof (log-depth hash chain), and each hash costs
hundreds of constraints. For a program with millions of memory accesses, this
becomes the dominant cost.

Ozdemir and others have demonstrated algebraic approaches that reduce memory
checking costs by 50x to 150x by replacing Merkle proofs with ``offline memory
checking'' -- a technique where the prover accumulates a fingerprint of all
reads and writes, and the verifier checks that the fingerprint is consistent at
the end. No per-access Merkle proofs, just a global consistency check. This is
the approach used by Jolt and by SP1's latest architecture. But even the
improved methods add substantial overhead compared to native memory access,
which costs exactly zero proof constraints.

\subsection{Source 3: Polynomial
Commitment}\label{source-3-polynomial-commitment}

After the constraints are constructed, the prover must commit to the polynomial
representations and prove their evaluations. The commitment step is where the
``zero-knowledge'' part happens: the prover seals the polynomials into
cryptographic commitments that reveal nothing about the underlying values, then
proves properties of the committed polynomials without opening them. This is
also, by far, the most computationally expensive step.

The commitment process typically involves Number Theoretic Transforms (NTTs) --
the finite-field analog of the Fast Fourier Transform -- which can consume up to
90\% of GPU proving time. Multi-scalar multiplications (MSMs) for group-based
commitments add further cost. For FRI-based systems (used with STARKs), the
commitment involves building Merkle trees over polynomial evaluations and
performing multiple rounds of degree-halving with random challenges. For
KZG-based systems (used with PLONK and Groth16), the commitment involves
computing elliptic curve group operations -- MSMs of size proportional to the
polynomial degree.

The three sources multiply. If field encoding costs 10x, constraint expansion
costs 50x, and polynomial commitment costs 10x, the total overhead is not
\(70\times\) but \(5{,}000\times\). This multiplicative composition is why the
overhead is so large and why improvements in any single source yield modest
overall gains. Reducing field encoding overhead by 10x (moving from 254-bit to
31-bit fields) and reducing constraint expansion by 2x (using lookup-based
approaches) yields a combined 20x improvement -- meaningful, but still leaving a
250x overhead. All three sources must be attacked simultaneously.

\subsection{Is the Overhead Fundamental?}\label{is-the-overhead-fundamental}

No.~Every source of overhead is under active attack by engineering and
mathematical innovation:

\begin{itemize}
\tightlist
\item
  \textbf{Field size}: The shift from 254-bit to 31-bit fields reduced
  per-operation cost by roughly 100x. A 31-bit field multiplication is a single
  machine word multiply followed by a modular reduction -- roughly 2 to 3 clock
  cycles versus the 50 to 100 cycles required for a 254-bit Montgomery
  multiplication.
\item
  \textbf{Memory checking}: Algebraic memory checking (Ozdemir et al.) reduces
  overhead by 50-150x versus Merkle-based approaches. Instead of hashing at
  every memory access, algebraic techniques use offline fingerprinting --
  accumulating a running product that can be checked at the end of the
  execution.
\item
  \textbf{Bit-level encoding}: Binius (Irreducible, 2025) reduces embedding
  overhead by 100x for bit-heavy workloads by working directly over binary tower
  fields, where a single bit \emph{is} a field element (no embedding required).
\item
  \textbf{Hardware acceleration}: GPU-based provers (BatchZK, ZKProphet) achieve
  throughput improvements of 10x to 100x through massive parallelism in NTT and
  MSM computations.
\item
  \textbf{Lookup-based architectures}: Jolt eliminates many constraint expansion
  costs by replacing polynomial constraints with table lookups. The 50-80
  constraints per instruction in a traditional zkVM drop to roughly 18 field
  element commitments per instruction.
\item
  \textbf{Folding schemes}: Nova, HyperNova, and Neo amortize the cost of
  proving many similar statements by ``folding'' them into a single accumulated
  instance. Instead of proving each step independently, the prover maintains a
  running accumulation that grows by a constant amount per step.
\end{itemize}

The cumulative effect is multiplicative. If small fields give 100x, algebraic
memory checking gives 50x, and GPU acceleration gives 10x, the combined
improvement is not 160x but potentially 50,000x -- enough to close much of the
gap between native and proven computation. The catch is that these improvements
compound only if they apply to the same bottleneck; in practice, eliminating one
bottleneck exposes the next. But the engineering trajectory points down.

\subsection{What the Overhead Feels Like in
Practice}\label{what-the-overhead-feels-like-in-practice}

Abstract multipliers are hard to internalize. Here are three concrete examples
that reveal the texture of the overhead tax.

\textbf{The single addition.} A function that adds two 64-bit integers takes
approximately 1 nanosecond on a modern CPU -- one clock cycle, one instruction,
done. The same addition, proven in zero knowledge, requires: encoding the
addition as a field operation (the 64-bit integer must be represented as one or
more field elements, with range checks to prove it fits in 64 bits), committing
to the input and output values (a polynomial commitment involving a multi-scalar
multiplication or hash-based Merkle path), and generating a proof that the
commitment is consistent with the constraint (running the full SNARK or STARK
prover pipeline). Total time: 10 to 100 microseconds, depending on the proof
system. Overhead: 10,000x to 100,000x. A single addition -- the simplest
possible computation -- pays the full fixed cost of the proof machinery.

\textbf{The Ethereum block.} An Ethereum block execution takes roughly 100
milliseconds of native computation: verifying signatures, executing smart
contract bytecode, updating the state trie. Proving the same block in a zkEVM
takes 6 to 35 seconds on GPU clusters (as reported by SP1, RISC Zero, and
Succinct in 2025 benchmarks). Overhead: 60x to 350x. This is far less than the
theoretical 10,000x to 50,000x because GPU parallelism and algorithmic
optimizations have eaten most of the overhead for large computations. The NTTs
and MSMs that dominate prover time are embarrassingly parallel -- they decompose
into millions of independent operations that map naturally onto GPU
architectures with thousands of cores.

\textbf{The Midnight transaction.} A Midnight shielded transfer involves a
Compact smart contract that reads and updates token balances behind
zero-knowledge proofs. The native computation -- checking a balance, subtracting
from one account, adding to another -- would take a few microseconds in any
programming language. Proving the same transaction in Midnight's ZK pipeline
takes approximately 20 seconds: the Compact compiler produces ZKIR instructions,
the backend lowers them to PLONKish constraints over BLS12-381, and the
Halo2-style prover generates the proof. Overhead: roughly 1,000,000x if measured
against the pure arithmetic of balance updates. But the comparison is
misleading, because the proof is doing far more than the arithmetic -- it is
proving that the balance update is consistent with the entire ledger state, that
the sender has sufficient funds, that the nullifier has not been previously
spent, and that the cryptographic commitments are correctly formed. The
``computation'' being proved is not the balance update itself but the entire
integrity argument surrounding it.

\textbf{The gap between the three.} A single addition suffers 10,000x to
100,000x overhead. An Ethereum block suffers 60x to 350x. A Midnight transaction
suffers a nominal 1,000,000x on the pure arithmetic but a more reasonable 1,000x
when measured against the full security computation it replaces. The difference
is not a measurement error. It reflects a fundamental asymmetry in the cost
structure: zero-knowledge proving has large fixed costs (setting up the
polynomial commitment, running the Fiat-Shamir transcript, computing the proof)
and relatively small marginal costs per additional constraint. A single addition
amortizes those fixed costs over one operation. An Ethereum block -- with
millions of constraints -- amortizes them over millions. The per-constraint
overhead might be identical, but the ratio of total proving time to native
execution time drops as the computation grows.

A table makes the pattern visible:

\begin{longtable}[]{@{}
  >{\raggedright\arraybackslash}p{(\linewidth - 8\tabcolsep) * \real{0.2453}}
  >{\raggedright\arraybackslash}p{(\linewidth - 8\tabcolsep) * \real{0.2453}}
  >{\raggedright\arraybackslash}p{(\linewidth - 8\tabcolsep) * \real{0.2264}}
  >{\raggedright\arraybackslash}p{(\linewidth - 8\tabcolsep) * \real{0.1887}}
  >{\raggedright\arraybackslash}p{(\linewidth - 8\tabcolsep) * \real{0.0943}}@{}}
\toprule\noalign{}
\begin{minipage}[b]{\linewidth}\raggedright
Computation
\end{minipage} & \begin{minipage}[b]{\linewidth}\raggedright
Native time
\end{minipage} & \begin{minipage}[b]{\linewidth}\raggedright
Proof time
\end{minipage} & \begin{minipage}[b]{\linewidth}\raggedright
Overhead
\end{minipage} & \begin{minipage}[b]{\linewidth}\raggedright
Why
\end{minipage} \\
\midrule\noalign{}
\arrayrulecolor{MidnightBlue}\midrule\arrayrulecolor{MidnightBorder}
\endhead
\bottomrule\noalign{}
\endlastfoot
Single 64-bit addition & \textasciitilde1 ns & 10-100 us & 10,000-100,000x &
Fixed costs dominate \\
SHA-256 hash (one block) & \textasciitilde300 ns & 1-10 ms & 3,000-30,000x &
Constraint expansion for bitwise ops \\
Ethereum block execution & \textasciitilde100 ms & 6-35 s & 60-350x & GPU
parallelism amortizes fixed costs \\
Midnight shielded transfer & \textasciitilde5 us (arithmetic) &
\textasciitilde20 s & \textasciitilde4,000,000x (arithmetic) & Large
cryptographic circuit, BLS12-381 field \\
\end{longtable}

The key insight: overhead is not uniform. Small computations suffer
disproportionately. Large computations amortize the fixed costs. Computations
over large fields (BLS12-381 at 254 bits) pay more per operation than those over
small fields (BabyBear at 31 bits). And computations that are inherently bitwise
(hashes, comparisons) pay more than those that are inherently arithmetic (field
operations, polynomial evaluations).

This is why zkVMs are viable for block-level proving (where the overhead is 100x
to 500x, manageable with GPU clusters) but impractical for individual function
calls (where the overhead is 10,000x to 100,000x, making a 1-microsecond
function take 100 milliseconds to prove). The economics of zero-knowledge
computation favor batching -- proving large computations in bulk rather than
small computations one at a time.

Where does the time actually go for a block-level proof? In a typical GPU-based
zkEVM prover (SP1, RISC Zero, or similar systems as measured in 2025
benchmarks), the breakdown looks roughly like this:

\begin{itemize}
\tightlist
\item
  \textbf{NTT (Number Theoretic Transform):} 40-60\% of total proving time. The
  finite-field analog of the FFT, used to convert polynomials between
  coefficient and evaluation representations. These are the workhorses of
  polynomial commitment. An NTT of size \(2^{24}\) involves roughly
  \(24 \cdot 2^{24} = 400\) million field multiplications -- each of which, even
  in a fast 31-bit field, takes a few nanoseconds.
\item
  \textbf{Polynomial commitment (MSM or hash-based):} 15-30\% of total proving
  time. For KZG-based systems, multi-scalar multiplications over elliptic
  curves. For FRI-based systems, Merkle tree construction over hash evaluations.
  FRI commitments require hashing the polynomial evaluations into a Merkle tree,
  then performing multiple rounds of folding (each requiring NTTs of decreasing
  size) and opening consistency proofs.
\item
  \textbf{Witness generation and constraint evaluation:} 10-20\% of total
  proving time. Filling in the execution trace and checking that all constraints
  are satisfied. This is the ``spreadsheet'' work: computing the values for
  every cell and verifying that every rule holds.
\item
  \textbf{Memory and communication overhead:} 5-15\%. Moving data between CPU
  and GPU, allocating buffers, serializing proof elements. For large proofs, the
  witness can be several gigabytes, and transferring it across the PCIe bus
  takes non-trivial time.
\end{itemize}

The dominance of NTT explains why GPU parallelism helps so much: NTTs decompose
into independent butterfly operations that map directly onto GPU warp-level
parallelism. A single NVIDIA A100 GPU can perform NTTs over \(2^{24}\) field
elements in under 100 milliseconds -- a task that would take several seconds on
a CPU. The algorithmic improvements from 2024 to 2026 (circle STARKs, WHIR,
lattice-based schemes that avoid NTTs entirely) are systematically attacking
this bottleneck.

One emerging approach avoids NTTs entirely. Sumcheck-based proof systems (like
Spartan, HyperNova, and Jolt) work with multilinear polynomials over the Boolean
hypercube rather than univariate polynomials over multiplicative subgroups.
Multilinear polynomials do not require NTTs for evaluation or commitment.
Instead, the prover performs structured summations -- which decompose into
independent, parallel operations without the butterfly dependency pattern of
NTTs. This is why the sumcheck-based architectures are gaining ground: they
eliminate the NTT bottleneck entirely, replacing it with a computation pattern
that is even more GPU-friendly.

This asymmetry also explains why recursive proof composition matters so much. If
you can batch thousands of small proofs into one large proof, and then prove the
large proof recursively, you move the computation into the regime where
amortization works in your favor. The fixed costs are paid once; the marginal
costs scale linearly. Recursion is an overhead-amortization strategy, not merely
a proof-size optimization.

The implications for system design are immediate. If you are building a zkVM for
general-purpose computation, you should optimize for large batch sizes: prove an
entire block at once, not individual transactions. If you are building a
privacy-preserving application (like Midnight), where each transaction requires
its own proof, you should invest in reducing the fixed costs: smaller fields,
faster commitment schemes, and more efficient constraint systems. The overhead
tax is not one number. It is a function of computation size, field choice,
constraint system, and proof system -- and the design space offers different
tradeoffs for different applications.

The 10,000-50,000x overhead of 2024 is not a permanent feature of provable
computation. It is the current state of a rapidly improving engineering
frontier. A reasonable projection is that overhead will decrease to 1,000-5,000x
within two to three years for general-purpose zkVMs, with application-specific
circuits already achieving lower ratios. Whether it can ever approach 100x or
below for general computation remains an open research question.

The trajectory is visible in the benchmarks. In 2022, proving a single Ethereum
block took minutes on specialized hardware. By 2024, it took 30 to 60 seconds on
GPU clusters. By early 2026, the fastest systems (SP1 Hypercube, RISC Zero 1.0)
demonstrate 6 to 15 seconds for the same workload, with further improvements
expected as circle STARKs, WHIR, and lattice-based commitments reach production
maturity. Each generation of improvements comes from a different source: the
move from 254-bit to 31-bit fields (2022-2023), the adoption of LogUp-GKR for
lookups (2023-2024), the shift to sumcheck-based architectures (2024-2025), and
GPU kernel optimization for NTTs and MSMs (ongoing). The overhead is falling not
because of one breakthrough but because of compounding engineering progress
across every layer of the stack.

For architects comparing systems, the following table normalizes the overhead by
system and field, using Ethereum block proving as the benchmark workload:

\begin{longtable}[]{@{}
  >{\raggedright\arraybackslash}p{(\linewidth - 10\tabcolsep) * \real{0.1081}}
  >{\raggedright\arraybackslash}p{(\linewidth - 10\tabcolsep) * \real{0.1486}}
  >{\raggedright\arraybackslash}p{(\linewidth - 10\tabcolsep) * \real{0.2162}}
  >{\raggedright\arraybackslash}p{(\linewidth - 10\tabcolsep) * \real{0.2297}}
  >{\raggedright\arraybackslash}p{(\linewidth - 10\tabcolsep) * \real{0.0811}}
  >{\raggedright\arraybackslash}p{(\linewidth - 10\tabcolsep) * \real{0.2162}}@{}}
\toprule\noalign{}
\begin{minipage}[b]{\linewidth}\raggedright
System
\end{minipage} & \begin{minipage}[b]{\linewidth}\raggedright
Base Field
\end{minipage} & \begin{minipage}[b]{\linewidth}\raggedright
Eth Block Time
\end{minipage} & \begin{minipage}[b]{\linewidth}\raggedright
Approx. Overhead
\end{minipage} & \begin{minipage}[b]{\linewidth}\raggedright
Year
\end{minipage} & \begin{minipage}[b]{\linewidth}\raggedright
Key Innovation
\end{minipage} \\
\midrule\noalign{}
\arrayrulecolor{MidnightBlue}\midrule\arrayrulecolor{MidnightBorder}
\endhead
\bottomrule\noalign{}
\endlastfoot
RISC Zero (v0.x) & BN254 (254-bit) & \textasciitilde60 s (GPU) &
\textasciitilde50,000x & 2023 & First general-purpose zkVM \\
SP1 (v1) & BabyBear (31-bit) & \textasciitilde15 s (16 GPU) &
\textasciitilde10,000x & 2024 & Small-field + multilinear STARK \\
SP1 Hypercube & BabyBear (31-bit) & 6.9 s (16 GPU) & \textasciitilde5,000x &
2025 & Sumcheck + precompiles \\
Stwo & Mersenne-31 (31-bit) & \textasciitilde10 s (cluster) &
\textasciitilde3,000-5,000x & 2025 & Circle STARK + 940x vs.~Stone \\
Airbender & BabyBear (31-bit) & \textasciitilde35 s (1 H100) &
\textasciitilde8,000x & 2025 & Single-GPU design \\
\end{longtable}

\section{Midnight's ZKIR: A Concrete Layer
4}\label{midnights-zkir-a-concrete-layer-4}

Abstract discussions of constraint systems benefit from a concrete example.
Midnight's ZKIR (Zero-Knowledge Intermediate Representation) provides one -- and
it reveals that real-world arithmetization carries more structure than the
mathematical formalism might suggest.

\subsection{The 24-Opcode DAG}\label{the-24-opcode-dag}

ZKIR is not itself a constraint system. It is a typed instruction-level
intermediate representation that sits \emph{above} the constraint system in the
compilation stack:

\begin{verbatim}
Compact (source language)
    |
    v
Compact IR (typed AST)
    |
    v
ZKIR (instruction DAG)  <-- This is what we are examining
    |
    v
PLONKish constraints (Halo2-style)
    |
    v
ZK proof (over BLS12-381)
\end{verbatim}

A ZKIR circuit is a directed acyclic graph of 24 base instructions organized
into eight categories:

\begin{itemize}
\tightlist
\item
  \textbf{Arithmetic} (3 opcodes): \texttt{add}, \texttt{mul}, \texttt{neg} --
  basic field arithmetic modulo the BLS12-381 scalar field (approximately
  \(2^{253}\)). There is no subtraction opcode; the compiler implements a - b as
  add(a, neg(b)).
\item
  \textbf{Constraints} (4 opcodes): \texttt{assert}, \texttt{constrain\_eq},
  \texttt{constrain\_bits}, \texttt{constrain\_to\_boolean} -- the enforcement
  mechanism.
\item
  \textbf{Comparison} (2 opcodes): \texttt{test\_eq}, \texttt{less\_than} --
  produce boolean results without enforcing them.
\item
  \textbf{Control flow} (2 opcodes): \texttt{cond\_select}, \texttt{copy} --
  conditional multiplexing and variable aliasing.
\item
  \textbf{Type encoding} (3 opcodes): \texttt{reconstitute\_field},
  \texttt{encode}, \texttt{decode} -- type-level serialization.
\item
  \textbf{Division} (1 opcode): \texttt{div\_mod\_power\_of\_two} --
  integer-style division for byte extraction.
\item
  \textbf{Cryptographic} (5 opcodes): \texttt{transient\_hash},
  \texttt{persistent\_hash}, \texttt{ec\_mul\_generator}, \texttt{ec\_mul},
  \texttt{hash\_to\_curve} -- elliptic curve and hash operations over the Jubjub
  curve embedded in BLS12-381.
\item
  \textbf{I/O} (4 opcodes): \texttt{private\_input}, \texttt{public\_input},
  \texttt{output}, \texttt{impact} -- the boundary between circuit and ledger
  state.
\end{itemize}

Each instruction consumes inputs (field elements or references to earlier
instruction outputs) and produces outputs. The DAG structure emerges from data
dependencies: instruction i depends on instruction j if it references j's
output. Variables are numbered sequentially (0, 1, 2, \ldots), and instructions
can only reference outputs of earlier instructions.

\subsection{constrain\_eq and constrain\_bits: The Enforcement
Backbone}\label{constrain_eq-and-constrain_bits-the-enforcement-backbone}

Two opcodes embody the fundamental challenge of arithmetization: ensuring that
abstract mathematical objects faithfully represent concrete computational
values.

\textbf{constrain\_eq} enforces that two field elements are identical. It
produces no output. If the values differ, the circuit rejects. This is the
fundamental correctness enforcement mechanism -- it appears after computations
to verify results, in transcript verification to bind circuit values to on-chain
state, and as the implicit check inside assertions.

There is a critical distinction between \texttt{constrain\_eq} (which
\emph{enforces} equality and fails the circuit if violated) and
\texttt{test\_eq} (which \emph{produces} a boolean result without enforcement).
The Compact compiler uses \texttt{test\_eq} for equality comparisons in program
logic and \texttt{constrain\_eq} for internal correctness checks. Confusing the
two is precisely the kind of constraint error that causes the under-constrained
vulnerabilities discussed in Chapter 3.

\textbf{constrain\_bits} enforces that a field element lies within a range
\([0, 2^N - 1]\). This is essential because ZKIR values are elements of the
BLS12-381 scalar field -- numbers up to approximately \(2^{253}\). But Compact
types often have bounded ranges: \texttt{Uint\textless{}8\textgreater{}} must be
in {[}0, 255{]}, \texttt{Uint\textless{}32\textgreater{}} in
\([0, 2^{32} - 1]\), \texttt{Boolean} must be exactly 0 or 1.

Without \texttt{constrain\_bits}, a malicious prover could substitute any
\(253\)-bit field element where an \(8\)-bit value was expected. If a circuit
adds two \texttt{Uint\textless{}8\textgreater{}} values, the honest result is at
most 510 -- but without range checking, a prover could claim the result is an
arbitrary 253-bit number, potentially extracting value or corrupting state.
Every \texttt{Uint\textless{}N\textgreater{}} value in compiled Compact code
includes a corresponding \texttt{constrain\_bits} to enforce the range
constraint.

A general principle is at work: in constraint systems, everything that is
\emph{not} explicitly constrained is implicitly \emph{allowed}. The prover will
satisfy exactly the constraints you write, and nothing more. If you forget a
constraint, the prover is free to exploit the gap. This is why under-constrained
circuits are the dominant failure mode in ZK systems.

\subsection{Where ZKIR Sits in the
Taxonomy}\label{where-zkir-sits-in-the-taxonomy}

ZKIR's relationship to the standard constraint system taxonomy is instructive:

\begin{longtable}[]{@{}
  >{\raggedright\arraybackslash}p{(\linewidth - 10\tabcolsep) * \real{0.2381}}
  >{\raggedright\arraybackslash}p{(\linewidth - 10\tabcolsep) * \real{0.1429}}
  >{\raggedright\arraybackslash}p{(\linewidth - 10\tabcolsep) * \real{0.1190}}
  >{\raggedright\arraybackslash}p{(\linewidth - 10\tabcolsep) * \real{0.2381}}
  >{\raggedright\arraybackslash}p{(\linewidth - 10\tabcolsep) * \real{0.1190}}
  >{\raggedright\arraybackslash}p{(\linewidth - 10\tabcolsep) * \real{0.1429}}@{}}
\toprule\noalign{}
\begin{minipage}[b]{\linewidth}\raggedright
Property
\end{minipage} & \begin{minipage}[b]{\linewidth}\raggedright
R1CS
\end{minipage} & \begin{minipage}[b]{\linewidth}\raggedright
AIR
\end{minipage} & \begin{minipage}[b]{\linewidth}\raggedright
PLONKish
\end{minipage} & \begin{minipage}[b]{\linewidth}\raggedright
CCS
\end{minipage} & \begin{minipage}[b]{\linewidth}\raggedright
ZKIR
\end{minipage} \\
\midrule\noalign{}
\arrayrulecolor{MidnightBlue}\midrule\arrayrulecolor{MidnightBorder}
\endhead
\bottomrule\noalign{}
\endlastfoot
Basic unit & Rank-1 constraint & Transition polynomial & Custom gate + wiring &
Matrix-vector product & Typed instruction \\
Structure & Flat constraint list & Uniform trace & Gate array + permutation &
Matrix equation & DAG of instructions \\
Abstraction level & Low & Low & Medium & Medium & \textbf{High} \\
Proof system binding & Groth16, Spartan & STARKs & Halo2, PLONK & Any IOP &
PLONKish (via backend) \\
\end{longtable}

ZKIR is not a competitor to R1CS, AIR, PLONKish, or CCS. It operates at a higher
abstraction level. Each ZKIR opcode \emph{generates} one or more underlying
PLONKish constraints: \texttt{add} generates an addition gate, \texttt{mul}
generates a multiplication gate, \texttt{constrain\_bits} generates range-check
constraints (potentially many gates for N-bit range), \texttt{ec\_mul} generates
a full scalar multiplication circuit (many internal gates),
\texttt{persistent\_hash} generates a hash circuit (many internal gates).

The ZKIR-to-PLONKish lowering is handled by the proof system backend. ZKIR
documents the \emph{semantic} layer -- what the circuit means. The actual
arithmetization into PLONK gates happens below, invisible to the Compact
developer.

Midnight's design philosophy becomes clear at this boundary. In Circom, the
developer writes constraints directly. In SP1, the zkVM generates constraints
automatically from the RISC-V execution trace. In Midnight, the Compact compiler
produces ZKIR instructions that carry type information and semantic meaning
(including blockchain-specific operations like ledger reads and writes), and the
backend translates these into PLONKish constraints. The developer never touches
the constraint system.

A ZKIR circuit could, in principle, be lowered to CCS instead of PLONKish. The
typed instruction set would need to be decomposed into the matrix-vector product
form that CCS requires. Whether Midnight's proof system will eventually migrate
from PLONKish to CCS depends on the maturity of CCS-based proof systems and the
availability of lattice-based folding schemes for production use -- a question
that connects Layer 4 directly to the post-quantum considerations at Layer 6.

\subsection{The BLS12-381 Field
Consequence}\label{the-bls12-381-field-consequence}

ZKIR operates over the BLS12-381 scalar field: a prime of approximately
\(2^{253}\), requiring 255 bits to represent. This is roughly 4x wider than the
Goldilocks field (64-bit) used by Neo and Plonky2, and approximately 8x wider
than the BabyBear (31-bit) or Mersenne-31 (31-bit) fields used by SP1, RISC
Zero, and Stwo.

The large field is necessary for Midnight's architectural choices. BLS12-381 is
a pairing-friendly curve, enabling KZG polynomial commitments and Groth16
verification. The Jubjub twisted Edwards curve embeds natively in BLS12-381's
scalar field, enabling in-circuit elliptic curve operations for Pedersen
commitments and key derivation. The mature ecosystem (Zcash, Ethereum 2.0)
provides audited tooling.

But the large field is also the primary performance cost. Each field operation
operates on 255-bit numbers using multi-precision arithmetic, while BabyBear or
M31 operations use single-register native arithmetic. This is why Midnight's
proof generation takes the order of 20 seconds per circuit (see Chapter 6 for
exact measurements) -- acceptable for privacy-preserving blockchain
transactions, but orders of magnitude slower than what small-field STARK systems
achieve.

The field choice at Layer 6 determines the arithmetic cost at Layer 4. There is
no escaping this dependency.

\section{Where the Layers Collapse}\label{where-the-layers-collapse}

This chapter has presented arithmetization as a distinct layer. But the evidence
from real systems shows that the boundary between Layer 4 and its neighbors is
porous.

\subsection{Layers 3 and 4: Jolt's Merger}\label{layers-3-and-4-jolts-merger}

In Jolt, witness generation \emph{is} the arithmetization. Every instruction in
the execution trace is decomposed into lookups on small subtables. The
decomposition happens simultaneously with trace generation -- there is no
meaningful step where ``first you generate the witness, then you arithmetize
it.'' The two processes are fused.

This has concrete implications. When the Feynman analysis asked ``if Layers 3
and 4 collapse into one in Jolt, does the seven-layer model actually work?'',
the answer is that the model is descriptive, not prescriptive. It identifies
conceptual concerns (witness generation, constraint encoding) that are always
present, even when the implementation fuses them.

\subsection{Layers 2 and 4: Cairo's
Co-Design}\label{layers-2-and-4-cairos-co-design}

Cairo, StarkWare's ZK-native language, was designed specifically so that its
instruction set would map efficiently to AIR constraints. The ISA \emph{is} the
constraint system. Language design (Layer 2) was dictated by arithmetization
efficiency (Layer 4).

The dependency runs opposite to the top-down model the book follows. In Cairo's
case, the constraint system came first, and the language was designed to match
it. Cairo's memory model is ``write-once'' -- once a value is written to an
address, it cannot be overwritten -- because write-once memory is much cheaper
to prove in AIR constraints than read-write memory. A conventional language
designer would never choose a write-once memory model. But the constraint system
designer knows that proving memory consistency for write-once memory requires a
simple sorted-access check (much cheaper than Merkle trees or fingerprinting),
so the language was shaped to match. The pedagogical order (language before
arithmetization) hides a real engineering dependency. Cairo was not ``compiled
to'' AIR; it was ``born from'' AIR.

\subsection{The Proof Core: Layers 4, 5, and
6}\label{the-proof-core-layers-4-5-and-6}

The most significant cross-layer dependency is the ``proof core'' -- the
inseparable triad of \{finite field, polynomial commitment scheme, polynomial
representation\} that straddles Layers 4, 5, and 6.

Choose a 31-bit field (Layer 6) and you get fast arithmetic but need FRI-based
commitments (Layer 5) and AIR or multilinear representations (Layer 4). Choose a
254-bit pairing-friendly field (Layer 6) and you can use KZG commitments (Layer
5) with univariate polynomials in Lagrange basis (Layer 4). Choose a
lattice-based commitment over a 64-bit Goldilocks field (Layer 6) and you get
post-quantum security with CCS-native constraints (Layer 4) and sumcheck-based
folding (Layer 5).

These are not three independent choices. They are one choice with three
manifestations. The seven-layer model usefully separates the \emph{concerns}
(what is being encoded? how is it committed? what field operations are
available?) even when the \emph{implementations} cannot be separated.

This is the collapse we warned about in Chapter 1. The seven-layer model is a
pedagogical map. The engineering territory has three layers at the proof core,
not seven, and the edges between them are bidirectional. Hold both models as you
read: the pedagogical stack (useful for learning each concern in isolation) and
the engineering DAG (useful for building real systems). Chapter 10 will draw the
honest map -- seven nodes, fourteen directed edges, no pretense of independence.

A concrete example of this coupling: RISC Zero originally used a 254-bit field
with KZG commitments and R1CS constraints. In 2023, they migrated to BabyBear
(31-bit field) with FRI commitments and AIR constraints. The migration was not
``swap out the field and keep everything else.'' It required simultaneously
changing the field (Layer 6), the commitment scheme (Layer 5), and the
constraint format (Layer 4) -- because none of the three could be changed
independently. BabyBear does not support KZG (which needs a pairing-friendly
curve), and FRI does not work naturally with R1CS (which lacks the
evaluation-domain structure that FRI requires). The three layers moved as a
unit, confirming that the ``proof core'' is a single design decision dressed up
as three.

\section{Where the Analogies Break}\label{where-the-analogies-break}

We promised at the beginning of this chapter to say where the analogies break
down.

Arithmetization is the hardest layer to explain because it is the layer where
computer science, algebra, and information theory collide in ways that resist
simplification. The ``spreadsheet with polynomial rules'' captures the structure
but not the mechanism. The ``Sudoku puzzle'' captures the
constraint-satisfaction flavor but misleads about uniqueness. The ``encoding''
metaphor captures the transformation but hides the overhead.

What actually happens at Layer 4 is a lossy translation. A computation in the
real world involves pointers, variable-length data, exceptions, floating-point
approximation, and timing. The arithmetized version strips all of that away and
replaces it with fixed-size field elements, fixed-structure polynomial
constraints, and deterministic evaluation. The gap between the two -- the
``abstraction tax'' of 10,000x to 50,000x -- is the price of making computation
mathematically verifiable.

Consider what is lost in the translation. A native C program uses 64-bit
integers with overflow semantics (values wrap around at \(2^{64}\)). The
arithmetized version uses field elements modulo a prime -- where overflow does
not exist, because field arithmetic is always exact. To faithfully represent
64-bit overflow behavior, the constraint system must decompose the values into
64 individual bits, check that each is boolean, compute the sum, and check that
only the low 64 bits are retained. The ``overflow'' that hardware handles in
zero cycles costs dozens of constraints. Similarly, a floating-point
multiplication that the CPU executes in one cycle using a dedicated FPU requires
hundreds of constraints to simulate in field arithmetic -- because there is no
floating-point hardware in a finite field, only integers. Every gap between
native computation and field arithmetic generates constraints. The overhead is
not laziness or bad engineering. It is the cost of bridging two incompatible
computational models.

That price is falling. It fell when AIR replaced R1CS for VM-style computations.
It fell when PLONKish introduced custom gates. It fell when CCS unified the
constraint systems. It fell when LogUp eliminated sorting from lookups. It fell
when Lasso made table sizes irrelevant. It fell when small fields replaced
254-bit primes. It falls every time a researcher finds a way to encode more
computation in fewer constraints.

But it has not fallen to zero, and it may never fall to zero. Provable
computation is inherently more expensive than unprovable computation. The
magician who performs backstage with no audience can cut corners. The magician
who must produce a sealed certificate -- one that any stranger can verify --
must record every step with mathematical precision. The overhead of
arithmetization is the cost of making the performance verifiable.

There is a theoretical lower bound that clarifies the situation. Any computation
that produces n bits of output requires at least n bits of communication to
verify (you need to at least read the output). The overhead above this
information-theoretic minimum comes from the cryptographic machinery: polynomial
commitments, random challenge generation, and the proof that the polynomial
identities hold. Whether this cryptographic overhead can be reduced to \(O(1)\)
multiplicative factor remains an open question. The current answer is: not yet,
but the constant factor is shrinking every year.

The honest summary of Layer 4 in 2026: arithmetization is hard, expensive, and
getting better fast. The constraint systems are converging toward CCS. The
lookup revolution is replacing hand-crafted constraints with table lookups. The
overhead is falling from 10,000x toward 1,000x and below. And the sumcheck
protocol -- invented in 1992, long before anyone imagined practical
zero-knowledge proofs -- has become the universal verification engine that makes
it all work.

From this point forward, the magician-and-audience framing will recede. Layers
5, 6, and 7 operate at a level of abstraction where the metaphor obscures more
than it reveals. The sealed certificate, the deep craft, the verdict -- these
section titles keep the theatrical frame alive as a mnemonic, but the
explanations will be increasingly technical. We will still speak of provers and
verifiers, because those are the real actors, but the stage curtains come down
here. When the metaphor returns in full force, it will be in Chapter 12, where
Midnight provides a concrete theater that makes the abstraction physical again.

\par\vspace{18pt}

\emph{The computation is encoded as mathematics. Every instruction, every memory
access, every comparison has been transformed into polynomial equations over a
finite field. The equations are organized -- as R1CS, as AIR, as PLONKish, or as
CCS -- and the lookup arguments have replaced the most expensive operations with
table references. The sumcheck protocol stands ready to verify that the
equations hold, without checking every cell in the spreadsheet.}

\emph{The constraint system is the scorecard. But a scorecard means nothing
until someone seals it into a certificate that cannot be tampered with. The next
chapter enters the proof system -- the cryptographic mechanism that seals the
certificate.}

\subsection{Reference Data}\label{reference-data}

\begin{itemize}
\tightlist
\item
  \textbf{R1CS} (2012): bilinear constraints (degree 2). One constraint per
  multiplication gate. Native format for Groth16 (128-byte proofs, constant-time
  verification) and Spartan.
\item
  \textbf{AIR} (2018): uniform transition constraints over execution traces.
  Native format for STARKs (transparent, post-quantum). Constraint description
  size independent of trace length.
\item
  \textbf{PLONKish} (2019): selector-gated custom gates with copy constraints
  via permutation arguments. Native format for Halo2, PLONK. Dominant in
  deployed systems (2020-2025).
\item
  \textbf{CCS} (Setty, 2023): unifies R1CS, AIR, and PLONKish without overhead.
  Native target for HyperNova, Neo, ProtoStar, ProtoGalaxy.
\item
  \textbf{Sumcheck protocol} (Lund et al., 1992): reduces verification of
  polynomial sums over \(2^n\) inputs to \(n\) rounds of interaction. Backbone
  of Spartan, HyperNova, Jolt, and SP1 Hypercube.
\item
  \textbf{Plookup} (2020): first practical lookup argument. Sorting-based,
  \(O(n \log n)\).
\item
  \textbf{LogUp} (2022): sorting-free lookup via logarithmic derivatives.
  \(O(n)\) prover cost.
\item
  \textbf{LogUp-GKR} (2023): logarithmic verifier cost for lookups. Used in SP1
  Hypercube and Stwo.
\item
  \textbf{Lasso} (2023): lookups into tables of size \(2^{128}\), prover cost
  independent of table size.
\item
  \textbf{Jolt} (2023): approximately 6x faster than RISC Zero in theoretical
  commitment cost analysis. Full RISC-V ISA via lookups.
\item
  \textbf{Overhead tax}: 10,000-50,000x versus native execution (2024-2025
  systems). Falling to 1,000-5,000x by 2027-2028.
\item
  \textbf{Overhead breakdown}: field encoding (10-100x), constraint expansion
  (50-100x), polynomial commitment (10-50x). Sources multiply.
\item
  \textbf{Ozdemir et al.}: 50-150x reduction in memory checking constraints via
  algebraic approaches.
\item
  \textbf{Binius} (2025): 100x reduction in bit-level embedding overhead via
  binary tower fields.
\item
  \textbf{Mersenne-31}: field modulus \(2^{31} - 1\). Fastest known modular
  reduction. Used by SP1 and Stwo.
\item
  \textbf{ZKIR}: 24 typed instructions, compiling Compact to PLONKish
  constraints over BLS12-381 (\(\sim 2^{253}\)).
\end{itemize}

\subsection{Sources}\label{sources}

\begin{itemize}
\tightlist
\item
  {[}R-L4-1{]} Gennaro, Gentry, Parno, Raykova. ``Quadratic Span Programs and
  Succinct NIZKs without PCPs.'' EUROCRYPT 2013. ePrint 2012/215.
\item
  {[}R-L4-2{]} Ben-Sasson, Bentov, Horesh, Riabzev. ``Scalable, Transparent, and
  Post-Quantum Secure Computational Integrity.'' ePrint 2018/046.
\item
  {[}R-L4-3{]} Gabizon, Williamson, Ciobotaru. ``PLONK.'' ePrint 2019/953.
\item
  {[}R-L4-4{]} Setty, Thaler, Wahby. ``Customizable Constraint Systems for
  Succinct Arguments.'' ePrint 2023/552.
\item
  {[}R-L4-5{]} Setty. ``Spartan: Efficient and General-Purpose zkSNARKs without
  Trusted Setup.'' ePrint 2019/550.
\item
  {[}R-L4-6{]} Gabizon, Williamson. ``Plookup.'' ePrint 2020/315.
\item
  {[}R-L4-7{]} Haboeck. ``LogUp.'' ePrint 2022/1530.
\item
  {[}R-L4-8{]} Papini, Shahar and Ulrich Haboeck. ``LogUp-GKR.'' ePrint
  2023/1284.
\item
  {[}R-L4-9{]} Setty, Thaler, Wahby. ``Lasso.'' ePrint 2023/1216.
\item
  {[}R-L4-10{]} Arun, Setty, Thaler. ``Jolt.'' ePrint 2023/1217.
\item
  Midnight ZKIR Reference (v2/v3), 119 oracle traces. Compact compiler v0.29.0.
\item
  Lund, Fortnow, Karloff, Nisan. ``Algebraic Methods for Interactive Proof
  Systems.'' JCSS 1992.
\end{itemize}

\par\vspace{18pt}

\emph{A note on the next three chapters.} Chapters 5, 6, and 7 cover
arithmetization, proof systems, and cryptographic primitives -- what this book
calls the ``proof core.'' In practice, these three layers are inseparable: the
choice of field (Layer 6) determines which arithmetization works (Layer 4),
which determines which proof system is viable (Layer 5). We present them
sequentially because a book must be linear, but they are best understood as a
single coupled design unit. If a choice in Chapter 7 seems to contradict a claim
in Chapter 5, it is because the dependency runs in both directions. Read all
three, then revisit.

\chapter{Chapter 6: Layer 5 -- The Sealed
Certificate}\label{chapter-6-layer-5-the-sealed-certificate}

\section{Sealing the Certificate}\label{sealing-the-certificate}

The magician has performed the trick backstage. Every step of the computation
has been recorded in a mathematical trace and encoded as a system of polynomial
equations. But a recording that never leaves the backstage is worthless. The
audience needs something it can hold in its hands, inspect, and trust -- without
ever seeing the performance itself.

That something is a sealed certificate.

The certificate attests that the trick was performed correctly -- ``this
transaction is valid,'' ``this person is over 18,'' ``this block was executed
correctly.'' It is small enough to carry in your pocket. It is unforgeable.
Anyone can verify it. Nobody needs to see the original performance. And a forged
certificate is not merely unlikely -- it is a mathematical impossibility, its
probability shrinking exponentially as the security parameter grows.

Think of it this way. In the previous two chapters, we watched the computation
get written down (the language), performed backstage (the witness), and
translated into a mathematical puzzle (the arithmetization). Now we reach the
moment the puzzle gets sealed into a certificate that will travel out into the
world, to be checked by strangers who have no reason to trust us and no access
to our private data. This is Layer 5: the proof system. It is the mechanism that
presses the wax seal.

\begin{quote}
\textbf{The Running Example: The Sudoku Proof}

Our 72 constraints over 16 witness variables are now sealed into a certificate.
The prover commits to the constraint polynomials using the chosen commitment
scheme (KZG, FRI, or Ajtai -- depending on the architectural path from Chapter
10). The verifier sends random challenge points (or they are derived via
Fiat-Shamir from the transcript hash). The prover evaluates the committed
polynomials at those points and returns the evaluations with opening proofs.

The verifier checks: do the evaluations satisfy the constraint relationships at
the challenged points? If yes, Schwartz-Zippel guarantees that the polynomials
themselves satisfy the constraints everywhere -- except with probability at most
\(d/|\mathbb{F}|\), where \(d\) is the polynomial degree and \(|\mathbb{F}|\) is
the field size. For our 4x4 Sudoku over a 256-bit field, \(d\) is roughly 4 and
\(|\mathbb{F}|\) is roughly \(2^{256}\), so the soundness error is vanishingly
small.

The result: a Groth16 proof of 192 bytes, or a STARK proof of around 50 KB, or a
folded lattice commitment -- depending on the path. The verifier learns that the
prover knows a valid solution. The verifier learns nothing about which numbers
go where. The sixteen secret values never leave the prover's machine. The 4x4
grid that was the witness has been compressed into a handful of group elements
and field evaluations -- the sealed certificate.
\end{quote}

Here is how the seal works. The proof system commits to a set of polynomial
evaluations. If the prover cheated anywhere in the computation, those
evaluations will be inconsistent -- the polynomial will disagree with itself at
a random point the verifier picks. The probability that the prover can guess
which point the verifier will pick, and cheat in exactly the right way to pass
that specific check, is negligible -- meaning it shrinks exponentially as the
security parameter grows.

A forged certificate does not look like a convincing imitation that might fool
someone. It looks like an impossible object. The proof system is designed so
that forged certificates simply cannot be produced in polynomial time, assuming
the underlying mathematical hardness assumptions hold. This chapter explains how
the sealing mechanism works, how it evolved from a single method into a family
of techniques with radically different properties, and why the engineering
choices within Layer 5 are changing the entire economics of blockchain
computation.

\section{The Three Families}\label{the-three-families}

Every zero-knowledge proof system in production today belongs to one of three
families. They differ in what they require before you start, how large a proof
they produce, and what mathematical assumptions keep them secure. Choosing among
them is not a philosophical exercise. It is an engineering decision with direct
consequences for cost, speed, security, and quantum resistance.

\subsection{Groth16: The Gold Standard for
Size}\label{groth16-the-gold-standard-for-size}

In 2016, Jens Groth published a proof system that produces the smallest possible
proofs: three elliptic curve group elements, totaling 192 bytes (two G1 points
and one G2 point on BLS12-381). Verification requires three pairing computations
(the Ethereum EVM implementation batches this as four pairings via EIP-1108) --
about 250,000 gas on Ethereum. Nearly a decade later, nothing comes close to
this combination of proof size and verification cost. Groth16 remains the final
wrapping target for almost every production ZK system that posts proofs
on-chain.

The catch is severe. Groth16 requires a per-circuit trusted setup ceremony. You
cannot change the circuit without re-running the ceremony. The ceremony produces
a structured reference string (SRS) that contains ``toxic waste'' -- secret
randomness that, if any single participant retains, would allow them to forge
proofs for that circuit forever. The 1-of-N trust model (discussed in Chapter 2)
mitigates this, but the ceremony is expensive, inflexible, and not
quantum-resistant.

Put differently: Groth16 seals the smallest possible certificate -- just 192
bytes, three elliptic curve points. But you need a custom seal for every circuit
you want to prove, and manufacturing that seal requires a ceremony involving
thousands of people, any one of whom could secretly keep a master key to forge
future certificates. If a quantum computer ever appears, every seal ever
manufactured becomes useless.

Despite these limitations, Groth16's on-chain economics are so favorable that it
persists as the outer shell of nearly every hybrid proving pipeline. The inner
proof system might be a STARK, a folding scheme, or something else entirely. But
the last step -- the proof that actually touches the blockchain -- is almost
always Groth16 over the BN254 curve, because Ethereum's EVM has precompiled
contracts that make BN254 pairings cheap.

\subsection{PLONK: The Universal Workhorse}\label{plonk-the-universal-workhorse}

In 2019, Gabizon, Williamson, and Ciobotaru published PLONK, which solved
Groth16's biggest practical problem: the per-circuit setup. PLONK uses a
universal structured reference string. Run one ceremony, and you can use the
resulting SRS for any circuit up to a fixed size. Change your program? No new
ceremony needed. Just compile a new circuit against the same SRS.

PLONK introduced ``PLONKish'' arithmetization -- a system of gate constraints
and copy constraints glued together by a permutation argument. This turned out
to be very flexible. Custom gates allow specialized operations (elliptic curve
arithmetic, hash functions, range checks) to be encoded efficiently. Lookup
tables (via Plookup and its successors) let the prover reference pre-computed
values instead of re-deriving them from scratch. The result was a proof system
that could be tuned for specific workloads while retaining universality.

Halo 2, developed by the Electric Coin Company for Zcash and later adopted by
other projects, extended PLONK with custom gates, lookup tables, and a flexible
backend that supports multiple commitment schemes. Note: the original Halo paper
(2019) used IPA commitments with no trusted setup. Halo 2, as deployed by Zcash
for Orchard and adopted by Midnight, uses KZG commitments requiring a ceremony
-- a different trust model despite the shared name. When instantiated with IPA
instead of KZG, Halo 2 eliminates the trusted setup entirely, and the proof size
grows from constant to logarithmic, but verification becomes somewhat more
expensive.

PLONK and its variants (UltraPlonk, TurboPlonk, Halo 2) are the workhorses of
production ZK today. They sit at the center of a design space between Groth16's
extreme compactness and STARKs' extreme transparency. Most ZK applications that
need flexibility -- circuits that change frequently, applications where a
per-circuit ceremony is impractical -- choose a PLONK-family system.

\subsection{STARKs: The Transparent Path}\label{starks-the-transparent-path}

STARKs (Scalable Transparent ARguments of Knowledge), introduced by Ben-Sasson,
Bentov, Horesh, and Riabzev in 2018, took a radically different approach. No
trusted setup at all. No pairings. No elliptic curves. The only cryptographic
assumption is the existence of collision-resistant hash functions -- a primitive
that is believed to be quantum-resistant.

The core idea is clean. STARKs prove statements about Algebraic Intermediate
Representations (AIR): polynomial constraints on an execution trace, where each
row represents one time step and the constraints enforce that consecutive rows
are consistent. The key mechanism is the FRI protocol (Fast Reed-Solomon
Interactive Oracle Proof), which verifies that a committed function is close to
a low-degree polynomial. If the computation was performed correctly, the trace
satisfies all constraints, and the polynomial is low-degree. If the prover
cheated, the polynomial's degree blows up, and FRI catches the inconsistency
with overwhelming probability.

The cost of this transparency is proof size. A STARK proof is hundreds of
kilobytes -- roughly 1,000 times larger than a Groth16 proof. On Ethereum, where
every byte costs gas, this difference translates directly into money. But STARKs
offer something the other families cannot: a path to post-quantum security, and
a proving architecture that scales almost linearly with computation size.

So the three families seal certificates with different properties. Groth16 seals
the smallest possible certificate -- just 192 bytes -- but requires a custom
seal for every circuit. PLONK seals a slightly larger certificate but can use
the same seal for any trick, eliminating the per-circuit ceremony. STARKs seal
certificates without any secret ingredient at all -- no ceremony, no toxic
waste, no trust assumptions beyond collision-resistant hashing -- but the
certificates are much bulkier: hundreds of kilobytes instead of hundreds of
bytes.

\subsection{Three Envelopes: An Intuition}\label{three-envelopes-an-intuition}

The technical differences between the three families are real, but their
\emph{character} is better grasped through analogy. Each family seals a
certificate. Think of each certificate as a letter placed inside an envelope.
The families differ not in what the letter says but in how the envelope is made,
what it costs, and what it reveals about the process of sealing.

\textbf{Groth16 is the smallest possible envelope.} Three numbers -- two points
on a curve and one more -- encode an entire computation. Consider the
compression ratio. A circuit proving that an Ethereum block was executed
correctly might involve billions of constraints, millions of intermediate
values, a witness consuming gigabytes of memory. The Groth16 proof of that
computation is 192 bytes. Three elliptic curve elements. It is as though someone
handed you a novel -- a thick, sprawling saga with hundreds of characters and
interlocking subplots -- and you compressed it into a haiku. Seventeen
syllables. And yet a reader who knows the rules of haiku composition can verify
that those seventeen syllables faithfully capture the plot. Not a summary. Not
an approximation. A mathematically exact encoding from which any deviation would
be detectable. The haiku either satisfies the verification equation or it does
not. There is no room for a ``pretty good'' forgery.

The cost of this extreme compression is inflexibility. The haiku form must be
custom-designed for each novel. You cannot take a haiku mold built for \emph{War
and Peace} and use it to compress \emph{Moby Dick}. In Groth16 terms, this means
a new trusted setup ceremony for every circuit. The ceremony produces the
specific algebraic structure -- the structured reference string -- that makes
the compression possible for that particular computation. Change the
computation, and you need a new ceremony. This is why Groth16 is almost never
used as the primary proof system. It is used as the \emph{final wrapper} -- the
outermost envelope -- because you only need one ceremony for the wrapping
circuit, and that circuit does not change.

\textbf{PLONK is the universal envelope.} One envelope fits all letters. The
envelope factory runs a single setup ceremony and produces a structured
reference string that works for any circuit up to a certain size. Write a new
smart contract? Same envelope. Update your program logic? Same envelope. The
setup cost is paid once. After that, any computation that fits within the size
bound can be sealed without a new ceremony.

The analogy is a standardized shipping container. Before containerization, every
type of cargo required its own packaging, its own crane, its own dockworker
expertise. A container does not care whether it holds televisions, bananas, or
machine parts. It is a universal interface between the thing being shipped and
the infrastructure that moves it. PLONK's universal SRS is the container. The
circuit -- whatever it computes -- is the cargo. The infrastructure -- the
verifier, the blockchain, the precompiled contract -- handles every circuit the
same way, because they all arrive in the same container.

The proofs are slightly larger than Groth16's (a few hundred bytes to a few
kilobytes, depending on the variant), and verification is slightly more
expensive. These are the costs of universality. You pay a modest premium for the
ability to change your message without manufacturing a new envelope. For most
applications, this tradeoff is overwhelmingly favorable, which is why
PLONK-family systems are the workhorses of production ZK.

\textbf{STARKs are glass envelopes.} Nothing is hidden in the construction.
There is no trusted setup, no toxic waste, no secret randomness that could
compromise the system if leaked. The envelope-making process is entirely public.
Anyone can inspect it, audit it, reproduce it. The cryptographic assumption is
minimal: collision-resistant hash functions exist. That is all.

The trade: glass envelopes are bulkier than paper ones. A STARK proof is
hundreds of kilobytes -- roughly a thousand times larger than a Groth16 proof.
On a blockchain where every byte costs gas, this bulk translates directly into
dollars. Glass is heavier than paper. You pay more to ship it. But glass has a
property that paper lacks: you can see through it. The transparency is not a
metaphor. It is a literal statement about the trust model. There is no ceremony
to trust, no participant to worry about, no toxic waste to dispose of. The
security rests on the hardest-to-break foundation in cryptography: the belief
that hash functions do not have secret backdoors.

And glass has another property that matters more each year: it does not shatter
under quantum impact. Hash-based cryptography is believed to resist quantum
attacks. Paper envelopes -- those built on elliptic curve assumptions -- will
dissolve the moment a sufficiently powerful quantum computer runs Shor's
algorithm. Glass envelopes will still be standing. The bulk that seems like a
disadvantage today may prove to be the price of survival.

But notice: the three envelope types are not competing products on a shelf. They
are components in a supply chain. The glass envelope (STARK) is manufactured
first, because it is transparent and quantum-resistant. Then the contents are
transferred into a paper envelope (Groth16) for shipping, because paper is
lighter and the postal system (the blockchain) charges by weight. The glass
envelope never reaches the destination. It does its job backstage -- proving the
computation with full transparency -- and then the result is repackaged into the
smallest, cheapest container for the final mile. Understanding this supply chain
is what the next section is about.

\section{The Hybrid Pipeline}\label{the-hybrid-pipeline}

Here is the secret that the field's own marketing has obscured: in production,
the three families are not competitors. They are components of a single
pipeline.

The dominant architecture in 2025 looks like this:

\begin{enumerate}
\def\labelenumi{\arabic{enumi}.}
\item
  \textbf{Prove the computation with a STARK.} The inner proof system generates
  a large, transparent proof using fast, small-field arithmetic. No trusted
  setup is required. The prover runs on GPUs.
\item
  \textbf{Recursively compress the STARK.} Apply one or more rounds of recursive
  STARK verification to shrink the proof from hundreds of kilobytes to tens of
  kilobytes.
\item
  \textbf{Wrap the compressed STARK in a Groth16 SNARK.} A small circuit
  verifies the STARK proof and produces a Groth16 proof: 192 bytes, 250K gas to
  verify on-chain. The wrapping circuit adapts the STARK's field arithmetic
  (say, Mersenne-31) to the BN254 field that Ethereum's precompiles support.
\item
  \textbf{Post the SNARK on-chain.} The Ethereum verifier contract checks the
  Groth16 proof. It has no idea that the inner computation used a STARK. From
  the chain's perspective, everything looks the same.
\end{enumerate}

Every major production system follows this architecture: SP1 (Succinct Labs),
Stwo (StarkWare), Polygon, ZKM, and most others. The STARK generates the raw
material; the SNARK seals it into a certificate the blockchain will accept. The
STARK provides transparency and prover efficiency. The SNARK provides on-chain
cost efficiency. Each family contributes what it does best.

The implication is worth spelling out. The ``SNARK vs.~STARK'' debate that
dominated conference panels from 2019 to 2023 was a false dichotomy. The field
converged on ``STARK inside, SNARK outside'' because the engineering tradeoffs
demanded it. Transparency for the prover. Compactness for the chain. The only
remaining question is whether the outer SNARK wrapper can itself become
post-quantum -- a problem that lattice-based proof systems (discussed later in
this chapter) are beginning to address.

\subsection{A Concrete Pipeline: From 1,000 Transactions to a 192-Byte
Proof}\label{a-concrete-pipeline-from-1000-transactions-to-a-192-byte-proof}

The hybrid architecture becomes vivid when you trace a specific workload through
it, step by step. Consider a ZK rollup operator -- running on Ethereum mainnet
-- who receives a batch of 1,000 transactions. Each transaction is a token
transfer, a swap, or a contract call. The operator must prove that executing all
1,000 transactions against the current state produces the claimed new state
root. Here is how the pipeline processes that batch in 2025.

\textbf{Step 1: Execute and generate the witness.} The operator's sequencer
replays all 1,000 transactions against a local copy of the rollup's state. Every
storage read, every storage write, every arithmetic operation is recorded in an
execution trace -- the giant spreadsheet from Chapter 4. The witness includes
all private data: account balances, nonces, intermediate computation values.
This step is ordinary software execution, no cryptography involved. On a modern
server, it takes 1 to 2 seconds. The output is a trace with millions of rows and
dozens of columns.

\textbf{Step 2: Generate a STARK proof over a small field.} The prover takes the
execution trace and produces a STARK proof using BabyBear (31-bit) or
Mersenne-31 arithmetic. This is where the heavy computation happens. The trace
is interpolated into polynomials, the polynomials are committed via FRI (Merkle
trees of evaluations), and the FRI protocol verifies low-degree proximity. On a
cluster of GPUs -- say, four NVIDIA A100s -- this step takes 3 to 5 seconds. The
output is a STARK proof: transparent, hash-based, quantum-resistant, and roughly
200 to 400 kilobytes in size.

\textbf{Step 3: Recursively compress the STARK.} The raw STARK proof is too
large to post on-chain economically. So the operator generates a second STARK
proof that verifies the first one. This is recursion: a proof about a proof. The
verifier circuit for a STARK is much smaller than the original computation
circuit, so the recursive proof is faster to generate and produces a smaller
output. One or two rounds of recursive compression shrink the proof from
hundreds of kilobytes to tens of kilobytes. This step takes 1 to 2 seconds.

\textbf{Step 4: Wrap in Groth16 over BN254.} The compressed STARK proof is now
small enough to verify inside a Groth16 circuit. A specialized wrapping circuit
takes the STARK verifier computation -- check the Merkle paths, verify the FRI
folding, confirm the constraint evaluations -- and expresses it as an R1CS
instance over the BN254 field. The Groth16 prover then seals this into a
192-byte proof: two G1 points and one G2 point. This is the field-crossing step,
where 31-bit STARK arithmetic is translated into 254-bit BN254 arithmetic. It is
computationally expensive per field operation, but the circuit is small (it is
only verifying a STARK, not re-executing the original computation). On a single
GPU, this step takes 5 to 10 seconds.

\textbf{Step 5: Post the proof to Ethereum.} The operator submits a transaction
to the rollup's on-chain verifier contract. The transaction contains the
192-byte Groth16 proof, the new state root, and a commitment to the batch of
transactions. The Ethereum verifier contract calls the BN254 pairing precompile
(EIP-1108), checks the three pairings, and accepts or rejects. Verification
costs approximately 250,000 gas -- roughly \$0.50 to \$1.00 at typical gas
prices. The state root is updated. The 1,000 transactions are finalized.

The audience sees only Step 5. A 192-byte proof appears on-chain. A smart
contract checks it in a few milliseconds. The state updates. Nobody knows -- or
needs to know -- that behind those 192 bytes lie four NVIDIA GPUs, two rounds of
recursive compression, a field-crossing circuit, and the execution traces of
1,000 individual transactions. The entire pipeline, from receiving the batch to
posting the proof, completes in under 15 seconds and costs under \$1.00.

That is the hybrid pipeline in concrete terms. The STARK did the heavy lifting:
proving the computation with transparency and quantum resistance. The Groth16
wrapper did the packaging: compressing everything into the smallest possible
on-chain footprint. Each proof system contributed what it does best. The
audience -- Ethereum's verifier contract -- received the finished product and
asked no questions about the manufacturing process.

Two years earlier, the same pipeline took minutes and cost tens of dollars. Two
years before that, it was a research prototype that could not process a full
Ethereum block at all. The economics shifted not because anyone invented a
fundamentally new proof system, but because each component in the pipeline got
faster -- small-field arithmetic, GPU parallelism, recursive compression,
optimized wrapping circuits -- and the improvements compounded multiplicatively
across stages. A 3x improvement in STARK proving, combined with a 2x improvement
in recursive compression, combined with a 2x improvement in Groth16 wrapping,
produces a 12x improvement end-to-end. This is why the cost curve has been
steeper than Moore's Law.

\section{Recursion vs.~Folding: Russian Dolls and
Snowballs}\label{recursion-vs.-folding-russian-dolls-and-snowballs}

Within any proof system architecture, there is a second design choice that
matters just as much: how do you handle computation that happens in steps?

A blockchain processes blocks sequentially. A zkVM executes instructions one at
a time. A rollup batches transactions and proves them in order. In every case,
the prover needs to demonstrate not just ``this single step is correct'' but
``every step from the beginning until now was correct.'' This is Incrementally
Verifiable Computation (IVC), first formalized by Valiant in 2008, and it sits
at the heart of how proof systems scale.

\subsection{Recursion: The Russian Doll}\label{recursion-the-russian-doll}

The first approach to IVC was recursive composition. The idea is deceptively
simple: to prove that steps 1 through N are all correct, you prove step N is
correct \emph{and} that you have a valid proof that steps 1 through N-1 are
correct. The verifier for steps 1 through N-1 is itself expressed as a circuit,
and the proof for step N includes a verification of the previous proof.

Think of Russian nesting dolls. Each doll contains a smaller doll inside it, and
that smaller doll contains an even smaller one. Opening the outermost doll and
confirming it contains a properly formed inner doll gives you confidence in the
entire chain. You never need to open all the dolls at once.

The problem is cost. The SNARK verifier is a complex circuit -- for
pairing-based systems like Groth16, it involves millions of gates. Embedding a
Groth16 verifier inside a Groth16 circuit means every step of computation must
pay the cost of a full SNARK verification on top of the actual computation. Zexe
(Bowe, Chiesa, Green, Miers, Mishra, Wu, 2018) demonstrated this approach using
depth-2 recursion over a 2-chain of pairing-friendly curves, achieving
constant-size 968-byte transactions regardless of offline computation. But the
proving cost per step was enormous.

\subsection{Folding: The Snowball}\label{folding-the-snowball}

In 2022, Abhiram Kothapalli, Srinath Setty, and Ioanna Tzialla published Nova,
and the field changed direction overnight.

Nova introduced a fundamentally different approach: instead of proving each step
and recursively verifying previous proofs, you \emph{fold} two claims into one.
Here is the intuition. In recursion, each step generates a complete proof, and
the next step verifies it. In folding, each step generates a \emph{claim} -- an
incomplete, unverified assertion -- and combines it with the running claim from
all previous steps. The combination is a random linear combination: the prover
takes the two claims, the verifier picks a random challenge, and the prover
produces a single new claim that is valid if and only if both original claims
were valid. No full proof is ever generated during the accumulation phase. The
certificate is not sealed along the way. Claims accumulate, and the expensive
seal is deferred to the very end. Only when all steps have been folded together
does the prover generate a single SNARK proof for the final accumulated claim.

The snowball analogy captures it precisely. Each new step is a handful of snow.
Instead of building a separate snowman for each step (recursion), you pack the
new snow onto the growing snowball (folding). The snowball gets slightly larger
with each step, but you never have to build an entire snowman until the very
end.

Nova's key technical insight was \emph{relaxed R1CS}. This is where the
mathematics earns its keep, and it is worth understanding why.

Standard R1CS says \(Az \circ Bz = Cz\) (where \(A\), \(B\), \(C\) are matrices
and \(z\) is the witness vector). Now suppose you try to combine two standard
R1CS instances by taking a random linear combination:
\((A(z_1 + r \cdot z_2)) \circ (B(z_1 + r \cdot z_2))\). When you expand this
product, cross-terms appear -- terms involving both \(z_1\) and \(z_2\)
multiplied together -- that do not fit the \(Az \circ Bz = Cz\) format. The
product of two linear combinations is quadratic, but R1CS demands bilinear
structure. The cross-terms break the format.

This is the kind of obstacle that looks fatal until someone sees through it.
Nova's insight was to relax the constraint. Instead of requiring
\(Az \circ Bz = Cz\) exactly, allow \(Az \circ Bz = u \cdot Cz + E\), where
\(u\) is a scalar and \(E\) is an error vector. When \(u = 1\) and \(E = 0\),
you recover standard R1CS. But the relaxed form is closed under random linear
combination: when you combine two relaxed instances, the cross-terms that would
have destroyed the standard format are absorbed into the error vector E. The
format survives. The scalar \(u\) tracks the linear combination's coefficient,
and \(E\) accumulates the algebraic debris that would otherwise break the
structure.

The folding verifier is tiny -- one scalar multiplication plus hashing, roughly
10,000 multiplication gates -- compared to millions for a full SNARK verifier.
This is the breakthrough: fold cheaply at every step, prove expensively only
once at the end. The per-step overhead of IVC dropped by orders of magnitude.

\subsection{The Snowball Does Not Fall
Apart}\label{the-snowball-does-not-fall-apart}

The natural worry about the snowball analogy is: snowballs fall apart. What is
the failure mode of folding?

The answer involves a subtlety that trips up even experienced cryptographers.
Folding does not provide soundness on its own. It provides a \emph{reduction}:
if the folded claim is valid, then both original claims were valid (with
overwhelming probability over the verifier's random challenge). But ``valid''
here means ``satisfies the relaxed constraint system.'' The final proof -- the
``decider'' SNARK at the end of the chain -- is what provides soundness. If the
underlying commitment scheme is binding and the random challenges are honestly
generated, then a cheating prover cannot produce a valid folded claim for an
invalid computation. The security proof works backward: a valid final claim
implies all intermediate claims were valid, which implies all computation steps
were correct.

The practical penalty comes from the size of the accumulated instance. In
classical Nova, the error vector E grows with each folding step. The 10,000-gate
figure for Nova's folding verifier is specifically the cost of the non-native
scalar multiplication required when working over a 2-cycle of curves. CycleFold
(Kothapalli and Setty, 2023) delegated this scalar multiplication to a
co-processor circuit on the secondary curve, where it can be computed natively
in approximately 1,500 gates. The folding verifier itself -- the hash and the
random linear combination -- is much cheaper. But the fundamental architecture
remains: fold cheaply, prove expensively but only once, at the end.

\section{The Folding Genealogy}\label{the-folding-genealogy}

Folding is not a single technique. It is a research program that has produced a
lineage of increasingly general schemes over the past four years. Understanding
this lineage matters for seeing where proof systems are headed -- including the
post-quantum frontier.

Before tracing each step, here is the map. Read it the way a naturalist reads a
field guide: not to memorize every species, but to understand the territory they
occupy. The folding genealogy advances along four independent axes:

\begin{itemize}
\tightlist
\item
  \textbf{Constraint generality:} R1CS (Nova) to CCS (HyperNova), covering all
  major arithmetizations.
\item
  \textbf{Instance generality:} Single-instance (Nova) to multi-instance
  (ProtoGalaxy), enabling parallel proving.
\item
  \textbf{Instruction generality:} Uniform IVC (Nova) to non-uniform IVC
  (SuperNova, SuperNeo), enabling VM execution.
\item
  \textbf{Cryptographic generality:} Elliptic curves (Nova through CycleFold) to
  lattices (LatticeFold, Neo, Symphony), enabling post-quantum security.
\end{itemize}

Every scheme in the genealogy that follows advances along at least one of these
axes. The reader can track the progress by asking a single question of each
paper: which axis did it push forward?

\subsection{Nova (2022): The Origin}\label{nova-2022-the-origin}

Kothapalli, Setty, and Tzialla. Folding for R1CS. The verifier performs one
scalar multiplication plus hashing -- roughly 10,000 R1CS gates of recursive
overhead. Prover cost is linear in the circuit size. This is the paper that made
folding practical. Every subsequent folding scheme either extends, generalizes,
or adapts Nova's core ideas.

\subsection{SuperNova (2022): Non-Uniform
IVC}\label{supernova-2022-non-uniform-ivc}

Kothapalli and Setty. Extended Nova to support multiple circuit types. Instead
of one step function that repeats, SuperNova allows different step functions at
each step -- exactly what you need for a virtual machine, where each instruction
has a different circuit. The scheme maintains one running instance per
instruction type and pays only for the circuit of the instruction actually
executed. This is the ``pay-per-instruction'' property that makes folding-based
zkVMs viable.

\subsection{HyperNova (2023): The Generalization
Point}\label{hypernova-2023-the-generalization-point}

Kothapalli and Setty again. This paper is the pivot of the entire genealogy.

HyperNova generalized folding from R1CS to CCS -- Customizable Constraint
Systems -- which subsume R1CS, PLONKish, and AIR in a single framework. The key
enabler was the sumcheck protocol, originally due to Lund, Fortnow, Karloff, and
Nisan in 1992, one of the most elegant tools in theoretical computer science. It
reduces the task of checking a sum over a large domain (say, \(2^{20}\)
evaluations of a multilinear polynomial) to checking a single evaluation at a
random point. The verifier sends random challenges in each round; the prover
responds with univariate polynomials. After \(\log(n)\) rounds, the verifier has
a single-point claim that can be checked directly.

HyperNova uses sumcheck to fold CCS instances without degree blowup. A CCS
constraint has the form: sum of products of matrix-vector multiplications equals
zero. This is inherently higher-degree than R1CS. Without sumcheck, folding such
constraints would require the verifier to handle cross-terms whose degree grows
with the constraint degree. Sumcheck reduces the multilinear CCS equation to a
single-point evaluation, allowing the folded instance to retain its structure.
The verifier cost remains \(O(\log m)\) field operations plus one scalar
multiplication -- essentially the same as Nova, despite handling a strictly more
general constraint system.

HyperNova is the generalization point because everything after it works in the
CCS framework. If Nova gave folding a body, HyperNova gave it a universal
language. CCS is to constraint systems what SQL is to databases: a common tongue
that lets you express any query regardless of the underlying storage engine.

\subsection{ProtoStar and ProtoGalaxy (2023): Alternative
Paths}\label{protostar-and-protogalaxy-2023-alternative-paths}

These two schemes took different approaches to generalizing folding. ProtoStar
(Bunz and Chen, ePrint 2023/620) built a generic accumulation framework for any
interactive argument satisfying ``special soundness'' -- a more abstract
starting point than HyperNova's CCS-specific approach. ProtoGalaxy (Eagen and
Gabizon) introduced multi-instance folding: folding k instances simultaneously
in a single round, rather than folding pairs sequentially. This enables
high-arity proof-carrying data trees, which are needed for parallel proof
generation.

Both contributed important ideas to the field, but the main trunk of the
genealogy runs through HyperNova because CCS became the dominant constraint
language.

\subsection{CycleFold (2023): The Practical
Fix}\label{cyclefold-2023-the-practical-fix}

Kothapalli and Setty. This is the engineering paper that made all the
theoretical folding schemes practical over elliptic curves. The problem: Nova's
folding verifier includes a scalar multiplication, which requires non-native
field arithmetic when working over a 2-cycle of curves. CycleFold delegates this
scalar multiplication to a co-processor circuit on the second curve, where it
can be computed natively. This reduced the second-curve circuit from roughly
10,000 gates to 1,500. Not glamorous, but essential. CycleFold is the standard
technique used in every practical implementation of Nova and HyperNova.

\subsection{LatticeFold (2024): Crossing Into Post-Quantum
Territory}\label{latticefold-2024-crossing-into-post-quantum-territory}

Boneh and Chen, published at ASIACRYPT 2025. LatticeFold was the first folding
scheme based on lattice assumptions rather than elliptic curve assumptions. This
matters enormously, because lattice problems (Module-SIS, Module-LWE) are
believed to resist quantum attacks, while elliptic curve discrete logarithm
falls to Shor's algorithm.

LatticeFold replaced elliptic curve commitments with Ajtai-style lattice
commitments: \(\text{Commit}(A, m, r) = A \cdot [m; r] \bmod q\), where \(A\) is
a public matrix, \(m\) is the message, and \(r\) is randomness. These
commitments are linearly homomorphic -- exactly the property that folding needs
to combine instances via random linear combination. But LatticeFold worked over
large fields (\(q \approx 2^{128}\)), which made arithmetic expensive.

\subsection{Neo (2025): Small Fields and
Pay-Per-Bit}\label{neo-2025-small-fields-and-pay-per-bit}

Wilson Nguyen and Srinath Setty. Neo adapted HyperNova's CCS folding to
lattice-based commitments over small fields -- specifically the Goldilocks field
(\(q = 2^{64} - 2^{32} + 1\)) with the cyclotomic ring
\(\mathbb{F}_q[X]/(\Phi_{81})\), where \(\Phi_{81}(X) = X^{54} + X^{27} + 1\).
Working over a 64-bit field instead of a 128-bit field makes arithmetic
dramatically faster, particularly on GPUs where 64-bit integer operations are
natively supported.

Neo introduced a property called ``pay-per-bit commitments.'' In an Ajtai
commitment, the cost of committing depends on the binary representation of the
witness. Committing to a single bit costs 32 times less than committing to a
32-bit value. This means the prover can commit to a binary witness at a fraction
of the cost of committing to field elements -- a property unique to the lattice
setting that has no analog in elliptic curve commitments.

SuperNeo, presented within the same paper, extended Neo to non-uniform IVC (the
lattice analog of SuperNova), supporting multiple instruction types for VM
execution.

\subsection{Symphony (2026): Production-Grade Lattice
Folding}\label{symphony-2026-production-grade-lattice-folding}

Independently extending lattice-based folding to high-arity settings, Symphony
refined the protocol for practical deployment. Key additions included optimized
Number Theoretic Transforms over the cyclotomic ring for fast polynomial
multiplication, a GPU-friendly architecture (lattice operations -- matrix-vector
products, NTTs -- parallelize naturally), and formal bridge theorems connecting
the folding scheme's knowledge soundness to the underlying Module-SIS and
Module-LWE assumptions.

Symphony demonstrated that lattice-based folding can be practical, not just
theoretical. With GPU acceleration, it achieves practical proving times, closing
the performance gap with elliptic-curve-based systems while maintaining
plausible post-quantum security.

Each step in the genealogy broadened the reach of folding without sacrificing
its fundamental advantage: logarithmic verifier cost and linear prover cost per
step. The entire genealogy, from Nova to Symphony, spans just three years. This
rate of progress is unusual even by the standards of a fast-moving field.

\section{Nightstream: What a Folding Engine Looks Like From the
Inside}\label{nightstream-what-a-folding-engine-looks-like-from-the-inside}

The genealogy reads like a family tree. But a family tree does not show you the
plumbing. It tells you who begat whom, not how the pipes connect, where the
pressure builds, or which joints leak under load.

Nightstream is a real implementation of the lattice-folding lineage described
above. Fifteen Rust crates, a Lean formal model, and a proving pipeline that
turns execution traces into folded obligations over the Goldilocks field. It is
a research prototype, not a production system -- its own README says so. But it
is the most complete public implementation of CCS-native lattice-based folding
available for study. And studying its engineering reveals something that theory
papers consistently leave out: the hardest problem in a folding system is not
any single algorithm. It is the alignment between all of them.

\subsection{A Pipeline, Not a Library}\label{a-pipeline-not-a-library}

The spine of Nightstream is a sequence of crates, each doing one thing, each
depending on the ones before it:

\texttt{neo-params} and \texttt{neo-math} fix the algebraic world -- Goldilocks
field, cyclotomic ring, norm bounds. \texttt{neo-ccs} defines the constraint and
evaluation relations. \texttt{neo-ajtai} provides the linear commitment backend.
\texttt{neo-transcript} enforces Fiat-Shamir sequencing through Poseidon2
hashing. \texttt{neo-reductions} implements the algebraic proof kernel.
\texttt{neo-memory} turns execution traces into per-step witness bundles.
\texttt{neo-fold} coordinates the shard-and-session folding runtime and emits
obligations.

This is Chapter 4's proof core triad -- constraint system, commitment scheme,
information-theoretic protocol -- made concrete in Rust. Each crate does one
thing. The system works only because every crate agrees with every other crate
on field choice, ring structure, witness layout, commitment semantics, and
transcript ordering. Break one agreement and the pipeline does not produce wrong
proofs. It produces no proofs at all.

That is the nature of a pipeline. A library offers tools. A pipeline imposes
discipline. You can use a library selectively. You must use a pipeline as
designed, or not at all.

\subsection{The Shared Bus}\label{the-shared-bus}

How does an execution trace become a foldable witness?

The answer lives in \texttt{neo-memory}, and the design is a shared CPU bus. The
execution trace -- every instruction, every memory access, every register state
-- is laid out as a single columnar spine. Each step in the computation produces
a row. The columns follow a fixed schema defined in the bus layout. Twist and
Shout memory arguments, which verify that the prover's claimed memory operations
are consistent, consume bus-extracted columns from the same spine rather than
maintaining separate committed witnesses.

This is what the backstage recording machinery from Chapter 4 looks like when
you unroll it into actual data structures. The magician's assistant is not just
scribbling notes -- she is filling in a spreadsheet with a fixed schema, one row
per step, and every downstream verification process reads from that same
spreadsheet.

The design works because it is unified. One trace spine means one source of
truth. No reconciliation between separate witness formats. No risk of the memory
argument seeing a different witness than the folding runtime.

The bus architecture connects six components in a fixed data flow:

\begin{longtable}[]{@{}
  >{\raggedright\arraybackslash}p{(\linewidth - 6\tabcolsep) * \real{0.2750}}
  >{\raggedright\arraybackslash}p{(\linewidth - 6\tabcolsep) * \real{0.1500}}
  >{\raggedright\arraybackslash}p{(\linewidth - 6\tabcolsep) * \real{0.3000}}
  >{\raggedright\arraybackslash}p{(\linewidth - 6\tabcolsep) * \real{0.2750}}@{}}
\toprule\noalign{}
\begin{minipage}[b]{\linewidth}\raggedright
Component
\end{minipage} & \begin{minipage}[b]{\linewidth}\raggedright
Role
\end{minipage} & \begin{minipage}[b]{\linewidth}\raggedright
Reads From
\end{minipage} & \begin{minipage}[b]{\linewidth}\raggedright
Writes To
\end{minipage} \\
\midrule\noalign{}
\arrayrulecolor{MidnightBlue}\midrule\arrayrulecolor{MidnightBorder}
\endhead
\bottomrule\noalign{}
\endlastfoot
Execution trace & Records each computation step as a row & Program input + state
& Columnar spine \\
Bus layout & Defines column schema (registers, flags, memory) & (architectural
constant) & All downstream consumers \\
neo-memory & Builds per-step witness bundles & Columnar spine & Folding
runtime \\
Twist memory args & Enforces memory consistency via permutation & Bus columns &
Obligation streams (val lane) \\
neo-fold & Accumulates shard/session folding obligations & Per-step bundles &
main + val obligation streams \\
Finalizer & Consumes both lanes and produces outer proof & main + val
obligations & Final proof artifact \\
\end{longtable}

The design is fragile because unified. Bus layout, witness builder, CPU
constraints, and oracle expectations must all agree on the exact column
semantics. If \texttt{neo-memory} produces a witness with columns in one order
and \texttt{neo-fold} expects them in another, the mismatch is not caught by
type checking. It is caught by a constraint that fails to satisfy, deep in the
pipeline, with an error message that does not point back to the layout
disagreement. This is the kind of bug that theory papers never discuss, because
in theory, witness layout is not a concept. In engineering, it is the concept.

\subsection{Three Reductions and Two
Lanes}\label{three-reductions-and-two-lanes}

The algebraic heart of Nightstream lives in \texttt{neo-reductions}, which
implements three reductions:

\textbf{Pi\_CCS} reduces CCS constraint satisfaction into evaluation claims.
This is the step that transforms ``does this witness satisfy these
constraints?'' into ``does this polynomial evaluate correctly at this random
point?'' -- the same transition from constraint checking to evaluation checking
that appears throughout modern proof systems.

\textbf{Pi\_RLC} performs random linear combination, batching multiple
evaluation claims into one. This is where the folding happens: two claims become
one, weighted by a verifier challenge.

\textbf{Pi\_DEC} handles decomposition, ensuring that the committed witness
values stay within the norm bounds required by the lattice commitment scheme.
Without this, the prover could fold claims using witness values that violate the
short-vector assumptions that make Ajtai commitments binding.

The reductions crate preserves a three-path architecture that deserves
attention. \texttt{Optimized} is the fast runtime path. \texttt{PaperExact}
mirrors the published protocol specification step for step, sacrificing
performance for auditability. \texttt{OptimizedWithCrosscheck} runs both paths
and compares their outputs. This is not paranoia. It is engineering discipline.
The optimized path inevitably diverges from the paper's notation, and a
reference implementation that can be run in parallel provides a semantic anchor.

After the reductions, the folding runtime in \texttt{neo-fold} introduces the
two-lane obligation model. The \texttt{main} lane carries the primary folded
claims. The \texttt{val} lane carries a separate obligation stream for the
Twist-related verification path, derived from a distinct random challenge. These
are not interchangeable. Verification reconstructs obligations in both lanes,
but verification alone does not finish the proof. Completion depends on a
finalizer that consumes both lanes and produces the outer proof. The system
emits a structured proof state, not a finished certificate.

What does a two-lane divergence look like in practice? Suppose the prover, at
folding step 3, manipulates a register value in the witness -- setting \(r_7\)
to 42 instead of the correct value 37. The \texttt{main} lane, using random
challenge \(\alpha\), accumulates an obligation that combines this register with
others:
\(\alpha \cdot r_1 + \alpha^2 \cdot r_2 + \cdots + \alpha^7 \cdot r_7 + \cdots\)
The wrong value shifts the accumulated sum, but with a single challenge, the
prover might get lucky -- if \(\alpha\) happens to land on a root of the error
polynomial, the corruption is invisible.

The \texttt{val} lane exists to prevent this. It uses an independent challenge
\(\beta\), drawn from a different Fiat-Shamir transcript fork. The probability
that the same error is invisible under \emph{both} \(\alpha\) and \(\beta\) is
at most \(d/|K|^2\), where \(d\) is the constraint degree and \(|K|\) is the
extension field size. For Neo's parameters (\(K = \mathbb{F}_{q^2}\) with
\(q \sim 2^{64}\)), this probability is less than \(2^{-127}\). The two lanes
provide soundness amplification: an error that survives one random challenge is
vanishingly unlikely to survive both. The finalizer checks that both lanes
produce consistent results -- and if the manipulated \(r_7\) corrupted the
\texttt{main} lane's obligation, the \texttt{val} lane catches the inconsistency
and the proof fails.

\subsection{The Lean Boundary}\label{the-lean-boundary}

Nightstream includes a formal subproject in Lean 4 that closes theorem surfaces
for the core protocol. The Lean model covers the algebraic reductions, the
commitment binding properties under Module-SIS, and the soundness chain from
folded claims back to original computation steps.

This is unusual for a research prototype. Most implementations at this stage
rely entirely on paper proofs and test suites. Nightstream's formal model
provides a higher standard of assurance for the mathematical core.

But there is a boundary that formal methods cannot cross on their own. The Lean
proofs verify that the \emph{mathematics} is correct -- that the reductions
preserve soundness, that the commitment scheme is binding under stated
assumptions. The Rust runtime must then consume those Lean-closed theorem
surfaces exactly as intended. The proofs live in one world; the code lives in
another. The residual risk is not that the theorems are wrong. It is that the
code, through a layout mismatch, an off-by-one index, or a misread parameter,
might inhabit a slightly different mathematical world than the one the theorems
describe. Mathematics proven correct in the abstract; the question is whether
the implementation faithfully instantiates it. This is a higher-quality problem
than most systems have. It is still the real remaining assurance boundary.

To be concrete about what Lean proves and what it does not: the Lean
formalization verifies that the mathematical reductions are sound -- that if the
folded claim passes the verifier's checks, then the original CCS instance was
satisfiable. It proves the \emph{if-then} chain from folded proof to original
computation. What Lean does \emph{not} prove is that the Rust code in
\texttt{neo-fold} correctly instantiates the reduction. An off-by-one index in a
matrix multiplication, a transposed loop bound in the commitment computation, a
misread parameter from \texttt{neo-params} -- any of these could cause the Rust
implementation to inhabit a slightly different mathematical world than the one
Lean verified. The residual risk is implementation fidelity, not mathematical
unsoundness. SP1's approach -- formal verification of opcode constraints against
the RISC-V Sail specification -- attacks the same gap from the opposite
direction: proving the code matches the spec rather than proving the spec is
sound. Neither project has closed the full loop from spec to code to hardware.
That loop is the frontier.

\subsection{Unfinished Scaffolding}\label{unfinished-scaffolding}

The core proving runtime -- reductions, memory, folding -- is maintained and
tested. The rest of the system is at varying stages of completion.

The finalizer that consumes both obligation lanes and produces an outer proof is
work in progress. \texttt{neo-spartan-bridge}, which would compress the folded
obligations into a Spartan-style SNARK, is explicitly experimental.
\texttt{neo-midnight-bridge}, which would connect Nightstream's output to
Midnight's PLONK/KZG verification layer, exists as a roadmap interoperability
path rather than a maintained product surface.

The project's README describes it as a research prototype. This is honest, and
the honesty matters. A system that accurately describes its own incompleteness
is more trustworthy than one that claims a completeness it has not achieved. The
core is real. The edges are still under construction. That distinction should be
preserved in any evaluation of the system.

\subsection{What the Plumbing Reveals}\label{what-the-plumbing-reveals}

The main lesson from Nightstream is not about any single cryptographic
primitive. It is about alignment.

The hardest engineering in a modern folding system is not inventing a new
commitment scheme or designing a new reduction. It is making the witness layout
match the fold expectations. Making the reductions agree with the commitment
semantics. Making the transcript bind every public input before challenges are
sampled. Making the bus columns line up between the builder that writes them and
the constraints that read them.

After 2023, the interesting work in proof systems moved from individual
algorithmic breakthroughs to pipeline coordination. Nova was a breakthrough.
HyperNova was a breakthrough. But getting fifteen crates to agree on column
ordering, norm bounds, transcript sequencing, and obligation semantics -- that
is not a breakthrough. It is the slow, patient, unglamorous work of making a
real system function. And it is where most of the time goes.

That lesson does not appear in any genealogy. It can only come from a codebase.

\section{Circle STARKs and Stwo: A Generational
Leap}\label{circle-starks-and-stwo-a-generational-leap}

While the folding lineage was evolving, a parallel revolution was happening in
the STARK world. In 2024, Haboeck (Polygon Labs), Levit, and Papini (StarkWare)
published Circle STARKs, and their production implementation -- Stwo -- became
the fastest proof system ever deployed.

The key insight is a change in the algebraic structure underlying the FFT, which
is the computational backbone of every STARK prover. Traditional STARKs use
multiplicative subgroups of finite fields for their FFT domains. These subgroups
exist in large fields (like BN254's 254-bit field), but finding smooth-order
subgroups in small fields is difficult. Circle STARKs replace the multiplicative
subgroup with the \emph{circle group} of a Mersenne prime field.

Why the circle group? The Mersenne-31 field (\(M31 = 2^{31} - 1\)) is a prime
with a special property: the circle group
\(C(\mathbb{F}_p) = \{(x, y) : x^2 + y^2 = 1 \text{ over } \mathbb{F}_p\}\) has
order \(p + 1 = 2^{31}\), which is a perfect power of 2. This enables a radix-2
FFT (the Circle FFT, or CFFT) analogous to the standard Cooley-Tukey FFT but
operating over circle points. Each FRI step halves the domain using the circle
group's ``squaring'' map, and the entire protocol adapts naturally to the circle
geometry.

Why does the field size matter so much? Because 31-bit numbers fit in a single
32-bit machine word. On modern CPUs with SIMD instructions, you can process 8 or
16 M31 elements in parallel per instruction. On GPUs, the advantage is even more
dramatic. Arithmetic over M31 is roughly 100 times faster than arithmetic over
BN254's 254-bit field. When your proof system spends most of its time doing
field arithmetic, a 100x speedup in the field operations translates almost
directly into a 100x speedup in proving.

The numbers confirm this. Stwo, StarkWare's production implementation of Circle
STARKs, went live on Starknet mainnet in 2025. Every Starknet block is now
proven by Stwo. The benchmarks:

\begin{itemize}
\tightlist
\item
  \textbf{940x throughput improvement} over Stone (StarkWare's previous prover)
\item
  \textbf{50x improvement} over ethSTARK (the academic reference implementation)
\item
  GPU acceleration via ICICLE-Stwo adds another 3.25x to 7x on top of the CPU
  SIMD backend
\item
  Sub-second recursive proving is within reach: roughly 20 milliseconds for
  10,000 Poseidon hash evaluations
\end{itemize}

These are not projections. This is production performance, running on mainnet,
proving real blocks with real transactions. The gap between ``academic proof
system'' and ``deployed infrastructure'' has closed.

\subsection{Why the Circle Group Matters}\label{why-the-circle-group-matters}

The phrase ``circle group'' sounds like an abstraction. It is not. It is a
geometric fact about numbers that makes everything else possible, and it
deserves a careful explanation.

Start with what came before. Traditional STARKs need a domain of points where
they can evaluate polynomials -- a set of evenly-spaced ``sampling points'' that
enable the FFT. In large fields like BN254, you find these points in the
multiplicative group: pick a generator \(g\), and the powers
\(g, g^2, g^3, \ldots, g^{2^k}\) form a subgroup of order \(2^k\). These powers
wrap around like hours on a clock face. The FFT works because the subgroup has
smooth order (a power of 2), so the Cooley-Tukey butterfly decomposition applies
perfectly.

But in small fields, this strategy collapses. The Mersenne-31 field has only
\(2^{31} - 2\) nonzero elements. Its multiplicative group has order
\(2^{31} - 2 = 2 \times 3 \times 357{,}913{,}941\). The largest power-of-2
subgroup has order just 2 -- it contains only \{1, -1\}. You cannot run an FFT
of size \(2^{24}\) over a group of size 2. The multiplicative group of M31 is,
for FFT purposes, useless.

This is where the circle enters. Consider the set of all pairs \((x, y)\) in M31
that satisfy \(x^2 + y^2 = 1\). This is the unit circle over the finite field --
not a continuous curve but a discrete set of points. The key fact: this set has
exactly \(p + 1 = 2^{31}\) points. Not \(2^{31} - 2\). Not some awkward
composite. Exactly \(2^{31}\). A perfect power of 2.

The coincidence is not a coincidence. It is a theorem. For any Mersenne prime
\(p = 2^n - 1\), the circle group \(C(\mathbb{F}_p)\) has order \(p + 1 = 2^n\).
This is because the circle group over \(\mathbb{F}_p\) is isomorphic to the
multiplicative group of \(\mathbb{F}_{p^2}\) modulo \(\mathbb{F}_p^*\) -- a
quotient that inherits the 2-adic structure of \(p + 1\). When \(p\) is a
Mersenne prime, that structure is maximally smooth: a pure power of 2.

This perfect power-of-2 structure gives you the cleanest possible FFT. The
Circle FFT (CFFT) decomposes a polynomial evaluation over \(2^{31}\) points into
layers of half-size evaluations, just as the standard Cooley-Tukey FFT does over
multiplicative subgroups. Each layer halves the domain. After 31 layers, you are
done. No padding, no awkward leftovers, no compromises.

\subsection{The Squaring Map: Folding a Circle in
Half}\label{the-squaring-map-folding-a-circle-in-half}

The FRI protocol -- the heart of every STARK's proof of low degree -- works by
repeatedly halving the evaluation domain. At each step, the prover combines
pairs of evaluations into single values, reducing the domain by a factor of 2.
After enough steps, the polynomial is so small that the verifier can check it
directly.

In a traditional STARK over a multiplicative group, the halving map is squaring:
the map \(x \mapsto x^2\) sends a subgroup of order \(2^k\) to a subgroup of
order \(2^{k-1}\). Each element and its ``twin'' (its negation in the group) map
to the same value, so the domain folds in half.

On the circle, the halving map is also a squaring -- but a \emph{geometric}
squaring. The circle has a group law: you can ``add'' two points by a formula
analogous to angle addition. The squaring map sends a point to its double under
this group law. Concretely, for a point \((x, y)\) on the circle
\(x^2 + y^2 = 1\), the doubling map sends \((x, y)\) to \((2x^2 - 1, 2xy)\).
This map is 2-to-1: each image point has exactly two preimages. So the circle of
\(2^k\) points folds onto a circle of \(2^{k-1}\) points.

The visual intuition is literal. Imagine folding a circle in half. The top half
maps onto the bottom half. Each point on the bottom half corresponds to two
points on the original circle -- one on top, one on bottom. The FRI protocol
does exactly this, algebraically. At each step, the prover commits to
evaluations on a circle, the verifier sends a random challenge, and the prover
uses the challenge to combine each pair of evaluations (the point and its ``fold
partner'') into a single value on the half-size circle. After \(\log(n)\) steps,
only a constant-size polynomial remains.

This geometric folding is why Circle STARKs achieve the same asymptotic
efficiency as traditional STARKs -- \(O(n \log n)\) prover time, \(O(\log^2 n)\)
verifier time, \(O(\log^2 n)\) proof size -- while operating over a field where
the elements are only 31 bits wide.

\subsection{Why 31-Bit Arithmetic Is So
Fast}\label{why-31-bit-arithmetic-is-so-fast}

The final piece of the Circle STARK advantage is not algebraic but
architectural. It concerns the physical reality of how modern processors handle
numbers.

A 64-bit CPU register holds 64 bits. A 254-bit BN254 field element requires four
64-bit ``limbs'' and careful carry propagation between them. A single field
multiplication over BN254 involves roughly 16 limb multiplications and a cascade
of additions and carries -- effectively 20 to 30 machine instructions per field
multiplication. This is multi-precision arithmetic, and it is inherently serial:
each carry depends on the result of the previous limb multiplication.

A 31-bit M31 field element fits in a single 32-bit word. A single field
multiplication is one 32-bit multiply followed by one modular reduction -- and
because \(M31 = 2^{31} - 1\) is a Mersenne prime, the reduction is a single
addition and a conditional subtraction. Two machine instructions. The ratio is
severe: 2 instructions for M31 versus 20-30 for BN254. A factor of 10 to 15, per
operation.

But the advantage compounds under parallelism. A 64-bit register can hold two
M31 elements side by side. An AVX-256 SIMD register holds eight. An AVX-512
register holds sixteen. A single SIMD instruction can multiply sixteen M31
elements simultaneously, producing sixteen field products in the time it takes
to perform one BN254 multiplication. The theoretical throughput ratio is not
10x. It is 100x or more.

On a GPU, the effect is even more dramatic. A GPU's streaming multiprocessors
are optimized for 32-bit integer and floating-point operations -- the native
word size of graphics workloads. M31 arithmetic maps directly onto the
hardware's sweet spot. BN254 arithmetic requires emulation using multiple 32-bit
operations per limb, with register pressure and carry chains that reduce
occupancy (the fraction of the GPU's compute units that are actively working).
In practice, Stwo's GPU backend -- implemented via ICICLE -- achieves 3.25x to
7x additional speedup on top of the already-fast CPU SIMD implementation. The
cumulative advantage of small-field arithmetic, from instruction-level to
chip-level, is why Circle STARKs proved to be not a modest improvement over
traditional STARKs but a qualitative jump.

The Mersenne prime M31 is not the only small field in production. BabyBear
(\(p = 2^{31} - 2^{27} + 1 = 15 \times 2^{27} + 1\)) offers similar 31-bit
arithmetic with a multiplicative group of smooth order (\(2^{27}\) divides
\(p - 1\)), enabling traditional multiplicative-subgroup FFTs rather than circle
FFTs. SP1 uses BabyBear for its inner proof system. The choice between M31 and
BabyBear is a choice between circle-group FFTs and multiplicative-group FFTs --
two paths to the same destination of fast, small-field proving. Both paths
converge on the same insight: the biggest optimization in proof system
engineering is not a cleverer algorithm. It is a smaller number.

\section{Real-Time Ethereum Proving}\label{real-time-ethereum-proving}

The performance revolution in proof systems has had a direct economic
consequence worth pausing to appreciate, because it changes the fundamental
calculus of what this technology is good for.

In December 2023, proving a single Ethereum block cost approximately \$80.21. By
December 2025, the cost had fallen to roughly \$0.04 -- a 2,000x reduction in 24
months. Airbender (ZKsync's prover) achieved \$0.0001 per transfer. These
numbers come from CastleLabs and Ethproofs benchmarks, and they represent a cost
curve steeper than Moore's Law.

But cost is only half the story. Speed matters equally, because Ethereum
produces a new block every 12 seconds. If proving takes longer than 12 seconds,
the prover cannot keep up with the chain. For years, this was a distant goal. In
late 2025, four teams crossed the threshold:

\begin{itemize}
\item
  \textbf{SP1 Hypercube} (Succinct Labs): Proved 99.7\% of Ethereum Layer 1
  blocks in under 12 seconds, using 16 NVIDIA RTX 5090 GPUs (hardware cost
  approximately \$32,000). SP1 Hypercube uses a multilinear polynomial stack
  built entirely on the sumcheck protocol, with a ``jagged'' polynomial
  commitment scheme that enables pay-per-use proving.
\item
  \textbf{ZKsync Airbender}: Achieved 21.8 million cycles per second on a single
  H100 GPU, proving Ethereum blocks in approximately 35 seconds. Open-source,
  moving toward formal verification.
\item
  \textbf{Two additional teams} from the Ethereum Foundation's proving ecosystem
  demonstrated sub-12-second proving under various hardware configurations.
\end{itemize}

The Ethereum Foundation responded by declaring the speed race ``operationally
viable'' in December 2025 and shifting its targets. The new requirements: less
than 10 seconds per block, less than \$100,000 in hardware, less than 10
kilowatts of power, 128-bit provable security, and proof sizes under 300
kilobytes. The pivot from speed to security signals that the raw performance
problem, which dominated proof system research for five years, has been
substantially solved. The next frontier is formal security guarantees -- a topic
we will revisit in Chapter 7.

This cost drop matters for reasons beyond Ethereum. When proving costs \$80, ZK
proofs are a luxury -- viable only for high-value transactions or well-funded
rollups. When proving costs four cents, ZK proofs become infrastructure -- cheap
enough to apply to every transaction, every block, every state transition. The
technology moves from ``expensive security upgrade'' to ``default operating
mode.'' That shift was enabled almost entirely by improvements at Layer 5.

\section{The Proof Core: Why Layers 4, 5, and 6 Are
Inseparable}\label{the-proof-core-why-layers-4-5-and-6-are-inseparable}

At this point in our tour of the seven-layer stack, a structural observation is
unavoidable. The boundaries between Layer 4 (arithmetization), Layer 5 (proof
system), and Layer 6 (cryptographic primitives) are not clean lines. They are
gradients.

Consider the design decisions involved in building a proof system:

\begin{itemize}
\item
  The \textbf{finite field} (Layer 6) determines which arithmetic is fast.
  Mersenne-31 enables SIMD-friendly 32-bit operations. Goldilocks enables
  efficient 64-bit operations on GPUs. BN254 requires expensive 254-bit
  multi-precision arithmetic. The field choice propagates upward through every
  layer.
\item
  The \textbf{commitment scheme} (Layer 6) determines the trust model and proof
  size. KZG commitments (from pairings, Layer 6) give constant-size proofs but
  require a trusted setup. FRI commitments (from hash functions, Layer 6) give
  logarithmic proofs with transparency. Ajtai commitments (from lattices, Layer
  6) give post-quantum security with larger proofs. The commitment scheme shapes
  the proof system architecture.
\item
  The \textbf{arithmetization} (Layer 4) determines which constraints the proof
  system must handle. R1CS, CCS, AIR, PLONKish -- each format has different
  properties, and the proof system must be designed to handle the chosen format
  efficiently. CCS folding (HyperNova) requires the sumcheck protocol. AIR
  proofs (STARKs) require FRI. The arithmetization and the proof system
  co-evolve.
\end{itemize}

These three choices -- field, commitment, arithmetization -- form a tightly
coupled triad. Change one, and the other two must adapt. This is why we call
them the ``proof core'': they function as a single design unit, even though our
seven-layer model places them in separate layers.

The layered model is still useful for understanding. It separates concerns that
are conceptually distinct: what the mathematics \emph{is} (Layer 6), how
computation is \emph{encoded} (Layer 4), and how the encoding is \emph{verified}
(Layer 5). But the reader should understand that in practice, these layers are
designed together, optimized together, and constrained by each other's choices.
A proof system is not assembled from independent components like bricks in a
wall. It is forged as a single alloy, where the properties of each ingredient
determine the properties of the whole.

\section{Fiat-Shamir Vulnerabilities}\label{fiat-shamir-vulnerabilities}

Every proof system we have discussed relies on a technique called the
Fiat-Shamir transform to convert an interactive protocol into a non-interactive
one. In the interactive version, the verifier sends random challenges to the
prover. In the non-interactive version, the prover generates the challenges by
hashing the protocol's transcript -- all previous messages -- into pseudorandom
values. This eliminates the need for the verifier to be online during proving.

The Fiat-Shamir transform is simple in principle but treacherous in practice.
The transcript that gets hashed must include \emph{everything} that the verifier
would have seen in the interactive version. If the prover omits anything -- a
public input, a commitment, a previous challenge -- then the resulting hash is
not a faithful simulation of the interactive protocol, and soundness can break.

\subsection{Frozen Heart and Fiat-Shamir
Vulnerabilities}\label{frozen-heart-and-fiat-shamir-vulnerabilities}

Why does this matter at Layer 5 specifically? Because the Fiat-Shamir transform
is the \emph{binding} mechanism between prover and verifier. In the interactive
version of any proof protocol, the verifier generates fresh random challenges
that force the prover to commit before seeing what will be checked. The
Fiat-Shamir transform replaces these live challenges with hash-derived
challenges -- but only if the hash input includes \emph{every} public value the
verifier would have seen. The transcript is the contract between the two
parties. Omit a single term, and the prover can retroactively choose commitments
that satisfy whatever challenges arise. The mathematical proof of soundness
assumes the transcript is complete. The implementation must deliver on that
assumption, term by term, or the proof of soundness is vacated.

This is not a theoretical concern. The ``Frozen Heart'' vulnerability class,
disclosed by Trail of Bits in 2022, affected six independent implementations
across three proof systems. The ``Last Challenge Attack'' of 2024 compromised
gnark's PLONK verifier, used by multiple Ethereum rollups. In 2025, Solana's ZK
ElGamal repeated the pattern. Chapter 8 catalogs these incidents in forensic
detail -- what was omitted, how the forgery was constructed, and what the
governance implications are for on-chain verification. Here at Layer 5, the
lesson is narrower but no less urgent: the gap between a proof system's
mathematical specification and its implementation is where real-world attacks
live.

Fiat-Shamir vulnerabilities are the ``SQL injection'' of zero-knowledge
cryptography: a well-understood class of bug that keeps recurring because it is
easy to get wrong and hard to detect by inspection. They remind us that the
security of a proof system is not just a property of its mathematical design. It
is a property of its implementation, its specification, and the gap between the
two. This is why formal verification of proof system implementations -- not just
their mathematical specifications -- is increasingly recognized as essential.
SP1 Hypercube's formal verification of all 62 RISC-V opcode constraints against
the official RISC-V Sail specification represents the state of the art in this
direction.

\section{Case Study: Midnight's Sealed
Certificate}\label{case-study-midnights-sealed-certificate}

To see how these abstract proof system choices play out in a real system,
consider Midnight -- the privacy-focused blockchain developed by Input Output
Global (IOG) for the Cardano ecosystem. Midnight's sealed certificate tells us
what Layer 5 looks like when theory meets production.

\subsection{The Proof System: Halo 2 / UltraPlonk over
BLS12-381}\label{the-proof-system-halo-2-ultraplonk-over-bls12-381}

Midnight chose a PLONK-family proof system: specifically Halo 2, an extension of
UltraPlonk with custom gates and lookup tables. The curve is BLS12-381, a
pairing-friendly curve with better security properties than BN254 (roughly
128-bit security vs.~BN254's \textasciitilde100 bits after Tower NFS advances).
Although the original Halo paper (2019) used inner product arguments (IPA) that
required no trusted setup, Midnight's deployment uses KZG polynomial commitments
over BLS12-381, which require a Powers-of-Tau ceremony -- a universal trusted
setup with a 1-of-N trust assumption.

This places Midnight firmly in the ``classical SNARK'' camp: pairing-based,
recursion-capable, not post-quantum. It is a mature, well-understood choice. The
PLONK arithmetization is flexible enough to support Midnight's
privacy-preserving smart contract model, where transactions carry zero-knowledge
proofs of valid state transitions.

\subsection{The Four-Phase Transaction
Pipeline}\label{the-four-phase-transaction-pipeline}

Midnight's proof generation follows a four-phase pipeline that illustrates how
the sealed certificate gets manufactured in practice:

\textbf{Phase 1: Circuit Execution (callTx).} The DApp calls a contract
function. The Compact compiler has already compiled this function into a ZKIR
(Zero-Knowledge Intermediate Representation) circuit. The circuit executes
locally -- on the user's machine -- producing an ``unproven transaction'' that
contains the execution trace but no proof.

\textbf{Phase 2: Proof Generation (proveTx).} The unproven transaction is sent
to a local proof server running on localhost:6300. This is a separate process,
launched via Docker, that generates the ZK proof. The witness (private inputs)
never leaves the client machine. The proof server runs Halo 2's prover and
returns a proven transaction. This step dominates latency. This is where the
certificate gets sealed.

\textbf{Phase 3: Fee Balancing (balanceTx).} The proven transaction is bound,
balanced (fee inputs are added from the user's DUST wallet), and signed. This
step is sub-second.

\textbf{Phase 4: Submission (submitTx).} The finalized transaction -- containing
the ZK proof and the public transcript of state reads and writes -- is submitted
to the blockchain. The node verifies the proof against the on-chain verifier
key, checks that the transcript is consistent with the current ledger state, and
applies the state transition. The sealed certificate has reached the audience.

\subsection{The Performance Reality}\label{the-performance-reality}

Measured performance on Midnight's devnet reveals the cost of sealing in
concrete terms:

\begin{longtable}[]{@{}lll@{}}
\toprule\noalign{}
Operation & Time & Bottleneck \\
\midrule\noalign{}
\arrayrulecolor{MidnightBlue}\midrule\arrayrulecolor{MidnightBorder}
\endhead
\bottomrule\noalign{}
\endlastfoot
Deploy (constructor circuit) & 17-27 seconds & Proof generation \\
Circuit call (simple increment) & 17-18 seconds & Proof generation \\
Circuit call (sealed bid) & 22-24 seconds & Proof generation \\
Circuit call (execute escrow) & 23.8 seconds & Proof generation \\
Balance + submit & \textless{} 1 second & Network \\
Failed assertion (local) & 0.1-0.5 seconds & No proof needed \\
\end{longtable}

The pattern is clear: proof generation accounts for 95\% or more of every
transaction's latency. A simple counter increment (the ``hello world'' of smart
contracts) takes 17 seconds because the Halo 2 prover must generate a full
zero-knowledge proof. A more complex circuit (escrow execution) takes 24
seconds. Everything else -- fee balancing, signing, network submission, on-chain
verification -- is negligible by comparison.

From the user's perspective, Layer 5 is a 17-to-28-second pause during which the
proof server is pressing the seal into wax. The cryptography is working. The
privacy is being maintained. But the user is waiting.

\subsection{Midnight vs.~the Frontier}\label{midnight-vs.-the-frontier}

Midnight's architecture makes different choices than the proof systems at the
performance frontier. It does not use STARKs. It does not use folding. It does
not use GPU acceleration (the proof server appears to be CPU-only). It does not
use small-field arithmetic. These are deliberate engineering choices that
prioritize maturity and correctness over raw speed.

A comparison with Ethereum is instructive:

\begin{longtable}[]{@{}
  >{\raggedright\arraybackslash}p{(\linewidth - 6\tabcolsep) * \real{0.1333}}
  >{\raggedright\arraybackslash}p{(\linewidth - 6\tabcolsep) * \real{0.2533}}
  >{\raggedright\arraybackslash}p{(\linewidth - 6\tabcolsep) * \real{0.4133}}
  >{\raggedright\arraybackslash}p{(\linewidth - 6\tabcolsep) * \real{0.2000}}@{}}
\toprule\noalign{}
\begin{minipage}[b]{\linewidth}\raggedright
Property
\end{minipage} & \begin{minipage}[b]{\linewidth}\raggedright
Midnight (Halo 2)
\end{minipage} & \begin{minipage}[b]{\linewidth}\raggedright
Neo/Symphony (Lattice folding)
\end{minipage} & \begin{minipage}[b]{\linewidth}\raggedright
SP1 Hypercube
\end{minipage} \\
\midrule\noalign{}
\arrayrulecolor{MidnightBlue}\midrule\arrayrulecolor{MidnightBorder}
\endhead
\bottomrule\noalign{}
\endlastfoot
Proof system family & PLONK & Folding (CCS) & STARK + sumcheck \\
Hardness assumption & Discrete log & Module-SIS/LWE & Collision-resistant
hashing \\
Post-quantum & No & Yes (plausible) & Yes (inner proof) \\
Field & BLS12-381 (255-bit) & Goldilocks (64-bit) & M31/BabyBear (31-bit) \\
Composition strategy & Recursion & Folding & STARK recursion + SNARK wrap \\
Proof time (simple circuit) & 17-18 seconds & GPU-accelerated NTT &
Sub-second \\
Proof size & \textasciitilde hundreds of bytes & Larger (lattice-based) &
\textasciitilde192 bytes (after Groth16 wrap) \\
Trust model & Trusted (universal SRS) & Transparent & Transparent inner, trusted
outer \\
\end{longtable}

Midnight's 17-second proof time for a simple increment reflects BLS12-381's
expensive 255-bit arithmetic, the absence of GPU acceleration, and the overhead
of a full Halo 2 proof per transaction. SP1 Hypercube's sub-second proving
reflects M31's cheap 31-bit arithmetic, GPU parallelism, and an architecture
optimized for throughput. These are not different implementations of the same
idea. They are different points in a design space where field size, hardware
utilization, and proof architecture interact multiplicatively.

None of this is meant to criticize Midnight. Its choices are appropriate for a
privacy-focused system where correctness and auditability matter more than raw
speed, and where the developer toolchain (Compact language, ZKIR, local proof
server) provides a coherent end-to-end experience. But it illustrates how the
choices made when sealing the certificate -- amplified by the field and
commitment choices at Layer 6 -- determine the user experience at the
application layer.

\section{SNARK Recursion vs.~Folding: The Full
Picture}\label{snark-recursion-vs.-folding-the-full-picture}

Now that we have seen both approaches in action -- Midnight's recursive Halo 2
architecture and the folding genealogy leading to Neo and Symphony -- we can
draw a more complete picture of how they compare.

\subsection{When Recursion Wins}\label{when-recursion-wins}

Recursive proof composition is the right choice when:

\begin{itemize}
\tightlist
\item
  The chain of computation is short (tens of steps, not millions)
\item
  The proof system already has a cheap verifier circuit (as Halo 2 does for its
  IPA-based verification)
\item
  The application needs per-step proofs (every transaction gets its own proof,
  as in Midnight)
\item
  The infrastructure is mature and the priority is correctness over performance
\end{itemize}

Recursion is conceptually simpler: each step produces a proof, and the next step
verifies it. There are no accumulated instances to manage, no decider to run at
the end, and no error vectors that grow with the chain length. For short chains
and mature tooling, recursion is the pragmatic choice.

\subsection{When Folding Wins}\label{when-folding-wins}

Folding is the right choice when:

\begin{itemize}
\tightlist
\item
  The chain of computation is long (millions of VM steps)
\item
  Per-step cost must be minimized (the overhead of a full SNARK verification per
  step is unacceptable)
\item
  The application is a zkVM or rollup that processes large batches of
  transactions
\item
  Post-quantum security is a requirement (lattice-based folding via
  Neo/Symphony)
\item
  Parallel proving is needed (multi-instance folding via ProtoGalaxy enables
  tree-structured PCD)
\end{itemize}

The asymptotic advantage of folding is clear: \(O(|F|)\) prover cost per step
(where \(|F|\) is the step circuit size) plus roughly 1,500 gates of folding
overhead (with CycleFold), compared to \(O(|F| + |V|)\) for recursion, where
\(|V|\) is the size of the verifier circuit (potentially millions of gates). For
a zkVM executing millions of RISC-V instructions, this difference is the gap
between practical and impractical.

\subsection{The Convergence}\label{the-convergence}

In practice, the distinction between recursion and folding is blurring. Most
production systems use a hybrid: folding for the inner loop (accumulate
computation steps cheaply), then a monolithic SNARK (Spartan, Groth16) as the
``decider'' that produces the final succinct proof. Nova uses Spartan as its
decider. Mangrove (Nguyen, Datta, Chen, Tyagi, Boneh, 2024) builds a k-arity PCD
tree where leaf nodes fold computation chunks and internal nodes merge folded
instances, with a final SNARK for the NP statement. The boundary between
``folding'' and ``recursion'' dissolves into a pipeline where both techniques
serve different stages.

The key decision variable is step count. For computations under approximately
1,000 steps, recursion with a fast inner proof system (Groth16 or PLONK) is
competitive in both prover time and implementation complexity -- the per-step
overhead of a full recursive proof is high but amortizes over few steps. For
computations exceeding 10,000 steps -- the regime of zkVMs proving full program
executions -- folding's \(O(|F|)\) per-step cost with \textasciitilde1,500 gates
of folding overhead dominates recursion's \(O(|F| + |V|)\) per-step cost, where
\(|V|\) is the verifier circuit size. The crossover region between 1,000 and
10,000 steps depends on the specific constraint system, field size, and
hardware: folding wins earlier on small fields (where \(|V|\) is relatively
large compared to \(|F|\)) and later on pairing-friendly fields (where recursive
verification is cheap). In practice, the distinction is blurring -- the dominant
architecture uses folding for the inner accumulation loop and a single recursive
SNARK compression as the final decider step.

\section{The Post-Quantum Horizon}\label{the-post-quantum-horizon}

Everything we have discussed in this chapter faces an existential question: what
happens when quantum computers arrive?

Shor's algorithm breaks the discrete logarithm problem in polynomial time. This
means every proof system built on elliptic curve cryptography -- Groth16, PLONK,
Halo 2, KZG commitments, all of Nova's original elliptic-curve-based
instantiations -- becomes insecure. Not ``might become insecure.'' Becomes
insecure. The mathematical fact is established; only the engineering timeline is
uncertain.

The NIST IR 8547 deprecation schedule targets 2035 for phasing out pre-quantum
cryptography. Conservative estimates for ``Q-Day'' -- the date when a
cryptographically relevant quantum computer exists -- range from 2032 to 2035.
For blockchain systems with 10+ year lifespans, this means systems deployed
today must either plan for migration or be built on quantum-resistant
foundations from the start.

The STARK family is already partially quantum-resistant: its security rests on
collision-resistant hash functions, which resist quantum attacks (though
Grover's algorithm reduces their effective security, SHA-256's collision
resistance drops from 128 bits to roughly 85 bits under the BHT algorithm,
though this attack requires quantum random-access memory, which is widely
considered physically impracticable with current technology). But the
STARK-to-SNARK wrapping pipeline reintroduces quantum vulnerability at the final
step, because the Groth16 wrapper uses BN254 pairings.

The lattice folding branch of the genealogy -- LatticeFold, Neo, Symphony --
represents the most direct path to post-quantum proof systems. Neo achieves
127-bit security under plausible lattice hardness assumptions (Module-SIS,
Module-LWE), operates over the GPU-friendly Goldilocks field, and supports the
full CCS constraint framework via sumcheck-based folding. Symphony extends this
to production-grade performance with GPU-optimized NTTs.

The remaining gap is on-chain verification. No post-quantum on-chain verifier
exists in production. Lattice-based proofs are larger than elliptic-curve-based
proofs (tens of kilobytes vs.~hundreds of bytes), and no blockchain has
precompiled contracts for lattice operations. Closing this gap -- either through
lattice-friendly L1 verification or through novel compression techniques -- is
one of the open problems at the frontier of the field.

\section{From Speed Race to Security
Race}\label{from-speed-race-to-security-race}

The story of Layer 5 over the last three years is a story of two races.

The first race was about speed. From 2022 to 2025, the question was: can you
prove computation fast enough for it to matter? Can you prove an Ethereum block
before the next block arrives? Can you bring the cost below a dollar, below a
dime, below a penny? This race has been substantially won. Real-time proving is
operational. Costs are in the single-digit cents. The hardware stack -- GPUs,
SIMD, and potentially FPGAs and ASICs -- has been mobilized.

The second race is about security. The Ethereum Foundation's December 2025
announcement marked the pivot: the target shifted from ``prove fast'' to ``prove
with 128-bit provable security.'' This means not just believing the proof system
is secure, but having a formal proof that a computationally bounded adversary
cannot forge proofs with probability better than \(2^{-128}\). It means formally
verifying the implementation against the specification. It means accounting for
the gap between the random oracle model (where Fiat-Shamir uses ideal hash
functions) and reality (where Fiat-Shamir uses SHA-256 or Poseidon).

SP1 Hypercube's formal verification of all 62 RISC-V opcode constraints against
the RISC-V Sail specification is a milestone in this second race. It
demonstrates that production proof systems can achieve the level of formal rigor
that was previously associated only with academic papers. But verifying opcodes
is only the beginning. The full stack -- from the Fiat-Shamir transform through
the polynomial commitment scheme through the field arithmetic -- must be
verified end-to-end. This is a harder problem, and it is the one that will
define Layer 5's trajectory over the next several years.

\section{The Sealed Certificate}\label{the-sealed-certificate}

Layer 5 is where the mathematics becomes a machine. The abstract polynomial
constraints from Layer 4 are sealed into a tamper-evident certificate that can
travel across networks, be verified by strangers, and survive adversarial
scrutiny. The certificate's properties -- its size, its verification cost, its
security assumptions, its quantum resistance -- are determined by the proof
system that seals it.

Three families of proof systems (Groth16, PLONK, STARKs) have converged into a
single hybrid pipeline where each contributes its strengths. Folding schemes
evolved from Nova's R1CS-specific innovation into a family of increasingly
general techniques spanning CCS, multi-instance proving, non-uniform IVC, and
lattice-based post-quantum security. Circle STARKs and Stwo delivered a 940x
throughput improvement through the simple insight that 31-bit field arithmetic
is 100 times faster than 254-bit arithmetic. Proving costs plummeted by three
orders of magnitude in two years, and proving speed crossed the real-time
threshold for Ethereum blocks.

But the vulnerabilities are equally real. Fiat-Shamir bugs (Frozen Heart, the
Last Challenge Attack) demonstrate that the gap between a proof system's
mathematical security and its implementation security is where real-world
attacks live. The proof core -- the inseparable triad of field, commitment
scheme, and arithmetization -- means that Layer 5 cannot be understood in
isolation from the layers above and below it.

And we have seen a real system, Midnight, deploy Halo 2 in production with a
four-phase transaction pipeline whose proof generation latency we measured. This
is the state of the art for privacy-preserving blockchain computation:
mathematically rigorous, practically functional, and waiting for the performance
frontier to catch up.

The certificate is sealed. But we have been treating the sealing mechanism as a
black box -- we know it produces trustworthy certificates, but we have not
examined what makes the seal unforgeable. What mathematical hardness assumptions
prevent a forger from creating a convincing fake? Why do we believe these
assumptions hold? And what happens if a quantum computer shatters them?

These are the questions of Layer 6. The seal works because certain mathematical
problems are hard. The next chapter examines those problems -- the foundations
that make the entire magic trick possible.

\chapter{Chapter 7: Layer 6 -- The Deep
Craft}\label{chapter-7-layer-6-the-deep-craft}

\section{The Laws That Break}\label{the-laws-that-break}

Richard Feynman liked to say that the laws of physics do not change. You can
test them in New York or on the moon, today or in a billion years, and you get
the same answers. The constants are constant. The symmetries hold. Nature does
not update its firmware.

Zero-knowledge proof systems have their own ``laws of physics'' -- mathematical
assumptions about which problems are hard to solve. These assumptions sit
beneath every layer we have examined so far. The setup ceremonies of Layer 1,
the constraint systems of Layer 4, the proof engines of Layer 5 -- all of them
rest on a handful of beliefs about the difficulty of certain computations. If
those beliefs are correct, the entire tower stands. If they are wrong, it
collapses -- not gracefully, not partially, but completely.

Here is the uncomfortable question that Feynman would have asked, leaning
forward with that half-grin that meant he had spotted something everyone else
was politely ignoring: \emph{What happens when quantum computers change these
``laws''?}

Physics does not change. Mathematics does not change either. But our
\emph{assumptions} about which mathematical problems are hard -- those change
whenever someone invents a better attack. And a quantum computer running Shor's
algorithm is not a better attack. It is not a faster way to pick the same lock.
It is a different kind of physics applied to the same mathematics, and it
renders certain problems trivially easy that we have spent fifty years assuming
were impossibly hard.

This chapter descends to the deepest layer of the zero-knowledge stack: the
cryptographic primitives that everything else is built upon. It is about
hardness assumptions, commitment schemes, finite fields, and the coming quantum
reckoning. It is also about a revolution in progress -- a shift from one family
of mathematical foundations to another that may dissolve what looked like
permanent tradeoffs.

Here the metaphor reaches its limit. We cannot avoid the mathematics. But we can
make it concrete. Every abstraction in this chapter corresponds to a specific
engineering choice made by real teams building real systems. When we say ``the
Goldilocks field,'' we mean a specific 64-bit prime number. When we say
``Module-SIS,'' we mean a specific problem involving short vectors in
high-dimensional lattices. The goal is not to teach the mathematics but to
explain why these choices matter and what they cost.

\section{Three Hardness Assumptions, Three
Worlds}\label{three-hardness-assumptions-three-worlds}

Every cryptographic system rests on a \emph{hardness assumption}: a belief that
a specific mathematical problem cannot be solved efficiently. The entire
security guarantee is conditional. ``This proof system is sound'' really means
``this proof system is sound \emph{assuming} that problem X is hard.'' If
someone finds a fast algorithm for problem X, the security guarantee vanishes.
Not slowly. Immediately.

There is a useful way to think about this. A hardness assumption is like a
combination lock. Classical computers try every combination one at a time -- for
a lock with a trillion trillion combinations, they will never finish. Quantum
computers do not try faster. They exploit the lock's internal structure to
narrow the possibilities. Shor's algorithm does exactly this to the discrete
logarithm problem: it uses quantum interference to find the answer directly. The
lock opens. For hash-based problems, quantum computers get a modest advantage
but the lock still holds if you make it big enough. For lattice problems
(Module-SIS), no one has found a quantum trick that exploits the lock's
structure at all. The tumblers do not vibrate. The lock holds.

Three hardness assumptions dominate zero-knowledge cryptography. Each creates a
different world of possibilities and constraints.

\subsection{The Discrete Logarithm
Problem}\label{the-discrete-logarithm-problem}

The oldest and most widely deployed assumption. Given a number \(g\) and a value
\(h = g^x\), find \(x\). On ordinary computers, the best known algorithms
require roughly \(2^{128}\) operations for carefully chosen groups --
effectively impossible. This assumption powers all elliptic curve cryptography,
which in turn powers KZG commitments, Groth16 proofs, PLONK, and every
pairing-based SNARK.

The DLP (Discrete Logarithm Problem -- the mathematical puzzle of figuring out
how many times a number was multiplied by itself to produce a given result,
which is easy to state but very hard to solve) world offers deep algebraic
richness. Elliptic curve groups have a bilinear pairing operation -- a special
function that takes two curve points and produces an element in a ``target
group.'' This pairing is the engine behind KZG polynomial commitments, which
produce constant-size proofs (a single curve point, about 48 bytes) and enable
constant-time verification (one pairing check). Nothing else in cryptography
achieves this combination of succinctness and speed.

The cost is existential. A quantum computer running Shor's algorithm solves the
discrete logarithm problem in polynomial time. Not ``might solve'' -- \emph{does
solve}, given enough qubits. The DLP world has an expiration date. We do not
know the day. But the clock is ticking, and the hands do not run backward.

\subsection{Collision-Resistant Hash
Functions}\label{collision-resistant-hash-functions}

A completely different kind of assumption. A hash function takes arbitrary input
and produces a fixed-size output. ``Collision resistance'' means it is hard to
find two different inputs that produce the same output. SHA-256, BLAKE3, and
Poseidon are all collision-resistant hash functions (we believe).

The CRHF world is simpler and more conservative. Hash functions require no
algebraic structure -- no groups, no pairings, no special number theory. This
simplicity is both a strength (fewer assumptions to break) and a weakness (fewer
mathematical tools to work with). FRI-based commitment schemes and STARKs live
in this world. They are transparent (no trusted setup) and plausibly
post-quantum, since hash functions are not broken by Shor's algorithm.

But ``plausibly post-quantum'' deserves scrutiny, and scrutiny reveals cracks.
Grover's algorithm gives a quantum computer a quadratic speedup for brute-force
search, halving the effective security level of hash preimage resistance: a
256-bit hash drops to 128-bit quantum security. More subtly, the BHT algorithm
(Brassard-Hoyer-Tapp) can reduce collision resistance by a factor of three:
SHA-256's 128-bit classical collision resistance becomes roughly 85-bit quantum
collision resistance, though this attack requires impractical amounts of quantum
random-access memory. And the FRI protocol's post-quantum security depends on
the soundness of the Fiat-Shamir transform in the quantum random oracle model --
a reduction that is known but carries non-tight security bounds.

The honest statement is that hash-based systems \emph{probably} survive quantum
computers with appropriate parameter adjustments, but the unqualified claim that
they are ``post-quantum secure'' gives false confidence. Intellectual honesty
demands we say: this is not yet fully understood.

\subsection{Module-SIS (Module Short Integer
Solution)}\label{module-sis-module-short-integer-solution}

The newest and most mathematically demanding assumption. Given a matrix \(M\)
over a polynomial ring, find a short nonzero vector \(z\) such that
\(M \cdot z = 0\). ``Short'' means the coefficients of \(z\) are small. The best
known algorithms (both classical and quantum) require exponential time for
properly chosen parameters.

Module-SIS is the foundation of lattice-based cryptography -- the family that
the post-quantum community has rallied around. NIST's post-quantum standards
(FIPS 203, 204, and 205, published August 2024) are built on lattice problems.
The assumption has been studied for over two decades, and no quantum algorithm
significantly outperforms classical ones against it.

The lattice world offers a distinctive property: \emph{module homomorphism}. An
Ajtai commitment (the lattice analogue of a Pedersen commitment -- a Pedersen
commitment is a cryptographic method for ``sealing'' a number using elliptic
curve arithmetic, so the committed value can be verified later but cannot be
changed after commitment) satisfies the equation
\(\rho \cdot \text{Com}(Z) = \text{Com}(\rho \cdot Z)\), where \(\rho\) is a
ring element. This is the algebraic structure that makes lattice-based folding
schemes possible. It is weaker than what pairings provide (no bilinear map to a
target group) but stronger than what hash functions provide (which have no
algebraic structure at all).

The upshot is that lattice-based schemes occupy a useful middle ground: they
have enough algebraic structure for folding and efficient composition, plus
post-quantum security, plus transparent setup. Whether they can match the
succinctness of pairing-based schemes is the open research question -- and the
answer is converging toward ``close enough.''

\section{Four Families of Commitment
Schemes}\label{four-families-of-commitment-schemes}

The hardness assumptions crystallize into concrete cryptographic tools called
\emph{polynomial commitment schemes}. A commitment scheme lets a prover ``seal''
a polynomial into a short commitment, then later prove that the polynomial
evaluates to a specific value at a specific point. This is the mechanism that
makes zero-knowledge proofs work -- it is how the prover demonstrates knowledge
without revealing the underlying data.

Four families dominate. Each inherits the properties and limitations of its
underlying assumption, the way a building inherits the geology of the ground it
stands on.

\subsection{KZG (Kate-Zaverucha-Goldberg)}\label{kzg-kate-zaverucha-goldberg}

Built on the DLP and the bilinear pairing. The prover commits to a polynomial
\(p(x)\) as a single elliptic curve point \(C = g^{p(s)}\), where \(s\) is a
secret from a trusted setup ceremony (the ``powers of tau''). To prove that
\(p(z) = y\), the prover produces a single group element -- an evaluation proof
-- that the verifier checks with one pairing operation.

KZG is the gold standard for succinctness. Proof size: constant, roughly 48
bytes on BLS12-381 regardless of polynomial degree. Verification time: constant,
one pairing check. These are numbers that no other scheme matches. They are, in
a precise mathematical sense, optimal.

The costs are equally clear. KZG requires a trusted setup -- a structured
reference string generated by a multi-party computation ceremony. At least one
participant must honestly destroy their secret contribution, or the entire
system can be broken. The ceremony for Ethereum's EIP-4844 drew the six-figure
participant count described in Chapter 2. And KZG is not post-quantum: Shor's
algorithm breaks the pairing assumption along with the DLP.

KZG powers Groth16, PLONK, Marlin, and virtually every pairing-based SNARK. It
is used by Midnight, Zcash, most Ethereum rollups (at the final verification
layer), and the EIP-4844 blob commitment scheme.

To understand what a polynomial commitment \emph{feels like}, consider what it
accomplishes. A polynomial of degree \(n\) encodes \(n + 1\) independent values
-- it is, in a precise sense, a compressed representation of an entire dataset.
A polynomial of degree one million contains a million pieces of information. KZG
seals all of that information into a single elliptic curve point: 48 bytes. One
point on a curve, and behind it, a million values, invisible but committed.
Later, anyone can ask ``what does the polynomial evaluate to at point \(z\)?''
and the prover produces a single additional curve point as proof. Not a proof
proportional to the polynomial's size. Not a proof that grows with the
complexity of the claim. A single point. Constant size. Whether the polynomial
has degree ten or degree ten million, the proof is 48 bytes.

This is, in a precise mathematical sense, miraculous. It is worth pausing to
feel the weight of that claim, because no amount of familiarity should make it
seem ordinary.

The miracle rests on the bilinear pairing -- a function \(e(P, Q)\) that takes
two elliptic curve points and produces an element in a target group, satisfying
\(e(aP, bQ) = e(P, Q)^{ab}\). This bilinearity allows the verifier to check
polynomial relationships without ever seeing the polynomial. The proof that
\(p(z) = y\) consists of a commitment to the quotient polynomial
\(q(x) = (p(x) - y) / (x - z)\). The verifier checks one equation:
\(e(C - yG, H) = e(\pi, H_s - zH)\), where \(C\) is the commitment, \(\pi\) is
the proof, and \(H_s\) is a point from the trusted setup encoding the secret
\(s\). If the equation holds, the polynomial evaluates to \(y\) at \(z\). If it
does not, the prover is lying. One equation. Two pairing evaluations. Done.

The trusted setup deserves its own intuition. Think of it as a ruler with
markings at positions \(s, s^2, s^3, \ldots, s^n\), where \(s\) is a secret that
no one knows. The ruler is published as a sequence of elliptic curve points:
\(g, g^s, g^{s^2}, \ldots, g^{s^n}\). Anyone can use these markings to
``measure'' polynomial evaluations -- to compute \(g^{p(s)}\) for any polynomial
\(p\) by combining the markings with the polynomial's coefficients. But no one
can recover \(s\) itself, because extracting \(s\) from \(g^s\) requires solving
the discrete logarithm.

The analogy is not casual. It captures the essential structure: the SRS is a set
of calibrated instruments whose internal workings are opaque but whose external
behavior is perfectly reliable. You do not need to know the secret to use the
ruler. You need to trust that someone built the ruler honestly -- that the
markings correspond to genuine powers of a single unknown \(s\), and that \(s\)
was destroyed after construction. This is the trust assumption that the
multi-party ceremony enforces. If even one participant destroys their
contribution, the ruler is sound. If all participants collude (or are
compromised), they can forge proofs. The ceremony scales trust across thousands
of independent parties, diluting the assumption to its practical vanishing
point.

The mathematical miracle, then, is this: pairings allow you to verify claims
about a polynomial you have never seen, committed in a form that reveals nothing
about its coefficients, using a ruler whose markings you cannot read but whose
geometry you can trust. Constant-size commitments. Constant-size proofs.
Constant-time verification. These three constants are the reason KZG has
dominated practical zero-knowledge for half a decade, and they are the benchmark
against which every alternative scheme is measured.

\subsection{FRI (Fast Reed-Solomon Interactive Oracle Proof of
Proximity)}\label{fri-fast-reed-solomon-interactive-oracle-proof-of-proximity}

Built on collision-resistant hashing alone. FRI tests whether a function
(represented as a table of evaluations) is close to the evaluations of a
low-degree polynomial. It works by recursive ``folding'' -- halving the domain
at each step and checking consistency via random linear combinations, with
Merkle trees providing the commitment structure.

FRI proofs are transparent (no trusted setup) and plausibly post-quantum
(security from hashing, not from algebraic structure). Proof sizes are
polylogarithmic -- \(O(\log^2 n)\) in theory, typically 50 to 200 kilobytes in
practice. This is orders of magnitude larger than KZG's 48 bytes, but the
absence of a trusted setup and the quantum resistance are compelling
compensations.

FRI is the commitment scheme inside every STARK: StarkWare's Stwo, Polygon's
Plonky2 and Plonky3, SP1 Hypercube, and RISC Zero. It works best over fields
with large multiplicative subgroups of order \(2^k\), which is why
STARK-friendly fields like Goldilocks and BabyBear exist.

The limitation: FRI has no algebraic homomorphism. You cannot add two FRI
commitments and get a commitment to the sum of the polynomials. This means FRI
cannot support folding directly, which is why STARK systems use recursion
(proving that a proof is valid) rather than folding (combining two claims into
one).

The intuition behind FRI is proximity testing, and it deserves a concrete
picture. Suppose you have a table of values -- say, \(2^{20}\) entries -- and
you claim these values are the evaluations of a polynomial of degree at most
\(2^{10}\). In other words, you claim there exists a smooth, low-degree curve
that passes through all million-odd points. FRI's job is to test that claim
without reading the entire table.

Here is how it works. Imagine the claimed polynomial plotted on a graph -- a
smooth curve undulating through a million points on the horizontal axis. FRI
asks: ``Is this really a smooth low-degree curve, or is it a jagged high-degree
impostor that happens to agree with a low-degree polynomial at most points?''
The test proceeds by \emph{folding}. The verifier sends a random challenge
\(\alpha\). The prover uses \(\alpha\) to combine pairs of evaluations: for each
point \(x\) in the domain, the prover computes a new value from \(f(x)\) and
\(f(-x)\), weighted by \(\alpha\). This produces a new table of half the size,
over a domain of half the width. If the original function was degree \(d\), the
folded function has degree \(d/2\).

Now repeat. Another random challenge. Another folding. The domain halves again.
The degree halves again. Each round is a ``zoom in'' -- the verifier is looking
at the function at finer and finer resolution, checking whether the smoothness
persists. If the original function was truly a low-degree polynomial, every
zoomed-in version remains smooth. The degree drops: \(2^{10}\), then \(2^9\),
then \(2^8\), all the way down to a constant. At the end, the prover reveals the
final polynomial directly -- it is small enough to check by inspection.

But if the original function was \emph{not} close to a low-degree polynomial --
if it was a high-degree impostor with hidden bumps -- then folding amplifies the
bumps. Each round of folding, guided by the random challenge, mixes pairs of
evaluations in a way that smooths genuine structure but destabilizes imposture.
The bumps do not cancel; they compound. By the time you have zoomed in far
enough, the remaining function is visibly not low-degree. The impostor is
caught.

The commitment mechanism is beautifully simple: Merkle trees. The prover commits
to each round's evaluation table by hashing it into a Merkle tree and publishing
the root. When the verifier wants to spot-check specific points, the prover
opens Merkle paths -- logarithmic-sized authentication paths from leaf to root.
The security rests entirely on collision resistance of the hash function. No
group structure, no pairings, no discrete logs. Just hashing.

This is why FRI replaces elliptic curves with hash functions. A Merkle tree
commitment to \(n\) evaluations costs \(O(n)\) hashes to build and \(O(\log n)\)
hashes to open a single leaf. The total proof consists of Merkle roots (one per
FRI round), Merkle paths (for the queried positions), and the final constant
polynomial. The result is polylogarithmic: \(O(\log^2 n)\) in total size,
typically 50 to 200 kilobytes. Enormously larger than KZG's 48 bytes. But
transparent. Plausibly post-quantum. And built from the simplest, most
conservative cryptographic primitive we have: the hash function.

The tradeoff crystallizes the philosophical divide in zero-knowledge
cryptography. KZG achieves its miracle of constant-size proofs by leveraging
deep algebraic structure -- bilinear pairings over elliptic curves, with all the
trust assumptions and quantum vulnerabilities that entails. FRI achieves
transparency and quantum plausibility by abandoning that structure entirely,
accepting larger proofs as the price. Neither choice is wrong. Each is a
coherent answer to a different question about what you are willing to assume and
what you are willing to pay.

\subsection{IPA / Bulletproofs (Inner Product
Argument)}\label{ipa-bulletproofs-inner-product-argument}

Built on the discrete logarithm problem without pairings. Bulletproofs use
Pedersen commitments -- a simpler construction than KZG that does not require
bilinear maps. The key insight is an inner product argument: prove that the
inner product of two committed vectors equals a claimed value, using a recursive
halving protocol that produces \(O(\log n)\) group elements.

IPA proofs are transparent (no trusted setup, just a random group element
generator). Proof sizes are logarithmic -- much smaller than FRI, but not
constant like KZG. Verification, however, requires \(O(n)\) work -- linear in
the statement size. This is the main drawback: the verifier is slow.

Halo (2019) proved that IPA-based schemes support recursion without pairings,
using a technique called ``nested amortization'' that defers expensive
verification across recursion steps. This was the conceptual breakthrough that
opened the door to transparent recursive proving. Halo2 powers Zcash's Orchard
protocol and influenced the design of the Pasta curves (Pallas and Vesta).

The limitation: IPA is not post-quantum. It still relies on the discrete
logarithm assumption. And the linear verification cost makes it impractical for
on-chain verification without wrapping in a more succinct outer proof.

The inner product argument deserves to be understood on its own terms, because
it is one of the most elegant constructions in modern cryptography. The problem
it solves is this: you have committed to two vectors \(\mathbf{a}\) and
\(\mathbf{b}\) of length \(n\), and you claim their inner product is some value
\(c\). You want to prove this claim without revealing the vectors. The naive
approach would require the verifier to reconstruct the entire inner product --
\(n\) multiplications, \(n - 1\) additions, linear work. The inner product
argument reduces this to logarithmic work.

The technique is recursive halving. Split both vectors into their left and right
halves: \(\mathbf{a} = (\mathbf{a_L}, \mathbf{a_R})\) and
\(\mathbf{b} = (\mathbf{b_L}, \mathbf{b_R})\). Compute two ``cross-terms'':
\(L = \text{Commit}(\mathbf{a_L}, \mathbf{b_R})\) and
\(R = \text{Commit}(\mathbf{a_R}, \mathbf{b_L})\). Publish \(L\) and \(R\). The
verifier sends a random challenge \(x\). Both parties compute new vectors:
\(\mathbf{a'} = x \cdot \mathbf{a_L} + x^{-1} \cdot \mathbf{a_R}\) and
\(\mathbf{b'} = x^{-1} \cdot \mathbf{b_L} + x \cdot \mathbf{b_R}\). The new
vectors have half the length. The new inner product claim can be derived from
the old one plus the cross-terms. Repeat.

Each halving adds exactly two group elements (\(L\) and \(R\)) to the proof. For
vectors of length \(n = 2^{20}\) -- roughly a million entries -- the protocol
runs for 20 rounds, producing 40 group elements. At 32 bytes per element on a
256-bit curve, that is 1,280 bytes. A proof that two million-entry vectors have
a specific dot product, in 1.3 kilobytes. Not constant like KZG. But
logarithmic, transparent, and no trusted setup.

This is why Halo was a watershed. Before Halo, recursive proof composition
required pairings -- you needed the bilinear map to verify a KZG proof inside a
circuit, and that meant you needed pairing-friendly curves, which meant you
needed a trusted setup. Halo demonstrated that IPA's logarithmic proofs could be
verified \emph{incrementally} across recursion steps, deferring the expensive
linear-time final check. The key was nested amortization: instead of verifying
the IPA proof fully at each recursion step (which would cost linear time and
destroy the efficiency), Halo accumulated the verification work and deferred it
to the end. This meant each recursion step added only constant overhead, and the
full linear verification happened once, at the very end of the recursion chain.

The result was the first recursive proof system with no trusted setup and no
pairings. It required a cycle of curves -- Pallas and Vesta, the ``Pasta''
curves, each one's scalar field being the other's base field -- but no ceremony,
no toxic waste, no trust assumptions beyond the hardness of the discrete
logarithm. Zcash adopted Halo2 for its Orchard shielded pool, replacing the
Sprout and Sapling ceremonies with transparent recursion. The ceremony was over.
The mathematics was enough.

\subsection{Lattice / Ajtai Commitments}\label{lattice-ajtai-commitments}

Built on Module-SIS. The Ajtai commitment scheme works over a cyclotomic ring
\(R_q = \mathbb{F}_q[X]/(\Phi(X))\). To commit to a vector \(Z\), compute
\(\text{Com}(Z) = M \cdot Z\) where \(M\) is a public random matrix over the
ring. The binding property reduces to the Module-SIS problem: forging a
commitment requires finding a short vector in a lattice.

Lattice commitments are transparent (the matrix \(M\) is generated from public
randomness) and post-quantum (Module-SIS resists quantum attacks). Proof sizes
are logarithmic -- \(O(\log n)\) ring elements, concretely 50 to 60 kilobytes
for current parameterizations.

The key algebraic property: Ajtai commitments are \emph{module-homomorphic} over
the ring. For a ring element \(\rho\) and a commitment \(\text{Com}(Z)\), you
get \(\rho \cdot \text{Com}(Z) = \text{Com}(\rho \cdot Z)\). This is strictly
richer than scalar homomorphism and is what enables lattice-based folding. The
challenges used in folding are ring elements from a carefully chosen ``strong
sampling set'' with small coefficients, which controls the norm growth of folded
witnesses.

Lattice commitments power Greyhound, LaBRADOR, LatticeFold, LatticeFold+, Neo,
and Symphony. They are the youngest family and the least battle-tested in
production, but they are the only family that simultaneously offers post-quantum
security, algebraic structure sufficient for folding, and transparent setup.

The geometry of the short vector problem rewards intuition. Imagine a forest of
trees planted in a perfectly regular grid -- rows and columns spaced exactly one
meter apart, stretching to the horizon in every direction. You are dropped at a
random point in this forest. Finding the nearest tree is trivial: you can see
the grid, feel its regularity, walk to the closest intersection. Now imagine the
same forest, but the grid has been distorted. The trees are no longer evenly
spaced. The rows curve. The columns tilt. The local pattern near any tree looks
regular, but the global structure is scrambled by a ``bad'' basis -- a set of
directions that are long, nearly parallel, and unhelpful for navigation. You can
still see trees around you, but finding the \emph{nearest} tree -- the closest
lattice point to your position -- has become exponentially hard. You wander
among trees that look close but are not closest. The geometry of the lattice
hides the answer in plain sight.

This is the Closest Vector Problem, and its computational cousin the Shortest
Vector Problem, on which Module-SIS rests. The matrix \(M\) in the Ajtai
commitment defines the lattice. The ``short vector'' condition -- finding \(z\)
such that \(M \cdot z = 0\) and each component of \(z\) is bounded by a
parameter \(\beta\) -- is the lattice analogue of finding that nearest tree. The
commitment to a message \(m\) uses randomness \(r\):
\(\text{commit} = M \cdot [m;\, r]\). Binding holds because finding two
different \((m, r)\) pairs that produce the same commitment requires finding a
short difference vector -- a short element in the kernel of \(M\). If the
lattice is properly parameterized (high dimension, appropriate modulus), no
efficient algorithm can find such a vector.

Why do lattices resist quantum computers? The answer is structural. Shor's
algorithm exploits \emph{periodicity}. The discrete logarithm problem has a
hidden periodic structure: given \(g^x\), the function
\(f(a, b) = g^a \cdot h^b\) is periodic with period related to \(x\). Shor's
quantum Fourier transform finds this period efficiently. Lattice problems have
no such periodicity. The Closest Vector Problem is not periodic -- it is
geometric. The difficulty comes from the high dimensionality and the scrambled
basis, not from any hidden algebraic cycle. Grover's algorithm provides a
generic quadratic speedup for unstructured search (turning \(2^{128}\) into
\(2^{64}\)), but lattice algorithms are not brute-force searches. The best
lattice algorithms (BKZ and its variants) operate by finding short vectors in
projected sublattices, and no quantum algorithm significantly accelerates this
process. The lattice estimator -- the standard tool for selecting parameters --
accounts for Grover's quadratic speedup and still produces comfortable security
margins. Neo's 127-bit post-quantum security, for instance, already incorporates
the best known quantum attacks.

The Module-SIS formulation adds one more layer. Instead of vectors over a field,
you work with vectors over a polynomial ring
\(R_q = \mathbb{F}_q[X]/(\Phi(X))\). Each ``entry'' in the vector is itself a
polynomial -- a ring element with \(d\) coefficients. This means a vector of
length \(\kappa\) over the ring actually contains \(\kappa \cdot d\) field
elements, giving the lattice its high dimension. The module structure is what
provides the algebraic homomorphism: because the ring has multiplication, the
commitment scheme inherits module-homomorphic properties that a plain integer
lattice would not provide. It is this marriage of lattice hardness and ring
algebra that makes the entire folding program possible.

\subsection{The Four Families at a Glance}\label{the-four-families-at-a-glance}

\begin{longtable}[]{@{}
  >{\raggedright\arraybackslash}p{(\linewidth - 8\tabcolsep) * \real{0.2000}}
  >{\raggedright\arraybackslash}p{(\linewidth - 8\tabcolsep) * \real{0.2000}}
  >{\raggedright\arraybackslash}p{(\linewidth - 8\tabcolsep) * \real{0.2000}}
  >{\raggedright\arraybackslash}p{(\linewidth - 8\tabcolsep) * \real{0.2000}}
  >{\raggedright\arraybackslash}p{(\linewidth - 8\tabcolsep) * \real{0.2000}}@{}}
\toprule\noalign{}
\begin{minipage}[b]{\linewidth}\raggedright
Property
\end{minipage} & \begin{minipage}[b]{\linewidth}\raggedright
KZG
\end{minipage} & \begin{minipage}[b]{\linewidth}\raggedright
FRI
\end{minipage} & \begin{minipage}[b]{\linewidth}\raggedright
IPA/Bulletproofs
\end{minipage} & \begin{minipage}[b]{\linewidth}\raggedright
Lattice/Ajtai
\end{minipage} \\
\midrule\noalign{}
\arrayrulecolor{MidnightBlue}\midrule\arrayrulecolor{MidnightBorder}
\endhead
\bottomrule\noalign{}
\endlastfoot
\textbf{Proof size} & \(O(1)\), \textasciitilde48 bytes & \(O(\log^2 n)\),
\textasciitilde50-200 KB & \(O(\log n)\), \textasciitilde1-5 KB & \(O(\log n)\),
\textasciitilde50-60 KB \\
\textbf{Verification} & \(O(1)\) pairings & \(O(\log^2 n)\) hashes & \(O(n)\)
group ops & \(O(\sqrt{n})\) or \(O(\log n)\) \\
\textbf{Trusted setup} & Yes (SRS) & No & No & No \\
\textbf{Post-quantum} & No & Plausible & No & Yes (MSIS) \\
\textbf{Homomorphic} & Additive & No & Additive & Module-homomorphic \\
\textbf{Algebraic structure} & Rich (pairings) & Minimal (hashes) & Moderate
(DLOG) & Rich (ring ops) \\
\end{longtable}

\section{The Trilemma -- And Its
Dissolution}\label{the-trilemma-and-its-dissolution}

The original paper presented a ``cryptographic primitives trilemma'': a claim
that any commitment scheme can achieve at most two of three desirable
properties.

\begin{enumerate}
\def\labelenumi{\arabic{enumi}.}
\tightlist
\item
  \textbf{Algebraic functionality} -- the homomorphic structure needed for
  folding, composition, and efficient recursive proving.
\item
  \textbf{Post-quantum security} -- resilience against quantum computers running
  Shor's and Grover's algorithms.
\item
  \textbf{Succinctness} -- small proofs and fast verification.
\end{enumerate}

The trilemma positioned the four families like this:

\begin{itemize}
\tightlist
\item
  \textbf{KZG} achieves algebraic functionality and succinctness but lacks
  post-quantum security.
\item
  \textbf{FRI} achieves post-quantum security but lacks algebraic functionality
  (no homomorphism, so no folding) and offers only moderate succinctness (large
  proofs).
\item
  \textbf{IPA} achieves moderate algebraic functionality but lacks both
  post-quantum security and full succinctness (linear verification).
\item
  \textbf{Lattice} achieves algebraic functionality and post-quantum security.
  Succinctness is the remaining gap.
\end{itemize}

This framing was useful as a historical snapshot. As a statement of permanent
truth, it is increasingly wrong.

The lattice revolution -- Greyhound (2024), LatticeFold (2024), LatticeFold+
(2025), Neo (2025), Symphony (2026) -- has been systematically closing the
succinctness gap. Greyhound demonstrated 50-kilobyte proofs with sublinear
verification. LaBRADOR achieved 58-kilobyte proofs for large constraint systems.
Symphony's high-arity folding can compress a final proof via a compact SNARK
that, if instantiated with a pairing-based scheme, produces constant-size output
-- and if instantiated with a lattice-based scheme, remains fully post-quantum.

The trilemma is better understood as a \emph{spectrum} that is being actively
compressed. The engineering challenge is real (lattice proofs are still 1000x
larger than KZG), but the trajectory is clear: lattice schemes are approaching
practical competitiveness, and the gap shrinks with each generation. What looked
like a permanent constraint on the geometry of the design space is turning out
to be an artifact of our current engineering, not a law of mathematical nature.

To see the trilemma clearly, state the three properties any polynomial
commitment scheme would ideally achieve:

\begin{enumerate}
\def\labelenumi{\arabic{enumi}.}
\tightlist
\item
  \textbf{Small proofs, constant size.} The proof should not grow with the size
  of the polynomial. Ideally, one group element or a fixed number of bytes,
  regardless of degree.
\item
  \textbf{Fast verification, constant time.} The verifier's work should not
  depend on the polynomial's complexity. Ideally, a single algebraic check.
\item
  \textbf{No trusted setup, transparent.} The scheme should require no ceremony,
  no toxic waste, no trust assumptions beyond the hardness of a mathematical
  problem.
\end{enumerate}

No known scheme achieves all three. KZG achieves the first two but requires a
trusted setup ceremony. FRI achieves the third (transparent) and offers
reasonable verification (polylogarithmic), but its proofs are polylogarithmic
rather than constant -- orders of magnitude larger. IPA achieves the third
(transparent) with logarithmic proofs (impressively small), but its verification
is linear -- the verifier must do work proportional to the polynomial's degree.
Lattice commitments achieve the third (transparent) with logarithmic proofs, but
verification is sublinear rather than constant.

The question that should keep a mathematician awake at night is: \emph{is this
trilemma fundamental?} Is there a theorem -- an impossibility result, an
information-theoretic lower bound -- proving that no commitment scheme can
simultaneously achieve constant-size proofs, constant-time verification, and
transparency?

The answer, as of 2026, is no. No one has proven that the ideal PCS is
impossible. The barriers are engineering barriers, not mathematical barriers.
The bilinear pairing that gives KZG its constant-size miracle is a specific
algebraic structure tied to elliptic curves, and elliptic curves require
structured reference strings to exploit pairings. But nothing in information
theory says that constant-size polynomial commitments \emph{require} pairings.
Nothing says that transparency \emph{requires} large proofs. The ideal scheme --
transparent, constant-size, constant-verification, post-quantum -- remains the
field's holy grail. It may not exist. But its impossibility has not been proven,
and the gap between what lattice schemes achieve today and what that grail
demands shrinks with every new construction. The trilemma may be less a law of
nature than a confession of our current ignorance.

\section{Small Fields}\label{small-fields}

If the commitment scheme is the ``which family'' decision, the finite field is
the ``which numbers'' decision. And this choice -- seemingly a detail buried
deep in the mathematics -- turns out to determine nearly everything about
performance.

A finite field is a set of numbers equipped with addition and multiplication
that ``wrap around'' at a fixed boundary, called the \emph{modulus} or
\emph{prime}. Every value in a zero-knowledge circuit is an element of some
finite field. Every arithmetic operation is performed modulo the field's prime.
The choice of prime cascades upward through every layer of the system. It is one
of those decisions that seems technical and narrow when you make it, and then
turns out to have been the most important decision you made.

\subsection{The Old World: BN254 and
BLS12-381}\label{the-old-world-bn254-and-bls12-381}

For most of the 2010s, the ZK world standardized on two large primes:

\textbf{BN254} (Barreto-Naehrig curve, 254-bit prime). This was the first widely
deployed pairing-friendly curve. Ethereum embedded BN254 pairing operations as
EVM precompiles in 2017, making it the de facto standard for on-chain Groth16
verification. Every deployed Groth16 verifier on Ethereum -- including those
used by the major rollups -- runs over BN254.

\textbf{BLS12-381} (Barreto-Lynn-Scott curve, \textasciitilde253-bit prime).
Introduced later with a higher security margin and better pairing efficiency.
Used by Zcash, Filecoin, Midnight, and the Ethereum KZG ceremony (EIP-4844).

Both are enormous primes -- 254 and 253 bits respectively. Arithmetic on 254-bit
numbers requires multiple machine words on any existing processor. A single
multiplication takes several CPU instructions and cannot exploit the native
32-bit or 64-bit arithmetic units that modern hardware is optimized for.

\subsection{The Security Erosion Problem}\label{the-security-erosion-problem}

BN254 was originally believed to provide 128-bit security -- meaning an attacker
would need roughly \(2^{128}\) operations to break the discrete logarithm. But
advances in the Tower Number Field Sieve (Tower NFS) {[}Kim and Barbulescu,
``Extended Tower Number Field Sieve,'' Mathematics of Computation, 2016;
Guillevic, ``Comparing the pairing efficiency over composite-order and
prime-order elliptic curves,'' ACNS 2013{]} have revised this estimate downward
to approximately 100 bits. This does not mean BN254 is broken -- \(2^{100}\)
operations is still astronomically expensive -- but it means the security margin
is significantly thinner than designed.

This matters because BN254 is embedded in Ethereum's EVM precompiles. Changing
the precompiled curves requires a hard fork. Every Groth16 verifier on Ethereum
depends on BN254. The security erosion is not academic -- it affects the most
widely deployed zero-knowledge infrastructure in the world.

BLS12-381 is not affected by Tower NFS and retains its 128-bit security
estimate. But it is also not immune to future cryptanalytic advances. And
neither curve survives a quantum computer.

\subsection{The New World: BabyBear, M31,
Goldilocks}\label{the-new-world-babybear-m31-goldilocks}

Starting around 2022, a radical idea took hold: \emph{use much smaller primes}.

\textbf{BabyBear} (\(p = 2^{31} - 2^{27} + 1\), a 31-bit prime). Fits in a
single 32-bit machine word. Arithmetic is native on every modern CPU and GPU.
Used by RISC Zero and Plonky3.

\textbf{Mersenne-31 / M31} (\(p = 2^{31} - 1\), a 31-bit Mersenne prime). The
simplest possible arithmetic -- reduction modulo a Mersenne prime is a single
addition. Used by StarkWare's Stwo and Circle STARKs.

\textbf{Goldilocks} (\(p = 2^{64} - 2^{32} + 1\), a 64-bit prime). Fits in a
single 64-bit machine word. Has high 2-adicity (\(2^{32}\) divides \(p - 1\)),
enabling efficient FFTs with domain sizes up to \(2^{32}\). Used by Polygon's
Plonky2 and by Neo/Nightstream.

The performance impact is not incremental. It is a factor of 100. Arithmetic on
31-bit numbers is roughly 100 times faster than arithmetic on 254-bit numbers
{[}Haboeck, Levit, and Papini, ``Circle STARKs,'' ePrint 2024/278; confirmed by
SP1 Hypercube benchmarks, Succinct Labs, 2025{]}. This is not algorithmic
improvement -- it is the raw physics of computer hardware. A 31-bit multiply is
one CPU instruction. A 254-bit multiply is an entire subroutine involving carry
propagation, multi-limb multiplication, and modular reduction.

This single insight -- that smaller fields make faster provers -- catalyzed the
performance explosion in zero-knowledge proving. Circle STARKs over M31 (Stwo)
achieve throughputs that were unimaginable with BN254-based systems. Plonky2
over Goldilocks enabled the first practical recursive STARKs.

But smaller fields introduce a subtlety that Penrose would appreciate for its
geometric elegance. A single 31-bit field element provides only 31 bits of
security against certain attacks. To achieve 128-bit security, systems use
\emph{extension fields}. An extension field is built by the same trick as
complex numbers: you take a small field and add extra ``dimensions'' to your
arithmetic, and the security grows with the dimension. The cost is slightly more
expensive arithmetic per operation -- but each operation now works in a larger,
more secure space -- enlarging \(\mathbb{F}_p\) to \(\mathbb{F}_{p^k}\) for some
small \(k\). In Stwo, the extension degree is 4, giving effectively 124 bits. In
Neo, the extension is \(\mathbb{F}_{q^2}\) over Goldilocks, giving 128 bits. The
extension adds complexity but the arithmetic is still vastly cheaper than native
254-bit operations.

\subsection{Why This Choice Is a One-Way
Door}\label{why-this-choice-is-a-one-way-door}

The field choice is perhaps the most consequential ``one-way door'' decision in
zero-knowledge system design. It cascades through every layer:

\begin{itemize}
\tightlist
\item
  The field determines which commitment schemes are efficient (pairing-friendly
  fields for KZG, STARK-friendly fields for FRI, cyclotomic fields for lattice
  schemes).
\item
  The commitment scheme determines which proof systems work (KZG enables
  Groth16/PLONK, FRI enables STARKs, Ajtai enables lattice folding).
\item
  The proof system determines the arithmetization format (PLONK gates, AIR,
  CCS).
\item
  The arithmetization determines which programs can be efficiently proved.
\end{itemize}

Changing the field after deployment means rewriting the compiler, the prover,
the verifier, the standard library, and every circuit. It is not a parameter
change. It is a complete system redesign.

\section{The Quantum Threat Horizon}\label{the-quantum-threat-horizon}

The question is not whether quantum computers will break DLP-based cryptography.
The question is when.

\subsection{What Shor's Algorithm Does}\label{what-shors-algorithm-does}

Shor's algorithm, published in 1994, solves the discrete logarithm problem in
polynomial time on a quantum computer. Given \(g\) and \(h = g^x\), it finds
\(x\) using roughly \(O(n^3)\) quantum gates, where \(n\) is the bit length of
the group order. For BLS12-381, this requires approximately 2,500 logical
qubits. Each logical qubit requires thousands of physical qubits for error
correction, meaning the actual hardware requirement is on the order of millions
of physical qubits -- roughly three orders of magnitude beyond current devices,
which have demonstrated only a few thousand physical qubits.

When a cryptographically relevant quantum computer (CRQC) exists, the
consequences are immediate and total:

\begin{itemize}
\tightlist
\item
  Every KZG commitment ever published becomes forgeable.
\item
  Every Groth16 proof ever verified becomes suspect.
\item
  Every elliptic curve signature (including BLS signatures and ECDSA) breaks.
\item
  Every Pedersen commitment loses its binding property.
\item
  Every system built on pairing-based cryptography fails simultaneously.
\end{itemize}

The failure mode is not gradual degradation. It is a cliff. One day the lock
holds. The next day every lock of that type, everywhere, opens at once.

\subsection{The Timeline}\label{the-timeline}

Estimates for when a CRQC will exist vary widely:

\begin{itemize}
\tightlist
\item
  \textbf{Optimistic (Forrester, 2024):} ``Q-Day'' by 2030.
\item
  \textbf{Conservative (academic consensus):} 2032-2035.
\item
  \textbf{NIST IR 8547 (November 2024):} Recommends that all federal agencies
  deprecate pre-quantum cryptographic algorithms by 2035.
\end{itemize}

NIST's position is the most policy-relevant. In August 2024, NIST published
three post-quantum cryptography standards: FIPS 203 (ML-KEM, a key encapsulation
mechanism based on Module-LWE), FIPS 204 (ML-DSA, a digital signature based on
Module-LWE), and FIPS 205 (SLH-DSA, a stateless hash-based signature). These are
not draft standards or proposals -- they are finalized, mandatory standards for
federal systems.

The IR 8547 guidance document sets a deadline: by 2035, federal systems must
have migrated away from RSA, ECDSA, and all DLP-based cryptography. The
reasoning is straightforward: if a CRQC arrives in 2035 and a migration takes
5-10 years, you need to start migrating now.

\subsection{The HNDL Threat}\label{the-hndl-threat}

The most insidious quantum threat is not future code-breaking but present data
collection. ``Harvest Now, Decrypt Later'' (HNDL) describes the strategy of
recording encrypted communications today for decryption by future quantum
computers. A Federal Reserve discussion paper (FEDS 2025-093) explicitly
identified HNDL as a risk to financial infrastructure.

For zero-knowledge systems, the HNDL analogue is this: an adversary records all
on-chain data -- commitments, proofs, public inputs -- today, waiting for a
quantum computer to extract the underlying secrets. In a system like Zcash or
Midnight, where commitments hide transaction amounts and sender/receiver
identities, a future CRQC could retroactively de-anonymize the entire
transaction history.

The concern is not theoretical. The blockchain is a permanent, public record.
Nothing posted to Ethereum or any other blockchain can be deleted. Every
BLS12-381 commitment, every BN254 proof, every elliptic curve public key is
preserved forever, waiting. The data is patient. A quantum computer needs only
to be built once.

\subsection{Deployed Systems at Risk}\label{deployed-systems-at-risk}

Any zero-knowledge system deployed today that relies on DLP-based cryptography
faces a choice:

\begin{enumerate}
\def\labelenumi{\arabic{enumi}.}
\item
  \textbf{Accept the expiration date.} Acknowledge that the system's security
  guarantees have a finite lifespan and plan accordingly. This is the pragmatic
  approach for systems that do not require long-term privacy (e.g., rollups
  where the transactions are already publicly visible).
\item
  \textbf{Migrate proactively.} Begin transitioning to post-quantum primitives
  before a CRQC exists. This is the only option for systems that provide
  long-term privacy guarantees (e.g., shielded transactions, confidential
  identity systems).
\item
  \textbf{Ignore the problem.} Hope that quantum computers take longer than
  expected, or that ``crypto-agility'' will allow a rapid migration when the
  time comes. This is the most common approach -- and the most dangerous,
  because ``crypto-agility'' in practice means ``complete architectural
  redesign.''
\end{enumerate}

\section{Lattice-Based Proving}\label{lattice-based-proving}

Against this backdrop -- the ticking clock, the harvest-now threat, the cliff
edge -- a research program has been steadily building an alternative:
lattice-based zero-knowledge proof systems that provide post-quantum security
without sacrificing the algebraic structure needed for efficient proving.

The progression has been fast, compressing a decade of typical cryptographic
development into two years.

\subsection{Stage 1: Greyhound (2024)}\label{stage-1-greyhound-2024}

The first demonstration that lattice-based SNARKs could be practical, not just
theoretical. Greyhound achieved approximately 50-kilobyte proofs with sublinear
(square root of N) verification, built entirely on Module-SIS. It was a
standalone SNARK, not a folding scheme -- a proof of concept that lattice proofs
could fit in the same order of magnitude as STARK proofs.

\subsection{Stage 2: LatticeFold (2024)}\label{stage-2-latticefold-2024}

The conceptual breakthrough. Dan Boneh and Binyi Chen adapted Nova-style folding
to work over cyclotomic rings with Ajtai commitments. The key insight: Ajtai's
module homomorphism is the lattice analogue of Pedersen's additive homomorphism,
and it is sufficient to enable the random-linear-combination technique that
makes folding work.

LatticeFold introduced three composable reductions -- \(\Pi_{\text{CCS}}\)
(constraint satisfaction to evaluation claims), \(\Pi_{\text{RLC}}\) (random
linear combination), and \(\Pi_{\text{DEC}}\) (decomposition to control norm
growth) -- that together form a complete folding scheme for CCS (Customizable
Constraint Systems).

But LatticeFold had a critical limitation. It was restricted to power-of-two
cyclotomic polynomials of the form \(X^d + 1\). Over the Goldilocks field, this
polynomial splits completely into degree-1 factors, meaning each NTT slot has
only 64-bit security -- insufficient for the 128-bit target.

\subsection{Stage 3: LatticeFold+ (2025)}\label{stage-3-latticefold-2025}

A comprehensive improvement: 5 to 10 times faster prover, simpler verification
circuit, shorter folding proofs, and a new purely algebraic range proof
replacing LaBRADOR's more complex approach. LatticeFold+ identified three
concrete parameterizations, including one over Goldilocks with the 81st
cyclotomic polynomial \(\Phi_{81} = X^{54} + X^{27} + 1\). This polynomial does
not split into linear factors over Goldilocks; instead, it factors into degree-2
irreducibles, giving the extension field \(K = \mathbb{F}_{q^2}\) with 128-bit
NTT slots.

LatticeFold+ also introduced a general composition theorem: if one reduction is
``phi-restricted'' with restricted knowledge soundness, and another is
``phi-relaxed'' knowledge sound, their composition is knowledge sound. This
provided the formal foundation for chaining reductions.

\subsection{Stage 4: Neo (2025)}\label{stage-4-neo-2025}

Neo overcame LatticeFold's power-of-two limitation directly. Its central
innovation is the \emph{rotation matrix encoding} -- the ``bar transform'' --
which represents ring elements as rotation matrices in a commutative subring of
the matrix ring. The map sends a ring element \(a\) in \(R_q\) to a
\(d \times d\) matrix \(\text{rot}(a)\), and this map is a ring isomorphism. The
payoff: the matrix commitment scheme becomes S-homomorphic:
\(\text{rot}(\rho) \cdot \text{Com}(Z) = \text{Com}(\text{rot}(\rho) \cdot Z)\).

Neo works natively over Goldilocks with \(\Phi_{81}\), giving \(d = 54\),
\(\kappa = 16\), \(m = 2^{24}\), and 127-bit post-quantum security (verified via
the standard lattice estimator). The guard condition
\((k+1) \cdot T \cdot (b-1) = 2{,}808 < 4{,}096 = B\) ensures that norm growth
remains bounded across arbitrarily many folding steps.

\subsection{Stage 5: Symphony (2026)}\label{stage-5-symphony-2026}

The most ambitious design. Symphony pushes folding to high arity -- folding
1,024 or more instances in a single step, rather than the standard two. This
eliminates the need for recursive IVC (which requires embedding hash
verification inside the SNARK circuit, a major overhead). Symphony folds
\(\text{poly}(\lambda)\) NP statements into two committed linear statements in
one shot, then proves these with a compact SNARK.

Symphony also introduces approximate range proofs (replacing the exact norm
proofs of LatticeFold+), reducing verification complexity further. Its concrete
instantiation can handle \(2^{32}\) R1CS constraints -- over four billion -- in
a single batch.

If Symphony's compact SNARK is instantiated with a pairing-based scheme
(Groth16), the final proof is constant-size. If instantiated with a
lattice-based scheme, the entire pipeline is post-quantum. This modularity is
the architectural insight: separate the bulk proving (which must be
post-quantum) from the final compression (which can optionally use classical
tools for maximum succinctness).

\subsection{The Key Algebraic Insight}\label{the-key-algebraic-insight}

The entire lattice folding line rests on a single algebraic fact that deserves
to be stated plainly, because it is the kind of fact that sounds narrow but
turns out to govern everything.

An Ajtai commitment over the ring \(R_q\) is \emph{module-homomorphic}. This
means that for any challenge element \(\rho\) drawn from a strong sampling set
with small coefficients, the equation
\(\rho \cdot \text{Com}(Z) = \text{Com}(\rho \cdot Z)\) holds. This is the
lattice analogue of the scalar homomorphism that makes Nova-style folding work
over elliptic curves.

But the lattice version is richer. The challenge \(\rho\) is not a scalar but a
\emph{ring element} -- equivalently, a \(d \times d\) rotation matrix. This
richer structure enables three things simultaneously:

\begin{enumerate}
\def\labelenumi{\arabic{enumi}.}
\tightlist
\item
  \textbf{Folding with norm control.} Challenges from the strong sampling set
  \(\mathcal{C}\) have small coefficients (in \(\{-2, -1, 0, 1, 2\}\) for Neo),
  so the folded witness grows slowly in norm.
\item
  \textbf{Sum-check compatibility.} Ring evaluation claims can be verified via
  the sum-check protocol over the base field, connecting the commitment layer to
  the constraint satisfaction layer.
\item
  \textbf{Decomposition.} The accumulated norm can be reduced back to the base
  bound \(b\) via bit-decomposition, enabling unbounded recursion without norm
  blowup.
\end{enumerate}

This algebraic trifecta -- homomorphism, sum-check compatibility, and
decomposition -- is what makes lattice-based folding possible. It is the ``deep
craft'' of Layer 6: not a single clever trick, but an interlocking set of
algebraic properties that together provide something no other family offers. The
geometry of the lattice (its distances, its short vectors, its algebraic
symmetries) does the work that pairings do in the elliptic curve world, but
without the quantum vulnerability.

\section{Case Study: Midnight}\label{case-study-midnight}

To see how Layer 6 choices play out in a real system, consider Midnight -- a
privacy-focused blockchain built by Input Output Global (IOG), the company
behind Cardano. Midnight makes every choice from the pairing-based, pre-quantum
playbook. The consequences cascade through every layer.

\subsection{The Stack}\label{the-stack}

\textbf{Scalar field:} BLS12-381, with a \textasciitilde253-bit prime modulus
\(r\). Every value in Midnight's zero-knowledge circuits -- inputs, outputs,
intermediate computations, token balances -- is an element of \(\mathbb{F}_r\).

\textbf{Commitment scheme:} KZG, implied by the choice of BLS12-381 (a
pairing-friendly curve). The wallet SDK caches BLS parameters locally (a
structured reference string from a trusted setup ceremony). Proof size is
constant. Verification is a single pairing check.

\textbf{Embedded curve:} Jubjub, a twisted Edwards curve whose order divides
\(r\). Jubjub lives ``inside'' BLS12-381's scalar field, enabling efficient
elliptic curve operations (point addition, scalar multiplication, hash-to-curve)
within zero-knowledge circuits without the overhead of non-native field
arithmetic.

\textbf{Hash functions:} Poseidon-family algebraic hashes, represented at the
ZKIR level as opaque opcodes (\texttt{transient\_hash},
\texttt{persistent\_hash}, \texttt{hash\_to\_curve}). Algebraic hash functions
are dramatically more efficient inside ZK circuits than traditional hash
functions like SHA-256, because their operations (field additions and
multiplications) are native to the circuit's arithmetic.

\textbf{Token model:} UTXO-based shielded tokens (similar to Zcash Sapling). A
coin is a triple (nonce, color, value) committed to a global Merkle tree via
\texttt{persistent\_hash}. Nullifiers prevent double-spending. Pedersen
commitments on Jubjub hide transaction values.

\subsection{What Midnight Gets}\label{what-midnight-gets}

Midnight occupies the ``high algebraic functionality + high succinctness''
corner of the design space. The pairing enables constant-size KZG proofs. Jubjub
enables rich in-circuit elliptic curve operations (key derivation, Pedersen
commitments, hash-to-curve) with native efficiency. The PLONK-like proof system
compiles from a purpose-built language (Compact) through a 24-opcode instruction
set (ZKIR). The standard library provides Merkle trees, shielded token circuits,
and a full Zswap protocol for private token transfers.

The result is maximum algebraic functionality. Every cryptographic primitive --
hashing, commitment, key derivation, signature verification -- operates natively
within the circuit's arithmetic. Nothing requires emulation or non-native field
arithmetic.

\subsection{What Midnight Gives Up}\label{what-midnight-gives-up}

Post-quantum resilience: none. Every component of Midnight's cryptographic stack
depends on either the discrete logarithm problem (Jubjub key derivation,
Pedersen commitments) or pairing-based assumptions (KZG proof verification).
Shor's algorithm breaks all of it.

The vulnerability assessment is total:

\begin{longtable}[]{@{}lll@{}}
\toprule\noalign{}
Component & Assumption & Post-Quantum Status \\
\midrule\noalign{}
\arrayrulecolor{MidnightBlue}\midrule\arrayrulecolor{MidnightBorder}
\endhead
\bottomrule\noalign{}
\endlastfoot
Proof verification (KZG) & q-SDH on BLS12-381 & Broken by Shor \\
Jubjub key derivation & ECDLP on Jubjub & Broken by Shor \\
Pedersen commitments & DLP on Jubjub & Broken by Shor (binding fails) \\
In-circuit hashing & CRHF (Poseidon) & Weakened but likely survivable \\
Merkle tree roots & CRHF & Likely survivable \\
Nullifiers & PRF & Likely survivable \\
\end{longtable}

The proof system is the deepest vulnerability. Even if Midnight replaced its
Pedersen commitments with hash-based constructions, and even if it switched from
Jubjub key derivation to a lattice-based signature scheme, the proof
verification mechanism is fundamentally tied to the BLS12-381 pairing. Changing
it would require replacing KZG with a different commitment scheme, which would
require changing the field, which would require rewriting the compiler, the
standard library, the prover, the verifier, and the wallet SDK.

The one-way-door property is in full force. Midnight's choice of BLS12-381 is
not a parameter that can be updated. It is the foundation on which every other
component is built. A post-quantum migration for Midnight would not be an
upgrade. It would be a new system.

\subsection{Midnight vs.~Neo: Opposite
Corners}\label{midnight-vs.-neo-opposite-corners}

The contrast with Neo/Nightstream makes the tradeoffs vivid:

\begin{longtable}[]{@{}lll@{}}
\toprule\noalign{}
Dimension & Midnight & Neo/Nightstream \\
\midrule\noalign{}
\arrayrulecolor{MidnightBlue}\midrule\arrayrulecolor{MidnightBorder}
\endhead
\bottomrule\noalign{}
\endlastfoot
Field & BLS12-381, \textasciitilde253 bits & Goldilocks, 64 bits \\
Ring & N/A (field-based) & \(\mathbb{F}_q[X]/(\Phi_{81})\), degree 54 \\
Commitment & KZG (pairing) & Ajtai (lattice) \\
Hash & Poseidon (algebraic) & Ring-SIS (lattice) \\
Proof size & \(O(1)\) curve points & \(O(\log n)\) ring elements \\
PQ secure & No & Yes (127-bit) \\
EC in-circuit & Yes (Jubjub, native) & No (would need circuit emulation) \\
Trusted setup & Yes (powers-of-tau) & No \\
\end{longtable}

Neo trades Midnight's in-circuit elliptic curve operations and constant-size
proofs for post-quantum security, transparent setup, and a simpler recursive
architecture (no curve cycles needed). Neither system dominates on every
dimension. The question is which tradeoffs matter more for your threat model and
time horizon.

For a system deployed today that needs maximum on-chain efficiency and whose
privacy guarantees are measured in years (not decades), Midnight's choices are
defensible. For a system that needs to protect sensitive data for 15 or more
years, or that must comply with NIST's 2035 deprecation timeline, Neo's choices
are the only viable path.

\section{The Cascade Effect}\label{the-cascade-effect}

The deeper lesson of Layer 6 is that it is not really a ``layer'' at all. It is
a foundation. Every choice made here propagates upward through the entire stack
with the force of mathematical necessity.

Consider Neo's parameter cascade:

\begin{verbatim}
Field: Goldilocks (q = 2^64 - 2^32 + 1)
  --> Ring: Phi_81, degree d = 54
    --> Commitment: kappa = 16 rows, m = 2^24 columns
      --> Folding: b = 2 (base), k = 12 (decomposition depth), B = 4096 (norm bound)
        --> Challenge: T = 216 (expansion factor), |C| ~ 2^125
          --> Security: 127-bit MSIS
            --> Guard: (k+1) * T * (b-1) = 2,808 < 4,096 = B
\end{verbatim}

Change any parameter and everything downstream shifts. Use a different prime and
the cyclotomic factorization changes, which changes the extension field, which
changes the security level, which changes the commitment parameters, which
changes the folding parameters. This is not optional coupling -- it is algebraic
necessity. The parameters are not chosen independently. They are derived from
each other, each one a consequence of the ones above it, the way the shape of a
crystal is a consequence of the geometry of its atoms.

The same cascade operates in the pairing world. BLS12-381's prime determines the
Jubjub embedding. Jubjub determines which in-circuit operations are efficient.
The pairing determines which commitment scheme works. The commitment scheme
determines the proof system. The proof system determines the arithmetization.

And so ``crypto-agility'' -- the ability to swap cryptographic primitives
without redesigning the system -- is largely a fiction for zero-knowledge
systems. You cannot change the field without changing everything. The choice at
Layer 6 is a one-way door, and once you walk through it, you are committed.

To make this concrete, here is the decision tree that every zero-knowledge
system architect walks, whether they realize it or not.

\textbf{If you choose BabyBear (31-bit prime, \(p = 2^{31} - 2^{27} + 1\)):} You
get SIMD-friendly arithmetic -- four field multiplications packed into a single
128-bit SIMD instruction, eight into a 256-bit AVX2 register. Your natural
commitment scheme is FRI, because BabyBear has a multiplicative subgroup of
order \(2^{27}\), large enough for practical NTT domains. Your constraint format
is AIR (Algebraic Intermediate Representation) or CCS, depending on your proof
system. Your setup is transparent -- no ceremony, no trust. Your proofs are
large (50 to 200 kilobytes) but your prover is \emph{fast}, because every
arithmetic operation is a single machine instruction. You achieve 128-bit
security via a degree-4 extension field (four BabyBear elements per extended
element, giving \textasciitilde124 bits). This is the path chosen by RISC Zero
and Plonky3. It optimizes for prover throughput at the cost of proof size.

\textbf{If you choose Goldilocks (64-bit prime, \(p = 2^{64} - 2^{32} + 1\)):}
You get native 64-bit arithmetic -- one multiplication per CPU instruction, no
multi-limb overhead. Your natural commitment scheme is FRI (exploiting 2-adicity
of \(2^{32}\) for large NTT domains) or lattice-based Ajtai commitments (using
the 81st cyclotomic polynomial for post-quantum security). Your constraint
format is CCS or R1CS. Your setup is transparent in either case. If you choose
FRI, your proofs are large but your prover leverages GPU-friendly 64-bit
arithmetic. If you choose Ajtai, you get post-quantum security and folding
capability, with proofs in the 50 to 60 kilobyte range. This is the path chosen
by Plonky2 (FRI) and Neo/Nightstream (Ajtai). It balances prover speed, proof
size, and -- if lattice-based -- quantum resilience.

\textbf{If you choose BLS12-381 (254-bit pairing-friendly curve):} You get the
full power of bilinear pairings -- KZG commitments with constant-size proofs (48
bytes), constant-time verification (one pairing check), and the richest
algebraic structure available. Your constraint format is PLONKish gates or R1CS.
Your setup requires a trusted ceremony (powers-of-tau). Your proofs are the
smallest in existence. But your arithmetic is the most expensive: a single
254-bit multiplication costs a multi-limb subroutine that is 100 times slower
than BabyBear's native operation. And you inherit an expiration date: Shor's
algorithm will break every pairing-based proof when a cryptographically relevant
quantum computer arrives. This is the path chosen by Midnight, Zcash
(pre-Orchard), every Ethereum rollup's final verification layer, and the
EIP-4844 blob scheme. It optimizes for proof succinctness and verifier
efficiency at the cost of prover performance and quantum resilience.

\textbf{If you choose Mersenne-31 (\(p = 2^{31} - 1\)):} You get the simplest
possible modular reduction -- subtraction of the carry bit, because
\(2^{31} \equiv 1 \pmod{p}\). Your commitment scheme is FRI, adapted via Circle
STARKs to work with M31's multiplicative group structure (which lacks large
2-adic subgroups but has a circle group of order \(2^{31}\)). Your prover is the
fastest in existence for STARK-based systems, because M31 arithmetic is cheaper
than any other field. Your proofs are transparent and plausibly post-quantum.
This is StarkWare's Stwo path -- maximum prover throughput, hash-based security,
no algebraic frills.

Each path is internally consistent. Each forecloses the others. You cannot start
down the BabyBear path and switch to KZG midstream -- the field does not support
pairings. You cannot start with BLS12-381 and add post-quantum security -- the
algebraic structure that gives you constant-size proofs is the same structure
that Shor's algorithm destroys. The decision tree is not a menu. It is a set of
branching tunnels, and once you enter one, the others seal behind you.

\section{Algebraic vs.~Traditional Hash
Functions}\label{algebraic-vs.-traditional-hash-functions}

One more Layer 6 choice deserves attention, because it illustrates how deeply
the primitive selection affects practical performance.

\textbf{Traditional hash functions} (SHA-256, BLAKE3, Keccak) are designed for
speed on general-purpose hardware. Their internal operations -- bitwise
rotations, XOR, addition with carry -- are cheap on CPUs but extremely expensive
inside zero-knowledge circuits, because the circuit's native operations are
field additions and multiplications. Proving a single SHA-256 computation inside
a SNARK requires tens of thousands of constraints.

\textbf{Algebraic hash functions} (Poseidon, Poseidon2, Rescue, Griffin) are
designed for the opposite environment. Their internal operations are field
multiplications and exponentiations -- exactly the operations that are native to
zero-knowledge circuits. A Poseidon hash inside a circuit costs hundreds of
constraints instead of tens of thousands.

The performance difference is 100x or more. This is why every system that does
significant hashing inside circuits (Merkle tree verification, Fiat-Shamir
challenges, commitment randomness) either uses algebraic hashes or pays an
enormous performance penalty.

But algebraic hashes are newer and less studied than SHA-256 or BLAKE3. Their
security rests on assumptions about the difficulty of algebraic attacks (Grobner
basis computations, interpolation attacks) that have not endured decades of
cryptanalysis. Poseidon, in particular, has seen several parameter revisions in
response to improved attacks.

The choice between algebraic and traditional hash functions is itself a Layer 6
decision that cascades upward. Midnight uses Poseidon-family hashes (maximizing
in-circuit efficiency at the cost of less mature security analysis). STARK-based
systems can use either, but algebraic hashes dramatically reduce the size of the
verification circuit when recursion or wrapping is needed.

\section{The Structural Advantage of
Lattices}\label{the-structural-advantage-of-lattices}

One insight from the lattice revolution deserves emphasis because it is easy to
miss amid the parameter details: \textbf{lattice-based schemes are
architecturally simpler} than their pairing-based predecessors, not just
quantum-resistant.

Pre-quantum recursive proof systems (exemplified by Zexe) require:

\begin{itemize}
\item
  \textbf{Cycles of elliptic curves.} To verify a proof inside another proof,
  the verifier's field arithmetic must be efficient in the prover's circuit.
  This requires two curves whose scalar fields are each other's base fields -- a
  ``cycle.'' Finding such cycles constrains parameter choices severely, and
  arithmetic on the second curve is typically 2x or more expensive.
\item
  \textbf{Non-native field arithmetic.} When the proof system operates over one
  field but the verifier checks computations in another, every operation must be
  emulated using multi-precision arithmetic inside the circuit. This is a major
  source of overhead.
\item
  \textbf{Multiple structured reference strings.} Each curve in the cycle needs
  its own trusted setup, doubling the ceremony burden.
\end{itemize}

Lattice-based folding eliminates all three requirements. Neo operates over a
single ring \(R_q\). The rotation matrix encoding makes everything native to one
algebraic structure. Recursion via folding requires no curve cycles and no
non-native arithmetic. The setup is transparent (public random matrix).

This simplification is not cosmetic. Fewer moving parts mean fewer places for
bugs, fewer parameters to choose and validate, fewer assumptions to audit. The
lattice path is not only quantum-resistant -- it is \emph{simpler}. And in
cryptographic engineering, simplicity is not a luxury. It is a security
property.

\section{Maturity and Readiness}\label{maturity-and-readiness}

As of early 2026, the picture looks like this:

\textbf{Deployed and battle-tested:} KZG (BN254 and BLS12-381), FRI/STARK
(Goldilocks, BabyBear, M31), IPA/Bulletproofs (Pasta curves). These power every
production ZK system -- Ethereum rollups, Zcash, Midnight, Starknet.

\textbf{Peer-reviewed and prototyped:} LatticeFold (ASIACRYPT 2025, presentation
by Boneh and Chen), LatticeFold+ (CRYPTO 2025). Neo has an active implementation
in Rust (the Nightstream repository, 15 crates). Concrete benchmarks are
emerging but sparse.

\textbf{Proposed and promising:} Symphony (ePrint 2025/1905, no implementation
yet). The high-arity folding concept is validated theoretically but awaits
engineering.

\textbf{Standards in place:} NIST FIPS 203/204/205 (August 2024) standardize
lattice-based key encapsulation and signatures. No standard yet exists for
lattice-based zero-knowledge proof systems, but the parameter selection
methodology (the lattice estimator) is well established.

The adoption trajectory suggests lattice-based proof systems will move from
research prototypes to production-ready systems in 2026-2027, with
Neo/Nightstream among the first to target production deployment.

\section{The One-Way Door}\label{the-one-way-door}

Layer 6 is unlike any other layer in the stack. Layers above it can be upgraded,
swapped, and optimized. A rollup can change its sequencer, rewrite its compiler,
switch arithmetization formats, even adopt a new proof system -- all without
changing the cryptographic foundations. But the foundations themselves are
effectively permanent.

For architects making this one-way decision, the following rubric distills the
tradeoffs:

\begin{longtable}[]{@{}
  >{\raggedright\arraybackslash}p{(\linewidth - 8\tabcolsep) * \real{0.1273}}
  >{\raggedright\arraybackslash}p{(\linewidth - 8\tabcolsep) * \real{0.1091}}
  >{\raggedright\arraybackslash}p{(\linewidth - 8\tabcolsep) * \real{0.2000}}
  >{\raggedright\arraybackslash}p{(\linewidth - 8\tabcolsep) * \real{0.3455}}
  >{\raggedright\arraybackslash}p{(\linewidth - 8\tabcolsep) * \real{0.2182}}@{}}
\toprule\noalign{}
\begin{minipage}[b]{\linewidth}\raggedright
Field
\end{minipage} & \begin{minipage}[b]{\linewidth}\raggedright
Size
\end{minipage} & \begin{minipage}[b]{\linewidth}\raggedright
PQ Status
\end{minipage} & \begin{minipage}[b]{\linewidth}\raggedright
Commitment Options
\end{minipage} & \begin{minipage}[b]{\linewidth}\raggedright
Sweet Spot
\end{minipage} \\
\midrule\noalign{}
\arrayrulecolor{MidnightBlue}\midrule\arrayrulecolor{MidnightBorder}
\endhead
\bottomrule\noalign{}
\endlastfoot
BN254 & 254-bit & Quantum-vulnerable & KZG (cheapest EVM verification) & Legacy
Ethereum rollups; Groth16 wrapper \\
BLS12-381 & 253-bit & Quantum-vulnerable & KZG (higher security margin) &
Privacy systems needing pairings (Midnight, Zcash) \\
BabyBear & 31-bit & Hash-PQ; lattice-PQ & FRI, Ajtai & Maximum prover speed;
RISC-V zkVMs (SP1, RISC Zero) \\
Mersenne-31 & 31-bit & Hash-PQ & FRI (Circle STARK) & Fastest arithmetic;
Stwo/Starknet ecosystem \\
Goldilocks & 64-bit & Hash-PQ; lattice-PQ & FRI, Ajtai & Balance of speed and
precision; Neo/Nightstream \\
\end{longtable}

The finite field determines the commitment scheme. The commitment scheme
determines the proof system family. The hardness assumption determines the
security lifespan. These choices are made once, at the beginning, and they
propagate upward through every component with the inexorability of mathematical
structure.

This is why the quantum threat is not a problem that can be deferred until
quantum computers actually exist. A system deployed in 2026 with BN254
foundations will still be running in 2036. If a CRQC arrives in 2035, that
system will have spent its final years accumulating a public record of
commitments and proofs that can be retroactively broken. The HNDL threat means
the privacy guarantees were never real -- they were deferred revelations,
secrets written in ink that merely required a light that had not yet been
invented.

The lattice revolution is a construction project, not an academic exercise. It
is building new foundations that can support the same architectural weight as
pairing-based cryptography -- folding, recursion, efficient composition --
without the quantum expiration date. The trilemma that seemed permanent is being
dissolved not by discovering new mathematics but by engineering better
constructions from mathematics that has existed for decades.

The laws of physics do not change. But our understanding of which mathematical
problems are hard does change, and a quantum computer represents a discontinuous
shift in that understanding. The systems that survive will be the ones whose
foundations were chosen with that shift in mind.

\par\vspace{18pt}

The physical laws are set. The field is chosen, the commitment scheme
determined, the hardness assumption staked. Everything from here upward -- the
proof system, the arithmetization, the language, the setup -- inherits the
possibilities and constraints of this foundation. But none of it matters until
someone checks the proof. Layer 7 is where the mathematics meets its audience:
on a blockchain, in a smart contract, through a governance structure that can
override everything we have built. The next chapter examines the verdict -- and
the uncomfortable truth that the audience's judgment is only as trustworthy as
the institution that seats them.

\chapter{Chapter 8: Layer 7 -- The Verdict}\label{chapter-8-layer-7-the-verdict}

\emph{``The audience can be deceived -- or worse, someone can replace the
audience entirely.''}

\section{The Social Layer}\label{the-social-layer}

Every layer of the zero-knowledge stack we have examined so far -- the ceremony,
the language, the witness, the arithmetization, the proof system, the
cryptographic primitives -- converges on a single moment: a piece of software
reads a proof and says \emph{yes} or \emph{no}.

That piece of software is the audience. We established this in Chapter 1: the
verifier is the audience, the entity that watches the trick and renders its
verdict. On Ethereum, the audience is a smart contract. On Midnight, it is a
node. On a private enterprise chain, it might be a service running in a data
center. Whatever form it takes, the verifier is the point where all the private
magic becomes a public verdict.

And here is the uncomfortable truth that most explanations of zero-knowledge
proofs prefer to gloss over: the audience can be replaced. Not by breaking the
cryptography. Not by forging a proof. By something much simpler.

By changing the software.

If three people on a governance multisig can upgrade the verifier contract to
one that accepts every proof -- or no proofs, or only their proofs -- then the
128-bit security of the proof system, the million-dollar ceremony, the carefully
audited circuits, all of it becomes decorative. The math does not protect you
from the admin key.

This chapter is about what happens after the proof is generated. It is about gas
costs, data availability, implementation bugs, governance attacks, and the
social structures that determine whether the cryptographic guarantees from
Layers 1 through 6 actually reach the people they are supposed to protect.

Layer 7 is where cryptography meets politics. And politics, as a rule, wins.

Layer 7 carries four distinct responsibilities, and this chapter treats each in
turn. First: the \emph{economics} of rendering a verdict --- what does
verification cost, and who pays? Second: \emph{implementation vulnerabilities}
that can corrupt the verdict --- Fiat-Shamir transcript bugs that enable proof
forgery. Third: \emph{governance structures} that can override the verdict ---
multisig attacks, upgrade mechanisms, and the social layer above the math.
Fourth: \emph{aggregation and data availability infrastructure} that sits
between the prover and the verifier --- SHARP, blob economics, and the emerging
DA marketplace. These four concerns are operationally convergent --- they all
determine whether the audience's verdict is trustworthy --- but they are
logically distinct. A system can have perfect verification economics and
catastrophic governance. Separating the concerns makes the trust analysis
sharper.

\section{The Price of a Verdict}\label{the-price-of-a-verdict}

Let us start with money, because money clarifies.

A Groth16 proof verification on Ethereum uses the BN254 elliptic curve pairing
precompiles introduced in the Byzantium hard fork (2017) and made cheaper by the
Istanbul upgrade's EIP-1108 (2019). The gas cost breaks down as follows:

\begin{longtable}[]{@{}ll@{}}
\toprule\noalign{}
Component & Gas Cost \\
\midrule\noalign{}
\arrayrulecolor{MidnightBlue}\midrule\arrayrulecolor{MidnightBorder}
\endhead
\bottomrule\noalign{}
\endlastfoot
Pairing check (4 pairings via EIP-1108) & 181,000 \\
Calldata (256-byte proof) & 4,096 \\
EVM scaffolding & \textasciitilde1,600 \\
Per public input & \textasciitilde7,160 each \\
\textbf{Total (fixed, no public inputs)} & \textbf{\textasciitilde207,700} \\
\end{longtable}

The formula is roughly \((181 + 6L) \times 1{,}000\) gas for \(L\) public
inputs. At typical Ethereum gas prices and ETH valuations, this works out to
somewhere between fifty cents and two dollars per verification. Call it a
dollar.

One dollar. To check a proof that summarizes thousands, or millions, of
computations. That is the economic engine of the entire zero-knowledge rollup
industry. It is also worth noting: the cost of \emph{rendering a verdict} on an
arbitrarily complex computation is effectively fixed. The computation can be ten
steps or ten billion. The verdict costs the same.

But notice what that dollar buys. It buys a \emph{Groth16} verification. Groth16
requires a trusted setup (Layer 1), uses elliptic curve pairings on BN254 (Layer
6), and produces the smallest proofs in the field -- three group elements that
fit in a tweet. The cheapness of the verdict is not free. It is subsidized by
decisions made five layers below.

What about STARKs? The paper being revised presents STARK verification as
expensive -- two to five million gas -- and contrasts this with SNARK cheapness.
This framing was arguably accurate in 2022. It is misleading in 2026. The reason
is simple: nobody posts raw STARKs to Ethereum.

The actual production pipeline looks like this:

\begin{enumerate}
\def\labelenumi{\arabic{enumi}.}
\tightlist
\item
  Generate a STARK proof (transparent, no trusted setup, large -- hundreds of
  kilobytes).
\item
  Recursively compress the STARK through multiple rounds.
\item
  Wrap the final compressed STARK inside a Groth16 proof.
\item
  Post the Groth16 proof to Ethereum.
\end{enumerate}

Starknet's SHARP (Shared Prover) does this. Succinct's SP1 does this. Polygon's
CDK does this. The verifier contract on Ethereum sees Groth16 in every case. The
``STARK path'' and the ``SNARK path'' converge at the courthouse door.

The actual cost differential between the two approaches, then, is not the 10-25x
that a naive comparison of raw STARK versus Groth16 verification would suggest.
It is closer to 2x -- the overhead of the wrapping step, amortized across many
proofs. The inner proof system matters enormously for prover economics (speed,
hardware requirements, parallelizability), but the on-chain verification cost is
nearly identical.

There is a throughput ceiling too. An Ethereum block has a 30-million-gas limit
(raised from the historical 15 million, with a 45-million effective target under
various proposals). At 207,700 gas per Groth16 verification, you can fit roughly
150 to 225 verifications per block. That sounds like a lot, until you realize
that each verification corresponds to a batch of rollup transactions. If
Ethereum hosts 50 rollups and each wants to verify once per block, they consume
less than a quarter of the block's capacity. But if we want real-time proving
(verification every L1 slot), with hundreds of rollups and bridges, the
verification gas budget starts to matter.

FFLONK, an alternative to Groth16, costs roughly 236,000 gas per verification --
slightly more, but with the advantage of a universal trusted setup (one ceremony
works for all circuits, unlike Groth16's per-circuit setup). The gas difference
is marginal. The governance and operational difference -- not needing a new
ceremony for each circuit -- is substantial.

\subsection{The Verification-Data Seesaw}\label{the-verification-data-seesaw}

Before March 2024, the dominant cost of running a ZK rollup on Ethereum was not
verification. It was data availability. Posting transaction data (or state
diffs) as calldata cost roughly 16 gas per byte. A typical rollup batch might
include hundreds of kilobytes of data, costing millions of gas -- dwarfing the
\textasciitilde200,000 gas for the proof check.

EIP-4844, deployed in the Dencun upgrade on March 13, 2024, changed this
calculus fundamentally. It introduced ``blob transactions'' -- a new data type
designed specifically for rollup data. Each blob contains 4,096 field elements
of 32 bytes (\textasciitilde128 KB), with a target of 3 blobs per block and a
maximum of 6. Critically, blobs have their own fee market, separate from
Ethereum's execution gas market, operating under a blob-specific EIP-1559
mechanism.

The result: rollup data costs dropped by 10-100x overnight. Blob fees settled
near zero because demand was well below the 3-blob target -- as of mid-2024,
only about 34\% of Ethereum blocks contained any blobs at all, and the average
was 1.33 blob transactions per block.

But Ethereum did not stop at EIP-4844. Two subsequent upgrades expanded DA
capacity further:

\begin{itemize}
\tightlist
\item
  \textbf{Pectra} (May 2025): Doubled blob targets from 3 to 6, and maximum from
  6 to approximately 9.
\item
  \textbf{Fusaka} (December 2025): Introduced PeerDAS (Peer Data Availability
  Sampling), implementing a distributed sampling scheme that raised the blob
  target to 14 and maximum to 21 -- an 8x increase in DA capacity over the
  original EIP-4844 specification.
\end{itemize}

The seesaw has tipped. With blob fees near zero and DA capacity expanding
rapidly, the \textasciitilde200,000 gas verification cost has become the
\emph{dominant} L1 settlement expense for many ZK rollups. This inversion
matters because it changes what is worth optimizing. Before EIP-4844, the
rational investment was in compression (minimizing data). After EIP-4844, the
rational investment is in proof aggregation (amortizing verification across more
transactions per batch) and in cheaper verification schemes.

\subsection{Beyond Ethereum: The DA
Marketplace}\label{beyond-ethereum-the-da-marketplace}

Ethereum is not the only source of data availability. A marketplace has emerged:

\textbf{Celestia} charges roughly \$0.07 per megabyte for data availability,
compared to Ethereum's blob cost of roughly \$3.83 per megabyte (when blobs are
priced above the floor). Celestia achieves this by being a purpose-built DA
layer -- it provides data ordering and availability guarantees without executing
any transactions. The intellectual lineage traces directly to Mustafa
Al-Bassam's LazyLedger (2019), which proposed a blockchain that does nothing but
guarantee data is available and ordered, leaving execution to sovereign rollups
that interpret their own transaction rules.

\textbf{EigenDA V2} targets 100 megabytes per second of throughput -- roughly
two orders of magnitude more than Ethereum's native DA capacity. It achieves
this by leveraging Ethereum's security through restaking (EigenLayer), where
validators stake ETH to back DA guarantees.

\textbf{Avail} offers a third alternative, with its own DAS-based light client
verification model.

The choice between these DA layers is not purely technical. A rollup that uses
Celestia for DA instead of Ethereum blobs trades Ethereum's full consensus
security for lower costs. This is a Layer 7 governance decision with Layer 6
security implications: the data availability guarantee is only as strong as the
weakest link in the DA provider's consensus mechanism.

\subsection{Data Availability}\label{data-availability}

The term ``data availability'' is one of those phrases that sounds
self-explanatory and is not. It does not mean ``the data exists somewhere.'' It
does not mean ``the data is stored on a server.'' It means something specific
and testable: if you send a transaction to a rollup, can any participant in the
world reconstruct the rollup's complete state using only publicly available
data?

If yes, the rollup has data availability. Anyone can verify that the rollup
operator is honest by replaying all transactions from genesis and checking that
the claimed state matches the computed state. If no -- if some of the data is
withheld, stored only on the operator's private servers, or available only to a
privileged set of participants -- then the operator could cheat and nobody would
know. The operator could include a transaction that steals every user's funds,
prove that the resulting state transition is ``valid'' (because the ZK proof
only proves that \emph{some} valid transition occurred), and nobody could
challenge it because nobody can see the inputs.

This is the critical subtlety that connects data availability to zero-knowledge
proofs. A ZK proof proves that a state transition was computed correctly. It
proves that if you start from state S and apply transactions T, you arrive at
state S'. What it does \emph{not} prove -- what it \emph{cannot} prove, by
design -- is what state S actually was. The proof attests to the correctness of
the computation, not the availability of the inputs. If the operator claims the
starting state was S but actually started from a fabricated state S\_fake, the
ZK proof will happily prove that the transition from S\_fake was computed
correctly. Without DA, nobody can verify the starting point.

Data availability is the anchor. The ZK proof is the chain. Without the anchor,
the chain secures nothing.

The three DA strategies represent different points on the cost-security
tradeoff:

\textbf{Ethereum calldata} is the oldest and most expensive approach.
Transaction data is posted directly as calldata in Ethereum L1 transactions,
stored permanently by every full node, and protected by Ethereum's full
consensus security. The cost is high -- 16 gas per byte of calldata -- but the
guarantee is absolute: if Ethereum's consensus is secure, the data is available.
This was the only option before March 2024, and it made rollup operations
expensive enough that most of the early rollup economics were dominated by DA
costs rather than verification costs.

\textbf{Ethereum blobs} (EIP-4844 and successors) are the middle ground. Blob
data is posted to Ethereum and protected by Ethereum's consensus during a
pruning window (currently approximately 18 days), after which nodes may discard
it. The data is available long enough for any challenge period to complete, and
it is significantly cheaper than calldata because blobs have their own fee
market and do not compete with execution gas. This is the default choice for
most production rollups in 2026.

\textbf{External DA layers} (Celestia, EigenDA, Avail) are the cheapest option
with a different trust model. The data is posted to a separate blockchain or
protocol that specializes in data ordering and availability guarantees. The cost
can be 10-100x lower than Ethereum blobs. The tradeoff is that the DA guarantee
depends on the external protocol's consensus and validator set, not Ethereum's.
A rollup using Celestia for DA inherits Celestia's security assumptions. If
Celestia's validator set colludes or fails, the rollup's data may become
unavailable even though Ethereum itself is functioning correctly.

The choice of DA strategy is, in practice, one of the most consequential
governance decisions a rollup team makes. It determines the rollup's operating
cost, its security model, its relationship to Ethereum's consensus, and its
vulnerability to the DA-saturation attacks discussed later in this chapter. It
is a Layer 7 decision with implications that cascade through every layer below.

\section{When the Transcript Lies: Fiat-Shamir
Vulnerabilities}\label{when-the-transcript-lies-fiat-shamir-vulnerabilities}

Chapter 6 introduced the Fiat-Shamir transform as the mechanism that seals Layer
5 proofs into non-interactive certificates, and flagged the binding requirement:
the hash must include every public value the verifier would have seen. This
section examines what happens when that requirement is violated in production --
and why the consequences are uniquely catastrophic at Layer 7, where on-chain
verifiers are permissionless and exploitation is automated.

The Fiat-Shamir heuristic is the mechanism that converts an interactive proof
(where the verifier asks random questions in real time) into a non-interactive
one (where the prover simulates the verifier's questions using a hash function).
Every non-interactive ZK proof deployed on a blockchain uses Fiat-Shamir. It is
the invisible thread that holds the entire verification model together.

And it is the single most dangerous implementation surface in the entire stack.

\subsection{Frozen Heart (2022)}\label{frozen-heart-2022}

In April 2022, Trail of Bits disclosed a vulnerability class they called
``Frozen Heart'' (a backronym: Forging Of Zero kNowledge proofs). The core error
was simple. Devastatingly simple. Multiple independent implementations of ZK
proof systems omitted public inputs from the Fiat-Shamir hash computation.

The implementations affected were not obscure academic prototypes. They were
production-grade libraries used by real projects:

\begin{itemize}
\tightlist
\item
  \textbf{Dusk Network} (PLONK implementation)
\item
  \textbf{Iden3/SnarkJS} (Groth16, used by Circom)
\item
  \textbf{ConsenSys/gnark} (PLONK implementation)
\item
  \textbf{ING Bank's zkrp} (Bulletproofs)
\item
  \textbf{SECBIT Labs' ckb-zkp} (Groth16)
\item
  \textbf{Adjoint Inc.'s bulletproofs} (Bulletproofs)
\end{itemize}

Six implementations. Three different proof systems. Four different
organizations. All made the same mistake: they left the public inputs out of the
hash that generates the verifier's challenges.

The consequence is total soundness failure. A malicious prover can forge proofs
for \emph{arbitrary false statements}. Not with some small probability. With
certainty. The ``proof'' passes verification because the challenges are no
longer bound to the specific statement being proved. The sealed certificate
attests to nothing. It is wax without an impression.

The rule that was violated is not subtle: the Fiat-Shamir hash must include
\emph{all} public values from the ZK statement and \emph{all} public values
computed during the proof. Every commitment, every public input, every piece of
data that the verifier would have seen in the interactive version must go into
the hash. Omit any of it, and the binding between challenge and statement
dissolves.

\subsection{The Last Challenge Attack
(2024)}\label{the-last-challenge-attack-2024}

The Last Challenge Attack, discovered during an audit of Linea's PLONK verifier
in the gnark library, is a more surgical variant of the same disease. In
KZG-based proof systems, the verifier often batches multiple polynomial
evaluations using a random ``batching challenge'' derived via Fiat-Shamir. The
Last Challenge Attack exploits the case where this batching challenge is
computed from a \emph{truncated} transcript -- one that excludes the evaluation
proofs themselves.

The attack is elegant in the way that a perfectly executed heist is elegant. The
malicious prover:

\begin{enumerate}
\def\labelenumi{\arabic{enumi}.}
\tightlist
\item
  Sets arbitrary (false) public inputs and proof components.
\item
  Computes the batching challenge from the truncated transcript.
\item
  Solves a linear system for the missing evaluation proofs.
\item
  The vulnerable verifier accepts the forged proof with probability 1.
\end{enumerate}

Not ``with high probability.'' With certainty. The forged proof is
deterministically constructed to pass verification. The audience has been
compromised not by force but by omission -- a single value left out of a hash,
and the entire edifice of mathematical certainty collapses.

The gnark advisory (GHSA-7p92-x423-vwj6) confirmed the vulnerability. The fix
was straightforward: compute the batching challenge only \emph{after} all
evaluation proofs are included in the transcript. But the vulnerability existed
in a production-quality library used by multiple rollup teams.

\subsection{Solana ZK ElGamal (2025)}\label{solana-zk-elgamal-2025}

The pattern repeated in 2025 on Solana, where the ZK ElGamal implementation was
found to have a Fiat-Shamir transcript that omitted the prover's challenge from
the hash computation. Same class of error. Same catastrophic consequence. The
same lesson, unlearned for the third time.

\subsection{The Pattern}\label{the-pattern}

These are not isolated incidents. They are symptoms of a structural problem: the
Fiat-Shamir heuristic is easy to describe (``hash everything the verifier would
see'') and remarkably easy to get wrong in implementation. The specification
says ``include all public values.'' The implementation omits one, because it
seemed redundant, or because it made the code cleaner, or because the developer
did not understand \emph{why} it needed to be there.

For on-chain verifiers, this class of vulnerability is uniquely dangerous. An
on-chain verifier is permissionless -- anyone can submit a proof. If the
verifier's Fiat-Shamir transcript is incomplete, exploitation is automated and
instantaneous. There is no human in the loop to notice that something looks
wrong. The forged proof passes the smart contract's checks, the state transition
is accepted, and the attacker drains whatever value the rollup is protecting.

Every on-chain SNARK verifier -- Groth16, PLONK, FFLONK, any KZG-based scheme --
must be audited specifically and primarily for Fiat-Shamir transcript
completeness. This should be the first check in any security review, not an item
buried in a general audit report.

\section{Governance: The Achilles Heel}\label{governance-the-achilles-heel}

If Fiat-Shamir bugs are the most exploited \emph{implementation} vulnerability
in zero-knowledge systems, governance is the most exploited \emph{architectural}
vulnerability. And unlike Fiat-Shamir bugs, governance vulnerabilities cannot be
fixed with better code review. They are features, not bugs.

Here the story changes genre. Until now, we have been watching a technical
narrative -- mathematicians and engineers building increasingly sophisticated
proof systems. Now the camera pulls back, and the audience discovers it has been
watching a different show than it thought. The threats at Layer 7 are not
mathematical. They are human.

\subsection{The Beanstalk Flash Loan Attack (\$182M, April
2022)}\label{the-beanstalk-flash-loan-attack-182m-april-2022}

Beanstalk was a permissionless stablecoin protocol -- a DeFi project built on
the idea that an algorithmic stablecoin (called Bean) could maintain its dollar
peg through a credit-based system of debt, deposits, and incentive cycles. The
protocol had attracted over \$100 million in total value locked. It had a
community. It had audits. It had a governance mechanism that allowed holders of
the protocol's internal Stalk token to propose and vote on changes to the
system's parameters, its contracts, its entire economic logic.

The mathematics were sound. The smart contracts were audited. The governance
mechanism was functioning exactly as designed.

That last sentence is the important one. Remember it.

Beanstalk's governance had one feature that seemed reasonable at the time: the
emergency commit threshold. If a proposal attracted more than two-thirds of
total Stalk voting power, it could be executed immediately -- no waiting period,
no time lock, no multi-day deliberation window. The designers' reasoning was
practical: if a supermajority of stakeholders agreed on something, why force
them to wait? Speed was a feature. In a fast-moving DeFi market, the ability to
respond quickly to exploits or market conditions was considered an advantage.

On April 17, 2022, at approximately 12:24 UTC, someone demonstrated why speed is
also a weapon.

The attacker -- whose identity remains unknown to this day -- began by taking
flash loans from three decentralized lending protocols: Aave, Uniswap V2, and
SushiSwap. A flash loan is a peculiar instrument unique to programmable
blockchains: it allows you to borrow any amount of money, provided you repay it
within the same transaction. If you cannot repay, the entire transaction reverts
as if it never happened. The borrowing cost is essentially zero -- just gas fees
and a small protocol fee. There is no credit check. There is no collateral.
There is no application form. You simply ask for the money, use it, and return
it, all within a single atomic operation that takes seconds.

On this day, the attacker borrowed approximately \$1 billion in assets. One
billion dollars, for thirteen seconds.

With the borrowed capital, the attacker swapped into Beanstalk's liquidity
pools, acquiring enough of the protocol's Stalk and Seed tokens to control over
67\% of total governance voting power. This was not a theoretical majority. It
was an absolute supermajority -- enough to clear the emergency commit threshold.

Then came the proposals. BIP-18 was the payload: a governance proposal whose
code, when executed, would transfer all of Beanstalk's protocol reserves --
every Bean, every LP token, every asset in the Silo -- to a wallet controlled by
the attacker. The code was not hidden. It was right there on the blockchain,
readable by anyone who looked. But governance proposals are submitted and voted
upon, not scrutinized line by line in the seconds between submission and
execution, and nobody was watching for a proposal backed by a billion dollars of
borrowed voting power.

BIP-19 was the other proposal: a donation of \$250,000 to the Ukraine war relief
wallet. Whether this was misdirection, moral compensation, ironic commentary, or
simply a way to make the governance transaction look routine is a question the
attacker left permanently unanswered. It remains one of the small, unsettling
details that elevate this from a theft to a performance.

The attacker voted on BIP-18 with the borrowed supermajority. The emergency
commit threshold was cleared. The governance system did what governance systems
do: it executed the will of the majority. The protocol's reserves flowed from
the Silo to the attacker's address. The attacker unwound the liquidity
positions, converted the assets, repaid the flash loans to Aave, Uniswap, and
SushiSwap -- in full, with fees -- and pocketed the difference.

Thirteen seconds. Borrow, vote, execute, extract, repay. The entire heist was a
single atomic transaction on Ethereum. If any step had failed -- if the flash
loan had been too small, if the voting power had been insufficient, if the
repayment had come up short -- every step would have reverted, and the
blockchain would have recorded nothing. But no step failed. The transaction
succeeded. The money was gone.

Total protocol loss: \$182 million in value destroyed. Net extraction by the
attacker: approximately \$77 million in non-Bean assets -- the portion that had
real market value independent of the now-collapsed protocol. The Bean stablecoin
depegged immediately and never recovered. The protocol's entire treasury was
emptied in a single block.

The root cause was not a code vulnerability. No contract was exploited in the
traditional sense. No buffer was overflowed. No access control was bypassed. No
reentrancy was triggered. The governance mechanism worked exactly as designed.
It accepted a vote from a stakeholder with a supermajority. It executed the
proposal that the supermajority approved. It transferred the funds that the
proposal specified. Every line of code behaved correctly.

The system was not broken. The system was \emph{used}.

The lesson lands differently depending on who you are. If you are a protocol
designer, the lesson is about time locks and minimum voting periods and the
danger of emergency execution without delay. If you are a governance theorist,
the lesson is about the difference between ownership and rental -- the attacker
did not own the voting power; he rented it for the cost of a flash loan fee. If
you are building a ZK rollup with token-weighted governance over an upgradeable
verifier contract, the lesson is existential: governance that can be rented by
the hour is governance that can be captured in seconds. The flash loan is the
instrument. The vulnerability is the assumption -- the assumption that token
holders are stakeholders, that voting power reflects long-term commitment, that
the people who hold the keys today will hold them tomorrow. Flash loans dissolve
that assumption into nothing. For the thirteen seconds that matter, anyone with
gas money is a supermajority stakeholder.

\subsection{The Tornado Cash Governance Attack (May
2023)}\label{the-tornado-cash-governance-attack-may-2023}

Tornado Cash was a privacy protocol built on zero-knowledge proofs -- the very
technology this book describes. It used ZK proofs to break the on-chain link
between depositors and withdrawers. You deposit ETH into a pool, receive a
cryptographic note, and later withdraw from the pool using a ZK proof that
demonstrates you possess a valid note without revealing which deposit was yours.
The cryptography was elegant, well-audited, and provably sound. The protocol's
privacy guarantees were genuine. Its governance was controlled by a DAO with
TORN token voting, and the governance was not.

The attack, when it came in May 2023, unfolded like a stage magic trick -- not
the kind where a rabbit appears from a hat, but the kind where the audience
watches the magician's right hand while the left hand replaces the entire stage.

To understand the trick, you need to understand two pieces of Ethereum
infrastructure that most users never think about.

The first is \texttt{CREATE2}. On Ethereum, when you deploy a contract, it gets
an address. Normally, this address is derived from the deployer's address and a
nonce (a sequential counter), so it is effectively unpredictable.
\texttt{CREATE2}, introduced in EIP-1014, changes the formula: the new
contract's address is derived from the deployer's address, a chosen salt, and
the \emph{hash of the bytecode being deployed}. This means you can calculate a
contract's address before deploying it. More importantly -- and this is the key
to the trick -- if you deploy an intermediary factory contract that itself uses
\texttt{CREATE} (the old opcode), and that factory deploys a child contract, and
then you destroy both the factory and the child via \texttt{selfdestruct}, and
then you redeploy the factory at its original \texttt{CREATE2} address, the
factory's nonce resets to zero, and it can deploy a \emph{completely different}
child contract at the \emph{same address} where the original child lived. The
address is reused. The code is not.

The second is the proxy pattern. Tornado Cash's governance system, like many DAO
governance contracts, used a proxy architecture (EIP-1967/UUPS). In a proxy
pattern, there is a permanent proxy contract at a fixed address that users
interact with. This proxy does not contain the actual governance logic. Instead,
it contains a pointer -- a storage slot at a specific, standardized location --
that holds the address of an \emph{implementation} contract. When you call a
function on the proxy, the proxy uses \texttt{delegatecall} to forward your call
to whatever implementation contract the pointer currently references. The
proxy's storage is used, but the implementation's code runs. This means whoever
can change the pointer controls what code executes when anyone interacts with
the governance system. Change the pointer, and you change the governance --
silently, without deploying a new visible contract, without changing the address
that everyone knows and trusts.

Now the trick.

The attacker submitted a governance proposal to the Tornado Cash DAO. The
proposal looked benign. Its description claimed it was identical to Proposal 16,
a previously approved and uncontroversial proposal that penalized certain
relayers for cheating. The voters did what voters do in a DAO with dozens of
proposals per month: they read the description, saw it matched something
familiar, and voted yes. They did not decompile and audit the proposal's
bytecode. Why would they? The description said it was the same proposal.
Reviewing raw EVM bytecode is not a skill most governance participants possess,
and the social norm in DAO governance is to review descriptions, not opcodes.

The vote passed. The proposal was approved by the DAO's governance process, with
legitimate TORN token holders casting legitimate votes through the legitimate
governance interface. Democracy had spoken.

Then the floor opened.

The proposal contract that the voters had approved contained a hidden
capability: \texttt{selfdestruct}. This EVM opcode does exactly what its name
suggests -- it destroys the contract at a given address, wiping its bytecode
from the blockchain state and sending any remaining ETH balance to a specified
recipient. After the vote passed and the proposal was executed, the attacker
triggered \texttt{selfdestruct} on the proposal contract. The code that the
voters had approved ceased to exist on the blockchain.

Then the attacker redeployed. Using the \texttt{CREATE2} intermediary trick
described above, the attacker deployed entirely new bytecode at the same address
where the original proposal contract had lived. The Tornado Cash governance
system still held a reference to that address. It still trusted that address.
But the code living there was now completely different from what the voters had
approved.

The new code did one thing: it gave the attacker the ability to mint TORN
governance tokens to themselves -- 10,000 TORN per iteration, repeatable, until
the attacker held 1.2 million votes. The entire legitimate DAO held roughly
700,000 votes. The attacker now controlled a permanent, unchallengeable
supermajority.

The misdirection was total. The malicious code was not present during the vote.
It did not exist when the voters examined the proposal. The voters approved code
A. The attacker destroyed code A and deployed code B at the same address. The
governance system, still pointing at that address, treated code B as if it had
the full authority of the vote that approved code A. The signed letter's text
changed after the seal was broken -- and the seal still looked intact.

Impact: complete control over Tornado Cash's governance. The attacker could
drain locked tokens, modify protocol parameters, brick the router contract, or
do anything else the governance system was authorized to do. Approximately
\$2.17 million was stolen directly. The TORN token price dropped 36\% as the
market priced in the total capture of the protocol's decision-making apparatus.
A privacy protocol whose zero-knowledge cryptography was unbroken -- whose
mathematical guarantees remained perfectly sound -- was nevertheless fully
compromised, because the human layer that governed it was exploitable through
misdirection and code replacement.

The root cause was two vulnerabilities woven together: a social one and a
technical one. The social vulnerability was that voters verified the proposal's
\emph{description} but not its \emph{code}. This is normal human behavior. It is
also, in hindsight, a systemic weakness of every DAO that presents proposals as
human-readable summaries rather than requiring formal verification of the
underlying bytecode. The technical vulnerability was the \texttt{selfdestruct} +
\texttt{CREATE2} pattern, which allowed post-approval code replacement at a
trusted address -- a capability that the governance system had no mechanism to
detect or prevent.

Neither vulnerability alone would have been sufficient. Together, they allowed
an attacker to go beyond exploiting the governance -- to \emph{become} the
governance. The Beanstalk attacker rented governance power for thirteen seconds.
The Tornado Cash attacker did something structurally worse: he permanently
replaced the governance with himself. Beanstalk was a heist. Tornado Cash was a
coup.

\subsection{ZK Rollup Governance Risk}\label{zk-rollup-governance-risk}

Both attacks targeted governance mechanisms that controlled upgradeable
contracts. And ZK rollup verifier contracts are almost always deployed behind
upgradeable proxy patterns -- the same patterns catalogued in the 2023 survey by
Meisami and Bodell, which documented EIP-1967 (OpenZeppelin transparent proxy),
EIP-1822 (UUPS), EIP-2535 (Diamonds), and Beacon proxies.

The proxy pattern introduces its own attack surface beyond governance: storage
layout corruption when state variables are reordered across upgrades, function
selector collisions between proxy admin and implementation functions, and the
fundamental risk that \texttt{delegatecall} means all storage operations in the
implementation affect the proxy's storage.

But the deepest risk is simpler than any of these. Whoever controls the proxy
admin controls the verifier. If the governance mechanism that controls the proxy
admin is vulnerable to flash loans (Beanstalk-style) or code replacement
(Tornado Cash-style), then the entire rollup's security reduces to the security
of its governance mechanism.

The cryptography could be perfect. The ceremony could have had a million
participants. The circuits could be formally verified. None of it matters if an
attacker can replace the verifier contract with one that returns \texttt{true}
for every proof. Six layers of mathematical elegance, and the seventh is a
multisig.

\subsection{L2Beat's Stages Framework}\label{l2beats-stages-framework}

L2Beat, the independent rollup monitoring organization, has formalized the
maturity of rollup decentralization into three stages:

\textbf{Stage 0 -- Full Training Wheels}: The rollup is effectively run by its
operators. It must have source-available software for state reconstruction from
L1 data, and it must have \emph{some} proof system to qualify. But governance
can override everything. Most rollup deployments begin here.

\textbf{Stage 1 -- Limited Training Wheels}: The proof system is fully
functional. Fraud proof submission (for optimistic rollups) or verification (for
ZK rollups) is permissionless. Users can exit without operator coordination
through forced inclusion or escape hatches. A Security Council may override the
proof system for bug fixes, but with constraints -- for example, a 6-of-8
multisig with a 7-day delay.

\textbf{Stage 2 -- No Training Wheels}: The rollup is fully managed by smart
contracts. The proof system is permissionless. Users get at least 30 days'
notice for unwanted upgrades. The Security Council is restricted to adjudicating
on-chain-provable soundness errors only. Users are fully protected from
governance attacks.

As of early 2026, most major ZK rollups are at Stage 0 or Stage 1. Achieving
Stage 2 requires either formally verified verifier contracts (so bugs are
unlikely enough that upgrade capability can be removed), multiple independent
implementations that cross-check each other, or bounded upgrade windows with
mandatory exit periods of 30 or more days.

The tension is real and irreducible. ZK verifier contracts are among the most
complex smart contracts ever deployed. They implement pairing checks, polynomial
evaluations, and Fiat-Shamir transcript verification in a language (Solidity, or
Yul) that was not designed for this kind of arithmetic. The probability of bugs
is non-trivial. But the ability to fix bugs via governance is the same ability
that allows governance to introduce them.

Stage 2 is where the cryptographic guarantees from Layers 1 through 6 actually
bind. Below Stage 2, they are advisory. A Stage 0 rollup with 256-bit proof
security is, from the user's perspective, only as secure as the governance
multisig's operational security.

\section{Proof Aggregation: The Missing
Layer}\label{proof-aggregation-the-missing-layer}

Between the prover (who generates proofs) and the on-chain verifier (who checks
them), a significant infrastructure layer has emerged that the seven-layer model
does not account for: proof aggregation services.

\textbf{SHARP (Shared Prover)}, built by StarkWare for Starknet, is the original
aggregation service. Multiple applications submit their execution traces to
SHARP, which generates a single STARK proof covering all of them, then wraps
that proof in Groth16 for on-chain verification. The verification gas cost is
amortized across all participating applications.

\textbf{Aligned Layer}, launched on mainnet with over \$11 billion in restaked
ETH (via EigenLayer), provides verification-as-a-service. Rollups and
applications submit proofs to Aligned Layer, which batches and verifies them,
posting the aggregated result to Ethereum.

\textbf{NEBRA}, live since August 2024, provides proof aggregation with a focus
on universal verification -- supporting multiple proof systems (Groth16, PLONK,
STARK) within a single aggregation layer.

The economic logic is straightforward. If a single Groth16 verification costs
\textasciitilde200,000 gas, and an aggregation service can batch 100 proofs into
a single on-chain verification, the per-proof verification cost drops from
\textasciitilde\$1 to \textasciitilde\$0.01. At sufficient volume, aggregation
makes proof verification nearly free.

But aggregation introduces a new trust assumption. Users must trust that the
aggregation service correctly includes their proof in the batch, and that the
aggregated proof faithfully represents all constituent proofs. If the
aggregation service is centralized (as SHARP is for Starknet), this is a single
point of failure at Layer 7 that can undermine the decentralization guarantees
of the underlying proof system.

\section{Case Study: Midnight and the Three-Token
Architecture}\label{case-study-midnight-and-the-three-token-architecture}

Midnight, developed by IOG (Input Output Global), provides an instructive case
study for Layer 7 because it makes different architectural choices than the
Ethereum rollup model. Where Ethereum rollups post proofs to a general-purpose
L1 and rely on upgradeable verifier contracts, Midnight integrates verification
into its consensus layer and uses a novel three-token economic model that
directly shapes the verification experience.

\subsection{The Verification Pipeline}\label{the-verification-pipeline}

Midnight uses a split execution model. Smart contracts are written in Compact, a
domain-specific ZK language, but they never execute on-chain in the traditional
sense:

\begin{enumerate}
\def\labelenumi{\arabic{enumi}.}
\tightlist
\item
  The user's SDK executes the Compact circuit locally, computing the new state.
\item
  A proof server generates a ZK proof that the state transition is valid.
\item
  The SDK packages the proof, fee inputs, and state delta into a transaction.
\item
  Every Midnight node verifies the proof against the circuit's verifier keys,
  which were deployed on-chain when the contract was created.
\item
  If valid, the blockchain updates the contract's ledger state.
\end{enumerate}

The critical difference from the Ethereum rollup model: verification is not
performed by a specialized verifier contract that can be upgraded via
governance. It is performed by every node as part of consensus. The verifier
keys are stored on-chain at the contract address and are immutable -- once
deployed, a contract's verification logic cannot be changed. A new contract must
be deployed for logic changes.

This is a strong answer to the governance-as-attack-surface problem. You cannot
upgrade what is immutable. But it creates a different problem: what happens when
there is a bug? The answer is migration -- deploy a corrected contract and
convince users to move to it. This is slower and messier than a governance
upgrade, but it cannot be exploited by an attacker with admin keys. The tradeoff
is explicit: Midnight accepts the inconvenience of immutability in exchange for
immunity to the Beanstalk-style attack.

\subsection{Three Tokens, Three Privacy
Levels}\label{three-tokens-three-privacy-levels}

Midnight's three-token model is not an arbitrary design choice. Each token
represents a different point on the privacy-transparency spectrum, and together
they create an economic system where verification costs, privacy guarantees, and
governance rights are separate and independently tunable.

\textbf{Night} is the unshielded native token. It is fully transparent -- all
operations are publicly visible on the ledger, stored as UTXOs with public-key
authentication. Night serves two purposes: staking (and therefore governance)
and backing for DUST generation. Every Night token registered for dust
generation produces DUST over time at a deterministic rate.

\textbf{Shielded tokens} are ZK-private custom tokens. Any Compact contract can
mint shielded tokens with unique type identifiers (``colors''). Balances and
transaction details are hidden via zero-knowledge proofs. Shielded tokens use
UTXOs with Pedersen commitments -- only the key holder can view or modify wallet
state. All shielded token transfers go through contracts, not direct
peer-to-peer.

\textbf{DUST} is the fee token. All transaction fees are denominated in DUST,
but DUST is not mined or minted in the traditional sense. It is a time-dependent
scalar computed from Night holdings. A user registers Night UTXOs for dust
generation, and DUST accumulates over time according to a deterministic formula
with parameters for rate, maximum cap, and creation time.

This creates a novel fee model with direct implications for verification
economics:

\begin{itemize}
\tightlist
\item
  \textbf{No fee market}: Fees are computed deterministically, not bid. There is
  no gas auction.
\item
  \textbf{Rate-limited spam}: Transaction throughput is naturally limited by
  DUST regeneration rate relative to Night holdings. An attacker who wants to
  spam the network must hold (or acquire) Night tokens and wait for DUST to
  regenerate.
\item
  \textbf{Staking alignment}: Only Night holders can generate DUST. Transaction
  capability is tied to network participation.
\item
  \textbf{Time-gated recovery}: A user who has spent all their DUST must wait
  for regeneration before transacting again. This is a natural circuit breaker
  against denial-of-service attacks.
\end{itemize}

\subsection{Disclosure Rules and Compiler-Enforced
Privacy}\label{disclosure-rules-and-compiler-enforced-privacy}

In Compact, everything is private by default. Values only appear on-chain when
the developer explicitly calls \texttt{disclose()}. The Compact compiler
enforces this through static disclosure analysis at compile time -- privacy is a
compiler guarantee, not developer discipline.

What must be disclosed includes contract ledger state (any
\texttt{export\ ledger} field), counter increments and decrements, state
transition deltas (so validators can verify the new state), nullifiers (to
prevent double-spending), and hash commitments (stored for future verification).
What stays private includes secret keys, witness values (circuit private
inputs), shielded balances, vote choices (only aggregate tallies change
on-chain), authorization preimages, and computation logic (ZK proof hides the
execution path).

This compiler-enforced privacy boundary is a significant Layer 7 innovation. In
the Ethereum model, what is public and what is private depends on the
developer's care in managing calldata, events, and storage. In Midnight, the
compiler draws the line, and crossing it requires an explicit annotation that is
visible in code review.

\subsection{Private Governance}\label{private-governance}

Midnight's DAO governance pattern demonstrates what private on-chain governance
can look like:

\begin{itemize}
\tightlist
\item
  \textbf{Anonymous identity}: Voters prove membership via hash commitments,
  never revealing their real identity.
\item
  \textbf{Weighted voting}: Contract-state token balances serve as anonymous
  voting weights via ZK circuit reads.
\item
  \textbf{Per-proposal nullifier domains}: Each proposal has its own nullifier
  space, preventing double-voting while allowing participation across multiple
  proposals.
\item
  \textbf{Vote privacy}: The voter's choice remains private -- only the
  aggregate tally changes on-chain.
\item
  \textbf{Irreversible state machines}: Proposal status is encoded as counter
  increments (0 unused, 1 open, 2 approved, 3 executed), and counters are
  monotonically increasing, so state transitions are irreversible by
  construction.
\end{itemize}

The multi-signature treasury adds M-of-N threshold approval with
propose/approve/execute circuits, where signer identity is verified via hash
commitment and double-vote prevention uses per-(proposal, signer) nullifiers.

Every design choice is a direct architectural response to the attacks we just
witnessed. Anonymous weighted voting means an attacker cannot flash-loan
governance tokens and vote -- they would need to know the secret key that
corresponds to a registered hash commitment, and flash-loaning tokens does not
give them that. Irreversible state machines mean a proposal cannot be rolled
back after execution. Per-proposal nullifier isolation means vote manipulation
in one proposal cannot leak to another.

The Beanstalk attacker borrowed a billion dollars of voting power for the
duration of a single transaction. Against Midnight's architecture, that
borrowing would be useless. You cannot vote with a key you do not possess.

\subsection{The Gaps}\label{the-gaps}

Midnight's approach is not without its own Layer 7 vulnerabilities:

\begin{itemize}
\tightlist
\item
  \textbf{Protocol upgrade governance}: The documentation does not describe how
  consensus-level parameters (fee rates, dust generation parameters, consensus
  rules) are governed. This is the most significant gap. Immutable contracts
  solve contract-level governance attacks, but someone must still govern the
  protocol itself.
\item
  \textbf{Oracle centralization}: All oracle patterns in the current
  documentation use single-party authorization via hash commitment. There is no
  multi-oracle or threshold-oracle pattern. A single compromised oracle can feed
  false data to every contract that depends on it.
\item
  \textbf{Fixed participant sets}: Current governance contracts hardcode 2-3
  participant slots. Production governance with dynamic participant sets would
  require Merkle-tree-based registration, which has been demonstrated (in the
  lending pool pattern) but not yet integrated into governance contracts.
\item
  \textbf{No emergency procedures}: There is no documented kill switch, pause
  mechanism, or emergency parameter override for deployed contracts.
  Immutability is a feature for preventing governance attacks, but it is a
  liability when a critical bug is discovered.
\end{itemize}

\section{The Deepest Symmetry}\label{the-deepest-symmetry}

There is a symmetry in the seven-layer model that becomes visible only at Layer
7, and it concerns the nature of trust.

At Layer 1, the setup ceremony, security rests on a social claim: ``At least one
of N participants honestly destroyed their toxic waste.'' This is not a
mathematical statement. It is a statement about human behavior. You trust the
ceremony because you trust that at least one person, out of thousands, did the
right thing.

At Layer 7, the verifier deployment, security rests on a parallel social claim:
``The governance mechanism that controls the verifier will not be captured by an
adversary.'' This is not a mathematical statement either. It is a statement
about institutional design, incentive alignment, and ultimately human judgment.

The layers in between -- the language, the witness, the arithmetization, the
proof system, the primitives -- are mathematical. They provide computational
guarantees that hold against any polynomial-time adversary. They are the part of
the trick that actually works by the laws of mathematics, not by the conventions
of human society. But they are sandwiched between two layers of social trust.

None of this represents a failure of the model. It is a description of reality.
Zero-knowledge proofs do not eliminate trust. They \emph{compress} it. Instead
of trusting a bank with your financial data every day, you trust that a ceremony
was run honestly once and that governance will not go rogue in the future. These
are weaker assumptions than trusting a single counterparty for every
transaction. But they are assumptions nonetheless.

The honest framing is not ``trustless.'' It is ``trust-minimized.'' And the
remaining trust assumptions -- ceremony integrity at the bottom, governance
integrity at the top -- are worth stating explicitly so that readers can
evaluate whether the trust reduction justifies the complexity.

Feynman, who had a gift for puncturing pretension, would probably say something
like this: ``You have built a beautiful machine that converts social trust into
mathematical certainty and back into social trust again. The mathematical part
in the middle is genuinely impressive. But do not pretend the social parts at
the ends do not exist.''

He would be right. And the fact that he would be right is itself the deepest
insight Layer 7 has to offer. The magic trick is real. The mathematics works.
But the trick is performed for an audience, and the audience is governed by
people, and people are not mathematical objects. The security of the whole
system is a chain, and the endpoints of that chain are anchored in human soil.

\section{Pricing Attacks}\label{pricing-attacks}

The relationship between verification costs and data availability costs creates
exploitable seams. A 2025 study by Chaliasos et al.~identified two novel attack
classes that exploit mismatches in how rollups price their three cost
dimensions: L2 execution, L1 data availability, and L1 settlement/proving.

\subsection{DA-Saturation Attacks}\label{da-saturation-attacks}

An attacker floods an L2 with data-heavy, compute-light transactions --
essentially random calldata followed by a STOP opcode. Each such transaction is
cheap in L2 gas (the computation is trivial) but expensive in L1 DA cost (it
fills blob space with incompressible junk). The study found that sustained
denial-of-service on Linea cost as little as 0.87 ETH per hour, and on Optimism
roughly 2 ETH per 30 minutes.

The effects cascade: the L2 is congested, finality delays increase by 1.45x to
2.73x compared to direct L1 blob stuffing, and the rollup operator hemorrhages
money because the fees collected from the spam transactions do not cover the L1
DA costs. All major rollups studied were found susceptible. Four bug bounties,
each worth tens of thousands of dollars, were paid.

\subsection{Prover-Killer Attacks}\label{prover-killer-attacks}

These exploit the mismatch between EVM gas metering and ZK proving costs. Not
all EVM opcodes are equally expensive to prove in zero knowledge. The study
measured ``cycles per gas'' ratios -- how many proving cycles are needed per
unit of gas cost:

\begin{longtable}[]{@{}llll@{}}
\toprule\noalign{}
Opcode / Precompile & EVM Gas & Proving Cycles/Gas & Attack Leverage \\
\midrule\noalign{}
\arrayrulecolor{MidnightBlue}\midrule\arrayrulecolor{MidnightBorder}
\endhead
\bottomrule\noalign{}
\endlastfoot
JUMPDEST & 1 & 1,039.79 & Very High \\
MODEXP & (varies) & 2,961.72 & Extreme \\
BN\_PAIRING & 45,000+ & 1,642.15 & High \\
SHA256 & 60+ & Moderate & Moderate \\
\end{longtable}

A MODEXP attack -- filling blocks with maximum-cost modular exponentiation
operations -- delayed finality by 94x (over 8 hours) and cost the rollup
operator \$42.26 per attack block. The rollup's proving system crashed after
10,266 seconds.

\subsection{The Concrete Scenarios}\label{the-concrete-scenarios}

These are not theoretical concerns. They are playbooks.

\textbf{DA-saturation in practice.} An attacker constructs transactions that
contain the maximum possible calldata -- random bytes, incompressible by design
-- followed by a STOP opcode. The EVM execution cost is negligible: STOP costs 0
gas, and the transaction is technically valid. But the data must be posted to L1
for data availability, and incompressible random data consumes maximum blob
space. The attacker submits these transactions continuously, filling every blob
slot in every Ethereum block with garbage.

The immediate effect: blob fees spike. Ethereum's blob fee market operates under
EIP-1559 dynamics -- when blobs are consistently full, the base fee increases
exponentially. Under sustained saturation, blob fees can increase by 100x or
more within minutes. Every legitimate rollup that needs to post data to Ethereum
sees its operating costs spike proportionally. The attacker pays blob fees too,
of course, but the attacker's cost is the cost of the attack. Every other
rollup's cost is collateral damage.

The secondary effect is subtler and worse: rollup operators, facing unexpectedly
high L1 costs, must choose between posting data at a loss (subsidizing
operations from their treasury), delaying batch submissions (increasing finality
time for users), or raising L2 fees (driving users to competitors). The attacker
does not need to sustain the attack indefinitely. A few hours of blob saturation
can cause lasting reputational and economic damage to rollups that depend on
predictable L1 costs.

The defense is multidimensional fee pricing on the L2 side: a separate DA fee
component that adjusts dynamically based on actual L1 blob costs, passed through
to users in real time rather than absorbed by the operator. Several rollups have
implemented this in response to the Chaliasos findings, but the adjustment is
inherently reactive -- the fee increases \emph{after} the attack begins, which
means the operator absorbs losses during the lag period.

\textbf{Prover-killer in practice.} An attacker submits transactions that are
cheap in EVM gas but catastrophically expensive to prove in zero knowledge. The
canonical example is MODEXP -- modular exponentiation with maximum-size inputs.
The EVM prices MODEXP based on the size of the operands and a formula that was
calibrated for native CPU execution, not for ZK circuit execution. A MODEXP
operation with 256-byte base, exponent, and modulus costs roughly 200 gas in the
EVM. Proving that same operation inside a ZK circuit requires the prover to
decompose the modular exponentiation into field arithmetic over the proving
system's native field, which involves thousands of multiplication gates per
limb, per exponentiation step. The ratio of proving cost to EVM gas cost -- the
``cycles per gas'' metric -- reaches nearly 3,000 for MODEXP. A single MODEXP
transaction can consume as much proving capacity as 3,000 normal transactions.

The attacker does not need exotic tools. They submit valid transactions --
MODEXP calls with legitimate inputs that any EVM will execute without complaint.
The transactions pass all validation checks. They pay standard gas fees. They
are included in blocks by the sequencer because the sequencer has no reason to
reject a valid, fee-paying transaction. But when those blocks reach the prover,
the prover must generate a ZK proof of the entire block's execution, including
the MODEXP operations. A single block stuffed with MODEXP calls can take the
prover 10,000 seconds to prove -- over 2.7 hours for a block that the sequencer
produced in seconds.

If the attacker sustains this for multiple blocks, the prover falls behind.
Proving latency grows. Finality -- the time before a rollup batch is verified on
L1 -- stretches from minutes to hours. The rollup is technically still
functioning, but its security guarantee degrades: until the proof is posted and
verified, the rollup's state transitions are unproven claims, not verified
facts. A sustained prover-killer attack can push a ZK rollup into a state where
it behaves, from the user's perspective, more like an optimistic rollup --
running on trust rather than proof, hoping nothing goes wrong during the gap.

The defense requires the sequencer to price transactions based on their
\emph{proving cost}, not just their EVM gas cost. This is the multidimensional
pricing problem: EVM gas, DA cost, and proving cost are three independent
resources, and a single gas price cannot accurately reflect all three. Until
rollups implement proving-cost-aware fee markets, the prover-killer attack
remains viable against any ZK rollup that prices transactions using only EVM gas
metering.

The root cause is fundamental: current rollup fee mechanisms use a
single-dimensional gas price that bundles L2 execution, L1 DA, and proving costs
into one number. When these three resources have different scarcity profiles,
the bundled price necessarily misprices at least one of them. The fix is
multidimensional pricing -- separate base fees for each resource type, each
following its own EIP-1559-style adjustment mechanism.

\section{Who Verifies the Verifier?}\label{who-verifies-the-verifier}

The verifier smart contract itself can have bugs. The FOOM Club exploit targeted
a misconfigured snarkjs deployment where the verification key parameter delta\_2
was set equal to gamma\_2, weakening the Groth16 verification equation. The
proof system was not broken -- the \emph{deployment configuration} was wrong.

The vulnerability is a supply-chain problem. The verifier contract depends on:

\begin{enumerate}
\def\labelenumi{\arabic{enumi}.}
\tightlist
\item
  The proof system specification (mathematical, usually correct).
\item
  The reference implementation (code, sometimes buggy -- see Frozen Heart).
\item
  The deployment configuration (operational, frequently wrong).
\item
  The Ethereum precompiles (hardware/protocol, generally reliable but not immune
  to bugs).
\item
  The compiler that compiled the verifier contract (Solidity, Vyper, or Yul --
  each with their own bug history).
\end{enumerate}

An analogy to the XZ Utils supply-chain attack (CVE-2024-3094) is apt. In that
case, a sophisticated attacker spent years contributing to an open-source
compression library, gained maintainer trust, and inserted a backdoor. The same
attack vector applies to ZK verifier libraries: snarkjs, gnark, arkworks, and
halo2 are open-source projects maintained by small teams. A compromised
maintainer could introduce a subtle verification bypass that passes all existing
tests.

Verifier ossification -- the strategy of deploying a verifier contract and
making it permanently immutable, treating it like a protocol-level constant
rather than upgradeable software -- is one defense. But it requires very high
confidence in the verifier's correctness, because bugs in an ossified verifier
cannot be fixed without deploying an entirely new contract and migrating all
dependent applications.

The tradeoff between immutable and upgradeable verifiers is one of the sharpest
architectural decisions at Layer 7:

\begin{longtable}[]{@{}
  >{\raggedright\arraybackslash}p{(\linewidth - 4\tabcolsep) * \real{0.2000}}
  >{\raggedright\arraybackslash}p{(\linewidth - 4\tabcolsep) * \real{0.3800}}
  >{\raggedright\arraybackslash}p{(\linewidth - 4\tabcolsep) * \real{0.4200}}@{}}
\toprule\noalign{}
\begin{minipage}[b]{\linewidth}\raggedright
Property
\end{minipage} & \begin{minipage}[b]{\linewidth}\raggedright
Immutable Verifier
\end{minipage} & \begin{minipage}[b]{\linewidth}\raggedright
Upgradeable Verifier
\end{minipage} \\
\midrule\noalign{}
\arrayrulecolor{MidnightBlue}\midrule\arrayrulecolor{MidnightBorder}
\endhead
\bottomrule\noalign{}
\endlastfoot
Bug patching & Impossible without contract migration & Possible via governance
vote or multisig \\
Governance capture & Immune --- no upgrade path to exploit & Vulnerable ---
Beanstalk/Tornado Cash-style attacks \\
Regulatory compliance & Fixed at deploy time; cannot adapt & Adaptable to
changing requirements \\
User trust model & Trust the code (audit once, rely forever) & Trust the
governance (ongoing vigilance) \\
L2Beat Stage & Stage 2 candidate (if verifier is correct) & Stage 0-1
(governance can override proofs) \\
Quantum migration & Requires full system replacement & Can upgrade to PQ
verifier via governance \\
Example & Midnight (immutable verifier keys) & Most Ethereum ZK rollups (proxy
pattern) \\
\end{longtable}

Neither choice dominates. Immutable verifiers maximize cryptographic integrity
at the cost of operational flexibility. Upgradeable verifiers maximize
adaptability at the cost of governance risk. The choice reflects a system's
threat model: does it fear bugs more, or governance capture more?

\section{On-Chain Verification in 2026}\label{on-chain-verification-in-2026}

The state of on-chain verification as of early 2026:

\textbf{Verification costs} have stabilized. Groth16 on BN254 remains the
dominant on-chain proof format, at roughly 200,000-250,000 gas per verification.
The verification cost floor is set by the pairing precompile gas schedule, which
is a protocol parameter that changes only through Ethereum governance (EIPs and
hard forks).

\textbf{Data availability} is abundant and cheap. Three Ethereum upgrades in two
years (Dencun, Pectra, Fusaka) have expanded DA capacity by roughly 16x.
Alternative DA layers (Celestia, EigenDA, Avail) provide even cheaper options at
the cost of different security assumptions.

\textbf{Governance maturity} lags. Most ZK rollups remain at Stage 0 or Stage 1
of L2Beat's framework. The path to Stage 2 -- where governance can no longer
override the proof system -- requires either formally verified verifier
contracts, multi-prover architectures, or long mandatory exit windows. No major
ZK rollup has achieved Stage 2 as of this writing.

\textbf{Fiat-Shamir security} is improving through hard experience. The Frozen
Heart disclosure, the Last Challenge Attack, and the Solana ZK ElGamal bug have
established Fiat-Shamir transcript completeness as a first-order security
property. Audit firms now check for it specifically. But new implementations
continue to be written, and the pattern will recur until proof system libraries
converge on a small number of battle-tested implementations.

\textbf{Proof aggregation} is maturing. SHARP, Aligned Layer, and NEBRA
demonstrate that aggregation can reduce per-proof verification costs by 10-100x.
But aggregation services are themselves centralization points that need their
own governance and security analysis.

The net picture: Layer 7 is the layer where the mathematical elegance of Layers
1 through 6 collides with the messy realities of software deployment, economic
incentives, governance design, and human judgment. The cryptography is strong.
The implementations are getting stronger. But the governance -- the social layer
that determines who can change the software that checks the math -- remains the
binding constraint on the security that zero-knowledge proofs can actually
deliver to end users.

Until the governance matures to Stage 2 -- until the smart contracts that verify
proofs are either immutable or governed by mechanisms that provably resist
capture -- the verdict remains provisional. The audience is competent. The math
checks out. But the audience serves at the pleasure of a committee that can
replace it at any time.

Layer 7 is the last layer. The seven-layer tour -- from setup ceremony to
on-chain verdict -- is complete. But zero-knowledge proofs do not operate in
isolation. They belong to a family of privacy-enhancing technologies -- MPC,
FHE, differential privacy, TEEs -- and understanding ZKPs without understanding
their siblings leads to architectures that reach for the right mathematics and
solve the wrong problem. Before we synthesize the seven layers in Part III, we
map the family.

\chapter{Chapter 9: Privacy-Enhancing
Technologies}\label{chapter-9-privacy-enhancing-technologies}

\emph{``Privacy is not a feature you bolt on. It is a property of the
architecture -- and if the architecture does not have it, no amount of
cryptography will give it to you.''}

\par\vspace{18pt}

This chapter makes a single argument: zero-knowledge proofs are necessary but
not sufficient for privacy. A system that deploys ZKPs alone will leak through
side channels, metadata, composition boundaries, and the gaps between what the
proof covers and what the system exposes. The architect who reaches for ZKPs
without understanding MPC, FHE, and differential privacy will build a house with
a vault door on the front and a screen door on the back. Each technology in the
PET family answers a different question. Understanding which question you are
actually asking -- and which tool answers it -- is the prerequisite for any
privacy architecture that works in practice.

We examine four technologies (ZKP, MPC, FHE, differential privacy), classify
their security guarantees into three tiers (information-theoretic,
computational, heuristic), study their composition patterns through five
real-world deployments, and end with a decision matrix for system architects
choosing between them.

\section{The Four Pillars}\label{the-four-pillars}

Zero-knowledge proofs are one member of a family. The family is called
privacy-enhancing technologies -- PETs -- and understanding ZKPs without
understanding their siblings is like watching a magician and concluding that all
performance is sleight of hand. It is not. Some performers sing. Some dance.
Some vanish entirely. And the best shows combine all of them.

There are four major PET categories that matter for system architects making
decisions in 2026. Each answers a different question. Each has a different
relationship to the stage.

\textbf{Zero-Knowledge Proofs (ZKPs)} answer: \emph{How do I prove a statement
about my private data without revealing the data itself?}

Think of ZKPs as the magician's core trick: the audience sees the result but not
the method. ``My balance exceeds the minimum.'' ``I am over 18.'' ``This
computation was performed correctly.'' The key insight, often missed, is that
ZKPs are a tool for \emph{selective disclosure}, not blanket privacy. The proof
reveals the \emph{truth of the statement} -- which is itself information.
Proving ``my balance exceeds \$1 million'' tells you something about the
balance, even though the exact figure stays hidden. The magician chooses which
cards to reveal. That choice is itself a disclosure.

If ZKPs are the magician's core trick -- proving truth without revealing method
-- then the three siblings each perform a different kind of magic.

\textbf{Secure Multi-Party Computation (MPC)} answers: \emph{How can multiple
parties jointly compute a function on their combined data without any party
revealing its input to the others?}

Picture three rival magicians who want to know whose trick is most popular --
but none will reveal their ticket sales to the others. MPC is the protocol that
lets them compute the answer as if a trusted accountant had all the books,
without any such accountant existing. The inputs stay private. Only the
agreed-upon output is revealed.

MPC is not a single protocol. It is a family, and the family members have very
different properties:

\begin{longtable}[]{@{}
  >{\raggedright\arraybackslash}p{(\linewidth - 6\tabcolsep) * \real{0.2963}}
  >{\raggedright\arraybackslash}p{(\linewidth - 6\tabcolsep) * \real{0.2407}}
  >{\raggedright\arraybackslash}p{(\linewidth - 6\tabcolsep) * \real{0.2778}}
  >{\raggedright\arraybackslash}p{(\linewidth - 6\tabcolsep) * \real{0.1852}}@{}}
\toprule\noalign{}
\begin{minipage}[b]{\linewidth}\raggedright
Protocol Family
\end{minipage} & \begin{minipage}[b]{\linewidth}\raggedright
Trust Model
\end{minipage} & \begin{minipage}[b]{\linewidth}\raggedright
Security Type
\end{minipage} & \begin{minipage}[b]{\linewidth}\raggedright
Best For
\end{minipage} \\
\midrule\noalign{}
\arrayrulecolor{MidnightBlue}\midrule\arrayrulecolor{MidnightBorder}
\endhead
\bottomrule\noalign{}
\endlastfoot
Shamir secret sharing & Honest majority (\textgreater50\% honest) &
Information-theoretic & Statistical computations, few parties \\
SPDZ (dishonest majority) & Any number of corruptions & Computational &
Adversarial settings, financial computation \\
Garbled circuits & Two parties, semi-honest or malicious & Computational &
Two-party computation \\
\end{longtable}

The distinction between honest-majority and dishonest-majority protocols matters
enormously. Shamir-based MPC with honest majority achieves
\emph{information-theoretic} security -- it remains secure even against an
adversary with unlimited computational power, including quantum computers. SPDZ
provides security against any number of corruptions, but relies on computational
hardness assumptions.

To understand what MPC actually does, consider the problem that started the
field. In 1982, Andrew Yao posed what is now called the Millionaires' Problem:
two millionaires want to determine who is richer without either revealing their
net worth. They are standing at a cocktail party. Neither will say a number.
Neither trusts the other to be honest. And no accountant is available whom both
would trust with the truth. How do they find out?

Here is the trick. Alice has a net worth of, say, \$7 million. Bob has \$5
million. Alice encodes her wealth into an encrypted lookup table -- a garbled
circuit -- that represents the comparison function ``is Alice's input greater
than Bob's input?'' She hands the garbled table to Bob. Bob, using a
sub-protocol called oblivious transfer, obtains the encryption key corresponding
to his own input (\$5 million) without Alice learning which key he selected. Bob
evaluates the garbled circuit with his key and obtains a single bit of output:
``Alice is richer.'' He announces the result. Neither party learned the other's
number. The function was computed. The inputs stayed private.

The garbled circuit deserves a moment of its own, because it is one of the most
counterintuitive constructions in all of cryptography. Alice takes the
computation she wants to perform and compiles it into a Boolean circuit -- AND
gates, OR gates, NOT gates, the same primitives that make up a physical
processor. She then encrypts every wire in the circuit with random labels: each
wire gets two labels, one for ``0'' and one for ``1,'' and Alice encrypts each
gate's truth table so that only the correct pair of input labels decrypts to the
correct output label. The result is a garbled mess -- a table of ciphertexts
that encodes the computation but reveals nothing about it. Bob can evaluate the
garbled circuit gate by gate, decrypting one entry per gate, following the
circuit from input to output. He sees the computation unfold, but the labels are
random strings. He learns the output and nothing else. Alice never sees Bob's
input. Bob never sees Alice's circuit internals. The computation happens in a
kind of cryptographic fog, visible only at the endpoints.

Scale this from cocktail-party curiosity to industrial infrastructure and the
applications multiply. Private auctions: bidders submit encrypted bids to an MPC
protocol that determines the winner and the clearing price without revealing any
losing bid. The auction house learns who won and at what price. It never learns
what the losers were willing to pay -- information that, in traditional
auctions, the house can exploit in future rounds. Dark pool matching in finance:
two investment banks want to match buy and sell orders for the same security
without revealing their order books to each other or to the market. MPC lets
them compute the intersection of their orders -- the trades that both sides want
-- without exposing the orders that did not match. The matched trades execute.
The unmatched orders remain invisible.

Private set intersection, or PSI, is the simplest and most widely deployed MPC
primitive. Two parties each hold a set of items. They want to learn which items
appear in both sets -- and nothing else. When COVID-19 contact tracing required
matching infected individuals against location databases, PSI offered a path
that did not require a central authority to hold everyone's location history.
Google and Apple's Exposure Notification framework used a related technique:
devices broadcast rotating pseudonymous identifiers, and the matching
computation happens locally on each device. The computation is distributed. The
data never congregates.

The cost of MPC is communication, not computation. Each gate in the circuit
requires the parties to exchange encrypted values. For garbled circuits, this
means transmitting roughly four ciphertexts per gate. A circuit with a billion
gates requires transmitting billions of ciphertexts -- tens of gigabytes over
the network. The computation itself is fast. The network is the bottleneck. This
is why MPC shines for problems where the function is simple but the privacy
requirement is absolute, and struggles for problems where the function is
complex and latency matters. The magician works quickly. The postal service does
not.

MPC is the ensemble performance: multiple magicians, each holding one card of a
shared secret, jointly computing a result that none of them could produce alone.
The audience sees the answer. No single performer ever sees the full hand.

\textbf{Fully Homomorphic Encryption (FHE)} answers: \emph{How can I outsource
computation on my data without the computing party learning anything about the
data?}

Craig Gentry, who invented FHE, gave it the perfect image: performing surgery on
a patient inside a sealed glovebox. The surgeon's hands are inside the gloves,
manipulating the patient, but the glovebox prevents any direct contact. The
surgeon can work -- she can cut, stitch, probe -- but she never touches the
patient directly, and nothing from the operating field crosses the barrier. It
is a vivid image from a different domain -- not our magician's stage but a
laboratory -- and we borrow it here because Gentry's metaphor has become
inseparable from the concept itself.

You encrypt your data. You send the ciphertext to a cloud provider. The cloud
provider performs operations on the ciphertext. You decrypt the result. The
cloud provider learns nothing -- not the data, not the result, not even which
operations were meaningful. The gloves never come off.

But the glovebox is thick, and the gloves are stiff. Current FHE computations
are 10,000x to 1,000,000x slower than their plaintext equivalents. Ciphertexts
accumulate noise with each operation, requiring periodic ``bootstrapping'' that
is computationally expensive. Not all operations are equally efficient --
additions and multiplications are the native operations, while comparisons and
divisions are far more costly. The surgeon can work, but she works very, very
slowly.

To understand why the gloves are so stiff, you need to see what an FHE
computation actually looks like from the inside. The dominant schemes -- BFV,
BGV, and CKKS -- all share a common mathematical structure. A plaintext value
(say, the number 42) is encoded as a polynomial in a ring, typically
\(\mathbb{Z}_q[X]/(X^n + 1)\) for \(n = 2^{13}\) or \(2^{14}\) and \(q\) a large
modulus, perhaps 200 to 800 bits long. Encryption adds a carefully sampled noise
term to this polynomial. The ciphertext is a pair of ring elements, each
thousands of bits wide. Where a plaintext integer fits in 64 bits, its
ciphertext might occupy 32 kilobytes. The glovebox is thick because the encoding
is thick.

Addition through the glovebox is relatively gentle. You add two ciphertexts
component-wise, and the noise terms add as well. The noise grows, but only
linearly. Encrypted addition is perhaps 100x slower than plaintext addition --
expensive, but not prohibitively so.

Multiplication is where the cost explodes. When you multiply two ciphertexts,
the underlying polynomial multiplication produces a result with noise that is
roughly the \emph{product} of the two input noise levels, not the sum. One
multiplication might double the noise budget. Two multiplications might
quadruple it. After a dozen consecutive multiplications without intervention,
the noise overwhelms the signal -- the ciphertext becomes a random-looking
polynomial that decrypts to garbage. The noise is not a flaw in the design. It
\emph{is} the security. Without noise, the lattice-based hardness assumptions
that make FHE secure would not hold. The glovebox is thick because thinning it
would make it transparent.

Gentry's breakthrough -- the idea that launched FHE from theoretical
impossibility to practical research program -- was bootstrapping. The concept is
recursive and almost paradoxical: to clean the noise from a ciphertext, you
decrypt it \emph{homomorphically}. That is, you take the noisy ciphertext,
encrypt the decryption key under a fresh public key, and then run the decryption
algorithm as a circuit \emph{inside the encryption}. The output is a fresh
ciphertext encrypting the same plaintext, but with reset noise -- as if you had
just encrypted the value for the first time. The surgeon, working through the
glovebox, performs a second surgery on the glovebox itself, replacing the foggy
glass with a clean pane, all without ever removing her hands.

The cost of this cleaning step is enormous. A single bootstrapping operation
might take 10 to 100 milliseconds on modern hardware -- which sounds fast until
you realize that the plaintext operation it replaces (a single multiplication)
takes about one nanosecond. The ratio is 10 million to one. And bootstrapping
must be performed after every few multiplications to keep the noise below the
fatal threshold. The 10,000x to 1,000,000x slowdown is not a single penalty
applied once. It is the accumulated cost of performing every arithmetic
operation on bloated polynomial ciphertexts and periodically cleaning the noise
through a decryption-inside-encryption cycle that is itself a complex
computation.

Recent optimizations have attacked every link in this chain. TFHE (Torus FHE)
reduces the bootstrapping cost for Boolean circuits by working over a different
algebraic structure. GPU acceleration parallelizes the polynomial ring
arithmetic. Hybrid schemes use leveled FHE (no bootstrapping) for shallow
circuits and switch to bootstrapping only when depth demands it. The overhead is
shrinking -- from a million-fold five years ago to ten-thousand-fold today --
but the fundamental structure remains: encrypted computation is expensive
because the noise that provides security must be managed, and managing it costs
orders of magnitude more than the computation itself.

And here is the connection that brings us back to our magician. FHE lets you
compute on encrypted data, but how do you know the computation was performed
\emph{correctly}? The cloud provider claims it evaluated your function honestly.
But the output is encrypted -- you cannot inspect the intermediate steps. The
provider could have computed a different function, or computed the right
function incorrectly, and you would not know until you decrypted the result and
found it nonsensical. Verifiable FHE -- zkFHE -- addresses this by having the
computing party produce a zero-knowledge proof that the homomorphic operations
were performed according to specification. The surgeon operates through the
glovebox, and a camera inside the glovebox records the procedure for the review
board. The patient stays sealed. The surgery is verified. This is where FHE and
ZKPs converge, and it is one of the most active research frontiers in applied
cryptography.

The field is improving rapidly. Zama, the leading FHE company, achieved a \$1
billion valuation in June 2025 and launched the Confidential Blockchain Protocol
testnet in July 2025. Their roadmap projects hundreds of transactions per second
with GPU acceleration. But a 10,000x overhead, even if it shrinks to 1,000x,
means FHE is practical only for computations where the privacy guarantee is
worth the performance cost. The glovebox is for surgery you cannot perform any
other way.

FHE is the trick performed inside a sealed box: the computation happens on
encrypted data, and even the magician who executes the computation never sees
the plaintext. The box opens to reveal only the result.

\textbf{Differential Privacy (DP)} answers: \emph{How can I release statistical
insights about a dataset while guaranteeing that no individual's record can be
reverse-engineered from the output?}

If the other PETs are stage tricks -- precise, targeted, visible to the audience
-- differential privacy is fog. It blurs the picture just enough that no
individual face can be identified, while the overall scene remains recognizable.
Apple uses it for iOS telemetry (since 2016, with \(\varepsilon = 2\) per day
for most data types). Google deployed RAPPOR (Randomized Aggregatable
Privacy-Preserving Ordinal Response) for Chrome usage monitoring. The US Census
Bureau used it for the 2020 Census -- the first-ever deployment at national
scale, motivated by the Dinur-Nissim database reconstruction theorem, which
proved that releasing too many exact statistics about a dataset inevitably leaks
individual records.

DP works by adding carefully calibrated noise to query results. The noise is
large enough to mask any individual's contribution but small enough to preserve
the statistical utility of the aggregate. The privacy guarantee is parameterized
by epsilon (lower epsilon = more privacy = more noise = less accuracy), and
composability is formalized by a composition theorem: sequential queries consume
a ``privacy budget,'' and once the budget is exhausted, no more queries can be
safely answered. The fog has a finite supply. Use it wisely.

The epsilon parameter deserves a longer look, because it is where the
mathematics of differential privacy meets the politics of data collection. Think
of epsilon as controlling the blur radius on a photograph. At
\(\varepsilon = 0.1\), you are looking at a Monet painting -- the water lilies
are recognizable as water lilies, but individual petals dissolve into
impression. The aggregate is preserved. The particular is lost. At
\(\varepsilon = 1\), you are looking at a photograph taken through frosted glass
-- shapes and proportions are clear, but faces are unreadable. At
\(\varepsilon = 10\), you are looking at a photograph with a slight smudge --
almost everything is visible, and a determined adversary with auxiliary
information might identify individuals. At \(\varepsilon = \infty\), there is no
blur at all. You are looking at the raw data.

The art of differential privacy is choosing the blur. Too much noise (low
epsilon) and the data is useless -- a census that cannot distinguish New York
from Nebraska serves no one. Too little noise (high epsilon) and the privacy
guarantee is hollow -- a medical database that lets researchers reconstruct
individual diagnoses is not private in any meaningful sense. The Dinur-Nissim
theorem makes the stakes precise: for any dataset of \(n\) individuals, if you
answer more than \(O(n)\) counting queries with accuracy better than
\(1/\sqrt{n}\), you can reconstruct the entire dataset. The blur is not
optional. Without it, the data eventually gives up everyone's secrets.

What does this look like in practice? Apple's deployment adds noise locally, on
each device, before data is transmitted -- a technique called local differential
privacy. When your iPhone wants to report which emoji you use most frequently,
it does not send ``thumbs up.'' It sends ``thumbs up'' with probability
\((e^\varepsilon)/(e^\varepsilon + 1)\) and a random emoji with probability
\(1/(e^\varepsilon + 1)\). Any individual report is plausibly random. But
aggregate millions of reports, and the noise cancels out, revealing the
population-level distribution. Apple uses \(\varepsilon = 2\) per day for most
data types and \(\varepsilon = 8\) for some health-related queries. Google's
RAPPOR uses a similar local model with a two-stage randomization that provides
both plausible deniability for individual responses and high accuracy for
aggregate statistics. You never see any of this. Your phone adds the noise
silently, the aggregation server receives randomized data, and the statistical
team extracts population trends from the collective fog.

The composition problem is the silent killer of differential privacy
deployments. Each query against a dataset consumes a portion of the privacy
budget. If you query a medical database once with \(\varepsilon = 1\), you get a
strong privacy guarantee. If you query it twice, the effective epsilon is (at
most) 2 -- weaker, but still meaningful. If you query it a thousand times, each
with \(\varepsilon = 1\), the effective epsilon is (at most) 1000 -- and at that
point, the privacy guarantee is essentially worthless. Advanced composition
theorems (Dwork, Rothblum, and Vadhan, 2010) give tighter bounds: k queries with
epsilon each compose to roughly \(\varepsilon \cdot \sqrt{k}\) rather than
\(\varepsilon \cdot k\). But the fundamental truth remains: privacy budgets are
finite, and every query spends them. A dataset that has been queried ten
thousand times is not the same, from a privacy standpoint, as one that has been
queried ten times. The fog dissipates with each question asked. The Census
Bureau's TopDown Algorithm was designed with this in mind: the total privacy
budget was fixed before any queries were defined, and the noise allocation was
optimized across all geographic levels simultaneously, from national aggregates
down to census blocks. The budget was spent once, carefully, and then the books
were closed.

\section{Three Kinds of Security}\label{three-kinds-of-security}

These four technologies provide fundamentally different \emph{types} of security
guarantees, and conflating them leads to bad architecture decisions. This is the
part where precision matters more than analogy.

\textbf{Information-theoretic security} means the guarantee holds even against
an adversary with unlimited computational power. No mathematical breakthrough,
no quantum computer, no advance in algorithms can break it. MPC with honest
majority (e.g., Shamir secret sharing where more than half the participants are
honest) achieves this. The security follows from information theory, not
computational hardness. The caveats: you need an honest majority, and the
communication cost scales with the number of parties and the complexity of the
computation.

\textbf{Computational security} means the guarantee holds against adversaries
bounded to polynomial-time computation. It rests on assumptions: ``discrete
logarithms are hard,'' ``the Learning With Errors problem is hard,'' ``the
Ring-LWE problem is hard.'' ZKPs and FHE provide computational security. If the
underlying hardness assumption falls -- as discrete logarithm will fall to
Shor's algorithm on a sufficiently large quantum computer -- the security
evaporates retroactively. Every proof ever generated under that assumption
becomes suspect. The lock does not weaken. It ceases to exist.

\textbf{Heuristic security} means the guarantee rests on practical observations
rather than formal proof. Trusted Execution Environments (TEEs) like Intel SGX,
AMD SEV, and ARM CCA provide heuristic security. The hardware manufacturer
attests that the enclave is isolated. But SGX has been broken by Spectre,
Meltdown, Foreshadow, Plundervolt, SGAxe, and AEPIC Leak. AMD SEV has shown
vulnerabilities (SEVered, CipherLeaks). Intel SGX was deprecated on consumer
processors in 2021. The 2025 attacks Battering RAM (\textasciitilde50 euros) and
Wiretap (\textasciitilde\$1,000) demonstrated physical attacks at commodity
prices. TEE security is real in practice against most adversaries, but it lacks
the mathematical foundation of cryptographic privacy and carries an expiration
date set by the next side-channel attack. It is a stage built from plywood
rather than steel: functional, but not what you want for the long run.

The metaphor of plywood is generous. To understand TEEs concretely, picture a
room within a room. Your CPU -- the physical chip on the motherboard -- creates
an isolated memory region called an enclave. Code and data inside the enclave
are encrypted in RAM. The operating system cannot read the enclave's memory. The
hypervisor cannot read it. Even a system administrator with root access and
physical possession of the machine cannot read it. The CPU itself enforces the
boundary: any attempt to access enclave memory from outside the enclave returns
encrypted garbage. Intel SGX (Software Guard Extensions) pioneered this
architecture. ARM TrustZone implements a similar concept at the processor level,
splitting the chip into a ``secure world'' and a ``normal world'' with
hardware-enforced isolation between them.

The promise is strong: you can run sensitive computation on an untrusted
machine, and the machine's owner cannot observe or tamper with it. Cloud
computing without trusting the cloud. The allure is obvious. The history is
cautionary.

Foreshadow (August 2018) broke SGX isolation through a speculative execution
attack. The CPU's branch predictor, trying to execute instructions ahead of time
for performance, would speculatively read enclave memory and leave traces in the
L1 cache. An attacker could measure cache timing to reconstruct the enclave's
secrets. The attack required no physical access -- it could be performed by a
process running alongside the enclave on the same machine. Intel patched it with
microcode updates, but the patch reduced performance and the fundamental
vulnerability -- that speculative execution can leak secrets across security
boundaries -- proved to be architectural, not incidental.

AEPIC Leak (August 2022) was worse. It exploited a bug in Intel's Advanced
Programmable Interrupt Controller to read stale data from the enclave's memory
hierarchy. Unlike Foreshadow, which required careful cache timing, AEPIC Leak
provided architecturally guaranteed data leakage -- the CPU would hand you the
enclave's data directly if you asked the right hardware register. No timing side
channel, no statistical analysis. A clean read.

Downfall (August 2023) exploited the Gather instruction, which loads data from
scattered memory locations into a vector register. The vulnerability allowed an
attacker to read data from other security domains -- including SGX enclaves --
by observing the contents of internal CPU buffers during Gather operations.
Intel's mitigation involved disabling the optimization that made Gather fast,
resulting in up to 50\% performance degradation for workloads that depended on
it.

Intel quietly deprecated SGX on consumer processors (12th generation and later)
beginning in 2021. The feature remains available on server-class Xeon
processors, where it is marketed for cloud confidential computing. But the
deprecation on consumer chips tells a story: Intel concluded that the attack
surface was too large and the performance cost of mitigations too high for a
feature intended to run on every laptop. The room within a room is still
available -- but only in the data center, where the threat model is different
and the economic calculus favors the convenience of hardware isolation despite
its known fragility.

For the system architect, this taxonomy matters because it determines \emph{what
you are actually trusting}. If your system uses MPC with honest majority for the
core computation and ZKPs for the verifiable output, you have
information-theoretic privacy for the computation and computational privacy for
the proof. If you then run the whole thing inside a TEE, the TEE adds
performance (fast) and convenience (no complex protocol choreography) but does
not strengthen the privacy beyond what the cryptography already provides -- and
may weaken it if the TEE is compromised.

The magician's guarantee depends on which lock protects the trick.
Information-theoretic security is a lock that cannot be picked, even with
infinite time. Computational security is a lock that cannot be picked in
practice -- but a quantum locksmith might change the calculus. Heuristic
security is a lock that has never been picked, without proof that it cannot be.

\section{Composability: When One PET Is Not
Enough}\label{composability-when-one-pet-is-not-enough}

The real power of PETs emerges when they are composed. No single instrument
plays the whole symphony. Consider a realistic healthcare scenario -- and notice
how each PET enters at the moment its particular strength is needed:

\begin{enumerate}
\def\labelenumi{\arabic{enumi}.}
\item
  \textbf{Step 1 (MPC)}: Five hospitals jointly compute aggregate statistics on
  a rare disease using their combined patient records. Each hospital contributes
  its data to an MPC protocol. No hospital sees any other hospital's records.
  The output is aggregate statistics: prevalence rates, treatment outcomes,
  demographic distributions.
\item
  \textbf{Step 2 (Differential Privacy)}: Before the aggregate statistics leave
  the MPC computation, differential privacy noise is added. This ensures that
  even the aggregate output cannot be used to infer individual patient records.
  The privacy budget (epsilon) is tracked across queries.
\item
  \textbf{Step 3 (FHE)}: The differentially private aggregate statistics are
  encrypted under FHE. An AI firm trains a predictive model on the encrypted
  data. The AI firm never sees the plaintext statistics. The hospitals never see
  the AI firm's model architecture (which may be proprietary). The glovebox,
  again.
\item
  \textbf{Step 4 (ZKP)}: The AI firm produces a zero-knowledge proof that the
  trained model meets accuracy and fairness criteria specified in a regulatory
  standard, without revealing the model's weights or the training data. A
  regulator verifies the proof and certifies the model for clinical use. The
  magician performs. The audience -- in this case, the regulator -- verifies.
\end{enumerate}

Each step uses the PET best suited to its specific trust problem: MPC for
multi-party data aggregation, DP for statistical disclosure control, FHE for
outsourced computation on sensitive data, and ZKP for verifiable compliance
without disclosure.

But the transitions between steps are not trivial. The MPC-to-FHE handoff
requires either the hospitals to encrypt the MPC output under the AI firm's FHE
public key (which means they see the plaintext), or a protocol that converts MPC
secret shares directly to FHE ciphertexts (an active research frontier). The
FHE-to-ZKP handoff requires verifiable FHE -- proving in zero knowledge that an
FHE computation was performed correctly -- which is emerging but not yet
production-ready.

The composability lesson: PETs compose in theory. In practice, each composition
point requires protocol engineering that is often harder than the individual PET
deployments. The system architect must understand not just what each PET does,
but how they hand off to each other. The orchestra sounds beautiful when
everyone enters on cue. Getting the cues right is the hard part.

Three composition patterns deserve particular attention, because they recur
across domains and will likely define the privacy architecture of the next
decade.

\textbf{ZKP + MPC: Verified Inputs to Joint Computation.} The healthcare
scenario above assumes each hospital contributes honest data to the MPC
protocol. But what if a hospital submits fabricated records -- inflating its
patient count to increase its share of research funding, or omitting records to
conceal a malpractice pattern? MPC computes correctly on whatever inputs it
receives. It does not verify that the inputs are truthful. This is where ZKPs
enter: each participant produces a zero-knowledge proof that its MPC input
satisfies agreed-upon constraints -- the records come from a certified database,
the patient count matches a signed attestation from the hospital's electronic
health record system, the data format conforms to the protocol specification.
The MPC protocol verifies these proofs before accepting the inputs. The joint
computation proceeds on data that is both private \emph{and} certified. No party
reveals its data. Every party proves its data is legitimate. The magician does
not merely perform behind a curtain -- she presents her credentials before
stepping onto the stage.

This pattern -- ZKP-verified inputs to MPC -- appears in private auctions (prove
your bid is backed by sufficient funds without revealing the bid amount), in
private voting (prove you are an eligible voter without revealing your
identity), and in collaborative machine learning (prove your training data meets
quality thresholds without revealing the data itself). In each case, MPC
provides the privacy during computation, and ZKPs provide the integrity of the
inputs. The two PETs are not redundant. They address orthogonal trust problems.
Privacy without integrity is a system that computes correctly on lies. Integrity
without privacy is a system that reveals everything it verifies.

\textbf{ZKP + FHE: Verifiable Encrypted Computation.} This is the zkFHE frontier
mentioned earlier, and it deserves a structural explanation. The problem: a
cloud provider performs FHE computation on your encrypted data and returns an
encrypted result. You decrypt and get an answer. But did the provider actually
compute the function you requested? Or did it compute a cheaper approximation,
or a different function entirely, or simply return a random ciphertext? FHE
guarantees confidentiality -- the provider cannot see your data. It does not
guarantee integrity -- the provider can lie about what it computed. ZKPs close
this gap. The provider produces a zero-knowledge proof that the sequence of
homomorphic operations it performed on the ciphertext corresponds exactly to the
function specification. The proof is verified against the input ciphertext, the
output ciphertext, and the function description. If it checks out, you know the
computation was honest. If it does not, you know to reject the result and find
another provider.

The difficulty is that proving FHE computations in zero knowledge is very
expensive. Each homomorphic operation involves polynomial arithmetic over large
rings, and the ZKP circuit must encode all of this arithmetic faithfully.
Current zkFHE prototypes achieve verification for small circuits -- a few
hundred multiplication gates -- and the proving overhead adds another order of
magnitude atop FHE's already steep costs. But the research trajectory is clear,
and the incentive is enormous: anyone who wants to outsource computation on
sensitive data to an untrusted cloud needs both confidentiality (FHE) and
integrity (ZKP). Neither alone is sufficient.

\textbf{The Privacy Stack.} The key insight about PET composition is
architectural: do not think of PETs as individual tools to be selected. Think of
them as layers in a protocol stack, analogous to the network stack that
separates TCP from IP from Ethernet. At the bottom, differential privacy
provides statistical-level guarantees for aggregate data releases -- the
coarsest and cheapest form of privacy, suitable for telemetry and census-scale
statistics. Above it, MPC provides computation-level privacy for multi-party
protocols -- stronger than DP (it protects individual inputs, not just
statistical aggregates), but more expensive and limited to specific interaction
patterns. Above that, FHE provides data-level privacy for outsourced computation
-- stronger still (the computing party learns nothing at all), but with the
highest performance cost. And at the top, ZKPs provide verification-level
privacy -- the ability to prove properties of private data or private
computation without revealing the underlying secrets.

Each layer addresses a different threat. Each has a different cost. And like
network layers, they compose vertically: a system might use DP for its
public-facing analytics dashboard, MPC for its inter-institutional data sharing,
FHE for its cloud-based model training, and ZKPs for its compliance proofs --
all within the same architecture, each operating at its appropriate level of the
stack. The system architect who treats PET selection as a single choice (``we
will use ZKPs'') is making the same mistake as the network engineer who treats
protocol selection as a single choice (``we will use TCP''). The answer is
almost always a stack, not a single layer.

\section{Real-World Deployments: Five Case
Studies}\label{real-world-deployments-five-case-studies}

\subsection{1. Decentriq and the Swiss National
Bank}\label{decentriq-and-the-swiss-national-bank}

In a federal pilot project beginning in 2021-2022, data clean room technology
enabled encrypted collaboration between the Swiss National Bank, SIX
(Switzerland's financial market infrastructure provider), and Zurich Cantonal
Bank. The goal was cybersecurity threat detection: analyzing patterns of
suspicious financial activity across institutions without any institution
revealing its transaction data to the others.

The architecture used MPC-style computation within Decentriq's confidential
computing platform, combining software-level privacy guarantees with hardware
TEEs. The result demonstrated that financial regulators can gain systemic risk
visibility without requiring banks to share raw transaction data -- a
significant precedent for privacy-preserving financial regulation. The regulator
sees the pattern. The banks keep the data. Everyone sleeps better.

\subsection{2. DTCC and the Canton Network}\label{dtcc-and-the-canton-network}

The Depository Trust and Clearing Corporation (DTCC), which processes virtually
all US securities transactions, partnered with Digital Asset's Canton Network in
December 2025 to tokenize US Treasuries on a permissioned blockchain with
privacy-preserving settlement. The architecture uses the Canton protocol's
built-in privacy model, where participants see only the portions of the ledger
relevant to them.

This deployment matters because of who is adopting. DTCC is not a startup
experimenting with privacy. It is the backbone of US securities infrastructure,
processing trillions of dollars annually. When DTCC chooses a privacy-preserving
architecture, it signals that privacy is not a nice-to-have feature but a
regulatory and competitive necessity. The largest financial plumbing system in
the world has decided it needs these tools. Pay attention.

\subsection{3. Partisia and Toppan Edge: Digital Student
IDs}\label{partisia-and-toppan-edge-digital-student-ids}

Toppan Edge and Partisia announced joint development of privacy-preserving
digital student IDs in 2025, with a proof-of-concept conducted at the Okinawa
Institute of Science and Technology from June to September 2025. The system
combines facial recognition for identity verification, decentralized identifiers
(DIDs) for credential management, smartphone NFC for physical access, and MPC
via Partisia's blockchain for privacy-preserving identity verification.

The key innovation: the student's biometric data is never stored in a single
location or revealed to a single party. MPC ensures that identity verification
can be performed without any single server holding the student's facial
template. The platform is targeted for students enrolling from April 2026. Your
face opens the door, but no one holds a copy of your face.

\subsection{4. Privacy Pools: Pragmatic On-Chain
Privacy}\label{privacy-pools-pragmatic-on-chain-privacy}

Privacy Pools, co-authored by Vitalik Buterin and implemented by 0xbow, launched
on Ethereum mainnet on April 1, 2025. Buterin was one of the first users,
depositing 1 ETH.

The design addresses the fundamental tension between on-chain privacy and
regulatory compliance -- a tension that destroyed Tornado Cash and haunts every
privacy protocol. Users deposit funds into a pool and can later withdraw them,
breaking the link between deposit and withdrawal addresses (similar to Tornado
Cash). But Privacy Pools add a compliance layer: an Association Set Provider
(ASP) screens deposits for connections to sanctioned or illicit addresses, and
the zero-knowledge proof used for withdrawal includes a proof that the user's
funds are drawn from a compliant ``association set.''

The result is ``pragmatic privacy'' -- transaction privacy for legitimate users,
with a built-in compliance mechanism that prevents sanctioned funds from mixing
with clean funds. As of early 2026, Privacy Pools has processed over \$6 million
in volume across more than 1,500 users. The broader ecosystem includes more than
35 teams pursuing approximately 13 distinct approaches to private transfers on
Ethereum. The magician proves she is not cheating -- and the regulator is
satisfied.

\subsection{5. Apple, Google, and the US Census Bureau: Differential Privacy at
Scale}\label{apple-google-and-the-us-census-bureau-differential-privacy-at-scale}

The largest PET deployments in the world are not blockchain systems. They are
not even close. They are differential privacy systems serving billions of users:

\begin{itemize}
\tightlist
\item
  \textbf{Apple} introduced DP in iOS 10 (2016) for emoji usage statistics,
  Safari search queries, HealthKit data, and keyboard autocorrect improvements.
  Each device adds local noise before transmitting data, with
  \(\varepsilon = 2\) per day for most data types.
\item
  \textbf{Google} deployed RAPPOR for Chrome settings monitoring, adding
  randomized responses to usage data before aggregation.
\item
  \textbf{US Census Bureau} used the TopDown Algorithm for the 2020 Census,
  adding calibrated noise to census statistics at every geographic level. The
  decision was motivated by a concrete threat: the Dinur-Nissim reconstruction
  theorem proved that releasing too many exact statistics from a dataset
  eventually allows full reconstruction of individual records.
\end{itemize}

These deployments demonstrate that differential privacy is the only PET to have
achieved planetary-scale adoption. ZKPs, MPC, and FHE remain orders of magnitude
smaller in deployment footprint. For the system architect, this suggests that DP
should be the first tool considered for statistical data release, with the other
PETs reserved for use cases that require computation on raw data or verifiable
individual claims. The fog machine is the most popular tool in the privacy
toolkit. The magic wand is catching up.

\section{Privacy Architectures for Smart Contracts: Kachina and
Zexe}\label{privacy-architectures-for-smart-contracts-kachina-and-zexe}

Two academic systems -- Kachina and Zexe -- represent the theoretical
foundations for how private smart contracts can be deployed on blockchains. They
take complementary approaches, and understanding both illuminates the design
space that every privacy-focused blockchain must work within.

\subsection{Kachina: Privacy as a
Parameter}\label{kachina-privacy-as-a-parameter}

Kachina, developed by Kerber, Kiayias, and Kohlweiss at the University of
Edinburgh and IOHK, provides a UC-secure (universally composable) framework for
privacy-preserving smart contracts. Its key innovation is treating privacy as a
\emph{parameter}, not a binary choice. Think of it as a dimmer switch rather
than an on/off toggle.

Contract state is split into \emph{shared public} state (on-chain) and
\emph{individual private} state (per party, off-chain). Users prove in zero
knowledge that their state transitions are valid given some private state and
input. The critical architectural insight is the \emph{state oracle transcript}:
instead of proving full state transitions (which would require locking shared
state), users capture oracle queries and responses as partial transcripts. These
transcripts are partial functions over state, enabling concurrent transactions
to succeed even when state changes between proof creation and proof submission.

The privacy leakage is formally captured by a \emph{leakage function} Lambda
that specifies exactly what information each transaction reveals. Lambda can be
tuned from ``full leakage'' (equivalent to Ethereum, where everything is
visible) to ``near-zero leakage'' (equivalent to Zerocash, where only nullifiers
and commitments are visible). This parameterization means the same framework can
model both transparent and private contracts, and everything in between. The
magician decides, contract by contract, how much of the trick to reveal.

Proving complexity is \(O(|T_\rho| + |T_\sigma|)\) -- proportional to the
\emph{transcript lengths}, not the full state size. This matters for
scalability: a contract with millions of state entries can support private
transactions that only touch a few entries, and the proving cost reflects only
the entries accessed.

\subsection{Zexe: Function Privacy}\label{zexe-function-privacy}

Zexe, developed by Bowe, Chiesa, Green, Miers, Mishra, and Wu (2018), takes a
UTXO-based approach with a stronger privacy guarantee: not only are the
transaction data hidden, but the \emph{function being computed} is hidden as
well. An observer cannot distinguish a token transfer from a governance vote
from a swap -- all transactions look identical on-chain. The audience sees
identical envelopes. Every envelope looks the same. The contents are unknowable.

The architecture uses a ``records nano-kernel'' (RNK) -- a minimalist shared
execution environment where records have birth and death predicates.
Transactions consume old records and create new ones by satisfying these
predicates in zero knowledge. The on-chain footprint is constant: 968 bytes for
a 2-input/2-output transaction, regardless of the complexity of the off-chain
computation.

Zexe uses recursive proof composition (bounded depth 2, not full recursion) with
a BLS-12 curve for inner SNARKs and a Cocks-Pinch curve for outer composition.
Proof generation takes roughly one minute plus computation-dependent time.
Verification takes tens of milliseconds.

\subsection{Comparison}\label{comparison}

\begin{longtable}[]{@{}
  >{\raggedright\arraybackslash}p{(\linewidth - 4\tabcolsep) * \real{0.4000}}
  >{\raggedright\arraybackslash}p{(\linewidth - 4\tabcolsep) * \real{0.3600}}
  >{\raggedright\arraybackslash}p{(\linewidth - 4\tabcolsep) * \real{0.2400}}@{}}
\toprule\noalign{}
\begin{minipage}[b]{\linewidth}\raggedright
Property
\end{minipage} & \begin{minipage}[b]{\linewidth}\raggedright
Kachina
\end{minipage} & \begin{minipage}[b]{\linewidth}\raggedright
Zexe
\end{minipage} \\
\midrule\noalign{}
\arrayrulecolor{MidnightBlue}\midrule\arrayrulecolor{MidnightBorder}
\endhead
\bottomrule\noalign{}
\endlastfoot
State model & Account-based (state machine) & UTXO-based (records) \\
Data privacy & Parameterizable (Lambda function) & Full \\
Function privacy & No (function identity visible) & Yes (all transactions
indistinguishable) \\
Concurrency & State oracle transcripts & UTXO model (naturally concurrent) \\
On-chain cost & \(O(\text{transcript length})\) & Constant (968 bytes) \\
Proving cost & \(O(\text{transcript length})\) & \textasciitilde1 minute +
computation \\
Security model & UC-secure (\(\mathcal{F}_\text{nizk}\),
\(\mathcal{G}_\text{ledger}\) hybrid) & Simulation-based \\
\end{longtable}

Midnight's architecture follows the Kachina model most closely --
parameterizable disclosure via \texttt{disclose()}, account-based state, and
compiler-enforced privacy boundaries. Aztec's design follows the Zexe model more
closely -- UTXO-based notes, client-side proving, and a Private Execution
Environment (PXE) that handles proof generation.

For the system architect choosing between these approaches, the key question is:
do you need function privacy? If yes (all transactions must be
indistinguishable), the Zexe/UTXO model is the natural choice. If no (you can
tolerate revealing which function was called, as long as the arguments are
private), the Kachina/account model offers simpler programming and easier state
management.

\section{The Regulatory Intersection}\label{the-regulatory-intersection}

Privacy-enhancing technologies do not operate in a regulatory vacuum. For the
system architect building in 2026, two regulatory developments demand attention
-- and both, in different ways, are pulling the same direction as the
technology.

\subsection{GDPR and the Blockchain Immutability
Paradox}\label{gdpr-and-the-blockchain-immutability-paradox}

The European Data Protection Board (EDPB) adopted Guidelines 02/2025 during its
April 2025 plenary, providing the most authoritative guidance to date on GDPR
compliance for blockchain systems. The guidelines address a fundamental tension:
blockchain's immutability directly conflicts with Article 17 of GDPR -- the
right to erasure. If personal data is stored on-chain, it cannot be deleted. The
regulation says it must be deletable. Something has to give.

The EDPB's answer is unambiguous: do not store personal data on-chain. The
recommended architecture is ``off-chain storage and hashing'' -- store
personally identifiable information (PII) in a mutable off-chain database, and
store only a cryptographic hash on-chain. If a data subject exercises their
right to erasure, delete the off-chain data and the cryptographic keys linking
it to the hash. The on-chain hash becomes an orphaned, meaningless string of
characters. The commitment remains. The secret it referenced is gone.

ZKPs play a natural role in this architecture. Instead of storing ``Alice is 25
years old and lives in Berlin'' on-chain, store a hash of Alice's credential and
allow Alice to generate ZK proofs about properties of that credential: ``I am
over 18'' (for age-gated access), ``I am a resident of the EU'' (for
jurisdictional compliance), ``I am not on a sanctions list'' (for regulatory
compliance). The on-chain system never learns Alice's age, address, or identity.
It learns only the truth of the specific claims she chooses to prove. The
magician reveals exactly what the audience needs to see. No more.

This pattern is called zKYC (zero-knowledge Know Your Customer), and it is
rapidly gaining traction. Galactica Network, zyphe, and hyli implement zKYC
systems that enable selective disclosure for regulatory compliance. The promise:
compliance without surveillance.

The tension between GDPR's right to erasure and blockchain's immutability is
worth dwelling on, because it illustrates a deeper architectural principle. The
naive response is to declare that blockchains and GDPR are incompatible -- that
you cannot have an append-only ledger and a right to delete. But the
off-chain-storage-with-on-chain-hash pattern resolves the tension elegantly, and
the resolution is instructive. The hash on-chain is not personal data. It is a
commitment -- a mathematical fingerprint that proves a piece of data existed at
a particular time, without revealing what the data was. Delete the off-chain
data, destroy the linking keys, and the hash is cryptographically orphaned. It
sits on the blockchain forever, a meaningless 32-byte string, pointing to
nothing. The right to erasure is satisfied not by deleting the blockchain entry
but by severing the link between the entry and the person it once referenced.
The commitment survives. The secret is gone. The regulation is satisfied. This
is not a workaround. It is good architecture -- the kind of architecture that
PETs make possible.

The zKYC pattern makes this concrete. Consider a bar that needs to verify a
customer is over 21. Today, the customer shows a driver's license, and the
bartender sees the customer's full name, date of birth, address, driver's
license number, organ donor status, and photograph. The bartender needs exactly
one bit of information: is this person at least 21? Instead, the customer
reveals a dozen pieces of personally identifiable information to a stranger.
With zKYC, the customer holds a verifiable credential in a digital wallet --
issued and signed by the government -- and generates a zero-knowledge proof:
``the date of birth in my credential, when compared to today's date, yields an
age of at least 21.'' The bartender's verification terminal checks the proof and
the government's signature. It learns exactly one fact: the customer is old
enough. The name, the address, the license number, the photograph -- none of it
crosses the bar. The trick reveals only what the audience needs to see.

\subsection{eIDAS 2.0 and the European Digital Identity
Wallet}\label{eidas-2.0-and-the-european-digital-identity-wallet}

The European Union's eIDAS 2.0 regulation, effective from 2024, mandates that
all EU member states offer citizens a European Digital Identity Wallet by 2026.
The wallet must support verifiable credentials and selective disclosure --
proving specific attributes (citizenship, age, professional qualifications)
without revealing the entire identity document.

ZKPs are a natural technical foundation for selective disclosure in identity
wallets. The Architecture and Reference Framework (ARF) for eIDAS 2.0 envisions
a credential ecosystem where issuers (governments, universities, professional
bodies) issue cryptographically signed credentials, holders store them in their
wallets, and verifiers check proofs of specific attributes without seeing the
full credential.

The scale of this mandate is easy to understate. By late 2026, every EU member
state must provide a digital identity wallet to every citizen who requests one.
That is a potential user base of 450 million people. The wallet must
interoperate across borders -- a Spanish wallet must be accepted by a German
verifier, a French credential must be verifiable in Italy. And the wallet must
support selective disclosure by design, not as an afterthought. When a Belgian
student presents her wallet to a Portuguese university, the university should be
able to verify her degree and her citizenship without learning her tax ID, her
medical history, or her home address. The credential is a bundle of attributes.
The wallet discloses only the attributes the verifier needs. The rest stays
sealed.

This is not a theoretical design. The EU's Large Scale Pilots (LSPs) --
POTENTIAL, EU Digital Identity Wallet Consortium, NOBID, and DC4EU -- have been
testing these architectures since 2023, with real users, real credentials, and
real cross-border verification. The technical challenge is not the ZKP itself
(the cryptography is well understood) but the credential format, the revocation
mechanism, the issuer trust framework, and the user interface that makes
selective disclosure comprehensible to a non-technical user. The magician's
trick is elegant. The stage production -- lighting, sound, audience management
-- is where the engineering budget goes.

The regulatory pull is significant: eIDAS 2.0 creates a legal mandate for the
privacy properties that ZKPs can provide. For the first time, a major regulatory
framework is not just permitting but \emph{requiring} selective disclosure. This
transforms ZKPs from a voluntary privacy choice to a compliance necessity for
any service that needs to verify European identities. The law now demands the
trick.

\subsection{Non-Compliance Is Expensive}\label{non-compliance-is-expensive}

GDPR violations can result in fines of up to 4\% of global annual revenue or 20
million euros, whichever is greater. For a technology company with \$10 billion
in revenue, a GDPR violation could cost \$400 million. This makes privacy
architecture decisions directly material to business risk.

The implication for system architects: the choice of PET is not merely a
technical decision. It is a risk management decision with quantifiable financial
exposure. An architecture that stores personal data on a public blockchain is
not just a privacy risk -- it is a potential nine-figure liability.

\section{The Decision Matrix}\label{the-decision-matrix}

For the system architect who needs to choose a PET (or a combination of PETs),
the decision depends on four questions:

\textbf{1. What is the trust model?} - If you need privacy against
computationally unbounded adversaries (including future quantum computers): MPC
with honest majority provides information-theoretic security. - If you trust
computational hardness assumptions: ZKPs and FHE provide computational security.
- If you trust hardware manufacturers: TEEs provide heuristic security with high
performance.

\textbf{2. Who has the data?} - If the data holder needs to prove a property:
ZKP (selective disclosure). - If multiple parties need to compute on their
combined data: MPC (collaborative computation). - If the data needs to be
processed by an untrusted third party: FHE (encrypted outsourcing). - If
aggregate statistics need to be released from a dataset: DP (statistical
disclosure control).

\textbf{3. What performance is acceptable?} - ZKP proof generation: seconds to
minutes. Verification: milliseconds. - MPC: communication rounds proportional to
circuit depth. Latency is the bottleneck. - FHE: 10,000-1,000,000x overhead over
plaintext. Improving rapidly, but still orders of magnitude slower. - DP:
negligible overhead (adding noise to query results is cheap). - TEE: near-native
performance (\textless5\% overhead for many workloads).

\textbf{4. What regulatory regime applies?} - GDPR/eIDAS 2.0: Selective
disclosure (ZKP), off-chain storage with on-chain hashing, right-to-erasure
compatibility. - Financial regulation (AML/KYC): zKYC (ZKP for compliance
proofs), Privacy Pools (ZKP for provenance), MPC for inter-institutional
analysis. - Healthcare (HIPAA, EU Clinical Trials Regulation): MPC for
multi-site computation, DP for statistical releases, FHE for outsourced AI model
training.

No single PET answers all four questions. The art is in composition -- and the
engineering is in the handoffs between them.

\section{Open Problems}\label{open-problems}

Three capabilities sit at the frontier of PET research and will likely shape the
next generation of privacy architectures:

\textbf{Verifiable FHE}: Proving in zero knowledge that an FHE computation was
performed correctly. This closes the loop in the healthcare scenario: the AI
firm not only computes on encrypted data but also proves that it computed
\emph{correctly} on the encrypted data. The surgeon not only operates through
the glovebox -- she provides a certificate that the operation was performed to
specification. The zkFHE project and SherLOCKED prototype (using RISC Zero's
Bonsai zkVM) are early implementations.

\textbf{Collaborative/threshold proving}: Distributing ZK proof generation
across multiple servers using MPC, so that no single server sees the full
witness. The work by Ozdemir and Boneh (USENIX Security 2022) and subsequent
improvements in 2024 demonstrate that proof generation itself can be
privacy-preserving. This creates a fifth proving model -- between client-side
(private but expensive) and delegated (cheap but witness-exposing) -- that
combines the privacy of the former with the performance of the latter. The
magician's backstage preparation is distributed across multiple locked rooms. No
single stagehand sees the whole act.

\textbf{Private Information Retrieval (PIR)}: Querying a database without
revealing which record you are accessing. A client-side prover in a private
rollup (like Aztec) needs to retrieve encrypted notes from the network without
revealing which notes belong to them. Recent advances at EUROCRYPT 2026 achieved
information-theoretic PIR with sublinear server time and quasilinear space,
moving PIR from theoretical curiosity toward practical deployment for
billion-entry databases.

\section{The Incomplete Stack}\label{the-incomplete-stack}

A thread runs through this chapter that is worth stating plainly.

Privacy is not a feature that you add to a system after it is built. It is a
property of the architecture, present from the first design document or absent
forever. You cannot retrofit privacy onto a transparent blockchain any more than
you can retrofit soundproofing onto a glass house. Or, to stay with our
metaphor: you cannot add trapdoors to a stage after the audience is already
seated.

The four PETs -- ZKPs, MPC, FHE, and DP -- are not competing technologies. They
are complementary tools in a single toolkit. The magician's wand, the glovebox,
the fog machine, and the locked vault where multiple parties contribute secrets
they never share. Each excels at a different trust problem. Each fails at
problems the others solve well. The system architect who understands all four,
and who understands how they compose, has a genuine advantage over one who knows
only ZKPs and treats every privacy problem as a nail to be hit with the
zero-knowledge hammer.

The regulatory environment is, for the first time, pulling in the same direction
as the technology. GDPR, eIDAS 2.0, and the global trend toward data sovereignty
create legal mandates for exactly the capabilities that PETs provide. The
question is no longer ``should we use privacy-enhancing technologies?'' but
``which ones, in what combination, and how do we prove to regulators that they
work?''

That last part -- proving to regulators that the privacy technology works -- is,
fittingly, itself a zero-knowledge problem. And we have the tools to solve it.

\partdivider{Part III: Synthesis and the Road Ahead}

\emph{The trick has been performed seven times, each time revealing a deeper
layer of the mechanism. Now we step back from the stage. What does the whole
show look like from the back of the theater? Where is the art heading? And does
the magic hold up outside the theater, in the harsh light of commerce,
regulation, and the passage of time?}

\chapter{Chapter 10: The Synthesis -- Three Paths, Not
Two}\label{chapter-10-the-synthesis-three-paths-not-two}

\section{The Binary That Broke}\label{the-binary-that-broke}

For three years, the zero-knowledge community organized its world along a single
axis: SNARK or STARK. Trusted setup or transparent. Small proofs or big ones.
Algebraic elegance or hash-based brute force. Every conference talk, every
investor pitch, every architectural decision began with this fork in the road.

That binary is dead.

It did not die because one side won. It died because the winning move turned out
to be using both sides simultaneously -- and then a third path appeared that
neither side had anticipated. This is the story of how a two-path map became a
three-path map, and why the seven-layer model from earlier in this book must
bend to accommodate what actually happened.

\section{The Map Redrawn}\label{the-map-redrawn}

Chapter 1 presented seven layers as a stack -- neat, linear, one resting on the
next. Nine chapters of evidence say otherwise. The seven layers are a directed
acyclic graph with fourteen causal edges, and we owe you the honest picture
before proceeding.

\begin{verbatim}
                    ┌──────────────────────────────────┐
                    │  THE SEVEN-LAYER CAUSAL WEB       │
                    │  (14 directed design-time edges)  │
                    └──────────────────────────────────┘

        ┌─────────┐         ┌─────────┐         ┌─────────┐
        │ Layer 1 │◄────────│ Layer 2 │◄────────│ Layer 3 │
        │  Setup  │         │Language │         │ Witness │
        └────┬────┘         └────┬────┘         └────┬────┘
             │                   │ ▲                  │
             │                   │ │                  │
             ▼                   ▼ │                  ▼
        ┌─────────┐         ┌─────────┐         ┌─────────┐
        │         │         │ Layer 4 │◄────────│         │
        │         │         │  Arith  │─ ─ ─ ─ ►│         │
        │         │         └────┬────┘         │         │
        │ Layer 5 │◄─────────────┘              │ Layer 7 │
        │  Proof  │◄─────────────┐              │ Verdict │
        │ System  │         ┌────┴────┐         │         │
        │         │◄────────│ Layer 6 │◄────────│         │
        │         │────────►│Primitive│────────►│         │
        └─────────┘         └─────────┘         └─────────┘

  ── ► Upward chain (6 edges): 6→5→4→3→2→1 (design flows up)
  ◄ ── Downward pressure (8 edges): 7→6, 7→5, 4→2, 3≡4,
       6→7, 1→5, 5→7, 2→3
\end{verbatim}

\textbf{Reading the diagram.} The upward chain (six edges) is what the book has
followed: field choice (Layer 6) determines commitment scheme (Layer 5), which
shapes arithmetization (Layer 4), which constrains witness layout (Layer 3),
which influences ISA design (Layer 2), which determines setup scope (Layer 1).
If this were the whole story, the stack metaphor would suffice.

It is not the whole story. Eight downward and cross-cutting edges complicate the
picture. Ethereum gas economics (Layer 7) force BN254 pairings (Layer 6) and
STARK-to-SNARK wrapping (Layer 5). Cairo's constraint system (Layer 4) dictated
its ISA design (Layer 2). Jolt fused witness generation (Layer 3) into
arithmetization (Layer 4). The graph has no cycles -- design-time constraints
are asymmetric -- but it has width: multiple independent paths connect the same
pair of nodes, and a single parameter change propagates through the web along
all of them simultaneously.

The three paths below are routes through this DAG, not floors in a building.
Each path makes different choices at each node, and the edges explain why those
choices are coupled.

\section{Path One: The Hybrid STARK-to-SNARK
Pipeline}\label{path-one-the-hybrid-stark-to-snark-pipeline}

The dominant production pattern in 2026 is not SNARK. It is not STARK. It is
STARK \emph{wrapped in} SNARK.

Here is how the trick works. A prover generates a STARK proof -- transparent, no
trusted setup, post-quantum by construction, but large (50-200 KB) and expensive
to verify on-chain (1-5 million gas on Ethereum). This STARK proof is then
compressed through recursive aggregation and finally wrapped in a Groth16 SNARK
proof -- tiny (192 bytes), cheap to verify on-chain (\textasciitilde250-300K
gas), but requiring a trusted setup ceremony (the Powers-of-Tau).

Think of it as a two-act show. In the first act, the magician performs behind a
glass wall -- everything is transparent, no trapdoors, no hidden compartments.
In the second act, the performance is photographed, and the photograph is sealed
in a tamper-proof envelope small enough to slip under a door. The audience in
the theater (the blockchain) only sees the envelope. But anyone who opens it can
verify that it faithfully captures the transparent performance.

Every major production system does this or plans to. SP1 Hypercube generates
multilinear STARK proofs over the BabyBear field, recursively compresses them,
then wraps the result in Groth16 over BN254 for Ethereum verification. Stwo
generates Circle STARK proofs over the Mersenne-31 field, aggregates them via
SHARP, then wraps to Groth16 via Herodotus for Ethereum settlement. RISC Zero,
Airbender, ZisK, Pico Prism -- all follow the same pattern. Even StarkWare, the
company that built its identity on transparent proving, wraps to Groth16 for
Ethereum L1 because the gas economics demand it.

The wrapping pipeline is not a Layer 5 phenomenon. It pierces three layers
simultaneously: Layer 5 (the proof system switches from STARK to SNARK), Layer 6
(the field transitions from BabyBear or M31 to BN254), and Layer 7 (the
verification target shifts from prover-internal consistency to EVM smart
contract). This vertical shaft through the stack is invisible in the original
framing, which treats STARKs and SNARKs as competing alternatives occupying the
same layer.

Why did this happen? Economics. Pure and simple. A raw STARK verification costs
1-5 million gas on Ethereum -- roughly \$5-\$25 at typical gas prices. A Groth16
verification costs \textasciitilde250K gas -- roughly \$0.50-\$1.00. When you
amortize this across thousands of transactions per batch, the difference is
enormous. The inner STARK gives you transparency and post-quantum readiness. The
outer SNARK gives you on-chain affordability. The combination gives you both.

The hybrid path has a structural weakness that deserves a name: the outer
Groth16 wrapper is not post-quantum. Even though the inner STARK is
quantum-resistant, the final on-chain verification depends on BN254 pairings,
which Shor's algorithm breaks in polynomial time. The chain of trust is only as
strong as its weakest link. Today's hybrid systems are transparent and fast on
the inside, but they inherit quantum vulnerability from their verification
wrapper. The glass wall is strong. The envelope has an expiration date.

\section{Path Two: Pure Transparent}\label{path-two-pure-transparent}

The Ethereum Foundation has staked a different position. For the L1 zkEVM -- the
project to prove every Ethereum block in zero knowledge -- the mandate is
explicit: no trusted setup, period. This is not a preference; it is a
requirement. The reasoning is straightforward: Ethereum's base layer cannot
depend on a trusted ceremony that quantum computers will eventually break. The
stage must be made of glass, all the way down.

This mandate forces a different engineering path. Pure transparent systems use
only hash-based commitments (FRI, Merkle trees) and avoid all pairing-based
cryptography. The cost is larger proofs and more expensive on-chain
verification. The benefit is that no ceremony is ever needed, no toxic waste
exists, and the security assumptions survive quantum computing (assuming hash
functions with sufficient output size).

The EF's December 2025 pivot is instructive. Having declared the speed race
``effectively won'' -- four teams proved \textgreater99\% of Ethereum blocks
within the 12-second slot time -- the Foundation shifted its 2026 targets to
security: 100-bit provable security by May 2026, 128-bit by December 2026. The
new primary metric is energy consumption per proof (kWh), not raw speed. The
Foundation explicitly rejects reliance on unproven mathematical conjectures
(proximity gap assumptions) for production soundness.

This matters because it validates a specific engineering philosophy: formal
provable security over empirical performance. SP1 Hypercube already eliminates
proximity gap conjectures in its multilinear STARK. The EF's security-first
stance means that systems with strong formal guarantees will be preferred for
the most security-critical application in the ecosystem -- proving the base
layer itself.

Brevis's Pico Prism occupies an interesting position on this path. Built on
Plonky3 with hash-based FRI commitments, it achieves 99\%+ real-time proving on
16 RTX 5090 GPUs while remaining fully transparent. For non-Ethereum-L1
applications that still need on-chain verification, it wraps to Groth16 -- but
the inner pipeline is transparent end to end.

\section{Path Three: Post-Quantum
Folding}\label{path-three-post-quantum-folding}

The third path is the newest and least traveled. It abandons both pairing-based
cryptography and hash-based commitments in favor of lattice-based constructions
that provide additive homomorphism (enabling folding), post-quantum security,
and increasingly competitive proof sizes.

The key systems are LatticeFold (Boneh \& Chen, ASIACRYPT 2025), LatticeFold+
(CRYPTO 2025), Neo/SuperNeo (Nguyen \& Setty, ePrint 2025/294), and Symphony
(Chen, ePrint 2025/1905). These systems use Module-SIS commitments --
lattice-based structures whose hardness NIST has validated through the FIPS
203/204 standardization process -- to achieve folding over polynomial rings
without relying on discrete logarithms or pairings.

The proof sizes are larger than Groth16 but rapidly improving: Greyhound (Nguyen
\& Seiler, CRYPTO 2024) achieves \textasciitilde50 KB proofs with lattice-based
commitments, and LaBRADOR (CRYPTO 2023) achieves \textasciitilde58 KB. These are
orders of magnitude larger than Groth16's 192 bytes but competitive with raw
STARK proofs. The performance trajectory suggests that lattice-based schemes may
reach practical production within 3-5 years.

Why does this path matter? Because it is the only path that survives quantum
computing without any caveats. The hybrid path (Path One) is quantum-vulnerable
at the verification wrapper. The pure transparent path (Path Two) relies on hash
functions whose collision resistance degrades under quantum attack (SHA-256
drops from 128-bit classical to \textasciitilde85-bit quantum collision
resistance via the BHT algorithm). The lattice path relies on Module-LWE/SIS,
which is believed to be quantum-resistant by design -- the same assumption
family that NIST chose for its post-quantum standards. This path does not merely
survive the quantum era. It was built for it.

The post-quantum folding path has a structural advantage that maps directly onto
the Layer 6 analysis from Chapter 7. That chapter presented a trilemma:
algebraic functionality, post-quantum security, and succinctness -- pick two.
Lattice-based commitments provide additive homomorphism (enabling folding, which
is a form of algebraic functionality), post-quantum security, and increasingly
competitive succinctness. The trilemma is not being sidestepped; it is being
actively compressed. Whether it can be fully dissolved remains an open question
-- KZG's \(O(1)\) proof size with full homomorphism has no post-quantum match
yet. But the gaps are narrowing with each paper.

\section{The Three-Path Table}\label{the-three-path-table}

\begin{longtable}[]{@{}
  >{\raggedright\arraybackslash}p{(\linewidth - 10\tabcolsep) * \real{0.0769}}
  >{\raggedright\arraybackslash}p{(\linewidth - 10\tabcolsep) * \real{0.0897}}
  >{\raggedright\arraybackslash}p{(\linewidth - 10\tabcolsep) * \real{0.2051}}
  >{\raggedright\arraybackslash}p{(\linewidth - 10\tabcolsep) * \real{0.2436}}
  >{\raggedright\arraybackslash}p{(\linewidth - 10\tabcolsep) * \real{0.1410}}
  >{\raggedright\arraybackslash}p{(\linewidth - 10\tabcolsep) * \real{0.2436}}@{}}
\toprule\noalign{}
\begin{minipage}[b]{\linewidth}\raggedright
Path
\end{minipage} & \begin{minipage}[b]{\linewidth}\raggedright
Setup
\end{minipage} & \begin{minipage}[b]{\linewidth}\raggedright
Inner Primitive
\end{minipage} & \begin{minipage}[b]{\linewidth}\raggedright
Outer Verification
\end{minipage} & \begin{minipage}[b]{\linewidth}\raggedright
PQ Status
\end{minipage} & \begin{minipage}[b]{\linewidth}\raggedright
Production Status
\end{minipage} \\
\midrule\noalign{}
\arrayrulecolor{MidnightBlue}\midrule\arrayrulecolor{MidnightBorder}
\endhead
\bottomrule\noalign{}
\endlastfoot
\textbf{Hybrid STARK-to-SNARK} & STARK inner (transparent) + Groth16 outer
(trusted) & Hash-based FRI / Merkle & BN254 pairing (\textasciitilde250K gas) &
Inner: quantum-safe; Outer: quantum-vulnerable & Dominant production default \\
\textbf{Pure Transparent} & Transparent only & Hash-based FRI, no pairings &
Large on-chain proof or alternative verification & Quantum-safe (with sufficient
hash output) & Ethereum L1 mandate; advancing \\
\textbf{Post-Quantum Folding} & Transparent (lattice-based) & Module-SIS
commitments & Lattice verification (higher cost) & Quantum-safe by design &
Research frontier; 3-5 year horizon \\
\end{longtable}

\section{The Causal Web: Why It Is a DAG, Not a
Stack}\label{the-causal-web-why-it-is-a-dag-not-a-stack}

The book has presented seven layers as floors in a building -- stacked, with
each resting on the one below. The evidence from Parts I and II tells a
different story. The layers are not a stack; they are a directed acyclic graph
(DAG) with bidirectional pressures. The building metaphor was useful for
learning. Now we must complicate it.

The most consequential causal chain runs \emph{upward} from Layer 6. Small-field
primitives (BabyBear, M31) at Layer 6 enabled Circle STARKs and efficient
multilinear proving at Layer 5, which enabled lookup-based and AIR-based
arithmetization at Layer 4, which shaped shard-based witness generation at Layer
3, which favored RISC-V ISAs at Layer 2, and universal zkVMs amortized setup
costs at Layer 1. The foundation shaped the building. That much the metaphor
gets right.

But there are equally important \emph{downward} pressures. The audience shapes
the show:

\textbf{Layer 7 forces Layer 6.} Ethereum gas economics demand Groth16
verification, which requires BN254 pairings, which constrains the outer proof
system. The verifier's economics shape the prover's cryptography. The audience's
ticket price dictates the magician's technique.

\textbf{Layer 7 forces Layer 5.} STARK-to-SNARK wrapping exists because
Ethereum's gas costs make raw STARK verification uneconomical. The verification
layer forces a compression step that the proof system layer would not otherwise
need.

\textbf{Layer 2 constrains Layer 4.} Cairo was designed so that its ISA
minimizes arithmetization cost. The language was shaped by the constraint
system, not the other way around. Layer 4 requirements propagated upward to
shape Layer 2 design. The choreography was rewritten to fit the stage machinery.

\textbf{Layer 3 collapses into Layer 4.} In Jolt, witness generation \emph{is}
the arithmetization -- every instruction is decomposed into lookups on
subtables. There is no meaningful boundary between ``generate the trace'' and
``encode the computation.''

The pedagogical ordering (Layer 1 first, Layer 7 last) follows the \emph{data
flow}: setup before language, language before witness, witness before proof.
This is how a user encounters the system. But the \emph{engineering causality}
is inverted: the field choice at Layer 6 determines the commitment scheme, which
determines the polynomial representation, which determines the arithmetization,
which shapes everything above it.

These four examples are not anomalies. They are the norm. Once you catalog every
causal arrow in the system, the picture that emerges is not a tower but a web --
and the web has a specific mathematical structure that is worth naming
precisely.

\subsection{The Shape of the Web}\label{the-shape-of-the-web}

Roger Penrose, in \emph{The Road to Reality}, draws a distinction between
structures that are merely complicated and structures that are \emph{irreducibly
entangled}. A stack is complicated: many parts, one ordering. A DAG is
entangled: many parts, many orderings, no cycles. The seven-layer model, once
you draw all the arrows, is a DAG with at least fourteen directed edges and zero
cycles. It has structure, but that structure is not linear.

Consider the full edge set. Layer 6 forces Layer 5 (field choice determines
commitment scheme). Layer 5 forces Layer 4 (commitment scheme shapes
arithmetization). Layer 4 shapes Layer 3 (constraint format determines witness
layout). Layer 3 shapes Layer 2 (witness cost influences ISA design). Layer 2
shapes Layer 1 (ISA scope determines setup complexity). That is the upward chain
-- six edges, roughly linear, roughly matching the pedagogical ordering. If this
were all, the stack metaphor would suffice.

But it is not all. Layer 7 forces Layer 6 (gas economics demand BN254). Layer 7
forces Layer 5 (verification cost forces wrapping). Layer 4 forces Layer 2
(constraint cost shapes ISA, as Cairo demonstrates). Layer 3 collapses into
Layer 4 (Jolt's lookup singularity). Layer 6 forces Layer 7 (field size
determines proof size, which determines verification cost). Layer 1 forces Layer
5 (setup type constrains which proof systems are available). Layer 5 forces
Layer 7 (proof format determines verifier contract design). Layer 2 forces Layer
3 (ISA instruction count determines trace width).

That is fourteen edges among seven nodes. The graph has no cycles -- you cannot
follow arrows from any node back to itself -- which is what makes it a DAG
rather than a general directed graph. But the graph has \emph{width}: multiple
independent paths connect the same pair of nodes. Layer 6 reaches Layer 7 both
directly (field determines proof size) and indirectly via Layer 5 (field
determines proof system, which determines verifier). Layer 7 reaches Layer 5
both directly (gas cost forces wrapping) and indirectly via Layer 6 (gas cost
forces BN254, which constrains proof system). These parallel paths are why
changing a single parameter -- say, the base field -- propagates unpredictably
through the stack. The change follows multiple routes, and those routes
interfere with each other.

Penrose would recognize this as a feature, not a bug. Physical theories have the
same structure: general relativity and quantum mechanics are not stacked but
entangled, each constraining the other through multiple channels. The seven
layers of zero-knowledge proofs exhibit the same irreducible entanglement. You
cannot understand Layer 5 (the proof system) without simultaneously
understanding Layer 6 (the field) and Layer 7 (the verifier). You cannot design
Layer 2 (the ISA) without understanding Layer 4 (the constraint system). The
system is not modular. It is coherent -- every part is connected to every other
part through at most two hops.

\subsection{Why No Cycles?}\label{why-no-cycles}

The absence of cycles is not obvious and deserves explanation. Why can you not
follow arrows from Layer 7 back to Layer 7?

The answer is that the arrows represent \emph{design-time constraints}, not
\emph{runtime data flow}. At runtime, data flows in a rough circle: the user
submits a transaction (Layer 2), which generates a witness (Layer 3), which is
arithmetized (Layer 4), which is proved (Layer 5), which is verified (Layer 7),
which triggers a state change that enables the next transaction (back to Layer
2). That loop is a cycle, and it is real. But the \emph{design} constraints --
which architectural choices force which other architectural choices -- are
acyclic. Choosing M31 at Layer 6 forces Circle STARKs at Layer 5, but choosing
Circle STARKs at Layer 5 does not force M31 at Layer 6 (you could use any
Mersenne prime). The arrows are asymmetric. The forcing goes one way.

This distinction -- cyclic runtime flow, acyclic design constraints -- explains
why the seven-layer model is useful despite being wrong. The model captures the
design-time DAG by projecting it onto a linear ordering. The projection loses
information (it hides the downward and cross-cutting arrows) but preserves the
acyclicity. A stack is the simplest DAG. The seven-layer stack is the simplest
correct projection of the seven-layer web. It is a useful lie that points toward
a more interesting truth.

A seven-layer model that acknowledges this bidirectionality is more accurate
than one that implies simple top-down dependency. The layers are aspects of a
single integrated system, not modules with clean interfaces. The magician, the
stage, and the audience are not separable. They are one show.

\section{Trust Decomposition: Seven Weaker
Assumptions}\label{trust-decomposition-seven-weaker-assumptions}

Now we arrive at the heart of the matter.

Everything in this book has been building toward this section. The seven layers,
the three paths, the proof core, the causal DAG -- all of it converges on a
single observation, and it is the most important thing this book has to say.

A traditional financial system requires trusting a single institution. The bank
holds your money, knows your balance, controls your transactions, and you trust
that it will behave honestly. You trust \emph{one entity} with
\emph{everything}.

Zero-knowledge proofs do not eliminate trust. They do something different. They
\emph{decompose} it. They shatter a single monolithic trust assumption into
seven independent, weaker assumptions. And because the assumptions are
independent, no single failure breaks the system. This is not trustlessness. It
is something better: trust that is distributed, auditable, and replaceable.

Here are the seven assumptions, and each one is a thread you can pull:

\textbf{Layer 1}: At least one of N ceremony participants was honest (for
trusted setups), or that hash functions are collision-resistant (for transparent
setups). This is the trust in the stage itself -- that the mathematical
parameters were generated fairly.

\textbf{Layer 2}: The circuit was correctly written and audited. The
under-constrained circuit epidemic catalogued in Chapter 3 remains the dominant
vulnerability class, and a bug here lets the prover prove false statements. This
is the trust in the choreography: that the script describes the trick
accurately.

\textbf{Layer 3}: The hardware running the prover does not leak the witness
through side channels. Timing attacks on Zcash's Groth16 prover leaked
transaction amounts with \(R = 0.57\) correlation. Cache timing attacks on
ZK-friendly hash functions (Poseidon, Reinforced Concrete) have been
demonstrated in cloud environments. This is the trust in the backstage -- that
no one can see behind the curtain through cracks in the wall.

\textbf{Layer 4}: The arithmetization correctly encodes the computation. If the
translation from program to polynomial constraints is wrong, the proof system
faithfully proves the wrong thing. This is the trust in the encoding -- that the
mathematical puzzle accurately represents the trick.

\textbf{Layer 5}: The proof system is sound -- no efficient adversary can forge
proofs. This depends on the Fiat-Shamir transform being correctly implemented
(the Frozen Heart bug of 2022 affected three independent implementations
simultaneously) and on the underlying interactive proof being sound. This is the
trust in the seal -- that the certificate cannot be forged.

\textbf{Layer 6}: The mathematical hardness assumptions hold. Discrete
logarithms are hard (or lattice problems are hard, or hash preimages are hard).
These are conjectures, not theorems. Tower NFS improvements already reduced
BN254's estimated security from \textasciitilde128 bits to \textasciitilde100
bits. This is the trust in mathematics itself -- the deepest assumption, and the
one we have the least power to verify.

\textbf{Layer 7}: The governance structure will not override the cryptography.
Most deployed ZK rollups are Stage 0 or Stage 1 on L2Beat's framework, meaning a
multisig can override the proof system. The Beanstalk flash-loan governance
attack (\$182M, April 2022) and the Tornado Cash CREATE2 contract replacement
(May 2023) demonstrate that governance can be exploited. This is the trust in
the theater management -- that the people who run the venue will not rig the
show.

\subsection{When Each Thread Snaps: Seven Failure
Scenarios}\label{when-each-thread-snaps-seven-failure-scenarios}

The seven assumptions above are not hypothetical. Every one of them has failed,
is failing, or will fail in a live system. Isaac Asimov once observed that the
most exciting phrase in science is not ``Eureka!'' but ``That's funny\ldots{}''
-- the moment when an assumption you did not know you were making turns out to
be wrong. In zero-knowledge systems, ``That's funny\ldots{}'' is the sound of
money disappearing. Here is what each failure looks like in practice, what
breaks, and -- critically -- what does \emph{not} break.

\textbf{If Layer 1 fails (compromised setup):} An attacker who knows the toxic
waste from a trusted setup ceremony can forge proofs of arbitrary statements.
They can mint tokens from nothing, approve transactions that never happened,
fabricate state transitions whole cloth. The terrifying property of this failure
is its \emph{invisibility}. A forged Groth16 proof is indistinguishable from a
legitimate one -- both are 192 bytes, both pass the verifier, both look
identical on-chain. No alarm fires. No anomaly appears in the logs. The
counterfeiting is perfect by construction. This is why Zcash's Powers-of-Tau
ceremony involved 87 independent participants across six continents -- the
assumption is that at least one of them destroyed their toxic waste. If all 87
were compromised (through coercion, incompetence, or a coordinated state-level
attack), the entire shielded pool would be silently forgeable. Note what does
\emph{not} break: the circuit logic (Layer 2) is still correct, the witness
privacy (Layer 3) is still intact, the math (Layer 6) still holds. The stage was
rigged, but the trick's choreography was genuine.

\textbf{If Layer 2 fails (buggy circuit):} The proof system faithfully proves a
false statement. This is the most common failure mode in production, and the
canonical example is Tornado Cash. A single missing constraint in Tornado Cash's
withdrawal circuit allowed an attacker to generate valid proofs for withdrawals
from deposits that never existed. The circuit was supposed to verify that a
nullifier corresponded to a real deposit commitment in the Merkle tree. A
missing range check meant the prover could satisfy the constraints with
fabricated values. The proof was \emph{valid} -- it passed the verifier --
because the verifier only checks that the proof matches the circuit, and the
circuit was wrong. The Zcash ``InternalH'' bug (CVE-2019-7167, February 2019) is
the same species: a missing check in the Sapling circuit would have allowed
unlimited counterfeiting of shielded ZEC. Found by a Zcash engineer during
routine review, it was quietly patched before exploitation. The gap between
``found by an auditor'' and ``found by an attacker'' was a matter of months.
Circom's under-constrained circuit epidemic -- where developers use the
\texttt{\textless{}-\/-} assignment operator instead of the
\texttt{\textless{}==} constraining operator -- produces the same failure at
industrial scale. The proof system is not broken. The program is broken. The
seal is perfect; the document it seals is a forgery.

\textbf{If Layer 3 fails (witness leakage):} The zero-knowledge property
evaporates. The proof is still valid, and the computation is still correct, but
the secret inputs are exposed. The system proves the truth and simultaneously
reveals what it was supposed to hide. Timing side-channel attacks on Zcash's
Groth16 prover demonstrated this concretely: by measuring how long proof
generation took, an observer could infer the transaction amount with
\(R = 0.57\) correlation. The proof said ``this transaction is valid'' without
revealing the amount -- but the \emph{time it took to generate the proof} leaked
the amount through a side channel. In cloud proving environments (AWS, GCP),
cache-timing attacks on ZK-friendly hash functions like Poseidon and Reinforced
Concrete are even more direct: a co-located VM can observe memory access
patterns and reconstruct the witness. This is a privacy catastrophe but not an
integrity catastrophe. The computations are still correct. The proofs are still
sound. But the magician's secrets are visible through cracks in the dressing
room wall.

\textbf{If Layer 4 fails (incorrect arithmetization):} The constraint system
does not faithfully represent the computation it claims to encode. This is
subtler than a Layer 2 bug because the \emph{program} may be correct -- the
error is in the \emph{translation} from program to polynomial constraints.
Consider a zkVM that claims to prove RISC-V execution. If the arithmetization
incorrectly encodes the behavior of, say, the \texttt{slt} (set-less-than)
instruction -- treating signed comparison as unsigned, or failing to handle the
overflow edge case at INT\_MIN -- then the zkVM produces valid proofs of
executions that never happened. The program was correct. The proof system was
sound. The encoding between them was wrong. This failure is especially dangerous
in hand-rolled constraint systems (pre-zkVM era), where a developer manually
translates each operation into R1CS or AIR constraints. SP1's formal
verification of all 62 RISC-V opcodes against the Sail specification exists
precisely to prevent this class of failure. The trust is not in the magician or
the seal, but in the translator standing between them -- and translators make
mistakes.

\textbf{If Layer 5 fails (broken proof system):} An attacker can forge proofs
without knowing the witness. The Frozen Heart vulnerability (2022) is the
textbook case. Three independent implementations of the Fiat-Shamir transform --
Bellman (Zcash), Gnark (ConsenSys), and an academic reference -- all made the
same mistake: they failed to bind the public inputs to the transcript hash. An
attacker could take a valid proof for one statement and re-use it for a
different statement. The proof said ``I know a witness for X'' but could be
replayed to claim ``I know a witness for Y.'' All three implementations were
broken simultaneously, because all three misunderstood the same subtle
requirement of the Fiat-Shamir transform. This is a soundness catastrophe. The
seal can be forged. Valid-looking proofs can be manufactured for false
statements. Unlike a Layer 2 failure (where the circuit is wrong but the proof
system is honest), a Layer 5 failure means the proof system itself is
compromised. Every proof it has ever generated becomes suspect.

\textbf{If Layer 6 fails (broken hardness assumption):} The mathematical
foundation dissolves. This has not happened catastrophically yet, but it has
happened incrementally, and incrementally is frightening enough. Tower NFS
improvements reduced BN254's estimated security from approximately 128 bits to
approximately 100 bits -- not a break, but an erosion. The 2023 lattice basis
reduction advances by Ducas and van Woerden tightened the known attacks on NTRU
lattices, and while they did not break any deployed scheme, they moved the
boundary closer. A full break of BN254's discrete logarithm problem would
compromise every Groth16 proof ever generated on that curve: past, present, and
future. Every wrapped STARK-to-SNARK proof on Ethereum would be forgeable. Every
ZK rollup using BN254 verification would lose its security guarantee
retroactively. Shor's algorithm achieves exactly this for all pairing-based and
elliptic-curve cryptography, given a sufficiently large quantum computer. The
question is not \emph{whether} this assumption will weaken but \emph{when} and
\emph{how fast}. This is the failure that the post-quantum folding path (Path
Three) exists to survive.

\textbf{If Layer 7 fails (governance override):} The cryptography is irrelevant
because the humans in charge simply bypass it. This is the most prosaic failure
and arguably the most dangerous, because it requires no mathematical
sophistication -- only social engineering, legal coercion, or economic
incentive. The Beanstalk flash-loan governance attack (\$182M, April 2022)
demonstrated the economic version: an attacker borrowed enough governance tokens
to pass a malicious proposal in a single transaction, draining the protocol's
treasury. The cryptography was never touched. The proofs were never forged. The
governance mechanism -- the human layer above the math -- was the attack
surface. The Tornado Cash CREATE2 replacement (May 2023) demonstrated the
supply-chain version: the deployer address used CREATE2 to replace the
governance contract with a malicious one, granting the attacker control over all
locked funds. Again, no cryptographic break was needed. Most deployed ZK rollups
operate with upgrade multisigs that can push new verifier contracts, effectively
overriding any proof system guarantee. If three of five multisig holders collude
(or are coerced by a state actor), they can deploy a verifier that accepts all
proofs, or no proofs, or only proofs from approved provers. The mathematical
fortress has a human door, and the door has a human lock.

\subsection{The Cascade Structure}\label{the-cascade-structure}

Not all failures are created equal, and not all are independent. The seven
assumptions form their own internal DAG of failure propagation.

A Layer 6 failure cascades into Layer 5 (if the hardness assumption breaks, the
proof system built on it is unsound) and Layer 1 (if discrete logs are easy,
trusted setup toxic waste can be reconstructed). But it does \emph{not} cascade
into Layer 2 (the circuit logic is still correct), Layer 3 (the witness
generation is still private against non-quantum adversaries), or Layer 4 (the
arithmetization is still faithful). A quantum computer that breaks BN254 makes
Groth16 proofs forgeable, but it does not introduce bugs into Circom circuits.
The choreography is fine. The seal is broken.

A Layer 2 failure is \emph{contained}. A buggy circuit produces provably wrong
results, but the proof system is still sound (it just proves the wrong thing),
the setup is still valid, the hardware does not leak, and the math still holds.
This containment is precisely what makes Layer 2 failures survivable: fix the
circuit, redeploy the verifier, and the system recovers. The Zcash InternalH bug
was patched in a single release. The stage machinery was fine; only the script
needed rewriting.

A Layer 7 failure is the most isolated and the most devastating. It requires no
interaction with any other layer. A governance override does not break the math,
corrupt the circuit, or leak the witness. It simply ignores all of them. This
isolation means that Layer 7 cannot be fixed by improving any other layer.
Better proof systems, stronger fields, more rigorous audits -- none of these
matter if a three-of-five multisig can replace the verifier contract. The only
defense is the same defense that human institutions have always relied on:
constitutional constraints, time-locks, social consensus, and the slow,
unglamorous work of governance design.

The cascade structure reveals a counterintuitive truth about the seven-layer
model. The \emph{deepest} failures (Layer 6: math breaks) are the most
catastrophic in scope but the least likely in practice. The \emph{shallowest}
failures (Layer 2: buggy circuit, Layer 7: governance override) are the most
common in practice but the most recoverable. The threat landscape is inverted:
the risks you encounter most often are the risks you can fix most easily. This
inversion is what makes the trust decomposition genuinely useful. A monolithic
trust model (trust the bank) gives you one failure mode: total. A decomposed
trust model gives you seven failure modes, most of which are partial and
recoverable. The system degrades gracefully, and graceful degradation is the
definition of resilient engineering.

Each of these is independently falsifiable, independently auditable, and
independently improvable. Breaking one does not necessarily break the others
(though some failures cascade -- a quantum computer breaks Layers 1, 5, and 6
simultaneously for pairing-based systems). This decomposition is the genuine
value proposition of zero-knowledge proofs: not trustlessness, but trust
minimization through distribution.

Remember the bank from Chapter 1? The institution that holds your money, knows
your balance, and controls your transactions? With zero-knowledge proofs, you no
longer trust one bank with everything. You trust that at least one ceremony
participant was honest. You trust that the circuit was correctly written. You
trust that the hardware does not leak. You trust that the math is hard. You
trust that the governance will not go rogue. Seven assumptions instead of one.
Each weaker. Each testable. Each replaceable.

That is not a marketing slogan. It is a structural transformation.

\section{``Trustless'' versus
``Trust-Minimized''}\label{trustless-versus-trust-minimized}

The analysis concludes with a question that should be printed on the wall of
every ZK team's office:

\begin{quote}
If zero-knowledge proofs provide ``trustless'' computation, but the setup
requires trusting ceremony participants (Layer 1), the program requires trusting
the developer not to make constraint errors (Layer 2), the witness requires
trusting the hardware not to leak side channels (Layer 3), the proof system
requires trusting that mathematical hardness assumptions hold (Layer 6), and the
verifier requires trusting that governance will not override the math (Layer 7)
-- then where, exactly, is the ``trustless'' part?
\end{quote}

The answer: nowhere. ``Trustless'' is a word that flatters the technology and
misleads the user. Zero-knowledge proofs do not eliminate trust. They minimize
and distribute it. Instead of trusting one bank with your financial data, you
trust that: (a) at least one ceremony participant was honest, (b) the circuit
was correctly written and audited, (c) the hardware is not leaking, (d) discrete
logarithms are hard, and (e) the governance multisig will not go rogue.

Each of these is a weaker assumption than trusting a single entity. The
combination is far more resilient than any single point of trust. But they are
assumptions nonetheless, and a responsible guide to zero-knowledge proofs should
catalog them explicitly.

We stated this thesis in the opening pages of Chapter 1: trust decomposition,
not trust elimination. Ten chapters later, the decomposition is precise. Seven
assumptions instead of one. Fourteen causal edges instead of a monolith. Three
architectural paths, each with different failure profiles. The word
``trustless'' obscures every one of these distinctions. The word
``trust-minimized'' preserves them.

The remaining chapters of this book will use ``trust-minimized'' rather than
``trustless.'' Not because it sounds better -- it sounds worse, deliberately --
but because it is accurate. And accuracy in naming things is where understanding
begins.

\chapter{Chapter 11: zkVMs -- The Universal
Stage}\label{chapter-11-zkvms-the-universal-stage}

\section{zkVMs Across the Stack}\label{zkvms-across-the-stack}

We need to talk about the thing that changed everything.

The zkVM is not a Layer 2 phenomenon. It is a technology that reaches into every
layer of the stack -- from the field choice at Layer 6, through the
arithmetization at Layer 4, to the verification economics at Layer 7. If the
seven-layer model is the map, the zkVM is the earthquake that reshaped the
terrain.

Before zkVMs, every zero-knowledge application required building a custom stage:
hand-crafting constraint systems, choosing a field, designing a witness format,
writing a custom verifier. Each application was a bespoke production -- a
one-night show with its own scenery, its own props, its own choreography. The
zkVM changed this by providing a \textbf{universal stage} that can host any
trick. A developer writes ordinary Rust, compiles it to RISC-V, and the zkVM
handles everything else: witness generation, arithmetization, proving,
compression, and verification. The stage is universal; the performances are
infinite.

The preceding layer-by-layer analysis reveals why this matters so deeply. The
most urgent findings at nearly every layer are not isolated gaps but
consequences of the same underlying shift. Polygon zkEVM's shutdown at Layer 2,
the Witness Gap's amplification at Layer 3, CCS and LogUp's emergence at Layer
4, folding's rise and the hybrid pipeline's dominance at Layer 5, the
small-field revolution at Layer 6, STARK-to-SNARK wrapping at Layer 7 -- these
are the seismic effects of a single tectonic event: the zero-knowledge virtual
machine reorganized the entire stack around itself.

\section{The Landscape Table (March 2026)}\label{the-landscape-table-march-2026}

The numbers tell the story concisely. Eight of ten major zkVMs now target
RISC-V. Only Stwo (Cairo) and zkWASM (WebAssembly) hold out -- and even
StarkWare's ecosystem hedges via Kakarot's EVM-on-Stwo path. The Ethereum
Foundation declared the speed race ``effectively won'' in December 2025 and
pivoted to 128-bit provable security by end of 2026.

\begin{longtable}[]{@{}
  >{\raggedright\arraybackslash}p{(\linewidth - 14\tabcolsep) * \real{0.1250}}
  >{\raggedright\arraybackslash}p{(\linewidth - 14\tabcolsep) * \real{0.1250}}
  >{\raggedright\arraybackslash}p{(\linewidth - 14\tabcolsep) * \real{0.1250}}
  >{\raggedright\arraybackslash}p{(\linewidth - 14\tabcolsep) * \real{0.1250}}
  >{\raggedright\arraybackslash}p{(\linewidth - 14\tabcolsep) * \real{0.1250}}
  >{\raggedright\arraybackslash}p{(\linewidth - 14\tabcolsep) * \real{0.1250}}
  >{\raggedright\arraybackslash}p{(\linewidth - 14\tabcolsep) * \real{0.1250}}
  >{\raggedright\arraybackslash}p{(\linewidth - 14\tabcolsep) * \real{0.1250}}@{}}
\toprule\noalign{}
\begin{minipage}[b]{\linewidth}\raggedright
\end{minipage} & \begin{minipage}[b]{\linewidth}\raggedright
SP1 Hypercube
\end{minipage} & \begin{minipage}[b]{\linewidth}\raggedright
RISC Zero
\end{minipage} & \begin{minipage}[b]{\linewidth}\raggedright
Jolt
\end{minipage} & \begin{minipage}[b]{\linewidth}\raggedright
Stwo
\end{minipage} & \begin{minipage}[b]{\linewidth}\raggedright
Airbender
\end{minipage} & \begin{minipage}[b]{\linewidth}\raggedright
ZisK
\end{minipage} & \begin{minipage}[b]{\linewidth}\raggedright
Pico Prism
\end{minipage} \\
\midrule\noalign{}
\arrayrulecolor{MidnightBlue}\midrule\arrayrulecolor{MidnightBorder}
\endhead
\bottomrule\noalign{}
\endlastfoot
\textbf{Org} & Succinct & RISC Zero & a16z & StarkWare & Matter Labs & SilentSig
& Brevis \\
\textbf{ISA} & RISC-V (RV32IM) & RISC-V (RV32IM) & RISC-V (RV32I) & Cairo
(custom) & RISC-V (RV32IM) & RISC-V 64 & RISC-V (RV32IM) \\
\textbf{Arithmetization} & Multilinear AIR + LogUp-GKR & AIR (DEEP-ALI) & R1CS +
Lasso lookups & Circle AIR + LogUp & AIR (degree-2) & AIR / PIL & AIR
(Plonky3) \\
\textbf{Proof system} & Multilinear STARK & FRI STARK & Spartan sumcheck &
Circle STARK & DEEP STARK & STARK (recursive) & STARK (Plonky3) \\
\textbf{Field} & BabyBear (31-bit) & BabyBear (31-bit) & BN254 (256-bit) & M31
(31-bit) & M31 (31-bit) & Goldilocks (64-bit) & BabyBear / M31 \\
\textbf{SNARK wrap} & Groth16 & Groth16 & Planned & None (native) & Groth16 &
Groth16 & Groth16 \\
\textbf{Eth block} & 6.9 s / 16 GPU & 44 s / cluster & N/A & N/A (Cairo) & 35 s
/ 1 GPU & 6.6 s / 24 GPU & 6.9 s / 16 GPU \\
\textbf{Maturity} & Production & Production & Beta & Production & Production &
Adv. testnet & Production \\
\end{longtable}

\textbf{Glossary.} \emph{ISA} = instruction set architecture, the fundamental
language a processor understands. \emph{Arithmetization} = the method of
translating computation into mathematical equations. \emph{AIR} = Algebraic
Intermediate Representation, encoding computation as polynomial constraints over
an execution trace. \emph{LogUp-GKR} = a sumcheck-based lookup argument.
\emph{Circle STARK} = a STARK adapted to work over the circle group of a
Mersenne prime. \emph{SNARK wrap} = compressing a large transparent proof into a
small proof for cheap on-chain verification. \emph{Field} = a set of numbers
with special arithmetic properties; ``31-bit'' and ``64-bit'' refer to element
size, with smaller fields enabling faster operations on modern hardware.

Three systems dropped from this table deserve brief mention. Nexus 3.0 abandoned
Nova-based folding for a Stwo backend after observing a 1000x speed penalty from
classical folding -- a telling result for the practical viability of folding in
production zkVMs. Valida uses a custom stack-based ISA designed from scratch for
ZK proving, with no independent large-scale benchmarks. zkWASM is the only
remaining PLONKish/KZG outlier, targeting WebAssembly rather than RISC-V.

For architects choosing a zkVM, the landscape table above describes \emph{what
exists}. The rubric below describes \emph{how to choose}:

\begin{longtable}[]{@{}
  >{\raggedright\arraybackslash}p{(\linewidth - 4\tabcolsep) * \real{0.2000}}
  >{\raggedright\arraybackslash}p{(\linewidth - 4\tabcolsep) * \real{0.4333}}
  >{\raggedright\arraybackslash}p{(\linewidth - 4\tabcolsep) * \real{0.3667}}@{}}
\toprule\noalign{}
\begin{minipage}[b]{\linewidth}\raggedright
Goal
\end{minipage} & \begin{minipage}[b]{\linewidth}\raggedright
Recommended
\end{minipage} & \begin{minipage}[b]{\linewidth}\raggedright
Rationale
\end{minipage} \\
\midrule\noalign{}
\arrayrulecolor{MidnightBlue}\midrule\arrayrulecolor{MidnightBorder}
\endhead
\bottomrule\noalign{}
\endlastfoot
General-purpose execution & SP1 Hypercube, RISC Zero & RISC-V, mature tooling,
broadest Rust ecosystem support \\
Maximum throughput & SP1 Hypercube, Airbender & Best Ethereum block proving
benchmarks (6.9s, 21.8M cycles/sec) \\
ZK-native efficiency & Stwo (Cairo) & Purpose-built ISA eliminates translation
overhead; Starknet-native \\
Post-quantum trajectory & Neo/SuperNeo (watch) & Lattice-based, CCS-native,
127-bit PQ security; 3-5 year horizon \\
Lookup-heavy design & Jolt & Lasso decomposable lookups; sumcheck-native; avoids
NTT bottleneck \\
Minimum trusted setup & Stwo, Pico Prism & Fully transparent (hash-based FRI);
no ceremony required \\
Formal verification & SP1 & 62 RISC-V opcodes verified against Sail spec;
strongest correctness guarantees \\
\end{longtable}

\section{Three zkVMs Through Seven
Layers}\label{three-zkvms-through-seven-layers}

The best way to understand why the seven-layer model bends under zkVM pressure
is to trace three representative systems through all seven layers. Watch how
they cross the same territory by different routes.

\subsection{SP1 Hypercube: The General-Purpose
Champion}\label{sp1-hypercube-the-general-purpose-champion}

\textbf{Layer 1 (Setup).} Hybrid -- transparent inner loop (hash-based Poseidon2
over BabyBear), trusted outer wrap (Groth16 over BN254). The KZG ceremony is for
the wrapper only. This demonstrates that the trusted-or-transparent choice is
actually ``both.''

\textbf{Layer 2 (ISA).} RISC-V (RV32IM). Developers write ordinary Rust; the
compiler handles the rest. SP1 implements 39+ RISC-V instructions, with all 62
core opcodes formally verified against the official RISC-V Sail specification.

\textbf{Layer 3 (Witness).} Shard-based execution traces with continuations.
Each shard is an independent proving unit; ``shared challenges'' enforce
consistency at shard boundaries. This is where the Witness Gap lives -- SP1's
witness generation is CPU-bound while its proving is GPU-accelerated, meaning
witness generation consumes an estimated 60-70\% of total proving time.

\textbf{Layer 4 (Arithmetization).} Multi-table AIR with LogUp-GKR cross-table
lookups over multilinear polynomials. Each RISC-V instruction type has its own
constraint table (``chip''). Precompiles (SHA-256, Keccak, secp256k1) are
independent STARK tables connected via LogUp.

\textbf{Layer 5 (Proof).} Four-stage pipeline: core STARK proof per shard,
recursive compression, shrink (field transition), Groth16 wrap. Proves 99.7\% of
Ethereum L1 blocks in under 12 seconds on 16 RTX 5090 GPUs.

\textbf{Layer 6 (Primitives).} BabyBear field (31-bit), Poseidon2 hash, Jagged
PCS (commits only to occupied trace rows, eliminating padding waste), FRI-based
commitment.

\textbf{Layer 7 (Verifier).} Groth16 on-chain verification at approximately
250-300K gas on any EVM chain. Live on Ethereum mainnet via the Succinct Prover
Network.

\textbf{What makes SP1 architecturally distinctive.} SP1 Hypercube is an
exercise in \emph{factoring a monolithic problem into independent pieces and
then reassembling them under a single algebraic umbrella}. The system's defining
architectural idea is not any single layer but the interaction between three
design choices that reinforce each other. First, the multi-table AIR
architecture assigns each RISC-V opcode its own constraint table -- a ``chip''
-- so that the constraint degree and column count for a SHA-256 precompile need
not compromise the constraint shape for a simple register-to-register add.
Second, LogUp-GKR cross-table lookups bind these independent chips together
using sumcheck over multilinear extensions, which avoids the quadratic blowup
that a naive permutation argument would impose as chip count grows. Third,
sharding with continuation challenges means that execution traces of arbitrary
length can be sliced into fixed-size proving units, each provable in parallel on
a separate GPU, with algebraic consistency enforced by shared random challenges
drawn after the shards are committed.

The result is a system that scales \emph{horizontally} in two independent
dimensions: more instruction types (more chips) and longer executions (more
shards). Adding a new precompile -- say, a BLS12-381 pairing for Ethereum
validator operations -- requires only a new chip table and new LogUp entries;
existing chips and the recursive compression pipeline remain unchanged. This
modular extensibility explains why SP1 has accumulated over 39 instruction types
and counting, whereas architectures with monolithic trace layouts face painful
refactoring when adding even one new opcode.

The Jagged PCS matters just as much. Standard polynomial commitment schemes
commit to a fixed-size domain, which means that a chip with 10,000 occupied rows
in a shard of capacity \(2^{20}\) wastes commitment work on a million empty
rows. Jagged PCS commits only to the non-trivial portion, eliminating padding
overhead. In practice, most shards have a few ``hot'' chips (ALU, memory) and
many ``cold'' chips (uncommon opcodes). Without Jagged PCS, the cold chips would
dominate commitment cost despite contributing almost nothing to the computation.
With it, proving cost tracks actual computation, not worst-case table size.

SP1's four-stage pipeline -- core STARK per shard, recursive compression,
field-transition shrink, Groth16 wrap -- is also a study in staged trust
assumptions. The inner stages are fully transparent, hash-based, and
post-quantum resilient. Only the final wrap introduces a trusted setup and
classical-security assumption. If and when a post-quantum on-chain verifier
becomes practical, SP1 can drop the Groth16 stage and expose the transparent
inner proof directly. The architecture is built to survive the quantum
transition, not merely to endure it.

\subsection{Stwo/Cairo: The ZK-Native ISA
Champion}\label{stwocairo-the-zk-native-isa-champion}

\textbf{Layer 1.} Fully transparent (hash-based). For Ethereum L1 settlement, a
Groth16 wrapper exists via Herodotus, but the primary pipeline remains
STARK-native.

\textbf{Layer 2.} Cairo -- a custom ISA designed to minimize arithmetization
cost. Here the model's top-down flow inverts. The ISA \emph{is} the constraint
system, by design. Layer 4 requirements shaped Layer 2. The choreography was
written to serve the stage machinery, not the other way around.

\textbf{Layer 3.} Cairo VM execution producing columnar M31 traces. The field
choice is baked into the VM itself, not merely the proving step -- a deeper
coupling than in RISC-V zkVMs.

\textbf{Layer 4.} Flat AIR with one component per instruction, LogUp
cross-component lookups, mixed-degree constraints. Circle STARKs use the circle
group over M31 (where \(p+1 = 2^{31}\) is a power of two) to enable FFTs over a
field that lacks large multiplicative subgroups.

\textbf{Layer 5.} Circle STARK with SHARP aggregation. Approximately 100x faster
than its predecessor Stone. Live on Starknet mainnet since November 2025.

\textbf{Layer 6.} M31 field (31-bit). M31 arithmetic is approximately 125x
faster than the 252-bit Stark field used by Stone. The field choice
\emph{created} the need for Circle STARKs -- the strongest example of Layer 6
forcing a Layer 5 invention.

\textbf{Layer 7.} Native STARK verification on Starknet (no wrapping needed).
For Ethereum L1: SHARP-aggregated Groth16.

\textbf{What makes Stwo architecturally distinctive.} Stwo is what happens when
you design the entire stack backward from a single mathematical insight:
\emph{the group of points on the circle \(x^2 + y^2 = 1\) over a Mersenne prime
has order \(p+1\), and when \(p = 2^{31} - 1\), that order is exactly \(2^{31}\)
-- a power of two}. This accident of number theory is the seed from which the
entire Stwo architecture grows.

Conventional STARKs require a field with large multiplicative subgroups for
FFT-based polynomial evaluation. The Mersenne prime \(M31 = 2^{31} - 1\) has no
such subgroups -- its multiplicative group has order \(2^{31} - 2\), which
factors badly. A naive approach would disqualify M31 from STARK construction
entirely. Circle STARKs solve this by abandoning the multiplicative group in
favor of the circle group, where the ``FFT'' becomes a circle-group analog (the
Circle Number Theoretic Transform). The evaluation domain is the set of points
on the unit circle over \(\mathbb{F}_{M31}\), not the powers of a generator in
\(\mathbb{F}_{M31}^*\). This is not a minor algebraic substitution; it requires
rethinking polynomial commitments, coset structures, and FRI queries from the
ground up.

The payoff is a 125x speedup. M31 arithmetic -- 31-bit integer addition and
multiplication with a single modular reduction -- maps directly onto 32-bit CPU
and GPU instructions with no multi-precision overhead. A single SIMD lane
processes one field element. Compare this with the 252-bit Stark field used by
Stone, where each field multiplication requires multiple 64-bit limb operations
and carry propagation. The measured 125x speedup over Stone is not a software
optimization; it is a consequence of matching the algebraic structure to the
hardware word size.

Cairo's role is just as distinctive. Where RISC-V zkVMs treat the ISA and the
constraint system as separate concerns -- the ISA defines computation, the
arithmetization encodes it -- Cairo \emph{collapses the two}. A Cairo
instruction is simultaneously a machine operation and a set of polynomial
constraints. The compiler does not translate programs into constraints; it emits
programs that \emph{are} constraints. This eliminates the ``arithmetization
tax'' that RISC-V zkVMs pay: the overhead of encoding a general-purpose
instruction (designed for silicon hardware) into an algebraic form (designed for
polynomial provers). Gassmann et al.'s finding that standard LLVM optimizations
yield 40\% improvement on RISC-V zkVMs -- because LLVM optimizes for caches and
branch predictors that do not exist in ZK execution -- quantifies exactly the
tax that Cairo avoids.

The SHARP (Shared Prover) aggregation layer adds a dimension absent from SP1 and
Jolt: \emph{amortization across applications}. SHARP batches proofs from
multiple independent Starknet applications into a single recursive proof, so
that the fixed cost of Ethereum L1 verification is shared among all applications
that submit proofs in the same batch window. This is economic aggregation, not
just cryptographic recursion. A small contract with ten transactions per hour
pays a fraction of the L1 verification cost that it would bear alone. The
theater shares its rent among all the acts on stage.

The trade-off is ecosystem lock-in. Cairo is not Rust, not C, not any language
with a pre-existing developer community of millions. Every Cairo developer is a
developer that StarkWare's ecosystem must recruit and train. The Kakarot project
-- an EVM interpreter written in Cairo, running on Stwo -- is the architectural
hedge: it lets Ethereum developers write Solidity while Cairo and Stwo handle
the proving underneath. Whether this bridge is sturdy enough to carry mainstream
adoption is the open strategic question.

\subsection{Jolt: The Lookup Singularity
Pioneer}\label{jolt-the-lookup-singularity-pioneer}

\textbf{Layer 1.} Trusted (Hyrax/Pedersen commitments), with transparent
alternatives planned (Basefold in Jolt-b).

\textbf{Layer 2.} RISC-V (RV32I). Same as SP1, but the divergence begins at
Layer 4.

\textbf{Layer 3.} Every instruction is decomposed into lookups on small
subtables. Witness generation \emph{is} the arithmetization -- no meaningful
boundary exists between Layers 3 and 4. This is the first crack in the
seven-layer model. The backstage preparation and the encoding of the trick are
the same act.

\textbf{Layer 4.} Lookup-based arithmetization via Lasso. Instead of encoding
each operation as constraints, the system verifies every step against
pre-approved entries in a comprehensive reference table. A thin R1CS wrapper
(\textasciitilde60 constraints per cycle) handles control flow via Spartan. This
is a genuinely distinct paradigm from AIR, PLONKish, or CCS.

\textbf{Layer 5.} Sumcheck-based proving. No recursion in production; no
STARK-to-SNARK wrapping pipeline. Approximately 6x faster than RISC Zero on
initial benchmarks, but the gap has narrowed with GPU acceleration.

\textbf{Layer 6.} BN254 scalar field (256-bit). 256-bit field operations are
approximately 100x more expensive per operation than 31-bit, but Jolt
compensates by performing far fewer operations per CPU step.

\textbf{Layer 7.} No production on-chain verifier. This is Jolt's most
significant gap relative to SP1 and Stwo. The trick is performed brilliantly,
but the theater has no box office.

\textbf{What makes Jolt architecturally distinctive.} Jolt is the most radical
of the three systems, because it asks a question that the other two never
consider: \emph{what if we stopped writing constraints entirely?} SP1 and Stwo
both encode computation as polynomial constraints -- they differ in how they
organize and evaluate those constraints, but both accept the premise that
proving a computation means constraining it. Jolt rejects this premise. In Jolt,
proving a computation means \emph{looking it up}.

The core insight, which Jolt inherits from the Lasso lookup argument, is that
any function on small inputs can be represented as a table. A 32-bit addition is
a function from two 16-bit operands to a 17-bit result -- a table with
\(2^{32}\) entries. Proving that ``a + b = c'' does not require writing a
constraint that the prover satisfies; it requires showing that the triple (a, b,
c) appears in the addition table. The prover's obligation shifts from ``solve
this system of equations'' to ``demonstrate membership in this pre-computed
set.'' This is not a small change in formalism. It is a different epistemology
of computation.

The practical problem is that \(2^{32}\)-entry tables are too large to
materialize. Lasso solves this by decomposing large lookups into combinations of
small subtable lookups, using the algebraic structure of the operations
themselves. Addition decomposes by limbs; bitwise operations decompose by
individual bits; shifts decompose by position. Each subtable is small enough to
commit to directly -- typically \(2^{16}\) entries or fewer. The Lasso sumcheck
protocol then proves that the decomposed lookups are consistent with the
original instruction. The result is a system where the per-instruction proving
cost scales with the \emph{decomposability} of the instruction, not with the
number of constraints needed to describe it.

This decomposition is why Jolt's Layer 3 and Layer 4 merge into a single act. In
SP1, the witness (Layer 3) is an execution trace -- a table of register states
and memory values at each cycle -- and the arithmetization (Layer 4) is a set of
polynomial constraints that the trace must satisfy. These are conceptually and
computationally distinct stages. In Jolt, the ``witness'' for each instruction
\emph{is} the lookup decomposition: the set of subtable indices and values that
reconstruct the instruction's behavior. Generating the witness and performing
the arithmetization are the same computation, executed in a single pass. There
is no moment where ``the trace is ready and now we constrain it.'' The trace is
the constraint.

The choice of BN254 as the base field is worth examining. SP1 and Stwo chose
31-bit fields for raw throughput, accepting hash-based commitments and
transparent proofs. Jolt chose a 256-bit elliptic-curve field, accepting 100x
slower per-operation arithmetic, because BN254 enables Hyrax commitments -- a
multi-scalar-multiplication-based polynomial commitment scheme with \emph{no
trusted setup} and \emph{logarithmic} verification time. Hyrax commitments over
BN254 are the reason Jolt avoids both a ceremony (unlike Groth16 wrappers) and
hash-chain verification (unlike FRI). The field choice is not a performance
concession; it is a commitment-scheme selection that happens to be expensive.

Jolt's current absence of a production on-chain verifier is not merely an
engineering gap waiting to be filled. It reflects a deeper tension: the
sumcheck-based proving paradigm produces proofs whose verification cost does not
compress as neatly into the constant-size Groth16 format that Ethereum L1
expects. The planned Jolt-b variant, which replaces Hyrax with Basefold (a
hash-based, transparent commitment scheme), would enable a more conventional
STARK-to-SNARK wrapping pipeline. But this replacement changes Jolt's security
model -- from discrete-log hardness to hash collision resistance -- and alters
the system's identity. Whether Jolt can ship a verifier without becoming a
different system is the question that will determine whether the lookup
singularity remains a research landmark or becomes a production architecture.

\section{The ``Proof Core'' Triad}\label{the-proof-core-triad}

The comparative analysis reveals a pattern that the seven-layer model obscures:
Layers 4, 5, and 6 form a tightly coupled triad -- the \textbf{proof core} -- in
every production zkVM. These three layers are not three choices. They are one
choice with three manifestations.

In Stwo: M31 (Layer 6) forces Circle groups (Layer 5) forces Circle AIR (Layer
4).

In SP1: BabyBear (Layer 6) enables multilinear PCS (Layer 6) enables LogUp-GKR
(Layer 4/5) enables Jagged multi-table AIR (Layer 4).

In Jolt: BN254 (Layer 6) enables Hyrax commitment (Layer 6) makes sumcheck
natural (Layer 5) makes Lasso lookups natural (Layer 4).

In each case, the field choice \emph{determines} the commitment scheme, which
\emph{determines} the polynomial representation, which \emph{determines} the
arithmetization. The proof core is the inseparable nucleus of \{field,
commitment scheme, polynomial representation\} that straddles Layers 4, 5, and
6. Acknowledging this concept does not require restructuring the seven-layer
model, but it does require acknowledging that the model's clean layer boundaries
are pedagogical simplifications, not architectural truths.

To see why this coupling is not merely incidental but structurally inevitable,
consider what happens when you try to swap one element of the triad while
holding the others fixed. Suppose you wanted to keep M31 arithmetic (for speed)
but use Hyrax commitments (for transparent elliptic-curve proofs without a
trusted setup). Hyrax requires multi-scalar multiplication over an elliptic
curve whose scalar field matches the proving field. No standard curve has a
31-bit scalar field -- the smallest curves used in practice have 254-bit or
256-bit scalar fields. You would need to embed M31 elements into a vastly larger
field for every commitment operation, destroying the throughput advantage that
motivated the M31 choice. The triad resists substitution because its elements
are not independent components connected by interfaces; they are facets of a
single algebraic structure viewed from different angles.

The same rigidity appears in the other direction. Suppose you wanted to keep
BN254 (for Hyrax compatibility) but switch to AIR-based arithmetization (for the
modular chip architecture that makes SP1 extensible). AIR evaluation requires
FFTs over the proving field, and BN254's scalar field has a multiplicative
subgroup of order \(2^{28}\) -- adequate but not generous. More critically,
BN254 field operations are 100x slower than BabyBear operations, so the
FFT-intensive AIR evaluation that runs in milliseconds over BabyBear would take
seconds over BN254. The AIR paradigm's viability depends on cheap field
arithmetic; cheap field arithmetic depends on small fields; small fields depend
on hash-based commitments. Pull one thread and the entire fabric moves.

This structural coupling explains an otherwise puzzling empirical observation:
despite the enormous design space of possible \{field, commitment,
arithmetization\} combinations, production systems cluster around a small number
of triads. The March 2026 landscape shows essentially three: \{small field,
hash-based FRI, AIR\} (SP1, Stwo, Airbender, RISC Zero, Pico Prism), \{small
field, hash-based FRI, AIR + lookup\} (ZisK, OpenVM), and \{large field,
MSM-based, sumcheck + lookup\} (Jolt). The clustering is not a failure of
imagination. It is the proof core exerting its gravitational pull: only certain
combinations are self-consistent, and the self-consistent combinations are few.

\textbf{The decisive fork is at Layer 6 (field choice).} SP1 and Stwo both chose
small fields (31-bit) optimized for hardware throughput, accepting extension
fields and hash-based commitments. Jolt chose a 256-bit field, accepting higher
per-operation cost for native elliptic curve compatibility. This single
parameter cascades through every other layer. The choice is not primarily about
arithmetic speed -- it is about which mathematical universe the entire proof
system will inhabit. A 31-bit field lives in the universe of hash trees, Merkle
commitments, and transparent proofs. A 256-bit elliptic-curve field lives in the
universe of pairings, discrete logarithms, and structured reference strings.
These universes have different physics, and a system that enters one cannot
easily visit the other.

\textbf{The second fork is at Layer 4 (arithmetization).} SP1 and Stwo both use
AIR-based constraint systems. Jolt abandons AIR entirely in favor of
lookup-based arithmetization. AIR systems \emph{describe} computation as
polynomial constraints; Jolt \emph{tabulates} computation as lookup entries. The
distinction is deeper than syntax. In an AIR system, the prover must
\emph{solve} the constraint system -- find a witness that satisfies all
polynomial equations simultaneously. The constraint system is a specification
that the witness must match. In a lookup system, the prover must
\emph{demonstrate membership} -- show that each instruction's input-output
behavior appears in a pre-computed table. The table is not a specification to be
satisfied but a reference to be consulted. One paradigm asks ``does this answer
satisfy the question?'' The other asks ``is this answer in the book?''

The practical consequence is that AIR systems scale with constraint complexity
(more constraints per instruction means more work per step), while lookup
systems scale with table size and decomposition depth (more subtables means more
sumcheck rounds). For simple arithmetic -- adds, multiplies, shifts -- the
lookup approach can be dramatically cheaper, because the subtables are small and
the decomposition is clean. For complex operations -- hash functions,
elliptic-curve arithmetic -- the lookup approach faces a combinatorial explosion
in table size, which is why even Jolt uses a thin R1CS layer for control flow
rather than attempting to tabulate branching logic.

\textbf{Where the model breaks.} Cairo shows Layer 4 shaping Layer 2
(bidirectional dependency). Jolt shows Layers 3 and 4 collapsing into one. The
STARK-to-SNARK wrapping pipeline is not a Layer 5 choice; it pierces Layers 5,
6, and 7 as a single vertical shaft. The seven layers are better understood as
seven \emph{aspects} of a single integrated system, not seven \emph{modules}
with clean interfaces. The proof core triad is the strongest evidence for this
view: three layers that the model presents as independent choices are in fact a
single crystalline structure, as rigid and as beautiful as any lattice in
mathematics. Change one axis and the crystal shatters. The map is useful. The
territory is more interesting.

\section{Performance: The Cost Collapse}\label{performance-the-cost-collapse}

The performance trajectory of zkVMs over the past 24 months is one of the
steepest cost collapses in applied cryptography. It deserves to be stated
plainly.

\textbf{Real-time Ethereum proving is solved.} Four independent teams proved
over 99\% of mainnet blocks within the 12-second slot time by late 2025:

\begin{longtable}[]{@{}llll@{}}
\toprule\noalign{}
System & Hardware & Avg Block Time & \% Blocks \textless{} 12s \\
\midrule\noalign{}
\arrayrulecolor{MidnightBlue}\midrule\arrayrulecolor{MidnightBorder}
\endhead
\bottomrule\noalign{}
\endlastfoot
SP1 Hypercube & 16x RTX 5090 & 6.9 s & 99.7\% \\
ZisK & 24x RTX 5090 & 6.6 s & 99.7\% \\
Pico Prism & 16x RTX 5090 & 6.9 s & 99\%+ \\
OpenVM 2.0 & 16x RTX 5090 & \textless{} 12 s (p99) & \textasciitilde99\%+ \\
\end{longtable}

Airbender proves a block in 35 seconds on a \emph{single} H100 GPU -- the
fastest single-GPU result, at 21.8 million cycles per second at the base STARK
layer.

\textbf{Cost trajectory.} The 2,000-fold cost collapse described in Chapter 6
continued to accelerate: within 2025 alone, a further 45x reduction (from \$1.69
to four cents per block). Roughly 10x per year, driven by algorithmic
improvements, GPU optimization, and competition. A real-time proving cluster
runs \$60,000-\$100,000: sixteen RTX 5090 GPUs (\textasciitilde\$32K),
dual-socket server, 512 GB RAM. For less than the price of a suburban house, you
can prove every Ethereum block in real time.

\textbf{The Witness Gap grows with acceleration.} As GPU provers drove
cryptographic proving time down 10-50x, witness generation -- still CPU-bound --
became the dominant bottleneck. The proportional shift described in Chapter 4 is
now the defining structural constraint of zkVM performance: witness generation
in a zkVM equals full VM emulation, which resists the parallelism that NTT and
MSM exploit so effectively. The magician's backstage preparation now takes
longer than sealing the proof.

Active optimization research is attacking this gap from multiple directions.
ZKPOG achieved up to 52x speedup by moving witness generation onto GPUs. OpenVM
2.0's SWIRL prover includes an ahead-of-time compiler executing at 3.8 GHz,
eliminating JIT overhead. Nexus 3.0 runs the program twice -- first to gather
memory statistics, then to produce an optimized trace.

\textbf{The EF security pivot (December 2025).} The Ethereum Foundation declared
the speed race won and shifted focus: - May 2026 target: 100-bit provable
security across all zkEVM teams - December 2026 target: 128-bit provable
security, sub-300 KB proofs - New primary metric: energy per proof (kWh),
replacing raw speed - The EF rejects unproven conjectures (proximity gap
assumptions) for production soundness

This pivot validates systems with strong formal security guarantees. It also
shifts the competitive axis from ``who can prove fastest'' to ``who can prove
most securely'' -- precisely where lattice-based and tensor-code approaches hold
structural advantages. The race for speed is over. The race for rigor has begun.

\section{RISC-V Convergence}\label{risc-v-convergence}

The numbers are unambiguous: eight of ten major zkVMs target RISC-V. The
original three-philosophy taxonomy -- EVM-Compatible, ZK-Native ISA,
General-Purpose ISA -- is accurate as a historical classification, but the
market has rendered its verdict. RISC-V won the general-purpose category
decisively. Even EVM-focused projects now build RISC-V backends and layer EVM
compatibility on top.

Why RISC-V? Three reasons converge.

First, RISC-V's register-transfer architecture maps cleanly onto tabular
execution traces, which are the native input format for AIR and lookup-based
arithmetization. A RISC-V instruction reads source registers, performs an
operation, and writes a destination register -- exactly one row in a trace
table.

Second, RISC-V's compiler ecosystem is decades deep. Any Rust, C, or C++ program
can be compiled to RISC-V using standard LLVM toolchains. This means millions of
existing programs become provable without modification. The universal stage
accepts any act, because any act can be translated into its language.

Third, RISC-V is open and royalty-free. Unlike ARM (proprietary) or x86
(legacy-encumbered), RISC-V has no licensing costs and no vendor lock-in. For an
open-source ecosystem, this matters.

The holdouts are instructive. Cairo (Stwo) is a ZK-native ISA designed to
minimize arithmetization cost -- the ISA \emph{is} the constraint system. This
gives Cairo a structural efficiency advantage: the compiler optimization study
(Gassmann et al., 2025) found that standard LLVM optimizations yield over 40\%
improvement on RISC-V zkVMs because they target hardware features (caches,
branch predictors) absent in ZK contexts. Cairo avoids this overhead by design.
Whether that advantage justifies a smaller developer ecosystem is the strategic
question StarkWare has answered with ``yes'' for Starknet and ``maybe not'' for
broader adoption (hence Kakarot's EVM-on-Stwo path).

zkWASM (Delphinus Lab) targets WebAssembly, offering the broadest language
support of any zkVM (any language that compiles to WASM). Its PLONKish/KZG
architecture is a generational outlier -- it uses pairing-based commitments
rather than hash-based -- and it has not demonstrated Ethereum block proving.
zkWASM appears to be a niche play for web-native applications rather than a
contender for the mainstream proving market.

\chapter{Chapter 12: Midnight -- The Privacy
Theater}\label{chapter-12-midnight-the-privacy-theater}

Every magic show needs a theater -- a physical space designed so that the
audience sees exactly what the magician intends and nothing more. The lighting,
the curtains, the sight lines, the trapdoors: all are engineered to control
information flow. What the audience sees is a choice, not an accident. What the
audience does not see is an architecture, not an oversight.

Midnight is such a theater. Every architectural choice -- from the language to
the proof system to the token model -- serves a single design goal: the audience
learns the truth of a claim, and nothing else. The question that animates this
chapter is whether the theater's engineering matches its ambition.

\section{Midnight as Test Case}\label{midnight-as-test-case}

Most ZK systems use zero-knowledge proofs as an optimization -- a way to
compress computation for cheaper on-chain verification. The magician's sealed
certificate is a convenience. Midnight is different. On Midnight, ZK proofs are
not an optimization layer bolted onto an existing execution model. They
\emph{are} the execution model. Every state transition is proven in zero
knowledge. Every contract deployment, every circuit call, every token transfer
passes through a local proof server before reaching the chain. The trick is not
incidental to the show. The trick \emph{is} the show.

This makes Midnight an unusually complete test of whether the seven-layer model
actually maps to a working system. With 473 pages of verified documentation
across five reference documents -- Developer Guide (191pp), Compact Language
Reference (75pp), ZKIR Specification (60pp), MidnightJS SDK Reference (87pp),
and Wallet SDK Reference (60pp) -- Midnight provides concrete evidence at every
layer. Where the earlier chapters' examples are necessarily abstract, Midnight
supplies specific opcodes, measured latencies, deployed contracts, and compiler
error messages.

The analysis that follows draws on findings from seven dedicated layer analysts.
Every claim references specific Midnight documentation or measured behavior.

\section{Midnight at a Glance}\label{midnight-at-a-glance}

Midnight is a Cardano sidechain that executes privacy-preserving smart contracts
written in \textbf{Compact}, a TypeScript-inspired domain-specific language. The
architecture follows a pipeline:

\begin{verbatim}
Compact source --> compactc compiler --> ZKIR circuits + TypeScript bindings + proving keys
                                              |                    |                |
                                         Proof server         DApp frontend    On-chain verifier
                                        (localhost:6300)      (browser/node)   (blockchain node)
\end{verbatim}

The compiler produces three artifacts from a single \texttt{.compact} file: ZKIR
circuit descriptions in JSON, TypeScript API bindings for the DApp frontend, and
cryptographic proving/verifier key pairs. This three-part output reflects the
fundamental architecture of privacy-preserving computation: what can be proven
(ZKIR), what runs privately (TypeScript witnesses), and what makes proofs
possible (keys).

Midnight's token model has three layers. \textbf{NIGHT} is the governance and
staking token, always unshielded (transparent). \textbf{DUST} is the fee token,
a public (unshielded) token per the wallet SDK, generated from staking NIGHT
over time, with a balance computed from generation parameters. \textbf{Custom
tokens} can be either shielded (encrypted, spent via ZK proofs and nullifiers)
or unshielded (transparent, spent via BIP-340 Schnorr signatures), at the
developer's choice per UTXO.

\section{Full Seven-Layer Mapping}\label{full-seven-layer-mapping}

\subsection{Layer 1: BLS12-381 and the Trusted
Ceremony}\label{layer-1-bls12-381-and-the-trusted-ceremony}

The book's Layer 1 asks: trusted or transparent? Midnight answers:
\textbf{trusted, with a universal SRS}.

Midnight uses the BLS12-381 elliptic curve with a PLONK-family proof system
(Halo 2). The Developer Guide's architecture diagram (p.8) shows a dedicated
\texttt{midnight-trusted-setup} repository for running a Powers-of-Tau ceremony.
The SRS is universal -- one ceremony serves all circuits -- but the compiler
generates per-circuit proving and verifier keys derived from that SRS. For a
simple counter contract, the compiler produces an \texttt{increment.prover} key
(13.7 KB) and an \texttt{increment.verifier} key (1.3 KB).

Midnight's cryptographic history reveals pragmatic trade-offs that echo the
theme from Chapter 2. The system originally used Pluto-Eris curves for recursive
proof composition but switched back to BLS12-381 for mainnet, citing faster
proof generation, wider ecosystem compatibility, and support for higher
transaction volumes. Theoretical optimality yielded to engineering reality -- a
pattern we have seen repeatedly throughout this book.

\textbf{Post-quantum implication}: BLS12-381 provides approximately 128-bit
classical security but zero post-quantum security. Shor's algorithm breaks the
discrete logarithm and pairing assumptions simultaneously. The stage is strong
today. Its quantum shelf life is finite.

\subsection{Layer 2: Compact as Fourth-Philosophy
DSL}\label{layer-2-compact-as-fourth-philosophy-dsl}

The three-philosophy taxonomy from earlier chapters -- EVM-Compatible, ZK-Native
ISA, General-Purpose ISA -- does not accommodate Compact. Compact represents a
\textbf{fourth philosophy}: the application-specific DSL.

Compact does not prove a processor. It compiles a domain-specific smart contract
language directly to chain-specific ZKIR, with the compiler enforcing
application-level invariants that no ISA-level approach can express. Its closest
relatives are Leo (Aleo) and Mina's o1js.

The critical differentiator is \textbf{disclosure analysis}. The Compact
compiler includes a \texttt{track-witness-data} analysis pass that traces
witness values through all program paths and rejects any program where private
data might reach public surfaces without explicit \texttt{disclose()} calls.
This is a hard compile-time error, not a warning. The Developer Guide documents
that a first attempt at private voting was rejected with 11 disclosure errors,
each tracing the path from witness to ledger operation.

No other ZK language prevents accidental privacy leaks at compile time. In
Circom, Noir, and Cairo, a disclosure mistake produces a privacy leak, not a
compiler error. Compact eliminates an entire vulnerability class -- accidental
disclosure -- at the language level. The stage manager locks the doors to
prevent accidental reveals. The magician must explicitly ask for the key.

\subsection{Layer 3: The disclose()
Boundary}\label{layer-3-the-disclose-boundary}

Witnesses in Midnight are arbitrary JavaScript functions that run off-chain. The
\texttt{disclose()} operator is the sole gateway from the witness world to the
circuit world -- the single controlled opening in the curtain between backstage
and audience.

Every ZKIR circuit has two transcript channels: the \texttt{publicTranscript}
(ledger operations visible to the verifier) and
\texttt{privateTranscriptOutputs} (witness values visible only to the prover).
Without \texttt{disclose()}, witness values remain entirely invisible to the
circuit and the chain. The ZKIR checker enforces transcript integrity with
specific diagnostic error messages -- tampering with either transcript causes
immediate rejection.

What does this boundary feel like to the developer who must work inside it every
day? Consider a programmer writing a private voting contract. She writes her
witness function in TypeScript -- ordinary, familiar, unexotic code that fetches
voter eligibility from a local database and computes a ballot. Every variable in
that function is invisible to the chain by default. She could write a hundred
lines of witness logic, and none of it would leave her machine. Then she reaches
the moment of commitment: she needs the circuit to know which candidate received
the vote, without revealing who cast it. She writes
\texttt{disclose(candidateId)}. That single call is the seam in the curtain, the
controlled opening where one piece of information -- and only that piece --
crosses from the private rehearsal room into the public theater. The compiler
has already analyzed every path through her code; if she accidentally wrote
\texttt{disclose(voterId)} three functions earlier, the build would have failed
with a traced error showing exactly how the private value reached the public
surface. The developer experience is not one of navigating cryptographic
abstractions. It is one of writing normal code inside a system that has opinions
-- strong, enforced, non-negotiable opinions -- about what leaves the room.
Asimov imagined robots governed by laws they could not violate. The Compact
developer works inside a compiler governed by disclosure laws it will not bend.
The fiction writer's dream of the incorruptible guardian is, in this narrow
domain, an engineering reality.

The theater analogy sharpens here. In a well-designed theater, the lighting grid
is the real enforcer of what the audience sees. A spotlight operator who
accidentally swings a beam toward the wings will reveal the stagehands, the
props not yet in play, the illusion's scaffolding. Midnight's disclosure
analysis is the lighting grid: it does not merely suggest where the light should
fall -- it physically prevents the spots from swinging toward the wings. The
developer sets the cues. The compiler locks the grid. And when the show runs,
only what was meant to be seen is seen.

\textbf{Side-channel gap}: None of the five reference documents address timing
attacks, cache attacks, or metadata leakage through the indexer's GraphQL API.
Proof generation dominates transaction time (\textasciitilde18-20 seconds),
providing natural but unintentional timing padding. The curtain is thick, but no
one has tested whether light leaks through the seams. A theater may control its
spotlights perfectly and still betray its secrets through the sound of trapdoors
opening, the vibration of machinery beneath the stage, the draft of air from a
hidden passage. Midnight's formal privacy model covers what the proofs reveal.
It does not yet cover what the infrastructure whispers.

\subsection{Layer 4: ZKIR as High-Level Constraint
IR}\label{layer-4-zkir-as-high-level-constraint-ir}

Midnight's ZKIR is a typed instruction-level intermediate representation with 24
base instructions organized into 8 categories: arithmetic (\texttt{add},
\texttt{mul}, \texttt{neg}), constraints (\texttt{assert},
\texttt{constrain\_eq}, \texttt{constrain\_bits},
\texttt{constrain\_to\_boolean}), comparison (\texttt{test\_eq},
\texttt{less\_than}), control flow (\texttt{cond\_select}, \texttt{copy}), type
encoding, cryptographic operations (targeting the Jubjub curve embedded in
BLS12-381), I/O, and division.

ZKIR sits above the mathematical constraint formalism but below the source
language. In the taxonomy of R1CS, AIR, PLONKish, and CCS, ZKIR is most
analogous to PLONKish but operates at a higher abstraction level -- the
developer sees typed operations with semantic meaning, not bare multiplication
gates. This is arithmetic with a human face.

\subsection{Layer 5: Halo 2 and the Four-Phase
Pipeline}\label{layer-5-halo-2-and-the-four-phase-pipeline}

Midnight uses Halo 2 (UltraPlonk) over BLS12-381. The proof server runs locally
at \texttt{localhost:6300}. The transaction lifecycle follows a four-phase
pipeline: \texttt{callTx()} (execute circuit locally) to \texttt{proveTx()}
(generate ZK proof, \textasciitilde17-24 seconds) to \texttt{balanceTx()} (bind,
sign, merge) to \texttt{submitTx()} (submit to chain, where the node verifies
the proof).

Proof generation dominates latency: deploy \textasciitilde17-28 seconds, circuit
call \textasciitilde17-24 seconds, balancing and submission sub-second. The
magician rehearses for twenty seconds. The audience's verdict takes a heartbeat.

\subsection{Layer 6: BLS12-381, Jubjub, and
Poseidon}\label{layer-6-bls12-381-jubjub-and-poseidon}

All ZKIR values are elements of the BLS12-381 scalar field (\(\sim 2^{253}\)).
The Jubjub twisted Edwards curve, embedded natively in BLS12-381, enables
in-circuit elliptic curve operations. Poseidon-based hashing via
\texttt{persistent\_hash} (two field elements for collision resistance) and
\texttt{transient\_hash} (single field element) provides ZK-friendly hash
computation.

The large field (253 bits vs.~31-64 bits for STARK-friendly fields) is the
primary performance cost -- 10-100x slower per operation -- but is necessary for
the pairing and embedded curve that enable KZG commitments and constant-size
proofs. Midnight pays for privacy with patience.

\subsection{Layer 7: Three Tokens and the Verifier Key
Lifecycle}\label{layer-7-three-tokens-and-the-verifier-key-lifecycle}

Midnight's three-token model (NIGHT/DUST/custom) and per-UTXO privacy choice
create a deployment architecture distinct from Ethereum-style gas economics.
DUST fees are paid in a public token generated by time-locked staking.

The economics deserve closer scrutiny, because they encode a philosophy. On
Ethereum, gas is purchased with ETH on the open market -- every fee payment is a
visible transaction, a data point for chain analysts, a signal. On Midnight,
DUST accrues silently from staked NIGHT over time, like interest accumulating in
an account the holder never visits. A developer deploying a contract pays
\textasciitilde490 trillion SPECK per circuit call, but that DUST was not
purchased in a transaction anyone can observe. It was generated by the passage
of time and the act of staking. The economic consequence is that fee payment
itself becomes a poor signal: an observer who watches DUST expenditures sees
activity, but cannot easily link that activity back to a market purchase, a
wallet funding event, or an exchange withdrawal. The three-token architecture is
not merely a governance convenience. It is a privacy mechanism at the economic
layer -- a recognition that in a system where computation is private, the
payment for computation must be private too, or else the ticket stub betrays the
show. Penrose once observed that the geometry of a space constrains what can
happen within it. Midnight's token geometry -- the separation of governance
(NIGHT), fees (DUST), and application value (custom tokens) -- constrains the
information that economic activity can leak. The shape of the money shapes the
privacy of the system. This is Layer 7 not as an afterthought but as
architecture.

The SDK API includes functions for dynamically managing verifier keys on-chain
(\texttt{submitInsertVerifierKeyTx}, \texttt{submitRemoveVerifierKeyTx}),
raising governance questions that parallel the book's discussion of upgradeable
proxy contracts. The theater management retains the ability to change the locks
-- to swap out the verifier that guards a contract's stage door, to retire a
circuit and replace it with another. In a traditional theater, this is the
producer's prerogative: the show can be recast, the set redesigned, the script
rewritten between seasons. But in a privacy theater, changing the verifier key
changes the terms under which secrets were originally committed. A user who
shielded tokens under one verifier's rules may find those rules altered by a
governance action she never approved. The trapdoor that was locked for the
performer's protection can be unlocked by the theater's owner. This tension --
between upgradeable infrastructure and the immutability that privacy demands --
is not resolved in Midnight's current design. It is, at most, acknowledged.

\section{Where Midnight Validates the
Model}\label{where-midnight-validates-the-model}

The seven-layer decomposition maps cleanly to Midnight's architecture in five
places:

\textbf{1. Layer 1 maps to \texttt{midnight-trusted-setup}.} The ceremony
produces the SRS; the compiler derives per-circuit keys. The capex/opex framing
applies directly: one-time ceremony cost amortized across all contracts, with
per-transaction proof costs of \textasciitilde18 seconds and \textasciitilde490
trillion SPECK.

\textbf{2. Layer 2 maps to Compact.} The language choice determines what
developers can express and what mistakes they cannot make. Disclosure analysis
validates the argument that language design has security implications beyond
expressiveness.

\textbf{3. Layer 4 maps to ZKIR.} The 24-opcode instruction set is a concrete
instance of the arithmetization layer, sitting above PLONKish constraints and
below the source language.

\textbf{4. Layer 5 maps to the proof server.} The four-phase pipeline
instantiates the proof generation and verification layer with measured latencies
and a clear component boundary.

\textbf{5. Layer 7 maps to the three-token model.} The verifier key deployment,
fee economics, and governance structure are concrete instances of deployment
concerns.

\section{Where Midnight Challenges the
Model}\label{where-midnight-challenges-the-model}

Three aspects of Midnight's architecture do not fit cleanly into seven layers:

\textbf{1. The compiler spans Layers 2, 3, and 4 simultaneously.} The Compact
compiler's nanopass architecture takes source code (Layer 2), performs
disclosure analysis (a Layer 3 concern -- who sees the witness?), and emits ZKIR
(Layer 4 arithmetization) in a single continuous pipeline of 26 intermediate
languages. There is no clean boundary where ``language'' ends and
``arithmetization'' begins. The magician's script, the backstage preparation,
and the mathematical encoding blur into a single creative act.

\textbf{2. The SDK is a cross-layer orchestrator.} The four-phase transaction
pipeline spans Layer 3 (witness construction in \texttt{callTx}), Layer 5 (proof
generation in \texttt{proveTx}), Layer 6 (BIP-340 signatures in
\texttt{balanceTx}), and Layer 7 (on-chain submission in \texttt{submitTx}). The
SDK is not ``at'' any single layer; it is the glue connecting all of them.

\textbf{3. Privacy is not a layer -- it is a cross-cutting concern.} The book
treats privacy primarily as a Layer 3 issue (witness secrecy). In Midnight,
privacy decisions propagate through every layer: BLS12-381's pairing structure
enables shielded commitments (Layer 1/6), disclosure analysis enforces privacy
at the language level (Layer 2), the two-transcript model separates public and
private data (Layer 3/4), the proof server keeps witnesses local (Layer 5), and
the shielded/unshielded UTXO model determines what the verifier sees (Layer 7).
Privacy is not a room in the theater. It is the architecture of the theater
itself.

The following table crystallizes what the Midnight case study proves and what it
leaves unresolved:

\begin{longtable}[]{@{}
  >{\raggedright\arraybackslash}p{(\linewidth - 4\tabcolsep) * \real{0.2000}}
  >{\raggedright\arraybackslash}p{(\linewidth - 4\tabcolsep) * \real{0.3636}}
  >{\raggedright\arraybackslash}p{(\linewidth - 4\tabcolsep) * \real{0.4364}}@{}}
\toprule\noalign{}
\begin{minipage}[b]{\linewidth}\raggedright
Dimension
\end{minipage} & \begin{minipage}[b]{\linewidth}\raggedright
Midnight Validates
\end{minipage} & \begin{minipage}[b]{\linewidth}\raggedright
Midnight Does Not Solve
\end{minipage} \\
\midrule\noalign{}
\arrayrulecolor{MidnightBlue}\midrule\arrayrulecolor{MidnightBorder}
\endhead
\bottomrule\noalign{}
\endlastfoot
Compiler-enforced privacy & Disclosure analysis catches 11 error types at
compile time & Side-channel leakage (timing), metadata via indexer/network
layer \\
Seven-layer decomposition & Clean mapping at 5 of 7 layers & Compact compiler
spans L2--L4; SDK spans L3--L7; layers not cleanly separable \\
Privacy as architecture & Every layer serves privacy -- UTXO model, local
proving, shielded state & Cross-contract token transfers between DApps remain
unsupported \\
Trust decomposition & Three-token model (Night/Shielded/DUST) is genuinely novel
& Governance key management and upgrade path remain centralized \\
Post-quantum readiness & Architecture acknowledged as non-PQ & No migration path
to lattice-based commitments without full redesign \\
Developer experience & Compact→ZKIR→proof pipeline is coherent end-to-end &
17-28s proof times create UX barrier; no GPU acceleration \\
\end{longtable}

\section{The Privacy Theater Analogy}\label{the-privacy-theater-analogy}

The magician metaphor that has sustained this book finds its fullest expression
in Midnight. But to see the analogy in its full depth, we must think not as the
audience watching the trick but as the architect who designed the theater -- the
one who decided where the walls would go, where the sight lines would converge,
where the trapdoors would open, and which doors would lock from the inside.

The \textbf{open stage} (unshielded) faces the audience directly. Performers
(transactions) are fully visible. The audience (validators) can see every
movement, verify every step. NIGHT tokens live here -- governance requires
transparency. The audience sees the magician's face. The lighting is full and
even, the equivalent of a bare bulb in a police interrogation room. Nothing
hides. Nothing can.

The \textbf{curtained stage} (shielded) is separated by a one-way mirror. The
audience cannot see the performers, but they can hear the music (the ZK proof)
and verify that it follows the score (the ZKIR circuit). Shielded tokens live
here. The performer's identity (UTXO owner) and movements (transaction values)
are hidden, but the proof guarantees the performance was legitimate. The trick
happens in the dark, and the sealed certificate emerges into the light. The
sight lines have been engineered so that from every seat in the house, the
audience sees only the proof -- a single sheet of paper slid under the curtain.
The geometry of the theater makes curiosity futile. There is simply no angle
from which the backstage is visible.

The \textbf{stage manager} (Compact compiler) enforces the rules. A performer
cannot accidentally step from behind the curtain onto the open stage -- the
compiler's disclosure analysis physically prevents the transition unless the
performer explicitly calls \texttt{disclose()}. This is not a best-practice
recommendation posted backstage; it is a locked door that requires a key. The
stage manager does not trust the performers to remember the blocking. She has
bolted the doors, wired the lighting grid to fixed positions, and removed the
handles from the wrong side. The performers can improvise within their space.
They cannot improvise their way into the audience's view.

The \textbf{proof server} is the rehearsal room. All practice (witness
computation) happens here, in private, at localhost:6300. Only the final
performance (the ZK proof) reaches the theater. The audience never sees the
rehearsal. The rehearsal room has no windows, no microphones, no cameras. It
exists on the performer's own machine, in a process that never opens a network
socket to anything but the local proof server. The separation is not a policy.
It is a wall.

The \textbf{trapdoors} are the governance mechanisms -- the verifier key
management functions, the upgrade paths, the administrative controls that the
SDK exposes. Every theater has trapdoors, and every trapdoor is a dual-use
technology: it enables the performer to appear and disappear as the trick
requires, but it also enables the theater owner to access spaces the audience
was promised would remain sealed. In Midnight's current architecture, the
trapdoors exist. They are documented. Whether they are adequately governed is a
question the documentation does not answer.

\textbf{DUST} is the ticket price. Every performance requires a ticket,
generated by staking NIGHT tokens over time. The ticket price is uniform for
same-type performances (\textasciitilde490 trillion SPECK per circuit call),
providing some metadata privacy -- you cannot tell what happened behind the
curtain by looking at the ticket stub. The tickets are not sold at a box office
window where a clerk might remember your face. They materialize in your wallet
through the silent mechanics of staking -- as if the theater rewarded loyal
patrons by slipping tickets under their doors in the night.

Penrose argued that the geometry of spacetime is not a backdrop against which
physics happens but the thing that \emph{is} physics -- that curvature and
matter are the same story told in two languages. Midnight's privacy theater
operates on a similar principle. The privacy is not a feature applied to a
blockchain. It is the shape of the blockchain itself. The curve choice, the
language design, the transcript separation, the token model, the proof locality
-- these are not independent decisions that happen to support privacy. They are
the curvature of Midnight's operational space, and privacy is what that
curvature produces. To remove privacy from Midnight would not be to disable a
feature. It would be to flatten the geometry, and the theater would cease to be
a theater at all.

\section{Five Lessons for ZK System
Design}\label{five-lessons-for-zk-system-design}

\subsection{Lesson 1: Privacy as Cross-Cutting
Concern}\label{lesson-1-privacy-as-cross-cutting-concern}

The book treats privacy primarily as a Layer 3 phenomenon. Midnight demonstrates
that privacy is an architectural decision at every layer: curve selection (L1),
language design (L2), witness boundaries (L3), transcript separation (L4), proof
locality (L5), commitment schemes (L6), and UTXO encryption (L7). A reader
applying the seven-layer model to Midnight would need to trace privacy through
all seven layers to understand the system's guarantees.

The OSI network model offers a precedent: security is not a layer but a property
that each layer must independently maintain. ZK privacy deserves the same
treatment.

\subsection{Lesson 2: The ``Compiler Protects You''
Philosophy}\label{lesson-2-the-compiler-protects-you-philosophy}

The under-constrained vulnerability epidemic described in Chapter 3 is the
dominant failure mode in ZK systems. Compact's disclosure analysis addresses an
equally severe class -- accidental disclosure -- at the language level. But
Compact's protection has limits: it prevents accidental leakage, but it cannot
prevent a developer from choosing to \texttt{disclose()} too much, storing
secrets in ledger state, or making application-logic errors. The locked door
keeps you from stumbling through. It does not stop you from handing someone the
key.

Compiler enforcement raises the floor of security but does not guarantee
correctness. The comparison between Circom's manual constraint authoring (where
under-constrained bugs thrive) and Compact's automatic constraint generation
(where disclosure bugs are caught) illustrates the Layer 2 security spectrum
concretely.

\subsection{Lesson 3: The Three-Token Economic
Model}\label{lesson-3-the-three-token-economic-model}

The book's Layer 7 focuses on Ethereum-style gas economics. Midnight's
three-token model represents a fundamentally different architecture: fees are
paid in a shielded token generated by time-locked staking, not purchased on the
open market. This means fee payment itself has privacy properties -- you cannot
determine a user's transaction volume by observing fee purchases. Layer 7
economics are not just about gas costs but about the information leakage of the
fee mechanism itself. Even the ticket stub can be a clue, and Midnight's design
tries to minimize what it reveals.

\subsection{Lesson 4: The Application-Specific DSL as Fourth
Philosophy}\label{lesson-4-the-application-specific-dsl-as-fourth-philosophy}

Compact demonstrates that a domain-specific ZK language can achieve properties
impossible for general-purpose approaches: compiler-enforced privacy boundaries,
first-class blockchain state, integrated token operations, and a unified DApp
development pipeline. The trade-off is vendor lock-in -- Compact contracts
cannot run on any chain except Midnight. The three-philosophy taxonomy should be
expanded to include this fourth philosophy.

\subsection{Lesson 5: ZKIR as Concrete Layer
4}\label{lesson-5-zkir-as-concrete-layer-4}

The book's Layer 4 discussion of R1CS, AIR, PLONKish, and CCS is necessarily
abstract. ZKIR provides a concrete example that sits above these mathematical
abstractions. Its 24 typed instructions with semantic meaning (not just
``multiplication gate'' but \texttt{persistent\_hash} and
\texttt{private\_input}) show that production arithmetization layers carry more
structure than the mathematical formalism suggests. The puzzle has names for its
pieces.

\subsection{Maturity Assessment}\label{maturity-assessment}

Midnight is best characterized as a late-stage testnet / early mainnet system.
The proof system works, the compiler catches real privacy bugs, and the devnet
supports end-to-end contract deployment and execution. However, cross-contract
token transfers fail with SDK errors, the \texttt{\textgreater{}} and
\texttt{\textless{}=} operators have a documented compiler bug, and deployment
latency (dominated by proof generation, as detailed in the Layer 5 section
above) indicates room for proving optimization. On the L2Beat Stages framework,
Midnight would sit at approximately Stage 0-1: operational with ZK proofs
providing validity guarantees, but with governance mechanisms retaining
significant centralized control.

The theater is built. The rehearsals are underway. The opening night has not yet
arrived.

Midnight is one theater. The zero-knowledge ecosystem has built dozens more --
each with different stages, different audiences, different trust bargains. The
next chapter surveys six market segments where the mathematics meets money, and
asks the question that every technology must eventually answer: who is buying
tickets, and what do they think they are paying for?

\chapter{Chapter 13: The Market
Landscape}\label{chapter-13-the-market-landscape}

The magician has performed. The audience has verified. The mathematics works.
But a magic show that no one attends is just a person alone in a room doing card
tricks. The most elegant proof system in the world, deployed on the most secure
stage, verified by the most rigorous audience, means nothing if no one pays for
a ticket.

The question that determines whether zero-knowledge proofs reshape industries or
remain a cryptographer's curiosity is not ``does the math work?'' but ``does
anyone pay for it?'' And the question that determines whether the money is
well-spent is sharper still: \emph{where is trust actually being minimized --
and where is it merely being moved?}

This chapter follows the money. But it watches where the trust goes.

\section{ZK Rollups: The Proving Grounds
(Production)}\label{zk-rollups-the-proving-grounds-production}

Zero-knowledge rollups are where the thesis of this chapter becomes immediately
testable. The mechanism is straightforward: execute transactions off-chain,
generate a ZK proof of correct execution, and post the proof plus compressed
state data to Ethereum for verification. Ethereum-grade security with 10-100x
lower transaction costs. At Layer 7, the audience can verify every proof. The
math is sound. The trick works.

But look at the governance. Every major rollup in production today operates at
Stage 0 or Stage 1 -- meaning a security council, a multisig, or a governance
committee can override the proof system. The audience can check the math, but a
committee can still replace the audience. Trust has been minimized at the
cryptographic layer and preserved, nearly intact, at the institutional layer.
This is not a criticism. It is a fact about where we are, and it should be
stated plainly.

The market has consolidated rapidly, and one of its most expensive lessons has
already been paid. Polygon zkEVM, once the flagship example of EVM-compatible ZK
rollups (Philosophy A in the original taxonomy), was shut down in 2025/2026
after approximately \$250 million in investment. The core team, led by
co-founder Jordi Baylina, spun off to found ZisK. A quarter-billion dollars
bought a lesson: in a field moving this fast, the first-mover advantage is often
a first-mover trap.

The current production leaders tell a more encouraging story -- but the trust
question follows them onto every stage:

\textbf{Scroll} (\$748M TVL as of early 2026) -- the largest zkEVM by total
value locked. Scroll uses a halo2-based proof system with KZG commitments,
generating PLONKish proofs that are verified directly on Ethereum L1. Seven
hundred forty-eight million dollars locked means users trust the math with real
money. Whether they have examined the governance is a different question.

\textbf{Linea} (\$2B TVL) -- ConsenSys's zkEVM, targeting full EVM equivalence.
Linea uses a custom prover with Fiat-Shamir-based PLONK and posts compressed
proofs to Ethereum. Two billion dollars locked. That is not a pilot program. It
is also not a trustless one.

\textbf{Starknet} -- powered by Stwo, the Circle STARK prover that went live on
mainnet in November 2025. Starknet stands apart: it uses the Cairo ISA rather
than EVM compatibility, and its STARK proofs are verified natively on L1 without
Groth16 wrapping (accepting larger proof sizes in exchange for transparency).
The only major rollup that performs the entire trick on a glass stage -- no
trusted setup, no opaque wrapping. The trust tradeoff is different here: you
trust the transparency of the mathematics, and you pay for it in proof size.

\textbf{ZKsync Era} -- Matter Labs' zkEVM, powered by the Airbender prover.
ZKsync achieved 21.8 million cycles per second on a single H100 GPU and deployed
via the Atlas upgrade. The 2026 roadmap targets formal verification and adoption
as a ``universal standard'' for ZK proving.

The economics have improved sharply. The 2,000-fold cost collapse described in
Chapter 6 means that at current rates, continuously proving every Ethereum block
costs roughly \$102,000 per year -- less than the GPU cluster required to do it.
Post-Pectra (May 2025) and Fusaka (December 2025), data availability costs have
also dropped, with blob capacity increasing 8x via PeerDAS. The trick has become
cheaper than the theater's electricity bill. But cost is not the same as trust.
A cheap proof verified by a governance council you cannot remove is still a
proof verified by a governance council you cannot remove.

\textbf{Trust relocated from:} L1 execution validators \textbf{to:} proof system
soundness + governance multisig. \textbf{Net:} genuine minimization for
computation integrity; governance remains the binding constraint.

\section{ZK Coprocessors: Off-Chain Computation, On-Chain Verification
(Growth)}\label{zk-coprocessors-off-chain-computation-on-chain-verification-growth}

ZK coprocessors represent a category that the original seven-layer framework
entirely missed -- and the trust dynamics are subtle enough to deserve careful
attention. A coprocessor allows a smart contract to verifiably read and compute
over historical blockchain data without the gas cost of executing that
computation on-chain. The contract sends a query; the coprocessor executes it
off-chain, generates a ZK proof of correct execution, and returns the result
with the proof for on-chain verification.

The magician, in this case, is not performing for a general audience. She is
performing for a smart contract -- an audience that cannot reason, only verify.
The proof guarantees that the computation was done correctly. But the
computation happened somewhere: on specific hardware, operated by a specific
entity, with access to the query inputs. Trust has not been eliminated. It has
been moved off-chain.

\textbf{Axiom} raised \$20 million in Series A funding and operates the leading
ZK coprocessor platform. Axiom's coprocessor can access any historical Ethereum
block header, account state, storage slot, transaction, or receipt, generate a
computation over that data, and deliver a verified result to a smart contract in
a single callback.

\textbf{Brevis} operates both a ZK coprocessor and the Pico Prism zkVM. Brevis's
ProverNet launched mainnet beta in December 2025, enabling decentralized proof
generation for coprocessor queries.

\textbf{Lagrange} provides ZK coprocessing for cross-chain state proofs and data
availability verification.

The coprocessor market matters because it extends ZK proofs beyond transaction
execution into data analytics. A DeFi protocol can use a coprocessor to compute
time-weighted average prices over 30 days of on-chain data, prove the
computation is correct, and use the result in a lending decision -- all without
trusting an oracle. The sealed certificate now covers not just ``the computation
was correct'' but ``the data was real.'' That is a genuine advance. But who
watches the prover's hardware? The proof says the computation was faithful. It
does not say the machine was honest about what it was asked to compute. The
trust has shifted, not vanished.

\textbf{Trust relocated from:} on-chain computation (expensive, public)
\textbf{to:} off-chain proving (cheap, private) + data availability assumptions.
\textbf{Net:} cost reduction with new trust in prover correctness and DA layer.

\section{ZKML: Provable Machine Learning
(Research)}\label{zkml-provable-machine-learning-research}

Here the trust question becomes genuinely uncomfortable. Zero-knowledge proofs
for machine learning inference -- proving that a specific model produced a
specific output on a specific input without revealing the model's weights -- can
prove that a neural network ran faithfully. They cannot prove the neural network
should have run at all.

Read that distinction carefully. ZKML proves the inference was correct
\emph{given the model}. It says nothing about whether the model is fair, whether
it was trained on representative data, whether its architecture is appropriate
for the task, or whether deploying it is wise. You can prove the magician
performed the trick exactly as rehearsed. You cannot prove the trick was worth
performing. Trust in the computation is not trust in the model's correctness,
and conflating the two is the central danger of this category.

The fundamental technical difficulty is that neural network operations (matrix
multiplications, activation functions, normalization) map poorly onto the finite
field arithmetic that ZK proof systems require.

\textbf{Lagrange's DeepProve} is the current leader, claiming 700x faster ZK
proofs for ML inference compared to previous approaches. DeepProve achieves this
through specialized arithmetization for neural network operations, avoiding the
general-purpose overhead of running ML models inside a generic zkVM.

\textbf{EZKL} provides an open-source toolkit for generating ZK proofs of neural
network inference, targeting the halo2 proof system. EZKL converts ONNX models
into ZK circuits.

To appreciate the difficulty, consider what proving a single inference of a
modest neural network (say, a 10-layer transformer with 100 million parameters)
actually requires. Each matrix multiplication involves billions of field
multiplications. Activation functions like ReLU require comparison operations
that are expensive in finite fields. Layer normalization involves division and
square roots, both of which must be decomposed into constraint-friendly
operations. The overhead tax from Chapter 5 -- already 10,000-50,000x for
general computation -- can be even steeper for ML workloads, because neural
network arithmetic is optimized for floating-point hardware that has no analogue
in finite field circuits. DeepProve's 700x improvement is impressive precisely
because the starting point was so far from practical.

The ZKML market is pre-revenue and research-heavy, but the applications are
clear: verifiable AI inference (prove that a content moderation decision was
made by a specific model), model privacy (prove model performance without
revealing weights), and auditable AI (regulatory compliance for model
decisions). In a world increasingly shaped by opaque AI systems, the ability to
prove properties of a model without revealing the model itself may become one of
the most valuable applications of zero-knowledge proofs. But only if the field
resists the temptation to let ``proven inference'' stand in for ``trustworthy
AI.'' The proof is one layer. The model is another. And no amount of
cryptographic elegance in the first layer compensates for negligence in the
second.

\textbf{Trust relocated from:} model operator's self-attestation \textbf{to:}
proof of inference correctness. \textbf{Net:} promising but pre-revenue; trust
in model integrity itself remains unaddressed.

\section{ZK Identity (Growth / Regulatory
Mandate)}\label{zk-identity-growth-regulatory-mandate}

This is the clearest win. In every other market segment, trust is being moved,
delegated, or repackaged. In identity, the magician's trick genuinely replaces a
gatekeeper. Before ZK identity, proving you were over 18 required handing your
driver's license to a stranger -- trusting them not to memorize your address,
your full name, your date of birth. The trust was total, and it was demanded by
a bouncer. ZK-based selective disclosure replaces that entire interaction with
mathematics: proving attributes about yourself (age over 18, citizenship of a
specific country, possession of a valid credential) without revealing the
underlying data. Trust shifted from institutions to proofs. Not trust minimized.
Trust \emph{replaced}.

\textbf{World} (formerly Worldcoin) uses iris-scanning orbs to generate unique
identity commitments, with ZK proofs enabling ``proof of personhood'' without
revealing biometric data. World's approach is controversial (biometric data
collection) but large-scale (millions of enrollments). The magician proves she
is human. The audience learns nothing else. But note the uncomfortable residue:
the orb that scans your iris is a piece of hardware operated by a company. The
proof is trustless. The enrollment is not. Even in the clearest case, trust has
a way of hiding in the infrastructure.

\textbf{EU eIDAS 2.0} mandates digital identity wallets for all EU citizens by
2026, with selective disclosure as a core requirement. ZK proofs are the leading
cryptographic mechanism for implementing the selective disclosure specified in
the regulation.

\textbf{Humanity Protocol} (\$1.1 billion valuation) focuses on palm-vein-based
biometric identity with ZK proofs for privacy-preserving verification.

The ZK identity market is projected to reach \$7.4 billion, driven by regulatory
mandates (eIDAS 2.0, GDPR's tension with blockchain transparency) and
institutional demand for verifiable credentials. When the law requires selective
disclosure and billions of citizens need identity wallets, the market does not
need to be created. It has been mandated. And of all the audiences the magician
now performs for, this one -- individual human beings trying to prove who they
are without surrendering who they are -- is the one that matters most.

\textbf{Trust relocated from:} credential presenter (showing full document)
\textbf{to:} credential issuer (signing the claim) + enrollment hardware.
\textbf{Net:} genuine minimization for disclosure; enrollment creates new trust
surface.

\section{Proving-as-a-Service: The Prover Market
(Production)}\label{proving-as-a-service-the-prover-market-production}

Here the privacy tradeoff from Chapter 4 returns with full force. Generating ZK
proofs is computationally expensive but highly parallelizable, so a natural
market has emerged: delegate your proving to someone with better hardware. The
efficiency gain is real. The trust consequence is this: the person who proves
for you sees your data.

Read that again. The architecture that protects your data the most --
client-side proving, where your secrets never leave your device -- requires
hardware that most people do not own. The architecture that works on any device
-- delegated proving -- requires trusting the prover with your secrets. Chapter
4 called this out as the economic structure of privacy in 2026. Here it becomes
a business model.

\textbf{Succinct} operates the leading prover network. By early 2026, the
Succinct Network had generated over 6 million proofs on mainnet, secured over
\$4 billion in value, and launched the \$PROVE token. The network's SP1
Hypercube zkVM powers multiple rollup deployments.

\textbf{RISC Zero Boundless} launched its open proof marketplace in September
2025. Boundless processed 542.7 trillion cycles by December 2025, then forced
migration from its centralized prover service to the decentralized marketplace.
This migration -- shutting down the centralized option to force adoption of the
decentralized alternative -- is a notable governance decision. Trust in a
centralized prover was replaced by trust in a marketplace mechanism. Whether
that is an improvement depends on what you think markets are.

\textbf{ZkCloud} (formerly Gevulot) provides production-ready proving
infrastructure with its Firestarter platform.

\textbf{Aligned Layer} takes a different approach: it uses EigenLayer restaking
(\$11 billion+ restaked ETH) to provide proof verification as a service,
allowing any application to submit proofs for verification without deploying its
own verifier contract.

The proving market is transitioning from an infrastructure cost (borne by rollup
operators) to a tradeable service (priced by market dynamics). The \$PROVE token
represents the first attempt to create a liquid market for computational
integrity. You can now buy and sell the ability to generate mathematical truth.
But the trust has been delegated to prover hardware, and the privacy question --
who sees the witness? -- has no market solution. It has only engineering
solutions, each with its own trust assumptions. TEEs, MPC-based proving,
client-side hardware acceleration: each moves the trust, none eliminates it.

The economics of proving-as-a-service deserve to be spelled out, because they
determine who actually pays for mathematical truth -- and at what margin. The
cost structure has three layers: hardware (GPU clusters, typically H100 or A100
machines running at \$2-3 per GPU-hour on cloud providers, less for owned
hardware with amortized capex), energy (a single H100 draws roughly 700W under
proving load; at industrial electricity rates of \$0.05-0.08/kWh, that is
\$0.035-0.056 per GPU-hour), and engineering (the prover software itself,
maintained by teams of 10-50 engineers at salaries that dwarf the hardware
costs). The first two are commodity inputs. The third is the moat.

Who pays? Today, the answer is straightforward: rollup operators pay. A ZK
rollup that generates proofs for every batch of transactions absorbs proving
costs as an operational expense, recouped through user transaction fees. At
current rates, proving costs for a major rollup run \$200,000-\$500,000 per year
-- significant, but a small fraction of the sequencer revenue for any rollup
with meaningful throughput. The user pays indirectly, through transaction fees
that embed a proving surcharge. That surcharge has been falling: the 2,000-fold
cost collapse described in Chapter 6 means the proving component of a typical L2
transaction fee has dropped from dollars to fractions of a cent. The user, in
most cases, does not notice.

But as ZK proofs extend beyond rollups into coprocessors, ZKML, identity, and
enterprise applications, the payment model diversifies. A DeFi protocol using a
ZK coprocessor to compute a verified TWAP pays per query -- perhaps
\$0.10-\$1.00 per coprocessor call, depending on the computation's complexity.
An enterprise using ZK proofs for compliance verification pays per attestation
-- perhaps \$5-50 per KYC proof, still far cheaper than the manual compliance
process it replaces. A ZKML application proving model inference pays per
inference -- and here the costs are still high enough to be prohibitive for all
but the most valuable use cases, because the computational overhead of proving
neural network operations remains steep.

The margin structure is revealing. Hardware and energy costs are transparent and
declining. The margin on proving-as-a-service therefore comes from three
sources: software differentiation (a faster prover generates more proofs per
GPU-hour, extracting more revenue from the same hardware), network effects (a
larger prover network attracts more proof requests, spreading fixed costs across
more revenue), and trust premiums (a prover network with a longer track record
and more stake at risk can charge more than a new entrant, because the customer
is paying not just for computation but for reliability). The first source
rewards engineering excellence. The second rewards early movers. The third --
and this is the one that connects to our thesis -- rewards \emph{visible
trustworthiness}. The margin on proving-as-a-service is, in part, a margin on
trust. The customer pays more to a prover they trust more. Trust has not been
eliminated. It has been priced.

Current gross margins for proving services are estimated at 40-60\% for
operators with owned hardware, dropping to 15-25\% for operators renting cloud
GPU capacity. As the market matures and competition intensifies, margins will
compress toward the cloud-rental floor -- unless operators differentiate on
software speed, privacy guarantees, or the economic security of their staking
mechanisms. The \$PROVE token's role in this dynamic is worth watching: it
attempts to align prover incentives with network reliability through slashing
(stake at risk if proofs are invalid or late), which is an economic mechanism
for trust. You do not trust the prover's goodwill. You trust the prover's
financial exposure. Whether that is trust-minimized or merely
trust-financialized is a question the market has not yet answered.

The long-term equilibrium may resemble cloud computing: a few large operators
(Succinct, RISC Zero, ZkCloud) competing on price and performance, with
specialized boutique provers serving niche markets (enterprise privacy, ZKML,
identity). The commodity layer -- raw proof generation -- will be a low-margin,
high-volume business. The value layer -- proof orchestration, privacy-preserving
proving, domain-specific optimization -- will be a higher-margin, lower-volume
business. In both layers, the customer pays for computation and receives a
proof. But in the value layer, the customer also pays for the \emph{conditions}
under which the proof was generated: was the witness kept private? Was the
hardware audited? Was the proving pipeline formally verified? These conditions
are trust assumptions, and they carry a price. The proving market does not sell
trustlessness. It sells \emph{degrees of trust}, at different price points, to
customers with different risk tolerances. This is an honest market. It is not
the market the cypherpunks imagined.

\textbf{Trust relocated from:} self-hosted proving infrastructure \textbf{to:}
proving marketplace operators + hardware availability. \textbf{Net:} cost
efficiency with new dependencies on prover liveness and correctness.

\section{Enterprise Pilots (Pilot)}\label{enterprise-pilots-pilot}

The magician has left the theater. She now performs in boardrooms, data centers,
and regulatory offices. The trick is the same. The audience has changed -- and
this audience does not care about trustlessness. It cares about compliance.

\textbf{Deutsche Bank, UBS, and Goldman Sachs} have all conducted pilots on
ZKsync, exploring tokenized assets with ZK-verified compliance. These pilots use
ZK proofs to demonstrate regulatory compliance (KYC/AML verification) without
revealing the underlying customer data to settlement counterparties. The bank
proves its customer is compliant. The counterparty learns nothing about who the
customer is. This is not trust-minimization in the way a cypherpunk would
recognize it. It is trust in compliance frameworks, verified by cryptographic
proofs but governed by regulatory requirements. The magician performs for
regulators now, and the trick she performs is: ``we followed the rules, and here
is a proof.''

The Deutsche Bank pilot reveals what enterprise adoption actually looks like
from the inside. Deutsche Bank's Project Guardian participation -- coordinated
through the Monetary Authority of Singapore -- tested whether
institutional-grade foreign exchange and government bond transactions could
settle on-chain with ZK-verified KYC attestations replacing the traditional
correspondent banking chain. The specific mechanism: Deutsche Bank's compliance
team verifies a client's KYC status through conventional channels, then issues a
ZK credential -- a proof that the client passed all required AML/CFT checks,
without encoding \emph{which} checks, \emph{what} documents, or \emph{whose}
name. The counterparty's smart contract verifies the proof on-chain and permits
settlement. The transaction completes in minutes rather than the T+2 standard.
The compliance data never crosses an institutional boundary.

Understand what has changed and what has not. The settlement speed improved. The
data exposure shrank. The compliance \emph{cost} dropped, because the
correspondent banking chain -- each intermediary performing redundant KYC on the
same client -- collapsed into a single proof verified by mathematics. But the
trust did not disappear. It migrated. Before: you trusted the correspondent
bank's compliance department. After: you trust Deutsche Bank's compliance
department \emph{and} the correctness of their ZK credential issuance system
\emph{and} the soundness of the proof system \emph{and} the smart contract that
verifies it. The number of trust assumptions increased. Their \emph{visibility}
increased too -- and visibility is not nothing. A trust assumption you can
examine is better than one you cannot. But it is still a trust assumption.

The pilot's economics are instructive. Correspondent banking fees for
cross-border institutional transactions typically run 0.1-0.3\% of notional
value. For a \$100 million FX trade, that is \$100,000-\$300,000 in intermediary
costs. The ZK-verified on-chain settlement reduced this to gas fees plus proving
costs -- roughly \$50-200 per transaction at current rates. The savings are not
marginal. They are structural. And they create an incentive so large that the
technology's imperfections become, from an institutional perspective, tolerable.
Banks do not adopt ZK proofs because they believe in trustlessness. They adopt
ZK proofs because the alternative costs a quarter of a million dollars per
transaction and takes two days. Trust-moved is acceptable when the cost savings
are three orders of magnitude.

\textbf{DTCC} (the Depository Trust \& Clearing Corporation) partnered with the
Canton Network for tokenized Treasury securities in December 2025, using ZK
proofs for privacy-preserving settlement. The DTCC clears approximately \$2.5
quadrillion in securities annually. If even a single-digit percentage of that
volume migrates to ZK-verified on-chain settlement within a decade, the proving
infrastructure required would dwarf every other ZK application combined.

The Canton Network partnership is structurally different. Canton is a
privacy-first blockchain designed by Digital Asset, built on the Daml smart
contract language, where every transaction is visible only to its direct
participants. ZK proofs enter not as a replacement for Canton's native privacy
(which is achieved through data partitioning) but as a \emph{bridge} between
partitions: proving to a regulator that aggregate settlement volumes across
private partitions satisfy capital adequacy requirements, without revealing the
individual transactions. The regulator sees a proof. The proof says: ``the sum
of all positions held by Institution X nets to Y, and Y satisfies threshold Z.''
The regulator learns Y and whether it exceeds Z. The regulator learns nothing
about the individual positions that produced Y.

This is trust architecture of a very specific kind. The DTCC is not minimizing
trust in itself -- it \emph{is} the trusted institution, and it intends to
remain so. It is minimizing the \emph{data exposure} required to maintain that
trust. The regulator still trusts the DTCC. But the regulator no longer needs to
see every transaction to verify compliance. The proof substitutes for the data.
The magician has not replaced the theater manager. She has given the theater
manager a way to check the books without reading every page.

\textbf{Partisia Blockchain} operates in a different corner of the enterprise
market, one that reveals how ZK proofs interact with multiparty computation in
government applications. Partisia's collaboration with the Danish government on
digital student credentials -- part of Denmark's broader digitization initiative
-- uses ZK proofs to enable students to prove enrollment status, degree
completion, or GPA thresholds to employers and other institutions without
revealing their full academic transcript. A student proves she graduated with
honors. The employer learns that fact and nothing else -- not the specific
courses, not the grades in individual subjects, not the institution's internal
student ID number.

The trust analysis here differs from the banking cases in a way that matters. In
the Deutsche Bank pilot, the institution issuing the credential (the bank) is
the same institution the counterparty already trusts. The ZK proof is an
efficiency gain, not a trust transformation. In the Partisia student credential
system, the institution issuing the credential (the university) and the
institution verifying it (the employer) have no pre-existing trust relationship.
The ZK proof is not making an existing trust relationship more efficient. It is
\emph{creating a new trust pathway} -- one that runs through mathematics rather
than through phone calls to a registrar's office. This is closer to genuine
trust minimization. The employer trusts the proof, not the university's
willingness to answer the phone.

But even here, the trust has not vanished. It has been moved to the credential
issuance step. The ZK proof guarantees that the credential was issued by the
university's signing key. It does not guarantee that the university's signing
key was not compromised, that the university's records were accurate, or that
the student did not commit academic fraud that went undetected. The mathematics
is honest about the student's credential. It says nothing about whether the
credential is honest about the student. Trust in the proof is not trust in the
underlying reality. This distinction -- between proving a document is authentic
and proving a document is \emph{true} -- is the gap that no amount of
cryptography can close. The magician can prove the card was in the deck. She
cannot prove the deck was not rigged.

The regulatory dimension of enterprise adoption inverts the usual narrative. In
every other context we have examined, ZK proofs are adopted \emph{despite}
regulatory uncertainty -- builders move fast, regulators catch up later. In the
enterprise space, the dynamic is reversed. EU eIDAS 2.0, the Markets in
Crypto-Assets Regulation (MiCA), Basel III's treatment of tokenized assets, and
the SEC's evolving stance on digital securities all create compliance
obligations that ZK proofs can satisfy. The regulation arrives first. The
technology follows. This means enterprise ZK adoption is not driven by
technological enthusiasm. It is driven by legal necessity. When a regulation
says ``you must verify compliance without sharing customer data across
jurisdictions,'' the list of technologies that satisfy that requirement is very
short. ZK proofs are at the top of it.

\textbf{Privacy Pools} (0xbow) launched on Ethereum mainnet in April 2025,
enabling users to prove the provenance of their funds -- that they did not
originate from sanctioned addresses -- without revealing their full transaction
history. This is a direct response to the regulatory challenges highlighted by
Tornado Cash sanctions. Privacy Pools is the most philosophically interesting
enterprise application: it uses ZK proofs to prove innocence without proving
identity. The trust assumption is not in a bank or a regulator but in the
mathematical definition of ``sanctioned address.'' That definition is still set
by a government. The mathematics enforces the policy. It does not choose the
policy.

\textbf{Trust relocated from:} centralized audit/compliance processes
\textbf{to:} ZK-verified attestations + institutional key management.
\textbf{Net:} mixed -- cryptographic integrity gains offset by institutional
adoption risk and key ceremony complexity.

\section{Market Sizing}\label{market-sizing}

The zero-knowledge proof market is growing fast from a small base:

\begin{longtable}[]{@{}lll@{}}
\toprule\noalign{}
Year & Market Size & Source \\
\midrule\noalign{}
\arrayrulecolor{MidnightBlue}\midrule\arrayrulecolor{MidnightBorder}
\endhead
\bottomrule\noalign{}
\endlastfoot
2024 & \$1.28 billion (total ZKP market) & Grand View Research {[}42{]} \\
2025 & \$1.54 billion (total ZKP market) & Grand View Research {[}42{]} \\
2025 & \$97 million (proving services only) & Chorus One \\
2030 & \$1.34 billion (proving services only) & Chorus One \\
2033 & \$7.59 billion (total ZKP market) & Grand View Research {[}42{]} \\
\end{longtable}

\begin{quote}
\textbf{Note on sources}: The total ZKP market figures (\$1.28B, \$1.54B,
\$7.59B) are from Grand View Research's ``Zero-Knowledge Proof Market Size
Report'' (2025), which covers all ZKP segments at a 22.1\% CAGR. The
proving-services sub-market figures (\$97M, \$1.34B) are from Chorus One's ``The
Economics of ZK-Proving: Market Size and Future Projections'' (2025), which
covers the narrower ZK proving services segment at a higher CAGR. Earlier
versions of this table conflated the two sources.
\end{quote}

The Grand View Research figures cover the full ZK ecosystem: proof generation
hardware and services, ZK rollup infrastructure, ZK identity systems, ZKML
tooling, and enterprise licensing. The Chorus One figures cover the narrower
proving-services sub-market. The largest segment of the total market is ZK
rollups (by transaction fee revenue), followed by proving-as-a-service (by
infrastructure investment), and ZK identity (by projected end-market size).

The cost trajectory suggests that ZK proving is following a classic deflationary
technology curve -- the same pattern that drove computing from mainframes to
smartphones. As costs approach zero, the binding constraint shifts from ``can we
afford to prove this?'' to ``what else can we prove?'' This shift is what drives
the expansion from rollups (proving transaction execution) into coprocessors
(proving data queries), ZKML (proving model inference), and identity (proving
personal attributes).

But the market numbers, by themselves, tell you nothing about trust. A \$7.59
billion market in which trust has merely been moved from one institution to
another is not the same as a \$97 million market in which trust has genuinely
been minimized. The thesis of this chapter -- that the market reveals where
trust is actually being minimized and where it is merely being relocated --
should make you read these projections with a specific question: in each
segment, what trust assumption remains that the buyer does not examine?

The enterprise market is the largest long-term opportunity but the slowest to
materialize. Financial institutions move on regulatory timelines (years, not
months), require extensive compliance review, and demand vendor stability. The
fact that major banks are conducting ZK pilots -- rather than dismissing the
technology -- suggests the institutional adoption curve has begun, but the
revenue impact will take 3-5 years to manifest at scale.

Six venues. Six audiences. In two of them -- rollups and proving-as-a-service --
the technology is production-grade and the trust reduction is measurable:
cryptographic proofs replace re-execution, and the economics are favorable. In
two -- identity and enterprise compliance -- the trust minimization is genuine
but the adoption path depends on regulatory mandates (eIDAS 2.0, DTCC pilots)
rather than pure market forces. In two -- coprocessors and ZKML -- the
technology works but the trust story is incomplete: new dependencies on data
availability, model integrity, and prover correctness introduce assumptions that
are not yet well-tested at scale. The market is honest about what it sells. The
question is whether buyers are honest about what they are buying.

The market is real, growing, and diversifying beyond its blockchain origins. The
magician now performs in six different venues -- rollups, coprocessors, ML
inference, identity wallets, proving marketplaces, and boardrooms -- for six
different audiences with six different trust requirements. In some venues, trust
is being genuinely minimized. In others, it is being moved to a new location and
given a new name. The numbers tell us where the technology is being adopted.
They do not tell us where it is honest. For that, we need to ask different
questions -- not ``how big is the market?'' but ``what remains unsolved?'' Those
questions are the subject of our final chapter.

The market is growing. The technology works. The money is real. But beneath
every market segment, open questions persist -- questions about governance,
quantum vulnerability, privacy guarantees, and whether the trust minimization
marketed to buyers matches the trust decomposition the mathematics actually
delivers. The final chapter names these questions, assesses which are solvable
and which may be permanent, and draws the line between what zero-knowledge
proofs have achieved and what remains to be built.

\chapter{Chapter 14: Open Questions and the Road
Ahead}\label{chapter-14-open-questions-and-the-road-ahead}

\section{The Seven Questions That Remain
Open}\label{the-seven-questions-that-remain-open}

Richard Feynman had a test for understanding. If you cannot explain something
simply, you do not understand it. If you can explain it simply and then find the
place where your explanation breaks down, you have found the frontier.

We have spent thirteen chapters explaining zero-knowledge proofs -- layer by
layer, system by system, from the stage to the audience. Now we arrive at the
places where the explanation breaks down. These are not rhetorical questions.
They are research problems with measurable progress and no current solution.
Each one is an open door. Consider this an invitation to walk through.

\subsection{Question 1: Can witness generation be made fully parallel on
GPUs?}\label{question-1-can-witness-generation-be-made-fully-parallel-on-gpus}

Witness generation involves dependency chains -- each intermediate value may
depend on previous ones in a directed acyclic graph that resists naive
parallelism. When cryptographic proving moves to GPUs (achieving 10-50x speedups
via NTT and MSM parallelism), witness generation remains CPU-bound, creating the
Witness Gap. ZKPOG (2025) demonstrated 22.8x speedups by restructuring witness
computation for GPU execution, but the general problem -- making arbitrary
computation graphs GPU-friendly -- remains open.

The Witness Gap grows with every proving speedup, as Chapter 4 documented:
accelerate the cryptographic step and the CPU-bound witness share swells to
dominate total latency. The asymptotic state of zkVM performance may be entirely
witness-bound. The backstage preparation may always take longer than the
performance. Or perhaps not -- but no one has proven otherwise.

The concrete stakes are worth spelling out. If witness generation stays
fundamentally sequential, then proving costs plateau regardless of how many GPUs
you throw at the problem. You could have a warehouse of A100s computing NTTs in
perfect lockstep, and the entire pipeline would still wait on a single CPU
thread chasing a dependency chain through the execution trace. The fastest GPU
prover in the world is only as fast as the slowest CPU witness generator feeding
it. This is Amdahl's Law wearing a cryptographic costume.

ZKPOG's approach -- restructuring the witness computation's dependency graph so
that independent sub-computations can be dispatched to GPU threads -- works when
the graph has enough independent width. Many real circuits do. A Merkle tree
hash has natural parallelism at each level. An elliptic curve multi-scalar
multiplication decomposes into independent scalar-point products. But general
computation does not cooperate this neatly. A loop where iteration \(n\) depends
on the result of iteration \(n-1\) is inherently sequential, and no amount of
clever scheduling changes that.

This creates a mismatch between two kinds of parallelism. Polynomial arithmetic
-- NTTs, MSMs, commitment evaluations -- exhibits \emph{regular} parallelism:
the same operation applied uniformly to millions of independent data points.
GPUs were designed for exactly this. Witness generation exhibits
\emph{irregular} parallelism: the computation graph's shape depends on the
specific program, the specific input, and the specific execution path taken.
GPUs were emphatically not designed for this. Regular parallelism maps to SIMT
execution. Irregular parallelism maps to task-graph scheduling, which is the
domain of CPU thread pools, not GPU warps.

The dependency graph approach points toward a middle path: analyze the witness
computation ahead of time, identify the parallel width at each stage, and
schedule GPU work only where the width justifies the dispatch overhead. But this
requires a compiler that understands both the circuit semantics and the GPU
memory hierarchy -- a compiler that does not yet exist in production. The gap
between ``this is possible in principle'' and ``this ships in a proving
pipeline'' is measured in years of compiler engineering.

What makes this question thrilling rather than depressing is that the bottleneck
is \emph{not} fundamental. Sequential witness generation is a property of
current architectures, not a property of the underlying mathematics. The witness
is a function of the public input and the private input. Nothing about that
function requires sequential evaluation. The sequentiality comes from how we
\emph{describe} the computation (as a step-by-step trace), not from the
computation itself. Whoever figures out how to describe the same computation in
a form that exposes its inherent parallelism -- without changing the circuit it
proves -- will unlock the next order-of-magnitude proving speedup. The NTT
revolution happened when people stopped thinking of polynomials as coefficient
lists and started thinking of them as evaluation vectors. The witness revolution
will happen when someone finds the equivalent change of representation for
execution traces.

\textbf{Executive Risk:} Proving costs plateau regardless of GPU investment.
Hardware budgets may hit diminishing returns before witness generation is
parallelized. \textbf{Timeline:} 2-4 years for resolution.

\subsection{Question 2: What is the proven lower bound on post-quantum proof
size?}\label{question-2-what-is-the-proven-lower-bound-on-post-quantum-proof-size}

Chapter 7 presented a trilemma: algebraic functionality, post-quantum security,
and succinctness -- pick two. But is this a fundamental limitation or a
reflection of current constructions? No paper proves that \(O(1)\)-size
post-quantum proofs are impossible. The lattice-based schemes are achieving
50-100 KB proofs with post-quantum security, approaching practical territory.
But whether there exists a post-quantum polynomial commitment scheme with
\(O(1)\) proof size and \(O(1)\) verification time -- matching KZG's properties
under lattice assumptions -- is unknown.

The ideal PCS problem comes into view: find a commitment scheme that is
transparent (no trusted setup), produces constant-size proofs, enables
constant-time verification, and relies only on post-quantum assumptions. If such
a scheme exists, it would collapse the three-path framework of Chapter 10 into a
single path. If it provably cannot exist, the three paths are permanent. Either
answer would reshape the field. Neither answer exists today.

To appreciate the gap, consider the numbers side by side. A KZG polynomial
commitment produces a proof that is a single elliptic curve point: 48 bytes. A
Groth16 proof -- the entire argument of knowledge, not just the commitment --
fits in 192 bytes. These are absurdly small. They are small because the pairing
structure of elliptic curves allows the verifier to check a polynomial identity
with a single bilinear map evaluation. No redundancy. No repetition. Just one
algebraic equation that either holds or doesn't.

Now look at the post-quantum world. The best lattice-based proofs run 50-100 KB.
Hash-based proofs (STARKs) run larger still -- hundreds of kilobytes to low
megabytes before recursion compresses them. The gap is three orders of
magnitude. A KZG proof fits in a tweet. A lattice proof fills a small PDF. Why?

The reason is structural, not incidental. Pairing-based schemes exploit a
specific algebraic trick: the pairing \(e(g, h)\) lets you check multiplicative
relations between committed values without revealing the values themselves. This
trick compresses verification into a constant number of pairing evaluations,
regardless of the polynomial's degree. Lattice-based schemes have no analogous
trick. Their security rests on the hardness of finding short vectors in
high-dimensional lattices -- a problem that is (probably) hard even for quantum
computers, but that does not offer the same algebraic leverage. Verification
requires checking that a vector is short, and the proof that a vector is short
inherently requires transmitting information proportional to the vector's
dimension.

The Wee-Wu line of results suggests that this gap may have deep roots. Their
work on compact functional encryption and related primitives shows that certain
cryptographic tasks require specific algebraic structure (pairings, or their
lattice analogues) to achieve compactness, and that this structure is hard to
instantiate from lattice assumptions without blowup. The barrier is not a proof
of impossibility -- nobody has shown that \(O(1)\)-size post-quantum proofs
\emph{cannot} exist. But the barrier results suggest that if such proofs exist,
they will require fundamentally new algebraic ideas, not incremental
improvements to existing lattice techniques.

What would a proof of impossibility mean? It would mean the ZK field permanently
bifurcates: fast-and-small proofs (pairing-based, quantum-vulnerable) versus
large-and-safe proofs (lattice-based, quantum-resistant), with no bridge between
them. Every system would need to choose a lane. The STARK-to-SNARK wrapping
strategy described in Chapter 10 -- use hash-based proofs for soundness, then
compress the proof with a pairing-based scheme for on-chain verification --
would be the permanent architecture, not a temporary compromise. The three paths
would harden into the three permanent roads.

Conversely, if someone constructs a post-quantum scheme with \(O(1)\) proof
size, it would be the most important result in applied cryptography since
Groth16. It would mean the entire field's current architecture -- the wrapping,
the recursion, the careful balancing of proof size against verification cost --
is a temporary artifact of our incomplete understanding, not a fundamental
constraint. Every ZK engineer is implicitly betting on which answer is correct.
The field has placed most of its chips on ``the gap is permanent'' and built its
infrastructure accordingly. If that bet is wrong, the infrastructure will need
to be rebuilt. If that bet is right, the infrastructure is the final form.
Either way, someone needs to settle the question. The mathematics is patient,
but the engineering decisions are being made now.

\textbf{Executive Risk:} Permanent bifurcation of the ZK ecosystem into
pre-quantum and post-quantum paths with no migration bridge. Systems deployed
today on pairing-based foundations face mandatory replacement, not upgrade.
\textbf{Timeline:} 5-10 years for theoretical resolution; practical deployment
3-5 years after.

\subsection{Question 3: When will Stage 2
bind?}\label{question-3-when-will-stage-2-bind}

L2Beat's Stages framework defines when ZK rollup governance can override
cryptographic guarantees: Stage 0 (centrally controlled), Stage 1 (limited
governance overrides with 7-day exit windows), Stage 2 (fully decentralized with
30-day exit windows and no governance override).

As of early 2026, most ZK rollups are Stage 0 or Stage 1. This means governance
can override the proof system. A Stage 0 rollup could have 256-bit security and
it would not matter if one person can push a verifier upgrade. The cryptographic
guarantees of Layers 1-6 do not actually bind until a system reaches Stage 2.
The magic trick is perfect, but the theater manager can rewrite the ending.

When will Stage 2 happen? Nobody knows, because Stage 2 requires resolving deep
engineering challenges (bug-free verifier contracts, escape hatches that work in
practice, proof submission that cannot be censored) and deep governance
challenges (who decides when a system is ready, and what happens when a critical
bug is found post-lock). The path forward likely involves multiple independent
verifier implementations that cross-check each other, formal verification of the
verifier contract's core logic, and mandatory exit windows long enough for users
to withdraw if they disagree with a governance decision. Each of these is
technically feasible; none has been fully achieved. This is the question where
mathematics ends and human institutions begin.

Consider what Stage 2 actually looks like in practice. Picture a ZK rollup where
no single entity -- not the founding team, not a multisig, not a governance
token majority -- can push a verifier upgrade without triggering a mandatory
30-day exit window. During that window, every user can inspect the proposed
change, compare the old and new verifier contracts, and withdraw their funds to
L1 if they disagree. Multiple independent verifier implementations -- written by
different teams, in different languages, tested against different fuzzing suites
-- cross-check each other's results. If any implementation disagrees with the
others, the system halts and the dispute resolution mechanism activates. The
verifier contract itself has been formally verified: not just audited by humans
who might miss an edge case, but proven correct by a theorem prover that has
checked every execution path.

That is the goal. Now consider the tensions that make it hard.

The first tension is between immutability and patchability. An immutable
verifier contract cannot be captured by governance -- but it also cannot be
patched when a critical bug is discovered. And critical bugs \emph{will} be
discovered. The history of smart contract security is unambiguous on this point:
every contract of sufficient complexity has bugs, and the bugs are found on
timelines measured in months to years, not days. A Stage 2 system that locks its
verifier contract is betting that the contract is bug-free. That bet has never
been won at scale. A Stage 2 system that retains upgrade capability, even with
time-locks, is betting that the governance process cannot be captured. That bet
has been lost repeatedly in traditional finance. The equilibrium, if it exists,
involves a narrow corridor: upgrades are possible, but only through a process so
slow and so transparent that capture is impractical.

The second tension is between safety and liveness. A system with 30-day exit
windows is safe -- users can always leave. But if a critical vulnerability is
discovered and exploited, 30 days is an eternity. An attacker who finds a
soundness bug can drain the rollup in a single transaction. The users' 30-day
exit window is worthless if the funds have already been stolen. This suggests
that Stage 2 needs not just exit windows but also emergency circuit-breakers --
but a circuit-breaker is a governance override, which is exactly what Stage 2 is
supposed to eliminate. The circle closes on itself.

The third tension is economic. Multiple independent verifier implementations are
expensive. Formal verification is expensive. Long exit windows impose
opportunity costs on users. Who pays? In a Stage 0 system, the founding team
pays for everything and recoups the cost through token appreciation or sequencer
fees. In a Stage 2 system, the costs must be borne by the protocol itself, which
means they are ultimately borne by users through fees or inflation. The economic
question is whether the security premium of Stage 2 is worth the cost -- and the
answer depends on how much value the rollup secures. A rollup holding \$10
billion has a very different cost-benefit calculus than a rollup holding \$10
million.

The honest answer may be that Stage 2 is not a destination but a spectrum.
Different systems will reach different points on that spectrum depending on
their value secured, their user base's risk tolerance, and the maturity of their
verifier contract. The first true Stage 2 system will likely be a rollup that
has been in production for years, has undergone multiple independent audits and
formal verification efforts, and secures enough value to justify the engineering
cost. It will be boring. It will be slow. It will be the most important thing in
the ecosystem. The question is not \emph{whether} it will happen, but whether
the incentives align for it to happen before a catastrophic failure forces it to
happen. The field should prefer the former.

\textbf{Executive Risk:} Governance override negates all cryptographic
guarantees. Over \$20B in TVL sits in systems where a multisig can replace the
verifier. This is the single largest unresolved risk in ZK infrastructure.
\textbf{Timeline:} 1-3 years for first Stage 2 rollups.

\subsection{Question 4: When will ``trustless'' become
real?}\label{question-4-when-will-trustless-become-real}

The trust decomposition in Chapter 10 identified seven independent assumptions
underlying ZK systems. Today, every deployed system requires trusting all seven
simultaneously. The trajectory points toward progressive reduction:

\begin{itemize}
\tightlist
\item
  Under-constrained circuits (Layer 2): formal verification tools (Picus, ZKAP,
  Coda) are reducing but not eliminating the vulnerability class catalogued in
  Chapter 3; the tools catch many but not all.
\item
  Fiat-Shamir bugs (Layer 5): standardization of transcript protocols and
  automated fuzzing (ARGUZZ found 11 bugs across 6 zkVMs) are addressing the
  implementation gap.
\item
  Governance override (Layer 7): Stage 2 maturation will eventually remove
  governance as an attack surface.
\item
  Quantum vulnerability (Layer 6): lattice-based and hash-based primitives
  provide migration paths.
\end{itemize}

But ``trustless'' -- meaning zero residual trust assumptions -- requires proving
that hardware does not leak (impossible without information-theoretic security),
proving that all software is correct (requiring formal verification of the
entire stack), and proving that mathematical hardness assumptions hold (which is
inherently impossible -- hardness is a conjecture, not a theorem).

The honest trajectory: trust will continue to decrease, asymptotically
approaching but never reaching zero. ``Trust-minimized, and getting better every
day'' is the accurate description. Like Zeno's arrow, the gap closes by half
with each advance. Zero is not the destination. The asymptote is the
destination, and it is a good one.

But the Zeno metaphor deserves to be pushed harder, because it reveals something
the optimistic narrative glosses over. Each layer's trust assumption is
shrinking independently. Layer 2 gets safer as formal verification improves.
Layer 5 gets safer as transcript standardization matures. Layer 6 gets safer as
lattice assumptions are studied. Layer 7 gets safer as governance decentralizes.
Viewed individually, each trend is encouraging. But the system's overall
trustworthiness is not the average of its layers -- it is the
\emph{conjunction}. All seven must hold simultaneously. If any single layer
fails, the system fails.

This is the Zeno's paradox of trust in its precise form. Suppose each layer
independently has a 99\% chance of holding its trust assumption. Seven
independent layers at 99\% each give a system-level confidence of
\(0.99^7 \approx 93.2\%\). Improve each layer to 99.9\% and the system reaches
99.3\%. Improve to 99.99\% and you get 99.93\%. The system-level confidence
always trails the weakest layer, and the gap between layer-level confidence and
system-level confidence grows with the number of layers. Adding layers -- which
is what happens as the stack matures and new trust assumptions are identified --
makes the individual layers stronger but the conjunction weaker.

This is not a reason for despair. It is a reason for precision. The trajectory
toward trust-minimization is real, but the conjunction effect means that
``almost trustless'' is a category error. A system with six perfect layers and
one compromised layer is not ``mostly trustless'' -- it is compromised. The
chain is as strong as its weakest link, and a seven-link chain has seven
opportunities for weakness.

The practical implication is that the field needs to track system-level trust,
not layer-level trust. A dashboard showing ``Layer 2: formally verified, Layer
5: standardized transcripts, Layer 6: post-quantum ready, Layer 7: Stage 1.5''
tells you about components. It does not tell you about the system. The
system-level question is: what is the probability that \emph{all seven} hold
simultaneously for \emph{this specific deployment} with \emph{this specific
configuration}? Nobody can answer that question today. Developing the
methodology to answer it -- a kind of actuarial science for cryptographic trust
-- is itself an open research problem.

The deepest version of this question is philosophical, not technical. Is there a
fundamental lower bound on residual trust, the way there is a fundamental lower
bound on measurement uncertainty in quantum mechanics? Can we prove that
\emph{any} computational system, no matter how carefully designed, must retain
some irreducible trust assumption? The answer is almost certainly yes -- you
always trust that the laws of physics haven't changed, that your hardware is
computing what you think it's computing, that the mathematical axioms underlying
the proofs are consistent. But articulating exactly where that floor sits, and
how close current systems are to reaching it, would transform
``trust-minimized'' from a vague aspiration into a measurable engineering
target. The person who defines that floor will have done for cryptographic trust
what Shannon did for information capacity: turned an intuition into a theorem
with a number attached to it.

\textbf{Executive Risk:} Marketing claims exceed technical reality.
Organizations deploying `trustless' systems inherit seven trust assumptions they
may not have assessed. Due diligence must enumerate all seven.
\textbf{Timeline:} Ongoing.

\subsection{Question 5: How do streaming witness approaches interact with
folding?}\label{question-5-how-do-streaming-witness-approaches-interact-with-folding}

Nair et al.~(2025) demonstrated streaming witness generation that uses space
proportional to trace width (registers) rather than trace length (steps). This
is essential for proving long computations without requiring hundreds of
gigabytes of RAM. But folding schemes (Nova, HyperNova, LatticeFold) accumulate
state across folding steps -- the accumulated instance grows with each fold. How
streaming witnesses interact with accumulating folding state is an open
architectural question.

This matters because folding and streaming attack different bottlenecks: folding
reduces prover time by avoiding redundant computation; streaming reduces prover
memory by generating witness chunks on-the-fly. An architecture that combines
both would be strictly superior to either alone, but the interaction is
non-trivial. Two good ideas that have not yet learned to dance together.

The architectural tension runs deeper than a scheduling problem. It is a
fundamental conflict between two strategies for managing state.

Folding \emph{accumulates} state. Each folding step takes the current
accumulated instance and a new instance and produces a new accumulated instance
that absorbs both. In Nova, this means the error vector \(\mathbf{E}\) grows --
or more precisely, the relaxed R1CS instance carries forward information about
all previous folding steps. In LatticeFold, the accumulated witness includes
norm-bounded vectors whose bounds increase with each fold. The accumulator is
the memory of the system. It remembers everything that has been folded into it.
That memory is what makes folding powerful: the verifier checks a single
accumulated instance at the end, rather than checking each step individually.
But that memory has a cost. The accumulated state must be stored, updated, and
eventually opened. It cannot be discarded mid-computation.

Streaming \emph{discards} state. The entire point of streaming witness
generation is to produce a chunk of the witness, use it, and then forget it. The
prover's memory footprint stays proportional to the width of the computation
(how many registers the program uses) rather than its length (how many steps the
program takes). A streaming prover for a billion-step computation uses the same
memory as a streaming prover for a million-step computation. That is the magic.
But it works only if the prover can forget the past. The moment you require the
prover to retain information about earlier chunks -- which is exactly what
folding demands -- the streaming property degrades.

The crux is this: folding needs the prover to remember the accumulated error
from all previous steps. Streaming needs the prover to forget the witness from
all previous steps. These are not the same thing. The accumulated error is a
compact summary (a single relaxed instance), while the witness is a massive
expansion (the full execution trace). In principle, you can stream the witness
while retaining only the accumulated instance. In practice, the folding step
itself requires access to both the accumulated witness \emph{and} the new
witness chunk simultaneously, because the folding combiner must compute
cross-terms between old and new. Those cross-terms are the technical barrier.

One possible resolution: a folding scheme where the cross-terms can be computed
incrementally, using only the compact accumulated instance and the current
witness chunk, without ever materializing the full accumulated witness. This
would require the cross-term computation to be decomposable into a
streaming-friendly form -- which is a non-trivial algebraic constraint on the
folding scheme's structure. Nova's cross-term involves an inner product between
the accumulated witness and a matrix applied to the new witness. That inner
product can, in principle, be computed in a streaming fashion if the
matrix-vector product is streamed. But the accumulated witness itself is not
compact; it is as large as the original witness. The circularity is apparent.

A different resolution: a layered architecture where streaming operates within
each folding step and folding operates across steps. Generate the witness for
step \(k\) using streaming (small memory), fold it into the accumulator (small
memory for the accumulator, large memory transiently for the fold), then discard
the witness for step \(k\) and move to step \(k+1\). The peak memory is
determined by the single-step witness size plus the accumulator size, not by the
total trace length. This is closer to what practical implementations will likely
achieve, but it requires that each folding step can be completed before the next
witness chunk is generated -- which imposes a sequential dependency between
folding and witness generation that limits parallelism.

Nobody has built an architecture that resolves all three constraints
simultaneously: streaming memory, folding time-savings, and parallel execution.
The system that does will dominate the proving market for the next generation of
zkVMs. It is the kind of problem where the first correct solution will seem
obvious in retrospect and impossible in prospect -- which is the surest sign
that it is worth working on.

\textbf{Executive Risk:} Memory requirements constrain prover hardware to
high-end data center configurations, limiting decentralization and increasing
vendor lock-in. \textbf{Timeline:} 2-3 years for streaming witness architectures
to mature.

\subsection{Question 6: Can constant-time ZK proving be made
practical?}\label{question-6-can-constant-time-zk-proving-be-made-practical}

Side-channel attacks on ZK implementations -- timing attacks on Zcash's Groth16
prover (\(R = 0.57\) correlation with transaction amounts), cache timing
leakages in Poseidon/Reinforced Concrete/Tip5, electromagnetic emanation from
field arithmetic -- demonstrate that ``zero-knowledge'' is a mathematical
property, not an implementation guarantee. The proof reveals nothing in theory.
The hardware whispers everything in practice.

Constant-time implementations exist (Monero's Bulletproofs prover achieves
\(R = 0.04\) correlation) but impose significant performance costs. GPU proving
complicates this further: GPUs execute in SIMT (Single Instruction, Multiple
Thread) mode, where thread divergence is observable. Constant-time code creates
artificial thread divergence, reducing GPU utilization. The tension between
side-channel resistance and hardware throughput has no known resolution. The
backstage walls need to be soundproofed, but soundproofing slows down the
preparation.

The GPU problem is worth understanding in detail, because it illustrates why
this is not simply a matter of writing more careful code.

A GPU does not execute threads independently. It executes them in \emph{warps}
-- groups of 32 threads (on NVIDIA hardware) that share a single instruction
pointer. When all 32 threads in a warp take the same branch, the warp executes
at full speed. When threads diverge -- some taking the \texttt{if} branch,
others taking the \texttt{else} branch -- the warp must execute \emph{both}
branches serially, masking out the threads that shouldn't participate in each
branch. This is thread divergence, and it is the fundamental mechanism by which
GPUs trade flexibility for throughput.

Constant-time code, by definition, must execute the same instructions regardless
of the input. In a CPU, this means taking both branches and selecting the result
with a conditional move -- a small overhead, typically 2x or less. In a GPU,
this means \emph{forcing} all 32 threads in a warp to execute both branches even
when all 32 threads would naturally take the same branch. The constant-time
requirement converts natural coherence (all threads agree) into artificial
divergence (all threads must pretend to disagree). On a CPU, constant-time code
is slower by a constant factor. On a GPU, constant-time code destroys the
execution model's fundamental assumption.

The numbers are stark. A well-optimized GPU NTT can achieve 80-90\% of
theoretical throughput because the computation is uniform: every thread does the
same butterfly operation on different data. A constant-time GPU implementation
of the same NTT -- one that pads all branches to equal length and eliminates
data-dependent memory access patterns -- drops to 30-50\% throughput. The field
arithmetic itself (modular addition, modular multiplication) can be made
constant-time without much overhead, because it is naturally branchless. But the
\emph{control flow} of the proving algorithm -- which polynomials to evaluate,
which commitment rounds to perform, how to handle edge cases in the MSM bucket
accumulation -- is where the branches live, and where constant-time enforcement
hurts.

This creates a stark engineering tradeoff. On one side: constant-time code on
CPU. Viable, proven (Monero does it), but slow. A CPU prover is 10-50x slower
than a GPU prover. On the other side: variable-time code on GPU. Fast, but the
timing variations leak information about the witness through observable
channels. The proving time itself becomes a side channel. An attacker who can
measure how long your proof took can infer something about your private input --
exactly the information the zero-knowledge property is supposed to hide.

No current system offers both. No current system even articulates a clear path
to offering both. The approaches being explored -- oblivious RAM for memory
access patterns, garbled circuits for branching logic, hardware enclaves for
execution isolation -- each solve part of the problem while introducing new
trust assumptions. Oblivious RAM has logarithmic overhead that compounds badly
at GPU scale. Garbled circuits are designed for two-party computation, not
single-prover execution. Hardware enclaves (SGX, TrustZone) are the very kind of
trusted hardware that ZK proofs were supposed to make unnecessary.

The resolution, when it comes, will likely involve a new hardware abstraction: a
proving accelerator designed from the ground up for constant-time cryptographic
computation, where the SIMT model is replaced by a model that treats uniform
execution time as a feature rather than a constraint. Such hardware does not
exist. Designing it is an open problem at the intersection of chip architecture,
cryptographic engineering, and side-channel analysis. The person who designs it
will need to understand all three. That is a rare combination, which is why the
problem remains open.

\textbf{Executive Risk:} Privacy leaks through side channels in production.
Proving time correlates with transaction characteristics (demonstrated
\(R = 0.57\) in Zcash). Systems marketed as `private' may leak information to
network observers. \textbf{Timeline:} 3-5 years for constant-time provers to
become standard.

\subsection{Question 7: Is ``seven'' the right number of
layers?}\label{question-7-is-seven-the-right-number-of-layers}

The evidence suggests that the seven-layer model is pedagogically useful but
architecturally approximate. Layers 3 and 4 collapse in Jolt. Layers 4, 5, and 6
form an inseparable ``proof core.'' STARK-to-SNARK wrapping pierces Layers 5, 6,
and 7. Data availability and proof aggregation are substantial infrastructure
layers not represented. Privacy is a cross-cutting concern, not a single-layer
property.

The model might be better as four macro-layers: (1) Trust Setup, (2) Programming
Model, (3) Proof Core \{witness + arithmetization + proof system + primitives\},
(4) On-Chain Settlement \{verification + data availability + governance\}. Or it
might be better as nine layers, adding data availability and proof aggregation
explicitly. Or seven might be exactly right for the pedagogical purpose it
serves, with explicit acknowledgment that production systems do not respect the
boundaries.

The question is not academic. The number of layers determines how people think
about the system, which determines which design decisions they consider
independent (and thus can optimize separately) versus coupled (and must
co-design). Getting the decomposition wrong leads to suboptimal architectures.
Getting it right enables the field to progress faster. The map shapes the
territory.

The practical cost of a wrong decomposition is already visible. Consider Layers
4 and 5 in this book's model -- the arithmetization layer and the proof system
layer. The model presents them as separate, and this separation is conceptually
clean: one translates the program into equations, the other proves the equations
are satisfied. Different teams can work on each. Different papers can advance
each. Different benchmarks can measure each. The separation enables a division
of labor that has served the field well.

But in practice, Layers 4 and 5 are deeply coupled. The arithmetization's
structure determines which proof system can operate on it efficiently. R1CS
works naturally with Groth16 and Nova. AIR works naturally with STARKs. CCS
works naturally with HyperNova. Choosing an arithmetization \emph{is} choosing a
family of proof systems, and choosing a proof system constrains the
arithmetization. Teams that optimize Layer 4 independently of Layer 5 --
designing a beautiful arithmetization and then discovering that no efficient
proof system can consume it -- have wasted months of work. This has happened. It
will happen again whenever the model's boundaries do not match reality's
boundaries.

The same coupling appears between Layers 3 and 4. Jolt demonstrated this
forcefully: in Jolt's architecture, the witness \emph{is} the arithmetization.
There is no separate ``witness generation'' step followed by a ``constraint
compilation'' step. The lookup table that defines the instruction semantics
simultaneously serves as the witness and the constraint. Trying to optimize
``witness generation'' and ``arithmetization'' separately in Jolt is like trying
to optimize a coin's heads and tails independently. They are the same object
viewed from two angles.

The four-macro-layer model -- Trust Setup, Programming Model, Proof Core,
On-Chain Settlement -- has the virtue of grouping the coupled components
together. The ``Proof Core'' macro-layer acknowledges that witness generation,
arithmetization, the proof system, and the underlying primitives are a single
coupled system that must be co-designed. But it loses granularity. It tells the
proof-core team ``this is all yours'' without helping them understand the
internal structure of their subsystem.

The nine-layer model -- adding data availability and proof aggregation as
explicit layers -- has the virtue of representing real infrastructure that the
seven-layer model ignores. Data availability is not a detail. It is a trust
assumption as fundamental as any in the stack: if the data is not available, the
proof is unverifiable in practice, regardless of its mathematical soundness.
Proof aggregation -- the recursive composition of proofs into proofs-of-proofs
-- is the mechanism that makes ZK rollups economically viable, and it introduces
its own trust assumptions (the aggregation circuit must be correct, the
recursive verifier must be sound). Omitting these from the model means omitting
trust assumptions from the analysis, which defeats the model's purpose.

The right answer may not be a fixed number at all. It may be that the useful
decomposition depends on the question you are asking. If you are asking ``how
does trust flow through a ZK system,'' seven layers (or nine, or four) provide
different but complementary views. If you are asking ``how should engineering
teams be organized,'' the coupling structure matters more than the layer count.
If you are asking ``what can go wrong,'' the answer is a threat model, not a
layer diagram, and the threat model cuts across every decomposition.

The OSI model survived not because seven was the right number of layers for
networking -- it was widely criticized for being both too many and too few,
depending on context -- but because it gave a generation of engineers a shared
vocabulary. This book's seven-layer model will succeed or fail on the same
terms. If it gives ZK engineers a shared vocabulary for discussing trust
assumptions, it will have served its purpose even if no production system has
exactly seven layers. If a better decomposition emerges -- one that captures the
couplings, represents the missing layers, and scales from pedagogy to production
-- it should replace this one without ceremony. The goal was never to be right
about the number. The goal was to be useful about the structure.

\textbf{Executive Risk:} Low. This is a pedagogical question, not a deployment
risk. The layer model is a communication tool; systems work regardless of how we
categorize their components.

\section{The Three Frontiers}\label{the-three-frontiers}

The ZK field is transitioning through three sequential frontiers, each building
on the previous:

\subsection{Frontier 1: Performance (2023-2025) -- Largely
Crossed}\label{frontier-1-performance-2023-2025-largely-crossed}

The performance frontier asked: can we prove computation fast enough and cheaply
enough for production use? The answer, as of late 2025, is yes. Real-time
Ethereum proving is achieved. The cost curve traced in Chapters 1 and 6 has
flattened at pennies per block. The EF declared the speed race won. The
remaining performance work is optimization (energy efficiency, memory reduction,
witness acceleration), not breakthrough. The magician can now perform in real
time. That question is settled.

\textbf{Evidence this frontier is active:} - SP1 Hypercube: 6.9s Ethereum block
proving on 16 GPUs (Dec 2025) - Airbender: 21.8M RISC-V cycles/sec on single
H100 (2025) - Proving cost: \$80 → \$0.04 in 24 months (2,000x reduction) - EF
declared speed race `effectively won,' pivoted to security (Dec 2025)

\subsection{Frontier 2: Security (2026-2028) -- Currently
Active}\label{frontier-2-security-2026-2028-currently-active}

The security frontier asks: can we prove that our proofs are actually sound,
with quantifiable security guarantees? The EF's 2026 targets (100-bit by May,
128-bit by December) define this frontier precisely. It includes formal
verification of zkVM implementations, provable soundness without conjectured
assumptions (proximity gaps), post-quantum readiness, and protection against
Fiat-Shamir vulnerabilities.

This frontier is harder than performance because security is a property you
prove, not a metric you optimize. You cannot benchmark soundness the way you
benchmark throughput. The tools (Picus, ZKAP, ARGUZZ, soundcalc) are emerging
but immature. The field's first formal verification efforts (SP1's RISC-V Sail
specification verification, Jolt's ACL2 lookup semantics verification) are
promising but cover only fragments of the full stack. Proving that the trick
works is easy. Proving that it \emph{cannot} be faked is the real challenge.

\textbf{Evidence this frontier is active:} - EF 2026 targets: 100-bit provable
security by May, 128-bit by December - SP1: formal verification of 62 RISC-V
opcodes against Sail specification - Arguzz (Hochrainer, Wustholz, Christakis,
2025): found 11 bugs across 6 major zkVMs - soundcalc: automated soundness
margin calculator for proof system parameters - Picus, ZKAP: emerging static
analysis tools for constraint under-specification

\subsection{Frontier 3: Privacy (2027+) --
Approaching}\label{frontier-3-privacy-2027-approaching}

The privacy frontier asks: can we build systems where users genuinely control
their data, where the ``zero-knowledge'' property holds not just mathematically
but in practice? This frontier includes constant-time implementations, metadata
privacy (timing, transaction structure, network-level information), and
application-layer privacy design (compiler-enforced disclosure boundaries,
selective disclosure for regulatory compliance).

Midnight represents an early attempt at this frontier. Its disclosure analysis
is the most sophisticated compiler-level privacy enforcement in any ZK system.
But the documentation's silence on side channels, the indexer's metadata
leakage, and the SDK-level timing correlations demonstrate how far the field has
to go.

The privacy frontier is the most difficult of the three because it requires
solving problems across all seven layers simultaneously. Performance is
primarily a Layers 3-5 problem. Security spans Layers 2-6. Privacy, as Midnight
demonstrates, is a property of the entire stack. The theater must be lightproof
at every joint.

\textbf{Evidence this frontier is active:} - Compact disclosure analysis: 11
compile-time privacy errors caught in single voting contract - Monero
Bulletproofs: constant-time implementation (\(R = 0.04\) timing correlation) -
eIDAS 2.0: EU mandate for selective disclosure in digital identity wallets
(2026) - Midnight testnet: end-to-end private smart contract execution with
local proving - Zcash timing vulnerability patched; field now aware of
side-channel class

\section{Convergence}\label{convergence}

The zero-knowledge proof ecosystem in March 2026 is in a state that would have
been difficult to predict three years ago. Real-time proving of arbitrary
computation is solved. Costs are negligible. Seven production zkVMs compete on a
standardized benchmark (Ethereum block proving). The Ethereum Foundation has
shifted its primary metric from speed to security. Lattice-based constructions
are approaching practical viability for post-quantum proving. Enterprise
financial institutions are conducting pilots. A \$97 million market is projected
to grow to \$7.59 billion in seven years.

The seven-layer model -- for all its imperfections, its outdated binary, its
missing primitives -- got the fundamental insight right: zero-knowledge proof
systems are stacks, not monoliths. Understanding them requires understanding
each layer's contribution and, critically, the interactions between layers. The
model bends where layers couple (the proof core) and breaks where layers
collapse (Jolt's witness-is-arithmetization), but it serves its pedagogical
purpose: giving a structured way to think about a technology that is
simultaneously simpler than it looks (the core math is elegant) and more complex
than it looks (the engineering is a web of coupled decisions).

The seven open questions remain open. The ideal PCS has not been found. Stage 2
governance has not bound. ``Trustless'' has not become real. But the trajectory
is unambiguous: the trust assumptions are getting weaker, the proofs are getting
cheaper, the security is getting stronger, and the privacy is getting more
principled.

Trust-minimized, not trustless -- and getting better every day.

\section{Coda}\label{coda}

The magician walks on stage. She faces an audience of strangers -- people who
have no reason to trust her and every reason to doubt. She makes a claim that
sounds impossible.

But the audience is no longer what it once was. Over thirteen chapters, the
stage has been built, inspected, and rebuilt. The choreography has been written
in languages that prevent the worst mistakes. The backstage recording has been
compressed into a mathematical certificate that no forger can reproduce. The
seal rests on problems that have resisted every attack for decades, and
increasingly on problems designed to resist quantum computers too. The verdict
is rendered by software that is slowly, painfully, being placed beyond the reach
of any committee or key-holder.

The trick still requires trust. It always will.

But the trust has been decomposed, distributed, and minimized until what remains
is not faith in a person or an institution but confidence in a conjunction of
mathematical facts -- each independently testable, each independently
replaceable, each weaker than trusting any single entity with everything.

Think about what that means for a real person.

Consider Alice. She is a small-business owner applying for a loan. In the world
before this technology, the bank demands everything: tax returns, account
balances, transaction histories, personal identification, credit reports. The
bank sees all of Alice's financial life, stores it on servers that get breached
with depressing regularity, and shares it with partners and regulators and data
brokers in ways Alice never consented to and cannot control.

Now consider Alice in the world this technology is building. She generates a
zero-knowledge proof that her income exceeds the lending threshold. She proves
her identity is verified and her funds are not sanctioned. She proves her credit
score falls within an acceptable range. The bank receives three proofs and a
loan application. It learns that Alice qualifies. It learns nothing else -- not
her exact income, not her account numbers, not her transaction history, not the
names of her customers. The sealed certificates say: she qualifies. The
mathematics guarantees it. The bank verifies in milliseconds.

Alice controls her data. The bank gets the answer it needs. The regulator can
verify that the process was followed correctly. And no database holds Alice's
financial life story, waiting to be breached.

That is what seven layers of mathematics make possible. Not trustlessness --
trust-minimization. Not perfection -- progress. Not magic -- engineering that
looks like magic until you understand every layer, at which point it looks like
something better: a system where telling the truth is easier than lying, and
where proving a fact does not require surrendering the context that makes it
private.

The magician performs. The audience verifies. And between them, seven layers of
mathematics ensure that the proof reveals nothing but the truth.

\chapter*{Complete Bibliography}\label{complete-bibliography}
\addcontentsline{toc}{chapter}{Complete Bibliography}

\subsection{Chapter 1: The Promise}\label{chapter-1-the-promise}

\begin{enumerate}
\def\labelenumi{\arabic{enumi}.}
\tightlist
\item
  Clarke, Arthur C. \emph{Profiles of the Future}. Harper \& Row, 1962.
\item
  Goldwasser, Shafi, Silvio Micali, and Charles Rackoff. ``The Knowledge
  Complexity of Interactive Proof Systems.'' \emph{SIAM Journal on Computing}
  18(1): 186--208, 1989.
\item
  Chaliasos, Stefanos, et al.~``SoK: What Don't We Know? Understanding Security
  Vulnerabilities in SNARKs.'' \emph{USENIX Security Symposium}, 2024.
\end{enumerate}

\subsection{Chapter 2: Layer 1 -- Building the
Stage}\label{chapter-2-layer-1-building-the-stage-1}

\begin{enumerate}
\def\labelenumi{\arabic{enumi}.}
\setcounter{enumi}{3}
\tightlist
\item
  Kate, Aniket, Gregory M. Zaverucha, and Ian Goldberg. ``Constant-Size
  Commitments to Polynomials and Their Applications.'' \emph{ASIACRYPT 2010}.
\item
  Bowe, Sean, Ariel Gabizon, and Ian Miers. ``Scalable Multi-party Computation
  for zk-SNARK Parameters in the Random Beacon Model.'' ePrint 2017/1050.
\item
  Groth, Jens. ``On the Size of Pairing-Based Non-interactive Arguments.''
  \emph{EUROCRYPT 2016}.
\item
  Gabizon, Ariel, Zachary J. Williamson, and Oana Ciobotaru. ``PLONK:
  Permutations over Lagrange-bases for Oecumenical Noninteractive arguments of
  Knowledge.'' ePrint 2019/953.
\item
  Ben-Sasson, Eli, et al.~``Scalable, Transparent, and Post-Quantum Secure
  Computational Integrity.'' ePrint 2018/046.
\item
  Bunz, Benedikt, et al.~``Bulletproofs: Short Proofs for Confidential
  Transactions and More.'' \emph{IEEE S\&P 2018}.
\item
  Wang, Faxing, Shaanan Cohney, and Joseph Bonneau. ``SoK: Trusted Setups for
  Powers-of-Tau Strings.'' \emph{FC 2025}. ePrint 2025/064.
\item
  Boneh, Dan and Binyi Chen. ``LatticeFold+: Faster, Simpler, Shorter
  Lattice-Based Folding for Succinct Proof Systems.'' \emph{CRYPTO 2025}. ePrint
  2025/247.
\end{enumerate}

\subsection{Chapters 3-5}\label{chapters-3-5}

\begin{enumerate}
\def\labelenumi{\arabic{enumi}.}
\setcounter{enumi}{11}
\tightlist
\item
  Pailoor, Shankara, et al.~``Automated Detection of Under-Constrained Circuits
  in Zero-Knowledge Proofs.'' \emph{PLDI 2023}. ePrint 2023/512. (See also ref
  {[}47{]}.)
\item
  Wen, Hongbo, et al.~``Practical Security Analysis of Zero-Knowledge Proof
  Circuits.'' \emph{USENIX Security 2024}. ePrint 2023/190. (See also ref
  {[}48{]}.)
\item
  Setty, Srinath, Justin Thaler, and Riad Wahby. ``Customizable Constraint
  Systems for Succinct Arguments.'' ePrint 2023/552.
\item
  Setty, Srinath, Justin Thaler, and Riad Wahby. ``Unlocking the Lookup
  Singularity with Lasso.'' ePrint 2023/1216.
\item
  Arun, Arasu, Srinath Setty, and Justin Thaler. ``Jolt: SNARKs for Virtual
  Machines via Lookups.'' ePrint 2023/1217.
\end{enumerate}

\subsection{Chapter 6: Layer 5 -- The Sealed
Certificate}\label{chapter-6-layer-5-the-sealed-certificate-1}

\begin{enumerate}
\def\labelenumi{\arabic{enumi}.}
\setcounter{enumi}{16}
\tightlist
\item
  Kothapalli, Abhiram, Srinath Setty, and Ioanna Tzialla. ``Nova: Recursive
  Zero-Knowledge Arguments from Folding Schemes.'' \emph{CRYPTO 2022}.
\item
  Kothapalli, Abhiram and Srinath Setty. ``HyperNova: Recursive Arguments for
  Customizable Constraint Systems.'' \emph{CRYPTO 2024}. ePrint 2023/573.
\item
  Bunz, Benedikt and Binyi Chen. ``ProtoStar: Generic Efficient
  Accumulation/Folding for Special Sound Protocols.'' \emph{ASIACRYPT 2023}.
  ePrint 2023/620.
\item
  Boneh, Dan and Binyi Chen. ``LatticeFold: A Lattice-based Folding Scheme and
  its Applications to Succinct Proof Systems.'' \emph{ASIACRYPT 2025}. ePrint
  2024/257.
\item
  Nguyen, Wilson and Srinath Setty. ``Neo: Lattice-based Folding Scheme for CCS
  over Small Fields and Pay-per-bit Commitments.'' ePrint 2025/294.
\item
  Trail of Bits. ``Frozen Heart: Forgery of Zero Knowledge Proofs.'' Blog post,
  April 2022.
\item
  Haboeck, Ulrich, David Levit, and Shahar Papini. ``Circle STARKs.'' ePrint
  2024/278.
\item
  LFDT-Nightstream. ``Nightstream: Lattice-Based Folding Implementation.''
  GitHub repository, 2025. https://github.com/LFDT-Nightstream/Nightstream
\end{enumerate}

\subsection{Chapter 7: Layer 6 -- The Deep
Craft}\label{chapter-7-layer-6-the-deep-craft-1}

\begin{enumerate}
\def\labelenumi{\arabic{enumi}.}
\setcounter{enumi}{23}
\tightlist
\item
  Shor, Peter W. ``Algorithms for Quantum Computation: Discrete Logarithms and
  Factoring.'' \emph{FOCS 1994}.
\item
  NIST. FIPS 203, 204, 205. August 2024.
\item
  NIST. ``Transition to Post-Quantum Cryptography Standards (IR 8547).''
  November 2024.
\end{enumerate}

\subsection{Chapter 8: Layer 7 -- The
Verdict}\label{chapter-8-layer-7-the-verdict-1}

\begin{enumerate}
\def\labelenumi{\arabic{enumi}.}
\setcounter{enumi}{26}
\tightlist
\item
  L2Beat. ``Stages Framework for L2 Maturity.'' https://l2beat.com/stages.
  Accessed March 2026.
\item
  Chaliasos, Stefanos, et al.~``Unaligned Incentives: Pricing Attacks Against
  Blockchain Rollups.'' arXiv 2509.17126, 2025.
\end{enumerate}

\subsection{Chapter 9: Privacy-Enhancing
Technologies}\label{chapter-9-privacy-enhancing-technologies-1}

\begin{enumerate}
\def\labelenumi{\arabic{enumi}.}
\setcounter{enumi}{28}
\tightlist
\item
  Gentry, Craig. ``Fully Homomorphic Encryption Using Ideal Lattices.''
  \emph{STOC 2009}.
\item
  Kerber, Thomas, Aggelos Kiayias, and Markulf Kohlweiss. ``Kachina --
  Foundations of Private Smart Contracts.'' \emph{IEEE CSF 2021}.
\item
  Buterin, Vitalik, et al.~``Blockchain Privacy and Regulatory Compliance:
  Towards a Practical Equilibrium.'' \emph{Blockchain: Research and
  Applications}, 2023. SSRN 4563364.
\end{enumerate}

\subsection{Chapters 10-11: Synthesis and
zkVMs}\label{chapters-10-11-synthesis-and-zkvms}

\begin{enumerate}
\def\labelenumi{\arabic{enumi}.}
\setcounter{enumi}{31}
\tightlist
\item
  Gassmann, Thomas, et al.~``Evaluating Compiler Optimization Impacts on zkVM
  Performance.'' arXiv 2508.17518, 2026.
\item
  Ozdemir, Alex, Fraser Brown, and Riad Wahby. ``CirC: Compiler Infrastructure
  for Proof Systems, Software Verification, and More.'' \emph{IEEE S\&P 2022}.
  ePrint 2020/1586.
\item
  Liu, Junrui, et al.~``Certifying Zero-Knowledge Circuits with Refinement Types
  (Coda).'' \emph{IEEE S\&P 2024}. ePrint 2023/547.
\item
  Maller, Mary, Sean Bowe, Markulf Kohlweiss, and Sarah Meiklejohn. ``Sonic:
  Zero-Knowledge SNARKs from Linear-Size Universal and Updatable Structured
  Reference Strings.'' CCS 2019. ePrint 2019/099.
\item
  Groth, Jens, Markulf Kohlweiss, Mary Maller, Sarah Meiklejohn, and Ian Miers.
  ``Updatable and Universal Common Reference Strings with Applications to
  zk-SNARKs.'' \emph{CRYPTO 2018}. ePrint 2018/280.
\item
  Kohlweiss, Markulf, Mary Maller, Janno Siim, and Mikhail Volkhov. ``Snarky
  Ceremonies.'' \emph{ASIACRYPT 2021}. ePrint 2021/219.
\item
  Kim, Taechan and Razvan Barbulescu. ``Extended Tower Number Field Sieve: A New
  Complexity for the Medium Prime Case.'' \emph{CRYPTO 2016} (LNCS 9814,
  pp.~543--571). ePrint 2015/1027.
\item
  Guillevic, Aurore. ``Comparing the Pairing Efficiency over Composite-Order and
  Prime-Order Elliptic Curves.'' \emph{ACNS 2013}. ePrint 2013/218.
\item
  Succinct Labs. ``SP1 Hypercube: Proving Ethereum in Real-Time.'' Blog post,
  May 2025. https://blog.succinct.xyz/sp1-hypercube/
\item
  ZKsync. ``Airbender: GPU-Accelerated RISC-V Proving.'' Product announcement,
  June 2025. https://www.zksync.io/airbender
\item
  Kadianakis, George. ``Shipping an L1 zkEVM \#2: The Security Foundations.''
  Ethereum Foundation Blog, December 2025.
  https://blog.ethereum.org/2025/12/18/zkevm-security-foundations
\item
  Chen, Binyi. ``Symphony: Scalable SNARKs in the Random Oracle Model from
  Lattice-Based High-Arity Folding.'' ePrint 2025/1905.
\end{enumerate}

\subsection{Chapter 12: Midnight}\label{chapter-12-midnight}

\begin{enumerate}
\def\labelenumi{\arabic{enumi}.}
\setcounter{enumi}{38}
\tightlist
\item
  Midnight Network. ``Compact Language Reference.'' 2025.
  https://docs.midnight.network/develop/reference/
\item
  Midnight Network. ``ZKIR Intermediate Representation Reference.'' 2025.
  https://docs.midnight.network/develop/reference/
\item
  Midnight Network. ``Developer Guide.'' 2025. https://docs.midnight.network/
\end{enumerate}

\subsection{Chapter 13: The Market
Landscape}\label{chapter-13-the-market-landscape-1}

\begin{enumerate}
\def\labelenumi{\arabic{enumi}.}
\setcounter{enumi}{41}
\tightlist
\item
  Grand View Research. ``Zero-Knowledge Proof Market Size, Share \& Trends
  Analysis Report.'' Report GVR-4-68040-808-5, 2025.
  https://www.grandviewresearch.com/industry-analysis/zero-knowledge-proof-market-report
\item
  CastleLabs. ``ZK Proofs: Is Privacy Cheap Enough to Be Mainstream?'' 2025.
  https://research.castlelabs.io
\item
  Ethproofs. ``ZK Proving Cost Tracker.'' Ethereum Foundation, 2025.
  https://ethproofs.org
\item
  Klich, Rafal (Chorus One). ``The Economics of ZK-Proving: Market Size and
  Future Projections.'' Research report, March 2025.
\item
  DTCC. ``DTCC and Digital Asset Tokenize US Treasuries on Canton Network.''
  Press release, December 2025. https://www.dtcc.com/digital-assets/tokenization
\item
  Tools for Humanity (World). ``World Whitepaper.'' 2024.
  https://whitepaper.world.org/
\item
  European Union. ``Regulation (EU) 2024/1183 -- European Digital Identity
  Framework (eIDAS 2.0).'' \emph{Official Journal of the European Union}, 2024.
\end{enumerate}

\subsection{Chapter 14: Open Questions}\label{chapter-14-open-questions}

\begin{enumerate}
\def\labelenumi{\arabic{enumi}.}
\setcounter{enumi}{44}
\tightlist
\item
  Nair, Vineet, Justin Thaler, and Michael Zhu. ``Proving CPU Executions in
  Small Space.'' ePrint 2025/611.
\item
  Ozdemir, Alex, Evan Laufer, and Dan Boneh. ``Volatile and Persistent Memory
  for zkSNARKs via Algebraic Interactive Proofs.'' \emph{IEEE S\&P 2025}. ePrint
  2024/979.
\item
  Pailoor, Shankara, et al.~``Automated Detection of Under-Constrained Circuits
  in Zero-Knowledge Proofs.'' \emph{PLDI 2023}. ePrint 2023/512.
\item
  Wen, Hongbo, et al.~``Practical Security Analysis of Zero-Knowledge Proof
  Circuits.'' \emph{USENIX Security 2024}. ePrint 2023/190.
\item
  Hochrainer, Christoph, Valentin Wustholz, and Maria Christakis. ``Arguzz:
  Testing zkVMs for Soundness and Completeness Bugs.'' arXiv 2509.10819, 2025.
\item
  Wee, Hoeteck and David J. Wu. ``Lattice-Based Functional Commitments: Fast
  Verification and Cryptanalysis.'' \emph{ASIACRYPT 2023}. ePrint 2024/028.
\item
  Mascelli, Jillian and Megan Rodden. ``\,`Harvest Now Decrypt Later': Examining
  Post-Quantum Cryptography and the Data Privacy Risks for Distributed Ledger
  Networks.'' FEDS 2025-093, Federal Reserve Board, 2025.
\item
  NIST. ``Module-Lattice-Based Digital Signature Standard (FIPS 204).'' August
  2024.
\end{enumerate}

\end{document}
